\documentclass[11pt,a4paper,oneside]{book}

% Preamble for Bosonize-Lean Pedagogical Monograph
% Plain LaTeX standard configuration without custom macro layers

\usepackage[utf8]{inputenc}
\usepackage[T1]{fontenc}
\usepackage{lmodern}

% Standard AMS mathematical typesetting
\usepackage{amsmath}
\usepackage{amssymb}
\usepackage{amsthm}
\usepackage{mathtools}

% Page layout and typography
\usepackage[a4paper, margin=2.5cm]{geometry}
\usepackage{microtype}
\usepackage{booktabs}
\usepackage{xcolor}
\usepackage{tikz}
\usetikzlibrary{decorations.pathreplacing}

% Hyperlinks
\usepackage[
    colorlinks=true,
    linkcolor=blue!75!black,
    citecolor=green!50!black,
    urlcolor=blue!75!black,
    pdfauthor={Bosonize-Lean Project},
    pdftitle={Constructive 1+1D Lattice Bosonization}
]{hyperref}

% Standard AMS theorem environments numbered within chapters
\theoremstyle{plain}
\newtheorem{theorem}{Theorem}[chapter]
\newtheorem{lemma}[theorem]{Lemma}
\newtheorem{proposition}[theorem]{Proposition}
\newtheorem{corollary}[theorem]{Corollary}

\theoremstyle{definition}
\newtheorem{definition}[theorem]{Definition}
\newtheorem{example}[theorem]{Example}

\theoremstyle{remark}
\newtheorem{remark}[theorem]{Remark}
\newtheorem{note}[theorem]{Note}

% CI supplies a revision stamp; local builds remain reproducible without it.
\InputIfFileExists{build-info.tex}{}{}
\providecommand{\BuildRevision}{local-development}


\title{\textbf{1+1D Lattice Bosonization}\\
\Large Explicit Finite-Lattice and Finite-Budget Models}
\author{\textbf{Bosonize-Lean Project}\\
\normalsize A Pedagogical Monograph}
\date{\today\\\small Source revision: \texttt{\BuildRevision}}

\begin{document}

\hypersetup{pageanchor=false}
\maketitle
\thispagestyle{empty}
\hypersetup{pageanchor=true}

\frontmatter

\chapter*{Preface}
\markboth{Preface}{Preface}

This monograph presents a pedagogical rederivation of selected results from the $1+1$-dimensional bosonization literature. We follow the standard bosonization program, keeping the lattice size and the energy restrictions explicit. At fixed size, the fermionic constructions use finite-dimensional CAR Fock spaces and finite sums. Some bosonic identities are most naturally expressed on an ambient polynomial algebra; we state separately when those identities apply to a finite energy budget.

This presentation is intended to make the algebra and its assumptions accessible. By maintaining explicit finite lattices and bounded cutoffs, we highlight algebraic structures such as boundary holonomies, Umklapp wrapping, and operator domain restrictions. However, this explicit finite framework also imposes limitations. It does not by itself establish a continuum limit, a thermodynamic limit, or replace the broader analytic treatments found in the established literature.

The material presented in Chapters~\ref{ch:lattice_band} through~\ref{ch:haldane_completeness} and Appendices~\ref{app:fourier_characters} through~\ref{app:density_partitions} corresponds to an implemented Lean 4 formalization. Note that Lean proof coverage applies strictly to the actual mathematical declarations and their explicit hypotheses, not automatically to every explanatory sentence in the text.

The later Chapters~\ref{ch:correlations}--\ref{ch:spinful} are manuscript drafts awaiting independent mathematical and editorial review. Chapter~19 uses the accepted Phase-B CH19/A08 development, whose Core promotion remains open. Chapters~20 and~21 do not yet have implemented chapter Lean modules. Their mathematical derivations, formal-series constructions, and finite comparisons must therefore be read with the scope and limitations stated in those chapters; compilation of this book does not certify their mathematics. Chapters~12--18 have separate private drafts and are not included in this edition, leaving an expository and prerequisite gap before reaching Chapter~19, although their implemented Lean formalizations are already promoted.

These notes are designed to be read by readers familiar with undergraduate quantum mechanics and linear algebra. Where our expository choices deviate from convention, or where specific structural insights emerge from the finite setting, dedicated \emph{Formalization Notes} discuss the exact definitions employed.

\tableofcontents

\chapter*{Notation and Conventions}
\addcontentsline{toc}{chapter}{Notation and Conventions}
\markboth{Notation and Conventions}{Notation and Conventions}

This monograph presents a completely self-contained, exact, and constructive formulation of $1+1$-dimensional lattice bosonization within algebraic quantum field theory (AQFT). Unlike traditional continuum heuristic treatments that rely on divergent operator products, formal distribution limits, or unbounded operators, all mathematical statements established here are exact theorems on strictly finite-dimensional carrier spaces, finite lattices, or explicitly graded polynomial algebras.

To maintain clarity and consistency across all chapters, this section establishes the unified mathematical notation used throughout the text.

\section*{Lattice Geometry and Discrete Kinematics}

\begin{itemize}
    \item \textbf{Lattice size:} $L \in \mathbb{N}$ denotes the number of discrete spatial lattice sites, strictly positive ($L \ge 1$). Throughout the fermionic and bosonization development, $L$ is assumed to be an even integer, parametrized by a half-length $h \ge 1$ such that $L = 2h$.
    \item \textbf{Spatial lattice:} $\Lambda = \mathbb{Z}/L\mathbb{Z}$ is the discrete 1D periodic chain of sites $x, y \in \Lambda$, endowed with the standard quotient group structure modeling periodic boundary conditions.
    \item \textbf{Momentum band:} $\Lambda^*$ denotes the centered discrete Brillouin zone of $L$ integers:
    \[
    \Lambda^* = \left\{ k \in \mathbb{Z} \;\middle|\; -L < 2k \le L \right\}.
    \]
    When $L = 2h$, this interval is explicitly $\Lambda^* = \{-h+1, -h+2, \dots, h-1, h\}$. This choice retains the positive Nyquist mode $k = +h = L/2$ and excludes $-h$, ensuring an exact bijection with $\Lambda$ while deliberately breaking reflection symmetry.
    \item \textbf{Projection and Section:}
    \begin{itemize}
        \item $\pi : \mathbb{Z} \to \Lambda$ denotes the canonical quotient projection $\pi(k) = k \bmod L$.
        \item $r : \Lambda \to \Lambda^*$ denotes the unique centered section of $\pi$ landing in $\Lambda^*$.
    \end{itemize}
    \item \textbf{Transported band arithmetic:}
    \[
    k \oplus p = r(\pi(k) + \pi(p)), \qquad \ominus k = r(-\pi(k)).
    \]
    The discrepancy between integer addition $k+p$ and band addition $k \oplus p$ is governed by the Umklapp integer $w \in \{-1, 0, 1\}$ such that $k+p = (k \oplus p) + wL$.
\end{itemize}

\section*{Difference Calculus and Fourier Analysis}

\begin{itemize}
    \item \textbf{Shift and differences:} For functions $f : \Lambda \to \mathbb{C}$, the shift operator is $(E f)(x) = f(x+1)$, the forward difference is $\Delta = E - I$, and the backward difference is $\nabla = I - E^{-1}$.
    \item \textbf{Fourier transform:} For $k \in \Lambda^*$ and $x \in \Lambda$, the character is $\chi(k,x) = \exp(2\pi i k x / L)$. The discrete Fourier transform $\mathcal{F}$ maps spatial functions to momentum space unitarily, satisfying exact Parseval and Plancherel identities over the finite cyclic group.
\end{itemize}

\section*{Fermionic Fock Spaces and Observable Nets}

\begin{itemize}
    \item \textbf{CAR Algebra:} The canonical anticommutation relations on the momentum band $\Lambda^*$ are generated by operators $c_k, c_k^\dagger$ satisfying:
    \[
    \{c_k, c_p^\dagger\} = \delta_{k,p} I, \qquad \{c_k, c_p\} = 0, \qquad \{c_k^\dagger, c_p^\dagger\} = 0, \quad \forall k, p \in \Lambda^*.
    \]
    \item \textbf{Spatial fermions:} Constructed via the finite discrete Fourier transform:
    \[
    \psi_x = \frac{1}{\sqrt{L}} \sum_{k \in \Lambda^*} \chi(-k, x)\, c_k, \qquad \psi_x^\dagger = \frac{1}{\sqrt{L}} \sum_{k \in \Lambda^*} \chi(k, x)\, c_k^\dagger.
    \]
    \item \textbf{Local nets:} For an interval $I \subset \Lambda$, $\mathfrak{A}(I)$ is the von Neumann / $C^*$-subalgebra generated by $\{\psi_x, \psi_x^\dagger \mid x \in I\}$.
\end{itemize}

\section*{Vacuum Structure, Charge Sectors, and Energy Budgets}

\begin{itemize}
    \item \textbf{Dirac sea vacuum:} $|\Omega\rangle$ denotes the half-filled reference sea configuration $S_0 = \{-h+1, \dots, 0\}$.
    \item \textbf{Admissible charge sectors:} For $N \in \mathbb{Z}$ with $-h \le N \le h$, the sector ground configuration is $S_N = \{-h+1, \dots, N\}$, with ground state $|N\rangle_0$ of total particle number $h + N$.
    \item \textbf{Excitation energy:} For an occupation configuration $\sigma \subset \Lambda^*$ of cardinality $h+N$, the excitation energy above the sector ground is:
    \[
    E(\sigma) = \sum_{k \in \sigma} k - \sum_{k \in S_N} k \ge 0.
    \]
    \item \textbf{Energy budget subspace:} For sector $N$ and energy cutoff $K \ge 0$:
    \[
    B(N, K) = \bigoplus_{E=0}^K \mathcal{H}(N, E),
    \]
    where $\mathcal{H}(N,E)$ is the subspace spanned by configurations with charge $N$ and excitation energy $E$.
\end{itemize}

\section*{Density Modes and Bosonic Representation}

\begin{itemize}
    \item \textbf{Normal-ordered density modes:} For mode integer $m \neq 0$:
    \[
    :\!\rho_m\!: = \sum_{k \in \Lambda^*} : \! c_{k \ominus m}^\dagger c_k \!:
    \]
    \item \textbf{Heisenberg / Kac--Moody algebra:} Under explicit margin conditions ($2M + K + |N| \le h$), the density modes satisfy the exact commutation relations on $B(N,K)$:
    \[
    \left[ :\!\rho_m\!:, :\!\rho_{-n}\!: \right] = m\, \delta_{m,n}\, I.
    \]
    \item \textbf{Bosonic partition states:} For an integer partition $\lambda = (\lambda_1 \ge \lambda_2 \ge \dots \ge \lambda_\ell > 0)$ of size $|\lambda| = \sum \lambda_i = E$:
    \[
    |\lambda; N\rangle = :\! \rho_{\lambda_1} \!: :\! \rho_{\lambda_2} \!: \dots :\! \rho_{\lambda_\ell} \!: |N\rangle_0.
    \]
    \item \textbf{Haldane completeness:} The states $\{|\lambda; N\rangle \mid |\lambda| = E\}$ form an exact basis of $\mathcal{H}(N, E)$, and $\bigcup_{E=0}^K \{|\lambda; N\rangle \mid |\lambda| = E\}$ forms a basis of $B(N, K)$.
\end{itemize}

\section*{Conventions on Constructive Realization}

Every mathematical theorem in this monograph corresponds to a fully verified, kernel-checked theorem in the Lean 4 formalization. Rather than polluting the mathematical discourse with code identifiers, tactic transcripts, or backticks, the text presents rigorous, complete human-readable proofs. 

Where constructive methods deviate from traditional non-constructive continuum habits (such as the choice of asymmetric Nyquist modes, finite decidable filters, quotient injectivity, or polynomial state carriers), the text highlights the underlying mathematical rationale in dedicated \emph{Formalization Note} remarks.


\mainmatter

\part{Discrete Foundations and the Fermionic Fock Net}

\chapter{Lattice and Band Geometry}
\label{ch:lattice_band}

\section{Introduction and Discrete Setting}

In $1+1$-dimensional field theory, continuum descriptions can produce ultraviolet and infrared divergences, as well as infinite mode sums that require separate regularization. Here we begin with a finite periodic chain and its exact Fourier-dual momentum band. The resulting operators and sums are finite, and the algebraic constraints imposed by periodicity can be handled before any continuum limit is considered.

Throughout this chapter, let $L \in \mathbb{N}$ denote the number of spatial lattice sites, assuming $L \ge 1$. While general difference calculus and Fourier transforms can be formulated for arbitrary positive integers $L$, the construction of half-filled Dirac seas and fermionic Fock spaces from Chapter~\ref{ch:fermions} onward strictly requires an even lattice length:
\begin{equation}
L = 2h, \quad h \in \mathbb{N}, \quad h \ge 1.
\end{equation}
The integer $h$ represents the half-filling particle number or the positive Nyquist cut-off.

\section{The Spatial Lattice}

The spatial geometry of the physical system is that of a discrete one-dimensional circular chain of $L$ sites with periodic boundary conditions.

\begin{definition}[Spatial Lattice]
The spatial lattice $\Lambda$ is defined as the cyclic quotient group of integers modulo $L$:
\begin{equation}
\Lambda \coloneqq \mathbb{Z}/L\mathbb{Z}.
\end{equation}
Elements $x, y \in \Lambda$ represent discrete spatial positions. Periodic boundary conditions are an intrinsic algebraic property of the quotient ring: translation by $L$ sites corresponds to the identity operation in $\Lambda$.
\end{definition}

Because $L \ge 1$, the spatial lattice $\Lambda$ is a finite commutative ring of cardinality $\#\Lambda = L$, endowed with canonical addition and multiplication inherited from $\mathbb{Z}$.

\begin{remark}[Finite-Lattice Remark]
The quotient description makes periodicity exact: adding $L$ to an integer representative leaves its site unchanged. Since $L$ is positive, this finite ring has exactly $L$ sites, and all position labels and translations are interpreted within it.
\end{remark}

\paragraph{Formal Correctness vs.\ Physical Implications: Discrete Chains vs.\ Continuum Circles.}
In condensed matter physics, a circular chain is often treated in the continuum limit by taking the lattice spacing $a \to 0$ and the number of sites $L \to \infty$ while holding the macroscopic circumference $\mathcal{L} = L a$ fixed. In that limit, spatial positions $x a$ tend to continuous coordinates on the circle $\mathbb{R}/\mathcal{L}\mathbb{Z}$, and fields become continuous functions or operator-valued distributions. 
In our exact algebraic framework, by contrast, the lattice spacing $a$ is set to $1$ without loss of generality, and all positions $x \in \Lambda$ are exact residue classes. This eliminates distribution-theoretic smearing and guarantees that every spatial operator algebra is strictly finite-dimensional.

\section{The Momentum Band and Nyquist Asymmetry}

In continuous quantum field theory, the dual momentum space is the real line $\mathbb{R}$ or the continuous circle $\mathbb{R}/(2\pi/a)\mathbb{Z}$. On a discrete spatial lattice of $L$ sites, the dual momentum band (the first Brillouin zone) must consist of exactly $L$ discrete momentum modes.

To center the physical momentum around zero while maintaining exact finite cardinality, we define the momentum band as an integer interval centered near the origin.

\begin{definition}[Momentum Band]
The momentum band $\Lambda^*$ is the subset of signed integers satisfying the two-sided inequalities:
\begin{equation}
\Lambda^* \coloneqq \left\{ k \in \mathbb{Z} \;\middle|\; -L < 2k \le L \right\}.
\end{equation}
Momenta are denoted by variables $k, p, q \in \Lambda^*$. The physical wavenumber corresponding to mode index $k$ is given by the scaled wavenumber:
\begin{equation}
p_{\mathrm{phys}} \coloneqq \frac{2\pi k}{L}.
\end{equation}
\end{definition}

The choice of inequalities $-L < 2k \le L$ is subtle and fundamental:
\begin{itemize}
    \item When $L = 2h$ is even, the condition $-2h < 2k \le 2h$ simplifies to $-h < k \le h$. Consequently, the momentum band consists explicitly of the $2h$ consecutive integers:
    \begin{equation}
    \Lambda^* = \{-h+1, -h+2, \dots, -1, 0, 1, \dots, h-1, h\}.
    \end{equation}
    \item The interval is intrinsically asymmetric: it includes the positive Nyquist mode $k = +h = +L/2$, but excludes the negative counterpart $k = -h = -L/2$.
    \item When $L$ is odd, say $L = 2m+1$, the condition becomes $-2m-1 < 2k \le 2m+1$, which yields the symmetric set $\{-m, -m+1, \dots, m-1, m\}$ of cardinality $2m+1$.
\end{itemize}

\paragraph{Formal Correctness vs.\ Physical Implications: Brillouin Zone Boundaries.}
In standard solid-state physics textbooks (such as Ashcroft and Mermin or Kittel), the first Brillouin zone is customarily drawn as the symmetric interval $[-\pi/a, \pi/a]$, or $[-h, h]$ in mode units. However, on an even lattice of $L = 2h$ sites, the modes $k = -h$ and $k = +h$ represent the \emph{same} physical wave on the lattice, since $e^{i (-h) (2\pi x / 2h)} = e^{-i \pi x} = (-1)^x = e^{i \pi x} = e^{i (+h) (2\pi x / 2h)}$.
Including both $-h$ and $+h$ yields $2h+1$ modes, which violates the fundamental dimension theorem $\dim(\ell^2(\Lambda)) = \dim(\ell^2(\Lambda^*)) = L$. Physics literature often circumvents this by assigning fractional weights $1/2$ to the boundary modes or passing implicitly to the thermodynamic limit. 
The discrete Fourier transform pairs spaces of the same finite dimension $L$. Retaining the positive Nyquist mode $+h$ while excluding $-h$ gives exactly $\#\Lambda^* = L$ distinct periodic modes.

\begin{remark}[Finite-Lattice Remark]
The band is an explicitly bounded set of integer modes, with exactly one representative of each residue class modulo $L$. Its asymmetric endpoints record the single positive Nyquist mode on an even lattice.
\end{remark}

\section{Quotient Projections and the Centered Section}

Because $\Lambda = \mathbb{Z}/L\mathbb{Z}$ is a quotient group, there is a natural projection from any integer to its equivalence class in the lattice.

\begin{definition}[Canonical Projection]
The canonical projection $\pi : \mathbb{Z} \to \Lambda$ maps an integer $k \in \mathbb{Z}$ to its periodic residue class:
\begin{equation}
\pi(k) \coloneqq k \bmod L.
\end{equation}
By restriction, $\pi$ maps any momentum mode $k \in \Lambda^*$ to an element $\pi(k) \in \Lambda$.
\end{definition}

To construct the inverse mapping, we seek a section of $\pi$ whose range lies strictly within the momentum band $\Lambda^*$. Every element $x \in \Lambda = \mathbb{Z}/L\mathbb{Z}$ possesses a unique canonical representative $x_{\mathrm{val}} \in \{0, 1, \dots, L-1\}$.

\begin{definition}[Centered Lift Section]
The centered lift map $r : \Lambda \to \Lambda^*$ is defined piecewise on the unsigned representative $x_{\mathrm{val}} \in \{0, 1, \dots, L-1\}$ by:
\begin{equation}
r(x) \coloneqq \begin{cases}
x_{\mathrm{val}}, & \text{if } 2 x_{\mathrm{val}} \le L, \\
x_{\mathrm{val}} - L, & \text{if } 2 x_{\mathrm{val}} > L.
\end{cases}
\end{equation}
\end{definition}

\paragraph{Physical Intuition for the Centered Lift.}
The map $r$ takes a non-negative residue $x_{\mathrm{val}} \in [0, L-1]$ and centers it around zero: values in the upper half of the range $[h+1, 2h-1]$ are translated downward by $L$ to become negative integers in $[-h+1, -1]$, corresponding to negative-frequency or left-moving momentum modes in physical terminology.

\begin{lemma}[Well-Definedness and Inversion]
\label{lem:section_identities}
For any positive integer $L \ge 1$:
\begin{enumerate}
    \item For every $x \in \Lambda$, the value $r(x)$ belongs to the momentum band $\Lambda^*$, and:
    \begin{equation}
    \pi(r(x)) = x.
    \end{equation}
    \item For every $k \in \Lambda^*$, the centered lift recovers the original mode:
    \begin{equation}
    r(\pi(k)) = k.
    \end{equation}
\end{enumerate}
\end{lemma}

\begin{proof}
We verify both statements directly from the definitions:
\begin{enumerate}
    \item Let $x \in \Lambda$ and let $a = x_{\mathrm{val}} \in \{0, 1, \dots, L-1\}$.
    \begin{itemize}
        \item If $2a \le L$, then $r(x) = a$. Since $a \ge 0$, we have $2a \ge 0 > -L$. Since $2a \le L$, $a$ satisfies $-L < 2a \le L$, so $r(x) \in \Lambda^*$. Furthermore, $\pi(r(x)) = \pi(a) = x$.
        \item If $2a > L$, then $r(x) = a - L$. Since $a \le L - 1$, we have $a - L \le -1$, so $2(a - L) \le -2 < 0 \le L$. Conversely, $2a > L$ implies $2(a - L) > -L$. Thus $-L < 2(a - L) \le L$, establishing that $r(x) \in \Lambda^*$. Moreover, $\pi(a - L) = \pi(a) - \pi(L) = x - 0 = x$.
    \end{itemize}
    In both branches, $r(x) \in \Lambda^*$ and $\pi(r(x)) = x$.
    
    \item Let $k\in\Lambda^*$. By part (1), $k$ and $r(\pi(k))$ have the same image under $\pi$, so their difference is an integer multiple of $L$. Since both lie in the band $(-L/2,L/2]$, their difference has absolute value strictly less than $L$. The only such multiple is zero, hence $r(\pi(k))=k$.
\end{enumerate}
\end{proof}

\begin{corollary}[Canonical Bijection]
The restriction of the canonical projection $\pi|_{\Lambda^*} : \Lambda^* \to \Lambda$ is a bijection, whose two-sided inverse is the centered lift section $r : \Lambda \to \Lambda^*$. In particular, the cardinality of the momentum band is strictly:
\begin{equation}
\#\Lambda^* = \#\Lambda = L.
\end{equation}
\end{corollary}

\begin{remark}[Finite-Lattice Remark]
For injectivity, take any $k_1,k_2\in\Lambda^*$ with $\pi(k_1)=\pi(k_2)$. Their difference is an integer multiple of $L$. Both lie in $(-L/2,L/2]$, so the diameter of the band is strictly less than $L$ and $|k_1-k_2|<L$. The difference must therefore be zero. Thus $k_1=k_2$, and the section identities give the asserted inversion.
\end{remark}

\section{Transported Band Arithmetic and Umklapp Wrapping}

While $\Lambda = \mathbb{Z}/L\mathbb{Z}$ is an abelian group under addition, the momentum band $\Lambda^*$ is an integer subset. The usual sum of two momenta $k, p \in \Lambda^*$ may lie outside the interval $(-L/2, L/2]$. 

To endow $\Lambda^*$ with a group structure isomorphic to $\Lambda$, we transport the addition and negation of the cyclic group through the bijection $(\pi, r)$.

\begin{definition}[Transported Band Arithmetic]
For any momenta $k, p \in \Lambda^*$, we define band addition $\oplus$ and band negation $\ominus$ by:
\begin{align}
k \oplus p &\coloneqq r(\pi(k) + \pi(p)), \\
\ominus k &\coloneqq r(-\pi(k)).
\end{align}
\end{definition}

\begin{remark}[Notation in later chapters]
The symbols $\oplus$ and $\ominus$ emphasize that these operations are transported from the quotient group to the chosen momentum band. In subsequent chapters, standard $+$ and $-$ symbols will denote these band operations when their domain is clear, reducing notational overhead.
\end{remark}

Because $(\pi, r)$ is a bijection between $\Lambda^*$ and the abelian group $\Lambda$, $(\Lambda^*, \oplus, \ominus, 0)$ forms a finite abelian group isomorphic to $\mathbb{Z}/L\mathbb{Z}$, with neutral element $0 = r(\pi(0))$.

\paragraph{Physical Intuition: Umklapp Scattering on a Lattice.}
In solid-state physics, when two quasi-particles with crystal momenta $k$ and $p$ scatter, their total momentum $k+p$ can exceed the boundary of the first Brillouin zone. The crystal lattice absorbs a reciprocal lattice vector $G = \pm L$, scattering the particles back into the first Brillouin zone with crystal momentum $k \oplus p = k + p \mp L$. This physical phenomenon is known as an \emph{Umklapp process}. Theorem~\ref{thm:wrap_lemma} makes this exact algebraically: the integer difference between arithmetic addition and band addition is always $w L$ with $w \in \{-1, 0, 1\}$.

\begin{theorem}[Umklapp Wrapping Lemma]
\label{thm:wrap_lemma}
For any $k, p \in \Lambda^*$, there exists a unique integer $w \in \{-1, 0, 1\}$ such that:
\begin{equation}
k + p = (k \oplus p) + w L.
\end{equation}
The integer $w$ is the Umklapp multiplier and is explicitly determined by:
\begin{equation}
w = \begin{cases}
+1, & \text{if } 2(k+p) > L, \\
0, & \text{if } -L < 2(k+p) \le L, \\
-1, & \text{if } 2(k+p) \le -L.
\end{cases}
\end{equation}
\end{theorem}

\begin{proof}
Let $s=k+p$. Since $k,p\in\Lambda^*$, their band bounds give $-L<s\le L$. By definition, $k\oplus p$ and $s$ have the same image under $\pi$, so $s-(k\oplus p)=wL$ for some integer $w$. The bounds $-L<s\le L$ and $-L/2<k\oplus p\le L/2$ imply $-3L/2<s-(k\oplus p)<3L/2$, hence $-3/2<w<3/2$. Therefore $w\in\{-1,0,1\}$.
By the defining property of the centered lift, $k\oplus p=r(\pi(s))$ is the unique element of $\Lambda^*$ with quotient image $\pi(s)$. Thus the integer $w$ above is uniquely determined by $s-(k\oplus p)=wL$. Its value must be one of $-1,0,1$, as established by the bounds.
Uniqueness of $w$ follows since $(w_1 - w_2) L = 0$ implies $w_1 = w_2$ when $L > 0$.
\end{proof}

\section{The Exceptional Nyquist Mode}

We now examine the behavior of band negation $\ominus k$ when the lattice length is even, $L = 2h$.

\begin{theorem}[Nyquist Mode Negation]
\label{thm:nyquist_neg}
Let $L = 2h$ with $h \ge 1$.
\begin{enumerate}
    \item The positive Nyquist mode $k = h = L/2$ is self-inverse under band negation:
    \begin{equation}
    \ominus \left(\frac{L}{2}\right) = \frac{L}{2}.
    \end{equation}
    \item For any mode $k \in \Lambda^*$ distinct from the Nyquist mode ($k \neq L/2$), band negation coincides with ordinary integer arithmetic negation:
    \begin{equation}
    \ominus k = -k.
    \end{equation}
\end{enumerate}
\end{theorem}

\begin{proof}
Let $L = 2h$.
\begin{enumerate}
    \item For the mode $k = L/2 = h$, ordinary negation yields $-k = -h = -L/2$. In the quotient ring $\Lambda = \mathbb{Z}/L\mathbb{Z}$:
    \begin{equation}
    -\pi(L/2) = \pi(-L/2) = \pi(-L/2 + L) = \pi(L/2).
    \end{equation}
    Applying the section $r$, we have:
    \begin{equation}
    \ominus (L/2) = r(-\pi(L/2)) = r(\pi(L/2)) = L/2,
    \end{equation}
    by Lemma~\ref{lem:section_identities}(2). Hence $\ominus(L/2) = L/2$.

    \item Suppose $k \in \Lambda^*$ with $k \neq L/2$. Since $k \in \{-h+1, \dots, h\}$ and $k \neq h$, we must have:
    \begin{equation}
    -h+1 \le k \le h-1 \implies -(h-1) \le -k \le h-1.
    \end{equation}
    Multiplying by $2$:
    \begin{equation}
    -2h < -2(h-1) \le 2(-k) \le 2(h-1) < 2h = L.
    \end{equation}
    Thus $-k$ satisfies the strict bounds $-L < 2(-k) \le L$, establishing that the integer $-k$ is already in $\Lambda^*$. By Lemma~\ref{lem:section_identities}(2), $r(\pi(-k)) = -k$. Therefore:
    \begin{equation}
    \ominus k = r(-\pi(k)) = r(\pi(-k)) = -k.
    \end{equation}
\end{enumerate}
\end{proof}

\begin{remark}[Finite-Lattice Remark]
The distinction between band negation and ordinary integer negation matters at the boundary: the positive Nyquist mode is self-inverse, while ordinary negation would leave the chosen band. This exact accounting is important in Fourier inversion and in transformations involving paired modes. Specifically:
\begin{itemize}
    \item The pair $(k, \ominus k)$ contains two distinct modes for all $k \notin \{0, L/2\}$.
    \item The fixed points of the negation involution $k \mapsto \ominus k$ on $\Lambda^*$ are precisely the two zero-energy/boundary modes: $k = 0$ and $k = L/2$.
\end{itemize}
This exact accounting ensures that no modes are double-counted or omitted when momentum sums are partitioned into pairs in later chapters.
\end{remark}

\section{Chapter Summary and Literature Comparison}

In this opening chapter, we established the discrete spatial and momentum geometry that underlies all subsequent field-theoretic constructions:
\begin{itemize}
    \item The spatial lattice is the finite cyclic group $\Lambda = \mathbb{Z}/L\mathbb{Z}$ ($L = 2h$).
    \item The momentum band is the asymmetric Brillouin zone $\Lambda^* = \{-h+1, \dots, h\}$, retaining $+h$ and excluding $-h$.
    \item The canonical projection $\pi$ and centered lift section $r$ form a two-sided bijection $\Lambda^* \cong \Lambda$.
    \item Group arithmetic on $\Lambda^*$ is transported via $(\pi, r)$, with integer discrepancies governed by the Umklapp wrapping multiplier $w \in \{-1, 0, 1\}$.
    \item The positive Nyquist mode $k = h$ is the unique non-zero self-inverse mode under band negation $\ominus$.
\end{itemize}

\paragraph{Comparison with Standard Continuum Literature.}
In standard condensed matter bosonization texts, such as Giamarchi~\cite{giamarchi2004}, Gogolin, Nersesyan, and Tsvelik~\cite{gogolin1998}, and von Delft and Schoeller~\cite{vondelft1998}, the spatial circle is typically taken as a continuous interval $[0, \mathcal{L}]$ with periodic boundary conditions, leading to an infinite set of discrete momentum modes $k_n = \frac{2\pi n}{\mathcal{L}}$ ($n \in \mathbb{Z}$). In that setting, there is no upper momentum cutoff, Umklapp scattering is treated as an extrinsic perturbation or ignored in low-energy Tomonaga--Luttinger models, and the Fermi sea has infinite depth.

By contrast, our finite-lattice formulation models the physical lattice from the outset. As emphasized by Haldane~\cite{haldane1981}, the lattice acts as an intrinsic ultraviolet regulator. The asymmetric Brillouin zone $\Lambda^*$ is the unique choice that simultaneously achieves:
\begin{enumerate}
    \item Exact finite dimension: $\dim(\ell^2(\Lambda)) = \dim(\ell^2(\Lambda^*)) = L$.
    \item Absence of duplicate physical modes.
    \item Unambiguous arithmetic at the band boundaries.
\end{enumerate}
In later chapters (specifically Chapter~\ref{ch:heisenberg_algebra} and Chapter~\ref{ch:haldane_completeness}), we will demonstrate that when particle excitations are restricted to an energy budget margin $2M + K + |N| \le h$, the Umklapp wrapping $w$ vanishes identically ($w = 0$), and the exact lattice dynamics reproduces the continuum current algebra of Giamarchi~\cite{giamarchi2004} and von Delft--Schoeller~\cite{vondelft1998} with zero error.

\chapter{Umbral Calculus Core}
\label{ch:umbral_calculus}

\section{Introduction and Algebraic Setting}
In continuum field theory, spatial derivatives $\partial_x$ act on smooth field operators and generate spatial momentum translations. On a discrete spatial lattice $\Lambda = \mathbb{Z}/L\mathbb{Z}$, continuous derivatives do not exist and must be replaced by exact difference operators. Umbral calculus provides the rigorous algebraic framework for replacing differential operators with finite differences without losing the core algebraic relations of Lie algebras and polynomial expansions.

\section{Core Definitions and Difference Operators}
Let $R$ be a commutative ring (such as $\mathbb{C}$ or $\mathbb{Z}$). The algebra of functions $\mathcal{F}(\Lambda, R)$ is acted upon by three fundamental operators:
\begin{itemize}
    \item \textbf{Shift Operator:} $(E f)(x) \coloneqq f(x+1)$.
    \item \textbf{Forward Difference:} $(\Delta f)(x) \coloneqq (E - I)f(x) = f(x+1) - f(x)$.
    \item \textbf{Backward Difference:} $(\nabla f)(x) \coloneqq (I - E^{-1})f(x) = f(x) - f(x-1)$.
\end{itemize}

\section{Key Results and Formal Discoveries}
\begin{itemize}
    \item \textbf{Discrete Leibniz Rule:} For any functions $f, g \in \mathcal{F}(\Lambda, R)$:
    \[
    \Delta(f g) = (\Delta f) g + E(f) (\Delta g) = (\Delta f) g + f (\Delta g) + (\Delta f)(\Delta g).
    \]
    \item \textbf{Discrete Summation by Parts:} Under periodic boundary conditions, the forward difference $\Delta$ and backward difference $-\nabla$ are exact formal adjoints under the spatial sum $\sum_{x \in \Lambda}$:
    \[
    \sum_{x \in \Lambda} f(x) (\Delta g)(x) = - \sum_{x \in \Lambda} (\nabla f)(x) g(x).
    \]
    \item \textbf{Umbral Commutator and Finite Carrier Obstruction:} While the position multiplier $X$ and derivative $\partial_x$ satisfy $[ \partial_x, X ] = I$ on infinite-dimensional polynomial spaces, no finite-dimensional carrier admits operators satisfying $[A, B] = I$ due to the trace identity $\mathrm{Tr}([A, B]) = 0$. The exact discrete formulation establishes how the finite spatial lattice regularizes this commutator.
\end{itemize}

\begin{remark}[Formalization Note]
In the verified Lean~4 core (Chapter 2, Umbral Calculus), difference operators are defined without division, operating over general commutative rings. Summation by parts is proved strictly constructively through index shifting in $\mathbb{Z}/L\mathbb{Z}$, without topological boundary terms.
\end{remark}

\chapter{The Finite Fourier Transform}
\label{ch:fourier}

\section{Introduction and Algebraic Setting}

The discrete Fourier transform is the central kinematic bridge of quantum field theory on a lattice. It connects the spatial coordinate representation on the periodic lattice $\Lambda = \mathbb{Z}/L\mathbb{Z}$ to the dual momentum representation on the centered band $\Lambda^* = \{-h+1, \dots, h\}$ (for even $L = 2h$).

In standard physics presentations, the Fourier transform is defined immediately on complex functions using transcendental exponentials $e^{2\pi i k x / L}$ and the unitary normalization factor $1/\sqrt{L}$. While physically standard, this conflates two distinct mathematical layers:
\begin{enumerate}
    \item \textbf{The algebraic character layer:} Character orthogonality and algebraic Fourier inversion require only a commutative field $K$ containing a primitive $L$-th root of unity $\zeta$, and the invertibility of the lattice size $(L : K) \neq 0$. Square roots and analytic limits are completely unnecessary.
    \item \textbf{The Hilbert unitary layer:} To make the Fourier transform an exact isometry that preserves Canonical Anticommutation Relations (CAR) between position and momentum fermions without spurious factors of $\sqrt{L}$, we instantiate $K = \mathbb{C}$ and introduce the metric normalization $a = 1/\sqrt{L}$.
\end{enumerate}

In this chapter, we develop both layers in full mathematical rigor, following the verified constructive formalization.

\section{Generic Character Algebra over Fields}

\begin{definition}[Primitive Root of Unity]
Let $K$ be a commutative field, and let $L \ge 1$ be a positive integer. An element $\zeta \in K$ is a primitive $L$-th root of unity if:
\begin{equation}
\zeta^L = 1, \qquad \text{and} \qquad \zeta^m \neq 1 \quad \text{for all } m \in \{1, 2, \dots, L-1\}.
\end{equation}
\end{definition}

When $K = \mathbb{C}$, the canonical primitive $L$-th root of unity is:
\begin{equation}
\zeta_{\mathrm{can}} \coloneqq \exp\left(\frac{2\pi i}{L}\right).
\end{equation}

\begin{definition}[Representative-Independent Character]
For any momentum mode $k \in \Lambda^*$ and spatial site $x \in \Lambda$, let $x_{\mathrm{val}} \in \{0, 1, \dots, L-1\}$ denote the unsigned residue of $x$. The discrete Fourier character is defined by:
\begin{equation}
\chi(k, x) \coloneqq \zeta^{k \cdot x_{\mathrm{val}}}.
\end{equation}
\end{definition}

\begin{lemma}[Representative Independence]
\label{lem:rep_independence}
For any integer $k \in \mathbb{Z}$ and any integer representative $\tilde{x} \in \mathbb{Z}$ of $x$ (so that $\tilde{x} \equiv x_{\mathrm{val}} \pmod L$):
\begin{equation}
\zeta^{k \tilde{x}} = \zeta^{k \cdot x_{\mathrm{val}}} = \chi(k, x).
\end{equation}
In particular, $\chi$ depends only on the quotient class $x \in \mathbb{Z}/L\mathbb{Z}$.
\end{lemma}

\begin{proof}
Since $\tilde{x} \equiv x_{\mathrm{val}} \pmod L$, there exists $m \in \mathbb{Z}$ such that $\tilde{x} = x_{\mathrm{val}} + m L$. Therefore:
\[
\zeta^{k \tilde{x}} = \zeta^{k (x_{\mathrm{val}} + m L)} = \zeta^{k x_{\mathrm{val}}} \cdot (\zeta^L)^{k m}.
\]
Because $\zeta$ is an $L$-th root of unity, $\zeta^L = 1$, so $(\zeta^L)^{k m} = 1^{k m} = 1$. Hence $\zeta^{k \tilde{x}} = \zeta^{k x_{\mathrm{val}}}$.
\end{proof}

\section{The Unscaled Analysis and Synthesis Transforms}

Because the spatial lattice $\Lambda$ and the momentum band $\Lambda^*$ are distinct sets, the Fourier transform consists of two directional linear maps between distinct function modules:
\begin{itemize}
    \item Spatial function space: $K^\Lambda \equiv (\Lambda \to K)$.
    \item Momentum function space: $K^{\Lambda^*} \equiv (\Lambda^* \to K)$.
\end{itemize}

\begin{definition}[Analysis and Synthesis Transforms]
We define the $K$-linear transforms:
\begin{enumerate}
    \item \textbf{Analysis Transform $S : K^\Lambda \to K^{\Lambda^*}$:}
    \begin{equation}
    (S f)(k) \coloneqq \sum_{x \in \Lambda} f(x)\, \chi(-k, x), \quad k \in \Lambda^*.
    \end{equation}
    \item \textbf{Synthesis Transform $T : K^{\Lambda^*} \to K^\Lambda$:}
    \begin{equation}
    (T g)(x) \coloneqq \sum_{k \in \Lambda^*} g(k)\, \chi(k, x), \quad x \in \Lambda.
    \end{equation}
\end{enumerate}
\end{definition}

Note that $\chi(-k, x) = \zeta^{-k x_{\mathrm{val}}} = (\zeta^{k x_{\mathrm{val}}})^{-1} = \chi(k, x)^{-1}$. By Lemma~\ref{lem:rep_independence}, evaluating at $-k \in \mathbb{Z}$ agrees with evaluating at the transported band negation $\ominus k \in \Lambda^*$, ensuring algebraic consistency even when $-k \notin \Lambda^*$.

\section{Dual Character Orthogonality and Algebraic Inversion}

The foundation of Fourier analysis on finite abelian groups is character orthogonality.

\begin{lemma}[Geometric Sum of Roots of Unity]
\label{lem:geom_sum}
Let $\zeta \in K$ be a primitive $L$-th root of unity. For any integer $m \in \mathbb{Z}$:
\begin{equation}
\sum_{x=0}^{L-1} \zeta^{m x} = \begin{cases}
L \cdot 1_K, & \text{if } m \equiv 0 \pmod L, \\
0, & \text{if } m \not\equiv 0 \pmod L.
\end{cases}
\end{equation}
\end{lemma}

\begin{proof}
If $m \equiv 0 \pmod L$, then $\zeta^m = 1$, so each of the $L$ terms equals $1$, giving $\sum_{x=0}^{L-1} 1 = L \cdot 1_K$.

If $m \not\equiv 0 \pmod L$, then $\zeta^m \neq 1$ because $\zeta$ is a primitive $L$-th root of unity. In any field, multiplying the sum by $(1 - \zeta^m)$ gives the telescoping sum:
\[
(1 - \zeta^m) \sum_{x=0}^{L-1} \zeta^{m x} = \sum_{x=0}^{L-1} \zeta^{m x} - \sum_{x=0}^{L-1} \zeta^{m (x+1)} = 1 - \zeta^{m L} = 1 - (\zeta^L)^m = 1 - 1 = 0.
\]
Because $K$ is a field and $(1 - \zeta^m) \neq 0$, dividing by $(1 - \zeta^m)$ proves $\sum_{x=0}^{L-1} \zeta^{m x} = 0$.
\end{proof}

\begin{theorem}[Dual Character Orthogonality Relations]
\label{thm:fourier_orthogonality}
Let $\zeta$ be a primitive $L$-th root of unity in $K$.
\begin{enumerate}
    \item \textbf{Spatial Orthogonality:} For any modes $k, p \in \Lambda^*$:
    \begin{equation}
    \sum_{x \in \Lambda} \chi(k - p, x) = L\, \delta_{k, p}.
    \end{equation}
    \item \textbf{Momentum (Dual) Orthogonality:} For any spatial sites $x, y \in \Lambda$:
    \begin{equation}
    \sum_{k \in \Lambda^*} \chi(k, x) \chi(-k, y) = L\, \delta_{x, y}.
    \end{equation}
\end{enumerate}
\end{theorem}

\begin{proof}
We prove each identity:
\begin{enumerate}
    \item The sum over $x \in \Lambda$ is explicitly the sum over unsigned residues $x_{\mathrm{val}} \in \{0, \dots, L-1\}$:
    \[
    \sum_{x \in \Lambda} \chi(k - p, x) = \sum_{x_{\mathrm{val}}=0}^{L-1} \zeta^{(k - p) x_{\mathrm{val}}}.
    \]
    By Lemma~\ref{lem:geom_sum}, this sum equals $L$ if $k - p \equiv 0 \pmod L$, and $0$ otherwise. Since both $k, p \in \Lambda^*$, their difference satisfies $-L < k - p < L$. The only multiple of $L$ in this range is $0$, meaning $k - p \equiv 0 \pmod L$ if and only if $k = p$. Thus the sum equals $L\, \delta_{k, p}$.

    \item By the canonical bijection $r : \Lambda \to \Lambda^*$ established in Chapter~\ref{ch:lattice_band}, summing over $k \in \Lambda^*$ is equivalent to summing over $k = r(z)$ for $z \in \Lambda$. Using representative independence (Lemma~\ref{lem:rep_independence}):
    \[
    \sum_{k \in \Lambda^*} \chi(k, x) \chi(-k, y) = \sum_{k \in \Lambda^*} \zeta^{k (x_{\mathrm{val}} - y_{\mathrm{val}})} = \sum_{z \in \Lambda} \zeta^{z_{\mathrm{val}} (x_{\mathrm{val}} - y_{\mathrm{val}})} = \sum_{z_{\mathrm{val}}=0}^{L-1} \zeta^{(x_{\mathrm{val}} - y_{\mathrm{val}}) z_{\mathrm{val}}}.
    \]
    By Lemma~\ref{lem:geom_sum}, this sum equals $L$ if $x_{\mathrm{val}} - y_{\mathrm{val}} \equiv 0 \pmod L$ (which means $x = y$ in $\Lambda$), and $0$ otherwise. This gives $L\, \delta_{x, y}$.
\end{enumerate}
\end{proof}

\begin{theorem}[Unscaled Inversion and Algebraic Isomorphism]
\label{thm:unscaled_inversion}
The compositions of analysis and synthesis satisfy:
\begin{equation}
T \circ S = L \cdot \mathrm{id}_{K^\Lambda}, \qquad S \circ T = L \cdot \mathrm{id}_{K^{\Lambda^*}}.
\end{equation}
Consequently, if $(L : K) \neq 0$ in the field $K$, $S$ and $T$ are linear isomorphisms, and the exact two-sided algebraic inverse of $S$ is:
\begin{equation}
S^{-1} = L^{-1} T.
\end{equation}
\end{theorem}

\begin{proof}
Let $f \in K^\Lambda$ and evaluate $(T(S f))(y)$ at site $y \in \Lambda$:
\begin{align*}
(T(S f))(y) &= \sum_{k \in \Lambda^*} (S f)(k)\, \chi(k, y) = \sum_{k \in \Lambda^*} \left( \sum_{x \in \Lambda} f(x)\, \chi(-k, x) \right) \chi(k, y) \\
&= \sum_{x \in \Lambda} f(x) \left( \sum_{k \in \Lambda^*} \chi(k, y) \chi(-k, x) \right).
\end{align*}
By Theorem~\ref{thm:fourier_orthogonality}(2), the inner sum equals $L\, \delta_{y, x} = L\, \delta_{x, y}$. Therefore:
\[
(T(S f))(y) = \sum_{x \in \Lambda} f(x) \cdot L\, \delta_{x, y} = L\, f(y).
\]
Hence $T \circ S = L \cdot \mathrm{id}_{K^\Lambda}$. Reversing the order and applying Theorem~\ref{thm:fourier_orthogonality}(1) similarly yields $S \circ T = L \cdot \mathrm{id}_{K^{\Lambda^*}}$.
\end{proof}

\begin{remark}[Formalization Note]
In the verified Lean~4 formalization (Chapter 3 and Appendix A01), all 61 Fourier character lemmas are proved constructively over roots of unity. The proof of orthogonality utilizes the geometric sum formula without analytical limits. This algebraic decoupling guarantees that the discrete Fourier transform is fully well-defined and bijective over any field of characteristic zero, without requiring metric completions or square roots of $L$.
\end{remark}

\section{Unitary Normalization on Complex Euclidean Space}

We now specialize to the field of complex numbers $K = \mathbb{C}$ equipped with the standard complex conjugation and the canonical root $\zeta = e^{2\pi i / L}$.

To define unitary operators, we endow the function spaces with their standard counting $\ell^2$ inner products:
\begin{itemize}
    \item Spatial Euclidean space $\mathcal{H}_{\mathrm{pos}} \coloneqq \ell^2(\Lambda)$:
    \begin{equation}
    \langle f_1, f_2 \rangle_{\mathrm{pos}} \coloneqq \sum_{x \in \Lambda} \overline{f_1(x)}\, f_2(x).
    \end{equation}
    \item Momentum Euclidean space $\mathcal{H}_{\mathrm{mom}} \coloneqq \ell^2(\Lambda^*)$:
    \begin{equation}
    \langle g_1, g_2 \rangle_{\mathrm{mom}} \coloneqq \sum_{k \in \Lambda^*} \overline{g_1(k)}\, g_2(k).
    \end{equation}
\end{itemize}

\begin{definition}[Unitary Fourier Transform]
Let $a \coloneqq \frac{1}{\sqrt{L}} \in \mathbb{R}_{>0}$, so that $L a^2 = 1$. The unitary discrete Fourier transform $U : \mathcal{H}_{\mathrm{pos}} \to \mathcal{H}_{\mathrm{mom}}$ and its inverse $U^{-1} : \mathcal{H}_{\mathrm{mom}} \to \mathcal{H}_{\mathrm{pos}}$ are defined by:
\begin{align}
(U f)(k) &\coloneqq \frac{1}{\sqrt{L}} (S f)(k) = \frac{1}{\sqrt{L}} \sum_{x \in \Lambda} f(x)\, \chi(-k, x), \\
(U^{-1} g)(x) &\coloneqq \frac{1}{\sqrt{L}} (T g)(x) = \frac{1}{\sqrt{L}} \sum_{k \in \Lambda^*} g(k)\, \chi(k, x).
\end{align}
\end{definition}

\begin{theorem}[Hilbert Adjoint and Unitarity]
\label{thm:unitary_fourier}
Under the counting inner products:
\begin{enumerate}
    \item The adjoint of the analysis map satisfies $S^\dagger = T$, and therefore:
    \begin{equation}
    U^\dagger = U^{-1}.
    \end{equation}
    \item $U$ is an exact Hilbert space isometry preserving inner products (Parseval / Plancherel identity):
    \begin{equation}
    \langle U f_1, U f_2 \rangle_{\mathrm{mom}} = \langle f_1, f_2 \rangle_{\mathrm{pos}} \quad \text{for all } f_1, f_2 \in \mathcal{H}_{\mathrm{pos}}.
    \end{equation}
\end{enumerate}
\end{theorem}

\begin{proof}
We evaluate the inner product directly:
\begin{align*}
\langle S f, g \rangle_{\mathrm{mom}} &= \sum_{k \in \Lambda^*} \overline{(S f)(k)}\, g(k) = \sum_{k \in \Lambda^*} \left( \sum_{x \in \Lambda} \overline{f(x)}\, \overline{\chi(-k, x)} \right) g(k) \\
&= \sum_{x \in \Lambda} \overline{f(x)} \left( \sum_{k \in \Lambda^*} g(k)\, \chi(k, x) \right) = \sum_{x \in \Lambda} \overline{f(x)}\, (T g)(x) = \langle f, T g \rangle_{\mathrm{pos}},
\end{align*}
using $\overline{\chi(-k, x)} = \overline{\zeta^{-k x_{\mathrm{val}}}} = \zeta^{k x_{\mathrm{val}}} = \chi(k, x)$. Thus $S^\dagger = T$.
Scaling by $a = 1/\sqrt{L} \in \mathbb{R}$:
\[
U^\dagger = (a S)^\dagger = a S^\dagger = a T = U^{-1}.
\]
Finally:
\[
\langle U f_1, U f_2 \rangle_{\mathrm{mom}} = \langle f_1, U^\dagger U f_2 \rangle_{\mathrm{pos}} = \langle f_1, U^{-1} U f_2 \rangle_{\mathrm{pos}} = \langle f_1, f_2 \rangle_{\mathrm{pos}}.
\]
This proves that $U$ is a unitary isomorphism.
\end{proof}

\section{Diagonalization of Difference Operators and Exact Dispersion}

We now connect the Fourier transform to the discrete difference operators of Chapter~\ref{ch:umbral_calculus}.

\begin{theorem}[Diagonal Action on Characters]
\label{thm:diagonal_dispersion}
Let $\chi_k(x) \coloneqq \chi(k, x)$ denote the Fourier character as a spatial function.
\begin{enumerate}
    \item The shift operator $E$ acts diagonally:
    \begin{equation}
    (E \chi_k)(x) = \zeta^k \chi_k(x).
    \end{equation}
    \item The forward and backward differences act diagonally:
    \begin{align}
    (\Delta \chi_k)(x) &= (\zeta^k - 1) \chi_k(x), \\
    (\nabla \chi_k)(x) &= (1 - \zeta^{-k}) \chi_k(x).
    \end{align}
    \item The discrete Laplacian acts diagonally with eigenvalue $\lambda_k = \zeta^k + \zeta^{-k} - 2$:
    \begin{equation}
    (\Delta\nabla \chi_k)(x) = (\zeta^k + \zeta^{-k} - 2) \chi_k(x).
    \end{equation}
\end{enumerate}
\end{theorem}

\begin{proof}
Evaluating the definitions:
\begin{enumerate}
    \item $(E \chi_k)(x) = \chi_k(x+1) = \zeta^{k (x_{\mathrm{val}}+1)} = \zeta^k \zeta^{k x_{\mathrm{val}}} = \zeta^k \chi_k(x)$ by Lemma~\ref{lem:rep_independence}.
    \item $(\Delta \chi_k)(x) = (E - I)\chi_k(x) = (\zeta^k - 1)\chi_k(x)$, and $(\nabla \chi_k)(x) = (I - E^{-1})\chi_k(x) = (1 - \zeta^{-k})\chi_k(x)$.
    \item Composing the two differences:
    \[
    \Delta(\nabla \chi_k) = (\zeta^k - 1)(1 - \zeta^{-k})\chi_k = (\zeta^k - 1 - 1 + \zeta^{-k})\chi_k = (\zeta^k + \zeta^{-k} - 2)\chi_k.
    \]
\end{enumerate}
\end{proof}

\begin{corollary}[Exact Lattice Dispersion in $\mathbb{C}$]
Under the canonical root $\zeta = e^{2\pi i / L}$, the Laplacian eigenvalue evaluates to:
\begin{equation}
\zeta^k + \zeta^{-k} - 2 = -4 \sin^2\left(\frac{\pi k}{L}\right).
\end{equation}
\end{corollary}

\begin{proof}
Using Euler's formula with angle $\theta = \frac{2\pi k}{L}$:
\[
\zeta^k + \zeta^{-k} = e^{i\theta} + e^{-i\theta} = 2 \cos\theta.
\]
Using the trigonometric half-angle identity $\cos\theta = 1 - 2 \sin^2(\theta/2)$:
\[
\zeta^k + \zeta^{-k} - 2 = 2 \cos\left(\frac{2\pi k}{L}\right) - 2 = 2 \left( 1 - 2 \sin^2\left(\frac{\pi k}{L}\right) \right) - 2 = -4 \sin^2\left(\frac{\pi k}{L}\right).
\]
\end{proof}

In the long-wavelength continuum limit where the mode index is small compared to the lattice size ($|k| \ll L$), we approximate $\sin(\pi k / L) \approx \frac{\pi k}{L}$. The dispersion becomes:
\begin{equation}
-4 \sin^2\left(\frac{\pi k}{L}\right) \approx -4 \left(\frac{\pi k}{L}\right)^2 = -\left(\frac{2\pi k}{L}\right)^2 = -p_{\mathrm{phys}}^2,
\end{equation}
recovering the continuous kinetic energy dispersion and Laplacian $-\partial_x^2$. On the discrete lattice, the trigonometric dispersion naturally truncates high-momentum modes, eliminating all ultraviolet divergences.

\chapter{CAR Representation and Fock Space}
\label{ch:car_fock}

\section{Introduction and Abstract Setting}

In quantum field theory, the fundamental kinematics of fermions is governed by the Canonical Anticommutation Relations (CAR). Before assigning spatial or momentum coordinates to fermionic modes, the algebraic theory of fermions must be established over a generic, abstract set of single-particle states.

Let $\iota$ be a non-empty, finite index set equipped with a linear order:
\begin{equation}
\iota = \{1, 2, \dots, d\}, \quad d \coloneqq \#\iota \ge 1.
\end{equation}
The choice of a total linear order $\le$ on $\iota$ is physically unobservable but mathematically essential: it fixes the sign conventions acquired during the permutation of fermionic creation operators, eliminating sign ambiguities.

Typical physical realizations of $\iota$ include:
\begin{enumerate}
    \item The discrete centered Brillouin zone $\iota = \Lambda^* = \{-h+1, \dots, h\}$, linearly ordered by the standard integer order $\le$ (Chapters~\ref{ch:lattice_band} and~\ref{ch:fermions}).
    \item Multi-species or spinful chiral fermions $\iota = \{R, L\} \times \{\uparrow, \downarrow\} \times \Lambda^*$, endowed with the lexicographic product order.
\end{enumerate}

\section{The Two-Tiered CAR Hierarchy: Algebraic vs. Hilbert Setting}

To separate pure algebraic identities from metric and inner-product properties, we distinguish between two tiers of CAR representations.

\begin{definition}[Algebraic CAR Representation — Tier 1]
Let $V$ be a complex vector space. An \emph{Algebraic CAR representation} over $\iota$ on $V$ consists of two families of linear operators $c, c^\dagger : \iota \to \mathrm{End}_{\mathbb{C}}(V)$ satisfying, for all $i, j \in \iota$:
\begin{align}
\{ c_i, c_j \} &\coloneqq c_i c_j + c_j c_i = 0, \label{eq:car_ann_ann} \\
\{ c_i^\dagger, c_j^\dagger \} &\coloneqq c_i^\dagger c_j^\dagger + c_j^\dagger c_i^\dagger = 0, \label{eq:car_cre_cre} \\
\{ c_i, c_j^\dagger \} &\coloneqq c_i c_j^\dagger + c_j^\dagger c_i = \delta_{i, j} I_V. \label{eq:car_ann_cre}
\end{align}
\end{definition}

Without assuming any inner product, adjoint operation, or topological structure on $V$, the algebraic relations alone imply fundamental commutation theorems.

\begin{theorem}[Purely Algebraic CAR Identities]
\label{thm:algebraic_car}
Let $(c, c^\dagger)$ be an Algebraic CAR representation on $V$.
\begin{enumerate}
    \item \textbf{Bilinear Commutator Identity:} For any modes $a, b, c, d \in \iota$:
    \begin{equation}
    \left[ c_a^\dagger c_b, c_c^\dagger c_d \right] = \delta_{b, c}\, c_a^\dagger c_d - \delta_{a, d}\, c_c^\dagger c_b.
    \end{equation}
    \item \textbf{Number Operator Commutators:} For $n_i \coloneqq c_i^\dagger c_i$ and any $j \in \iota$:
    \begin{align}
    [n_i, c_j^\dagger] &= \delta_{i, j} c_j^\dagger, \\
    [n_i, c_j] &= -\delta_{i, j} c_j.
    \end{align}
    \item \textbf{Idempotence and Number Commutativity:}
    \begin{equation}
    n_i^2 = n_i, \qquad [n_i, n_j] = 0 \quad \text{for all } i, j \in \iota.
    \end{equation}
\end{enumerate}
\end{theorem}

\begin{proof}
We verify each algebraic identity:
\begin{enumerate}
    \item We expand the commutator $[c_a^\dagger c_b, c_c^\dagger c_d] = c_a^\dagger c_b c_c^\dagger c_d - c_c^\dagger c_d c_a^\dagger c_b$. 
    Using $c_b c_c^\dagger = \delta_{b, c} I - c_c^\dagger c_b$ on the first term:
    \[
    c_a^\dagger c_b c_c^\dagger c_d = c_a^\dagger (\delta_{b, c} I - c_c^\dagger c_b) c_d = \delta_{b, c} c_a^\dagger c_d - c_a^\dagger c_c^\dagger c_b c_d.
    \]
    Using $c_c^\dagger c_a^\dagger = - c_a^\dagger c_c^\dagger$ and $c_d c_b = - c_b c_d$ on the second term:
    \[
    c_c^\dagger c_d c_a^\dagger c_b = c_c^\dagger (\delta_{a, d} I - c_a^\dagger c_d) c_b = \delta_{a, d} c_c^\dagger c_b - c_c^\dagger c_a^\dagger c_d c_b = \delta_{a, d} c_c^\dagger c_b - (-c_a^\dagger c_c^\dagger)(-c_b c_d).
    \]
    Subtracting the two expressions, the quartic products $-c_a^\dagger c_c^\dagger c_b c_d$ cancel identically, leaving:
    \[
    [c_a^\dagger c_b, c_c^\dagger c_d] = \delta_{b, c} c_a^\dagger c_d - \delta_{a, d} c_c^\dagger c_b.
    \]
    
    \item Setting $a = b = i$ and $c = j$ in the bilinear identity gives $[n_i, c_j^\dagger c_d] = \delta_{i, j} c_i^\dagger c_d - \delta_{i, d} c_j^\dagger c_i$. Contracting with a test state or using $\{c_i, c_j^\dagger\} = \delta_{i,j}$:
    \[
    n_i c_j^\dagger = c_i^\dagger c_i c_j^\dagger = c_i^\dagger (\delta_{i, j} I - c_j^\dagger c_i) = \delta_{i, j} c_i^\dagger - c_i^\dagger c_j^\dagger c_i = \delta_{i, j} c_j^\dagger + c_j^\dagger c_i^\dagger c_i = \delta_{i, j} c_j^\dagger + c_j^\dagger n_i,
    \]
    which gives $[n_i, c_j^\dagger] = \delta_{i, j} c_j^\dagger$. The annihilation commutator follows analogously.
    
    \item For idempotence: $n_i^2 = c_i^\dagger c_i c_i^\dagger c_i = c_i^\dagger (I - c_i^\dagger c_i) c_i = c_i^\dagger c_i - (c_i^\dagger)^2 c_i^2$. Since $(c_i^\dagger)^2 = 0$, $n_i^2 = n_i$. Finally, $[n_i, n_j] = [c_i^\dagger c_i, c_j^\dagger c_j] = \delta_{i, j} c_i^\dagger c_j - \delta_{i, j} c_j^\dagger c_i = 0$.
\end{enumerate}
\end{proof}

\begin{definition}[Hilbert Space CAR Representation — Tier 2]
Let $V$ be a finite-dimensional complex Hilbert space equipped with positive-definite inner product $\langle \cdot, \cdot \rangle$. A \emph{Hilbert CAR representation} over $\iota$ on $V$ is an Algebraic CAR representation $(c, c^\dagger)$ that additionally satisfies \emph{Hilbert Adjoint Compatibility}:
\begin{equation}
c_i^\dagger = (c_i)^\dagger \quad \text{for all } i \in \iota,
\end{equation}
where $(A)^\dagger$ denotes the unique Hilbert space adjoint satisfying $\langle u, A v \rangle = \langle A^\dagger u, v \rangle$.
\end{definition}

\begin{remark}[Formalization Note]
In the verified Lean formalization (\texttt{Bosonize.A02}), the separation between \texttt{AlgebraicCAR} and \texttt{CAR} is strictly maintained. Adjoint compatibility is an independent geometric axiom: if $(c, c^\dagger)$ is an algebraic CAR pair and $S$ is an invertible non-unitary operator, the similarity transform $\tilde{c}_i \coloneqq S c_i S^{-1}$ and $\tilde{c}_i^\dagger \coloneqq S c_i^\dagger S^{-1}$ preserves all anticommutators, but $\tilde{c}_i^\dagger \neq (\tilde{c}_i)^\dagger$. Thus, adjointness is an essential condition that fixes the physical Hilbert space metric.
\end{remark}

\section{The Concrete Occupation Fock Space}

We now construct the canonical Fock space representation of the CAR algebra.

\begin{definition}[Occupation Basis and Fock Space Carrier]
Let $\iota$ be a finite set. The fermionic Fock space $\mathcal{F}(\iota)$ is the complex Euclidean space indexed by the power set $\mathcal{P}(\iota)$ of all subsets of $\iota$:
\begin{equation}
\mathcal{F}(\iota) \coloneqq \ell^2(\mathcal{P}(\iota)) \equiv \bigoplus_{S \subseteq \iota} \mathbb{C} |S\rangle.
\end{equation}
For each subset $S \subseteq \iota$, the basis vector $|S\rangle$ is the indicator state representing the configuration where exactly the single-particle modes in $S$ are occupied.
\end{definition}

\begin{proposition}[Dimension and Inner Product]
\begin{enumerate}
    \item The complex dimension of the Fock space is:
    \begin{equation}
    \dim_{\mathbb{C}} \mathcal{F}(\iota) = 2^{\#\iota} = 2^d.
    \end{equation}
    \item Basis states form an orthonormal basis under the standard Euclidean inner product:
    \begin{equation}
    \langle S | T \rangle = \delta_{S, T} \quad \text{for all } S, T \subseteq \iota.
    \end{equation}
    \item The state $|\emptyset\rangle$ is the unique normalized fermionic vacuum state.
\end{enumerate}
\end{proposition}

\section{Exact Fermionic Sign Mechanics}

To define creation and annihilation operators on $\mathcal{F}(\iota)$ that satisfy anticommutation, we must account for the permutation phase acquired when inserting or removing a fermion.

\begin{definition}[Preceding Count and Fermionic Sign]
Let $\iota$ be linearly ordered by $\le$. For any mode $i \in \iota$ and any occupation subset $S \subseteq \iota$:
\begin{enumerate}
    \item The \emph{preceding count} $\sigma(i, S)$ is the number of occupied modes in $S$ strictly preceding $i$:
    \begin{equation}
    \sigma(i, S) \coloneqq \#\{j \in S \mid j < i\}.
    \end{equation}
    \item The \emph{fermionic sign} $\epsilon(i, S)$ is the sign factor:
    \begin{equation}
    \epsilon(i, S) \coloneqq (-1)^{\sigma(i, S)} \in \{+1, -1\}.
    \end{equation}
\end{enumerate}
\end{definition}

In constructive mathematics, to avoid underflow errors when computing counts in $\mathbb{N}$ (where $0 - 1 = 0$), all set operations are formulated additively.

\begin{lemma}[Additive Atomic Counts]
\label{lem:atomic_counts}
Let $i, j \in \iota$ and $S \subseteq \iota$.
\begin{enumerate}
    \item If $j \notin S$, then:
    \begin{equation}
    \sigma(i, S \cup \{j\}) = \sigma(i, S) + \begin{cases} 1, & \text{if } j < i, \\ 0, & \text{if } j \ge i. \end{cases}
    \end{equation}
    \item If $j \in S$, then:
    \begin{equation}
    \sigma(i, S) = \sigma(i, S \setminus \{j\}) + \begin{cases} 1, & \text{if } j < i, \\ 0, & \text{if } j \ge i. \end{cases}
    \end{equation}
    \item Self-invariance: $\sigma(i, S \cup \{i\}) = \sigma(i, S)$ and $\sigma(i, S \setminus \{i\}) = \sigma(i, S)$.
\end{enumerate}
\end{lemma}

\begin{proof}
The preceding set $\{k \in S \cup \{j\} \mid k < i\}$ is the disjoint union of $\{k \in S \mid k < i\}$ and $\{j\}$ if $j < i$, or simply $\{k \in S \mid k < i\}$ if $j \ge i$. Taking set cardinalities yields the additive formula. Part (2) is the inverse statement, and part (3) follows because $i < i$ is false.
\end{proof}

\begin{theorem}[Fermionic Sign Exchange Identities]
\label{thm:sign_exchange}
For any distinct modes $i \neq j$ in $\iota$:
\begin{enumerate}
    \item \textbf{Creation Pair:} If $i, j \notin S$:
    \begin{equation}
    \epsilon(j, S)\, \epsilon(i, S \cup \{j\}) = - \epsilon(i, S)\, \epsilon(j, S \cup \{i\}).
    \end{equation}
    \item \textbf{Annihilation Pair:} If $i, j \in S$:
    \begin{equation}
    \epsilon(j, S)\, \epsilon(i, S \setminus \{j\}) = - \epsilon(i, S)\, \epsilon(j, S \setminus \{i\}).
    \end{equation}
    \item \textbf{Mixed Pair:} If $i \notin S$ and $j \in S$:
    \begin{equation}
    \epsilon(j, S)\, \epsilon(i, S \setminus \{j\}) = - \epsilon(i, S)\, \epsilon(j, S \cup \{i\}).
    \end{equation}
\end{enumerate}
\end{theorem}

\begin{proof}
Because $\iota$ is linearly ordered and $i \neq j$, either $j < i$ or $i < j$.
\begin{enumerate}
    \item Suppose $j < i$. By Lemma~\ref{lem:atomic_counts}(1):
    \[
    \sigma(i, S \cup \{j\}) = \sigma(i, S) + 1 \implies \epsilon(i, S \cup \{j\}) = (-1)^{\sigma(i, S)+1} = - \epsilon(i, S).
    \]
    On the other hand, since $i > j$, inserting $i$ does not affect the count of elements preceding $j$:
    \[
    \sigma(j, S \cup \{i\}) = \sigma(j, S) \implies \epsilon(j, S \cup \{i\}) = \epsilon(j, S).
    \]
    Therefore:
    \[
    \epsilon(j, S) \epsilon(i, S \cup \{j\}) = \epsilon(j, S) (-\epsilon(i, S)) = - \epsilon(i, S) \epsilon(j, S \cup \{i\}).
    \]
    If $i < j$, the exact same argument applies with the roles of $i$ and $j$ interchanged.
    \item Parts (2) and (3) follow by the same parity discrepancy using Lemma~\ref{lem:atomic_counts}(2).
\end{enumerate}
\end{proof}

\section{Creation and Annihilation Operators on Fock Space}

\begin{definition}[Fermionic Operators on Basis States]
For each mode $i \in \iota$, the annihilation operator $c_i$ and creation operator $c_i^\dagger$ are defined on basis states $|S\rangle$ by:
\begin{align}
c_i |S\rangle &\coloneqq \begin{cases}
\epsilon(i, S) |S \setminus \{i\}\rangle, & \text{if } i \in S, \\
0, & \text{if } i \notin S,
\end{cases} \label{eq:fock_c} \\
c_i^\dagger |S\rangle &\coloneqq \begin{cases}
\epsilon(i, S) |S \cup \{i\}\rangle, & \text{if } i \notin S, \\
0, & \text{if } i \in S,
\end{cases} \label{eq:fock_cdag}
\end{align}
and extended linearly to all states in $\mathcal{F}(\iota)$.
\end{definition}

\begin{theorem}[Verification of the Concrete CAR Representation]
\label{thm:fock_car_verification}
The operators $\{c_i, c_i^\dagger \mid i \in \iota\}$ satisfy:
\begin{enumerate}
    \item \textbf{Adjoint Compatibility:} $c_i^\dagger = (c_i)^\dagger$ under the Euclidean inner product.
    \item \textbf{Anticommutation Relations:}
    \begin{equation}
    \{c_i, c_j\} = 0, \qquad \{c_i^\dagger, c_j^\dagger\} = 0, \qquad \{c_i, c_j^\dagger\} = \delta_{i, j} I.
    \end{equation}
\end{enumerate}
In particular, $\mathcal{F}(\iota)$ is a concrete Hilbert CAR representation.
\end{theorem}

\begin{proof}
We verify both properties on arbitrary basis states $|S\rangle, |T\rangle$:
\begin{enumerate}
    \item For adjointness, we compute the matrix elements $\langle S | c_i | T \rangle$ and $\langle c_i^\dagger S | T \rangle$:
    \begin{itemize}
        \item If $i \in T$ and $S = T \setminus \{i\}$, then $c_i |T\rangle = \epsilon(i, T) |S\rangle$, so $\langle S | c_i | T \rangle = \epsilon(i, T)$. Also, $i \notin S$ and $S \cup \{i\} = T$, so $c_i^\dagger |S\rangle = \epsilon(i, S) |T\rangle$. By Lemma~\ref{lem:atomic_counts}(3), $\sigma(i, S) = \sigma(i, T)$, so $\epsilon(i, S) = \epsilon(i, T)$. Thus $\langle c_i^\dagger S | T \rangle = \epsilon(i, S) = \epsilon(i, T) = \langle S | c_i | T \rangle$.
        \item In all other cases, both matrix elements vanish identically.
    \end{itemize}
    Since this holds for all basis states, $c_i^\dagger = (c_i)^\dagger$.

    \item For anticommutators:
    \begin{itemize}
        \item If $i = j$: $c_i^2 |S\rangle = 0$ because if $i \in S$, $c_i |S\rangle \propto |S \setminus \{i\}\rangle$, which no longer contains $i$, so applying $c_i$ again yields $0$. Similarly, $(c_i^\dagger)^2 = 0$. For the mixed term: if $i \in S$, $c_i^\dagger c_i |S\rangle = |S\rangle$ and $c_i c_i^\dagger |S\rangle = 0$; if $i \notin S$, $c_i^\dagger c_i |S\rangle = 0$ and $c_i c_i^\dagger |S\rangle = |S\rangle$. In both cases, $(c_i c_i^\dagger + c_i^\dagger c_i) |S\rangle = |S\rangle = \delta_{i, i} |S\rangle$.
        \item If $i \neq j$: when $i, j \notin S$, $(c_i^\dagger c_j^\dagger + c_j^\dagger c_i^\dagger)|S\rangle = \big( \epsilon(j, S) \epsilon(i, S \cup \{j\}) + \epsilon(i, S) \epsilon(j, S \cup \{i\}) \big) |S \cup \{i, j\}\rangle$. By Theorem~\ref{thm:sign_exchange}(1), the sum of signs is $0$. The same holds for $\{c_i, c_j\} = 0$ and $\{c_i, c_j^\dagger\} = 0$ using parts (2) and (3) of Theorem~\ref{thm:sign_exchange}.
    \end{itemize}
\end{enumerate}
\end{proof}

\section{Observables, Parity, and the Hopping Algebra}

\begin{definition}[Number and Parity Observables]
\begin{enumerate}
    \item \textbf{Local Number Operator:} $n_i \coloneqq c_i^\dagger c_i$. On basis states:
    \begin{equation}
    n_i |S\rangle = \begin{cases} |S\rangle, & \text{if } i \in S, \\ 0, & \text{if } i \notin S. \end{cases}
    \end{equation}
    \item \textbf{Total Particle Number Operator:} $\hat{N}_{\mathrm{tot}} \coloneqq \sum_{i \in \iota} n_i$. On basis states:
    \begin{equation}
    \hat{N}_{\mathrm{tot}} |S\rangle = (\#S)\, |S\rangle.
    \end{equation}
    \item \textbf{Global Fermion Parity Operator $\Gamma$:} Defined as the ordered product over the sorted index set $\iota = \{i_1 < i_2 < \dots < i_d\}$:
    \begin{equation}
    \Gamma \coloneqq \prod_{k=1}^d (I - 2 n_{i_k}).
    \end{equation}
\end{enumerate}
\end{definition}

\begin{proposition}[Properties of Global Fermion Parity]
\label{prop:parity}
The parity operator $\Gamma$ satisfies:
\begin{enumerate}
    \item Action on basis states:
    \begin{equation}
    \Gamma |S\rangle = (-1)^{\#S} |S\rangle.
    \end{equation}
    \item Involution and self-adjointness: $\Gamma^2 = I$ and $\Gamma^\dagger = \Gamma$.
    \item Anticommutation with fermionic generators:
    \begin{equation}
    \{\Gamma, c_i\} = 0, \qquad \{\Gamma, c_i^\dagger\} = 0 \quad \text{for all } i \in \iota.
    \end{equation}
\end{enumerate}
\end{proposition}

\begin{proof}
For any $i \in \iota$, $(I - 2n_i)|S\rangle = -|S\rangle$ if $i \in S$, and $+|S\rangle$ if $i \notin S$. Since all $n_i$ commute with each other, the product evaluates to $\prod_{i \in S} (-1) = (-1)^{\#S}$.
Since $(-1)^{\#S} \in \{+1, -1\}$, $\Gamma^2 |S\rangle = (-1)^{2\#S} |S\rangle = |S\rangle$, and $\Gamma$ is self-adjoint.
Finally, $c_i$ reduces $\#S$ by $1$ if $i \in S$, so $\Gamma c_i |S\rangle = (-1)^{\#S - 1} c_i |S\rangle = - c_i (-1)^{\#S} |S\rangle = - c_i \Gamma |S\rangle$.
\end{proof}

\begin{definition}[Hopping Operator]
For any modes $i, j \in \iota$, the hopping operator $T_{i, j}$ is defined by:
\begin{equation}
T_{i, j} \coloneqq c_i^\dagger c_j.
\end{equation}
\end{definition}

\begin{theorem}[Hopping Lie Algebra]
The hopping operators satisfy:
\begin{enumerate}
    \item Adjoint relation: $T_{i, j}^\dagger = T_{j, i}$.
    \item Commutation relations forming the Lie algebra $\mathfrak{u}(d)$:
    \begin{equation}
    [T_{a, b}, T_{c, d}] = \delta_{b, c} T_{a, d} - \delta_{a, d} T_{c, b}.
    \end{equation}
\end{enumerate}
\end{theorem}

\begin{proof}
This is an immediate consequence of Theorem~\ref{thm:algebraic_car}(1) and adjoint compatibility $c_i^\dagger = (c_i)^\dagger$.
\end{proof}

The hopping operators $T_{i, j}$ form the exact algebraic building blocks for the kinetic Hamiltonian and the density operators in the chapters that follow.

\chapter{Position and Momentum Fermions}
\label{ch:fermions}

\section{Introduction and Physical Framework}

Having established the abstract Canonical Anticommutation Relations (CAR) and Fock space representations in Chapter~\ref{ch:car_fock}, and the discrete Fourier transform in Chapter~\ref{ch:fourier}, we now unite these structures to describe physical fermions on the one-dimensional lattice.

Throughout this chapter and the remainder of the monograph, the lattice length is assumed to be an even positive integer:
\begin{equation}
L = 2h, \quad h \in \mathbb{N}, \quad h \ge 1.
\end{equation}
The single-particle index set $\iota$ is instantiated as the centered momentum band:
\begin{equation}
\iota = \Lambda^* = \left\{ k \in \mathbb{Z} \;\middle|\; -L < 2k \le L \right\} = \{-h+1, -h+2, \dots, h-1, h\},
\end{equation}
endowed with the standard linear order inherited from $\mathbb{Z}$. 

The fermionic state space is the Fock space over the momentum band:
\begin{equation}
\mathcal{F}(L) \coloneqq \ell^2(\mathcal{P}(\Lambda^*)) \equiv \bigoplus_{S \subseteq \Lambda^*} \mathbb{C} |S\rangle,
\end{equation}
which has exact complex dimension $\dim_{\mathbb{C}} \mathcal{F}(L) = 2^L$. The momentum creation and annihilation operators $\{c_k, c_k^\dagger \mid k \in \Lambda^*\}$ satisfy the standard CAR algebra:
\begin{equation}
\{c_k, c_p^\dagger\} = \delta_{k, p} I, \qquad \{c_k, c_p\} = 0, \qquad \{c_k^\dagger, c_p^\dagger\} = 0, \qquad c_k^\dagger = (c_k)^\dagger.
\end{equation}

\section{Spatial Fermions via the Unitary Fourier Transform}

In continuous quantum field theory, field operators $\psi(x)$ are operator-valued distributions that must be smeared with test functions. On our discrete spatial circle $\Lambda = \mathbb{Z}/L\mathbb{Z}$, spatial fermions are well-defined bounded operators constructed as finite linear combinations of momentum modes.

Recall from Chapter~\ref{ch:fourier} the real normalization factor $a \coloneqq 1/\sqrt{L} \in \mathbb{R}_{>0}$ satisfying $L a^2 = 1$, and the canonical primitive $L$-th root $\zeta = e^{2\pi i / L}$.

\begin{definition}[Spatial Fermionic Operators]
For every spatial site $x \in \Lambda$, the spatial annihilation operator $c_x$ and spatial creation operator $c_x^\dagger$ are defined by:
\begin{align}
c_x &\coloneqq \frac{1}{\sqrt{L}} \sum_{k \in \Lambda^*} \chi(k, x)\, c_k = \frac{1}{\sqrt{L}} \sum_{k \in \Lambda^*} \zeta^{k x_{\mathrm{val}}}\, c_k, \label{eq:spatial_c} \\
c_x^\dagger &\coloneqq \frac{1}{\sqrt{L}} \sum_{k \in \Lambda^*} \chi(-k, x)\, c_k^\dagger = \frac{1}{\sqrt{L}} \sum_{k \in \Lambda^*} \zeta^{-k x_{\mathrm{val}}}\, c_k^\dagger. \label{eq:spatial_cdag}
\end{align}
The local spatial fermion number operator is defined by:
\begin{equation}
n_x \coloneqq c_x^\dagger c_x.
\end{equation}
\end{definition}

\begin{lemma}[Hilbert Adjoint Compatibility]
The spatial creation operator is the exact Hilbert adjoint of the spatial annihilation operator:
\begin{equation}
(c_x)^\dagger = c_x^\dagger \quad \text{for all } x \in \Lambda.
\end{equation}
\end{lemma}

\begin{proof}
Taking the Hilbert adjoint of the linear combination~\eqref{eq:spatial_c}:
\[
(c_x)^\dagger = \left( \frac{1}{\sqrt{L}} \sum_{k \in \Lambda^*} \zeta^{k x_{\mathrm{val}}} c_k \right)^\dagger = \frac{1}{\sqrt{L}} \sum_{k \in \Lambda^*} \overline{\zeta^{k x_{\mathrm{val}}}} (c_k)^\dagger.
\]
Because $(c_k)^\dagger = c_k^\dagger$ and $\overline{\zeta^{k x_{\mathrm{val}}}} = \zeta^{-k x_{\mathrm{val}}} = \chi(-k, x)$, this sum coincides with the definition of $c_x^\dagger$ in~\eqref{eq:spatial_cdag}.
\end{proof}

\section{Exact Position CAR Inheritance}

The fundamental theorem of discrete spatial fermions is that they inherit the local canonical anticommutation relations without distortion.

\begin{theorem}[Exact Spatial CAR]
\label{thm:spatial_car}
For all spatial sites $x, y \in \Lambda$, the spatial operators satisfy:
\begin{align}
\{c_x, c_y^\dagger\} &= \delta_{x, y} I, \label{eq:pos_car_mixed} \\
\{c_x, c_y\} &= 0, \label{eq:pos_car_ann} \\
\{c_x^\dagger, c_y^\dagger\} &= 0. \label{eq:pos_car_cre}
\end{align}
\end{theorem}

\begin{proof}
We evaluate each anticommutator using the bilinearity of $\{\cdot, \cdot\}$:
\begin{enumerate}
    \item For the mixed anticommutator:
    \begin{align*}
    \{c_x, c_y^\dagger\} &= \left\{ \frac{1}{\sqrt{L}} \sum_{k \in \Lambda^*} \chi(k, x) c_k, \; \frac{1}{\sqrt{L}} \sum_{p \in \Lambda^*} \chi(-p, y) c_p^\dagger \right\} \\
    &= \frac{1}{L} \sum_{k \in \Lambda^*} \sum_{p \in \Lambda^*} \chi(k, x) \chi(-p, y) \{c_k, c_p^\dagger\}.
    \end{align*}
    Substituting the momentum CAR $\{c_k, c_p^\dagger\} = \delta_{k, p} I$:
    \[
    \{c_x, c_y^\dagger\} = \frac{1}{L} \sum_{k \in \Lambda^*} \chi(k, x) \chi(-k, y) I = \frac{1}{L} \left( \sum_{k \in \Lambda^*} \zeta^{k (x_{\mathrm{val}} - y_{\mathrm{val}})} \right) I.
    \]
    By Theorem~\ref{thm:fourier_orthogonality}(2) (dual character orthogonality), the momentum sum equals $L\, \delta_{x, y}$. Therefore:
    \[
    \{c_x, c_y^\dagger\} = \frac{1}{L} (L\, \delta_{x, y}) I = \delta_{x, y} I.
    \]

    \item For the annihilation anticommutator:
    \[
    \{c_x, c_y\} = \frac{1}{L} \sum_{k, p \in \Lambda^*} \chi(k, x) \chi(p, y) \{c_k, c_p\} = 0,
    \]
    because $\{c_k, c_p\} = 0$ for all $k, p \in \Lambda^*$. Taking adjoints yields $\{c_x^\dagger, c_y^\dagger\} = 0$.
\end{enumerate}
\end{proof}

\begin{remark}[Formalization Note]
In the verified Lean~4 core (Chapter 5, Position and Momentum Fermions), this result is proved purely through bilinear expansion and character orthogonality. Notice an essential conceptual feature of this construction: spatial fermions are obtained directly as global Fourier transforms of momentum fermions, rather than through local Jordan--Wigner transformations from spin operators. Consequently, there are no non-local string operators connecting spatial modes to momentum modes.
\end{remark}

\section{Operator Fourier Inversion}

The relationship between position and momentum fermions is completely symmetric under Fourier inversion.

\begin{theorem}[Fourier Inversion for Operators]
\label{thm:op_fourier_inversion}
The momentum operators are recovered from spatial fermions via the inverse discrete Fourier transform:
\begin{align}
c_k &= \frac{1}{\sqrt{L}} \sum_{x \in \Lambda} \chi(-k, x)\, c_x = \frac{1}{\sqrt{L}} \sum_{x \in \Lambda} \zeta^{-k x_{\mathrm{val}}}\, c_x, \label{eq:inv_ck} \\
c_k^\dagger &= \frac{1}{\sqrt{L}} \sum_{x \in \Lambda} \chi(k, x)\, c_x^\dagger = \frac{1}{\sqrt{L}} \sum_{x \in \Lambda} \zeta^{k x_{\mathrm{val}}}\, c_x^\dagger. \label{eq:inv_ckdag}
\end{align}
\end{theorem}

\begin{proof}
Substitute the definition of $c_x$ from~\eqref{eq:spatial_c} into the right-hand side of~\eqref{eq:inv_ck}:
\[
\frac{1}{\sqrt{L}} \sum_{x \in \Lambda} \chi(-k, x) \left( \frac{1}{\sqrt{L}} \sum_{p \in \Lambda^*} \chi(p, x) c_p \right) = \frac{1}{L} \sum_{p \in \Lambda^*} \left( \sum_{x \in \Lambda} \chi(p - k, x) \right) c_p.
\]
By Theorem~\ref{thm:fourier_orthogonality}(1) (spatial character orthogonality), $\sum_{x \in \Lambda} \chi(p - k, x) = L\, \delta_{p, k}$. Thus:
\[
\frac{1}{L} \sum_{p \in \Lambda^*} (L\, \delta_{p, k}) c_p = c_k.
\]
Taking Hilbert adjoints yields the corresponding identity for $c_k^\dagger$.
\end{proof}

\section{Conservation of Total Particle Number}

\begin{theorem}[Particle Number Invariance]
\label{thm:number_invariance}
The total particle number operator $\hat{N}_{\mathrm{tot}}$ satisfies:
\begin{equation}
\hat{N}_{\mathrm{tot}} = \sum_{k \in \Lambda^*} c_k^\dagger c_k = \sum_{x \in \Lambda} c_x^\dagger c_x = \sum_{x \in \Lambda} n_x.
\end{equation}
\end{theorem}

\begin{proof}
Expand the sum of spatial number operators using~\eqref{eq:spatial_c} and~\eqref{eq:spatial_cdag}:
\begin{align*}
\sum_{x \in \Lambda} c_x^\dagger c_x &= \sum_{x \in \Lambda} \left( \frac{1}{\sqrt{L}} \sum_{k \in \Lambda^*} \chi(-k, x) c_k^\dagger \right) \left( \frac{1}{\sqrt{L}} \sum_{p \in \Lambda^*} \chi(p, x) c_p \right) \\
&= \frac{1}{L} \sum_{k, p \in \Lambda^*} c_k^\dagger c_p \left( \sum_{x \in \Lambda} \chi(p - k, x) \right).
\end{align*}
Applying spatial character orthogonality $\sum_{x \in \Lambda} \chi(p - k, x) = L\, \delta_{p, k}$:
\[
\sum_{x \in \Lambda} c_x^\dagger c_x = \frac{1}{L} \sum_{k, p \in \Lambda^*} c_k^\dagger c_p (L\, \delta_{p, k}) = \sum_{k \in \Lambda^*} c_k^\dagger c_k = \hat{N}_{\mathrm{tot}}.
\]
\end{proof}

\section{The Dirac Fermi Sea and Linearized Hamiltonians}

In relativistic field theory, the physical vacuum is not the bare state $|\emptyset\rangle$ devoid of particles, but the filled Dirac sea where all negative energy states are occupied.

\begin{definition}[Dirac Sea Configuration and Vacuum]
For even lattice length $L = 2h$, the half-filled Dirac Fermi sea configuration $S_\Omega \subset \Lambda^*$ consists of all non-positive momentum modes:
\begin{equation}
S_\Omega \coloneqq \{k \in \Lambda^* \mid k \le 0\} = \{-h+1, -h+2, \dots, -1, 0\}.
\end{equation}
The cardinality of the Fermi sea is $\#S_\Omega = h = L/2$. The Dirac sea vacuum state is the corresponding occupation state:
\begin{equation}
|\Omega\rangle \coloneqq |S_\Omega\rangle \in \mathcal{F}(L).
\end{equation}
\end{definition}

\begin{definition}[Bare and Physical Momentum Hamiltonians]
\begin{enumerate}
    \item The bare kinetic Hamiltonian $H_0$ is defined by:
    \begin{equation}
    H_0 \coloneqq \sum_{k \in \Lambda^*} k\, c_k^\dagger c_k = \sum_{k \in \Lambda^*} k\, n_k.
    \end{equation}
    \item The ground state sea energy $E_\Omega$ is the eigenvalue of $H_0$ on $|\Omega\rangle$:
    \begin{equation}
    E_\Omega \coloneqq \sum_{k \in S_\Omega} k.
    \end{equation}
    \item The physical, normal-ordered Hamiltonian $\hat{P}$ is defined by:
    \begin{equation}
    \hat{P} \coloneqq H_0 - E_\Omega \cdot I,
    \end{equation}
    so that $\hat{P} |\Omega\rangle = 0$.
\end{enumerate}
\end{definition}

\begin{lemma}[Evaluation of Ground State Energy]
For $L = 2h$ with $h \ge 1$, the bare ground state energy evaluates to:
\begin{equation}
E_\Omega = -\frac{h(h-1)}{2}.
\end{equation}
\end{lemma}

\begin{proof}
By definition, $S_\Omega = \{-h+1, \dots, 0\}$. Summing the integers:
\[
E_\Omega = \sum_{j=-h+1}^0 j = - \sum_{m=1}^{h-1} m = - \frac{(h-1) h}{2} = -\frac{h(h-1)}{2}.
\]
\end{proof}

\begin{theorem}[Spectra and Commutation Relations]
\label{thm:h0_commutators}
\begin{enumerate}
    \item On any occupation basis state $|S\rangle$ ($S \subseteq \Lambda^*$):
    \begin{equation}
    H_0 |S\rangle = \left( \sum_{k \in S} k \right) |S\rangle, \qquad \hat{P} |S\rangle = \left( \sum_{k \in S} k - E_\Omega \right) |S\rangle.
    \end{equation}
    \item For any mode $p \in \Lambda^*$, the creation and annihilation operators satisfy:
    \begin{align}
    [H_0, c_p^\dagger] &= [\hat{P}, c_p^\dagger] = p\, c_p^\dagger, \\
    [H_0, c_p] &= [\hat{P}, c_p] = -p\, c_p.
    \end{align}
\end{enumerate}
\end{theorem}

\begin{proof}
We prove each statement:
\begin{enumerate}
    \item Each $n_k$ is diagonal on basis states: $n_k |S\rangle = |S\rangle$ if $k \in S$, and $0$ otherwise. Summing over $k \in \Lambda^*$ gives $H_0 |S\rangle = (\sum_{k \in S} k) |S\rangle$. Subtracting $E_\Omega I$ gives the eigenvalue for $\hat{P}$.
    \item Using the bilinear CAR commutator $[n_k, c_p^\dagger] = \delta_{k, p} c_p^\dagger$ from Theorem~\ref{thm:algebraic_car}(2):
    \[
    [H_0, c_p^\dagger] = \left[ \sum_{k \in \Lambda^*} k\, n_k, \; c_p^\dagger \right] = \sum_{k \in \Lambda^*} k\, [n_k, c_p^\dagger] = \sum_{k \in \Lambda^*} k (\delta_{k, p} c_p^\dagger) = p\, c_p^\dagger.
    \]
    Similarly, $[n_k, c_p] = -\delta_{k, p} c_p$ yields $[H_0, c_p] = -p\, c_p$.
    Because $\hat{P} = H_0 - E_\Omega I$ differs from $H_0$ only by a central scalar multiple of the identity, the commutators with $\hat{P}$ are identical.
\end{enumerate}
\end{proof}

Theorem~\ref{thm:h0_commutators} confirms that $c_p^\dagger$ raises the total physical momentum energy by $p$, while $c_p$ lowers it by $p$. This provides the kinematic foundation for constructing particle-hole excitations and density modes in the subsequent chapters.

\chapter{The Lattice AQFT Net}
\label{ch:lattice_aqft_net}

\section{Introduction to Algebraic Quantum Field Theory on a Lattice}
Algebraic Quantum Field Theory (AQFT), following the Haag--Kastler axioms, assigns to each localized spacetime region a unital $*$-algebra of local observables. On the discrete spatial circle $\Lambda = \mathbb{Z}/L\mathbb{Z}$, a localized region is a discrete interval $I \subset \Lambda$.

\section{Local Net of Observable Algebras}
For any non-empty interval $I \subset \Lambda$, we define the local field algebra $\mathfrak{F}(I)$ as the unital $*$-subalgebra generated by $\{\psi_x, \psi_x^\dagger \mid x \in I\}$.
The observable algebra $\mathfrak{A}(I)$ is the even subalgebra consisting of operators invariant under the global fermion parity automorphism $\alpha$:
\begin{equation}
\mathfrak{A}(I) \coloneqq \{A \in \mathfrak{F}(I) \mid \alpha(A) = A\}.
\end{equation}

\section{Core Axiomatic Theorems Proved on the Lattice}
\begin{itemize}
    \item \textbf{Isotony:} If $I_1 \subset I_2$, then $\mathfrak{A}(I_1) \subset \mathfrak{A}(I_2)$.
    \item \textbf{Twisted / Bosonic Locality:} If $I_1$ and $I_2$ are disjoint intervals ($I_1 \cap I_2 = \emptyset$), then for any $A_1 \in \mathfrak{A}(I_1)$ and $A_2 \in \mathfrak{A}(I_2)$, they commute:
    \[
    [A_1, A_2] = 0.
    \]
    Even observables commute at spacelike separation.
    \item \textbf{Irreducibility and Factor Properties:} The global algebra $\mathfrak{A}(\Lambda)$ acts irreducibly on the even/odd fermionic sectors.
    \item \textbf{Boundary Extension (Ch06Ext):} Extension to twisted boundary conditions (antiperiodic fermions) and holonomy phases, showing that the local net remains gauge-invariant and causal under boundary parameter variations.
\end{itemize}

\begin{remark}[Formalization Note]
In the verified Lean~4 core (Chapter 6 and the Boundary Extension module), the Haag--Kastler axioms are proved constructively as theorems rather than postulated as abstract axioms.
\end{remark}

\chapter{Vacuum, Sectors, Energy, and the Budget Space}
\label{ch:vacuum_budget}

\section{Introduction and the Energy Budget Philosophy}

In continuum quantum field theory, the Dirac sea has infinite negative depth. Attempting to define particle-hole operators, charges, and Hamiltonians directly in the continuum leads immediately to divergent vacuum expectations and unbounded operators. 

To develop the finite-lattice bosonization without topological approximations or divergent series, we introduce the \emph{Energy Budget regularization scheme}. The entire state space is decomposed into finite-dimensional subspaces governed by two discrete parameters:
\begin{enumerate}
    \item Net charge sector deviation $N \in \mathbb{Z}$.
    \item Maximum excitation energy cutoff $K \in \mathbb{N}$.
\end{enumerate}

Throughout this chapter, let $L = 2h$ with $h \ge 1$ be an even positive integer. The momentum band is $\Lambda^* = \{-h+1, \dots, h\}$, and $\mathcal{F}(L)$ is the $2^L$-dimensional Fock space.

\section{Admissible Charge Sectors and Sector Grounds}

Because the single-particle band $\Lambda^*$ contains exactly $2h$ modes, an occupation configuration $S \subseteq \Lambda^*$ can have particle number $\#S$ between $0$ and $2h$. 

\begin{definition}[Admissible Charge Sectors]
The net charge deviation $N \in \mathbb{Z}$ of an occupation configuration $S \subseteq \Lambda^*$ relative to half-filling is defined by:
\begin{equation}
N(S) \coloneqq \#S - h \in \mathbb{Z}.
\end{equation}
A charge sector $N$ is \emph{admissible} if it can be realized by a subset of $\Lambda^*$:
\begin{equation}
-h \le N \le h.
\end{equation}
For any admissible sector $N$, the unique \emph{sector ground configuration} $\mathcal{S}_N \subset \Lambda^*$ is the set of the lowest $h + N$ momentum modes:
\begin{equation}
\mathcal{S}_N \coloneqq \{ k \in \Lambda^* \mid k \le N \} = \{-h+1, -h+2, \dots, N\}.
\end{equation}
The \emph{sector ground state} $|N\rangle_0 \in \mathcal{F}(L)$ is the corresponding Fock basis vector:
\begin{equation}
|N\rangle_0 \coloneqq |\mathcal{S}_N\rangle.
\end{equation}
When $N = 0$, $\mathcal{S}_0 = \{-h+1, \dots, 0\}$ is the reference half-filled Dirac vacuum, and $|0\rangle_0 = |\Omega\rangle$.
\end{definition}

\begin{lemma}[Cardinality of Sector Grounds]
For any admissible charge $N \in \{-h, \dots, h\}$:
\begin{equation}
\#\mathcal{S}_N = h + N, \qquad N(\mathcal{S}_N) = N.
\end{equation}
\end{lemma}

\begin{proof}
The modes in $\mathcal{S}_N$ are consecutive integers from $-h+1$ to $N$. The number of elements is $N - (-h+1) + 1 = N + h$. Subtracting $h$ yields $N(\mathcal{S}_N) = (h+N) - h = N$.
\end{proof}

\section{Integer Energy Grading and Observables}

To maintain strict mathematical rigor, all combinatorial quantities are defined as exact signed integers in $\mathbb{Z}$.

\begin{definition}[Integer Charge, Momentum, and Ground Energy Offset]
For any configuration $S \subseteq \Lambda^*$:
\begin{enumerate}
    \item The half-filled reference sea has integer momentum energy $E_\Omega$; the total bare momentum relative to it is:
    \begin{equation}
    E_\Omega \coloneqq \sum_{k\in\mathcal{S}_0}k=-\frac{h(h-1)}2,\qquad
    P(S) \coloneqq \sum_{k\in S}k-E_\Omega\in\mathbb Z.
    \end{equation}
    \item The \emph{sector ground reference energy offset} $E_0(N)$ is defined for $N \in \mathbb{Z}$ by:
    \begin{equation}
    E_0(N) \coloneqq \frac{N(N+1)}{2} \in \mathbb{Z}.
    \end{equation}
    Because one of $N$ or $N+1$ is always even, the product $N(N+1)$ is divisible by $2$ in $\mathbb{Z}$, so $E_0(N)$ is an exact integer for every $N \in \mathbb{Z}$.
    \item The \emph{excitation energy} $e(S)$ of configuration $S$ is defined by:
    \begin{equation}
    e(S) \coloneqq P(S) - E_0(N(S)).
    \end{equation}
\end{enumerate}
\end{definition}

\begin{quote}
\item[]\paragraph{Physical Intuition for the Ground Energy Offset $E_0(N)$.}
Why does the ground state of charge sector $N$ carry energy $E_0(N) = N(N+1)/2$?
Consider adding $N > 0$ particles to the half-filled Dirac sea $\mathcal{S}_0 = \{-h+1, \dots, 0\}$. To minimize total kinetic energy, the particles must occupy the lowest available empty modes immediately above the Fermi surface, which are the consecutive positive modes $k = 1, 2, \dots, N$. The total momentum added is:
\begin{equation}
\sum_{k=1}^N k = 1 + 2 + \dots + N = \frac{N(N+1)}{2}.
\end{equation}
Similarly, if $N < 0$ particles are removed, the lowest-energy configuration removes the highest occupied modes $k = 0, -1, \dots, N+1$ from the sea. The net momentum change is $-\sum_{k=N+1}^0 k = \frac{|N|(|N|-1)}{2} + |N| = \frac{N(N+1)}{2}$.
In 1D physics, this quadratic term $E_0(N) = \frac{1}{2} N^2 + \frac{1}{2} N$ reflects the macroscopic charging energy of a 1D Luttinger liquid: adding particles increases the Fermi momentum $k_F \propto N$, so the total kinetic energy increases quadratically as $\int_0^N k \, dk \sim \frac{1}{2} N^2$.
\end{quote}

\begin{definition}[Operator Observables]
The relative charge operator $N$ and excitation energy operator $E$ belong to $\operatorname{End}_{\mathbb C}(\mathcal F(L))$ and are defined on basis states by:
\begin{align}
N |S\rangle &\coloneqq N(S) |S\rangle = (\#S - h) |S\rangle, \\
E |S\rangle &\coloneqq e(S) |S\rangle.
\end{align}
\end{definition}

\begin{lemma}[Vacuum Ground State Evaluation]
The ground states satisfy:
\begin{equation}
P(\mathcal{S}_N) = E_0(N), \qquad e(\mathcal{S}_N) = 0 \quad \text{for all } N \in \{-h, \dots, h\}.
\end{equation}
In particular, $E |N\rangle_0 = 0$.
\end{lemma}

\begin{proof}
For $N\ge0$, summing the added modes gives
\begin{equation}
P(\mathcal S_N)=\sum_{k=1}^{N}k=\frac{N(N+1)}2=E_0(N).
\end{equation}
For $N<0$, reindexing the removed modes gives
\begin{equation}
P(\mathcal S_N)=-\sum_{j=0}^{-N-1}(N+1+j)=\frac{N(N+1)}2=E_0(N).
\end{equation}
Thus $e(\mathcal S_N)=P(\mathcal S_N)-E_0(N)=0$ in either case.
\end{proof}

\section{The Rank Formula and Ground State Uniqueness}

A central theorem of the energy budget is that the excitation energy $e(S)$ is strictly non-negative for all configurations, and vanishes if and only if $S$ is the sector ground configuration $\mathcal{S}_N$.

\begin{definition}[Ordered Modes, Ground Positions, and Particle Displacements]
Let $S \subseteq \Lambda^*$ be an occupation configuration with cardinality $\#S = h + N$.
\begin{enumerate}
    \item The \emph{ordered occupied modes} of $S$ are the elements of $S$ written in strictly increasing order:
    \begin{equation}
    s_0 < s_1 < s_2 < \dots < s_{h+N-1}.
    \end{equation}
    \item The \emph{ground reference positions} are the consecutive elements of the sector ground $\mathcal{S}_N$:
    \begin{equation}
    g_i \coloneqq -h + 1 + i \quad \text{for } i \in \{0, 1, \dots, h+N-1\}.
    \end{equation}
    \item For each particle index $i \in \{0, 1, \dots, h+N-1\}$, the \emph{particle displacement} $d_i$ is defined by:
    \begin{equation}
    d_i \coloneqq s_i - g_i \in \mathbb{Z}.
    \end{equation}
\end{enumerate}
\end{definition}

\begin{quote}
\item[]\paragraph{Physical Intuition for Particle Displacements.}
The ground state $\mathcal{S}_N$ is a compact Fermi droplet where particles occupy consecutive lowest modes $g_i$. Any excited state $S$ is created by displacing particles upward from their ground positions. The integer $d_i = s_i - g_i$ measures how far the $i$-th fermion has been promoted above its ground state reference position.
\end{quote}

\begin{theorem}[Rank Formula and Non-Negativity of Excitation Energy]
\label{thm:rank_formula}
Let $S \subseteq \Lambda^*$ be any configuration with cardinality $\#S = h + N$ (so $N(S) = N$). 
\begin{enumerate}
    \item The displacements are non-negative, non-decreasing, and bounded above by the band edge:
    \begin{equation}
    h+N=0\;:\;S=\emptyset,\qquad
    h+N>0\;:\;0\le d_0\le\cdots\le d_{h+N-1}\le 2h-(h+N)=h-N.
    \end{equation}
    \item The excitation energy is the sum of particle displacements:
    \begin{equation}
    e(S) = \sum_{i=0}^{h+N-1} d_i.
    \end{equation}
    \item Non-negativity: $e(S) \ge 0$ for all $S \subseteq \Lambda^*$.
    \item Uniqueness: $e(S) = 0$ if and only if $S = \mathcal{S}_N$.
\end{enumerate}
\end{theorem}

\begin{proof}
We prove each assertion:
\begin{enumerate}
    \item If $S=\emptyset$, the displacement assertions are vacuous. Otherwise $s_0\ge g_0$, so $d_0\ge0$; for consecutive indices, $s_{i+1}-s_i\ge1$ and $g_{i+1}-g_i=1$, hence $d_{i+1}-d_i\ge0$. Finally, $s_{h+N-1}\le h$ and $g_{h+N-1}=N$ give $d_{h+N-1}\le h-N$.

    \item Using $P(S)=\sum_{k\in S}k-E_\Omega$ and $P(\mathcal S_N)=E_0(N)$, the sorted occupations and ground positions give the telescoping identity
    \begin{equation}
    e(S)=P(S)-E_0(N)=\sum_{k\in S}k-\sum_{k\in\mathcal S_N}k
    =\sum_{i=0}^{h+N-1}(s_i-g_i)=\sum_{i=0}^{h+N-1}d_i.
    \end{equation}
    \item Since each $d_i \ge 0$, the sum of non-negative integers is non-negative: $e(S) \ge 0$.
    
    \item Since $d_i \ge 0$, the sum $\sum d_i = 0$ holds if and only if $d_i = 0$ for every $i \in \{0, \dots, h+N-1\}$.
    This means $s_i = g_i$ for all $i$, which implies $S = \mathcal{S}_N$.
\end{enumerate}
\end{proof}

\begin{corollary}[One-Dimensional Ground Subspace]
\label{cor:ground_subspace}
In each admissible sector $N$, the zero-excitation energy subspace is strictly one-dimensional:
\begin{equation}
\mathcal{H}(N, 0) \coloneqq \mathrm{span}_{\mathbb{C}} \{ |S\rangle \mid N(S) = N, \; e(S) = 0 \} = \mathbb{C} \cdot |N\rangle_0.
\end{equation}
\end{corollary}

\begin{quote}
\item[]\paragraph{Ground-state uniqueness and the gapless spectrum.}
In condensed matter physics, Corollary~\ref{cor:ground_subspace} carries profound physical meaning:
\begin{enumerate}
    \item \textbf{Non-Degenerate Ground State:} For any fixed particle number $h + N$, the ground state is unique (non-degenerate). There are no spontaneous symmetry-breaking degeneracies at zero temperature in 1D.
    \item \textbf{The Energy Gap:} The lowest possible non-zero excitation energy is obtained by setting $d_{h+N-1} = 1$ and all other $d_i = 0$, which corresponds to promoting the highest fermion across the Fermi level into mode $N+1$. The excitation energy of this state is $e(S) = 1$. In physical energy units, the excitation gap is $\Delta \epsilon = \frac{2\pi v_F}{L}$. As $L \to \infty$, the gap vanishes as $\mathcal{O}(1/L)$, confirming that the 1D free Fermi gas is a \emph{gapless (critical) system} with a continuous low-energy spectrum.
\end{enumerate}
\end{quote}

\section{The Frozen Margins Theorem}

Because the displacements $d_i$ are non-negative and non-decreasing, an upper bound on the total excitation energy $e(S) \le K$ strongly constrains the positions of particles at the bottom and top of the band.

\begin{theorem}[Frozen Margins]
\label{thm:frozen_margins}
Let $S \subseteq \Lambda^*$ be a configuration with $N(S) = N$ and excitation energy $e(S) \le K$.
\begin{enumerate}
    \item \textbf{Deep Sea Rigidity:} Every momentum mode deep below the Fermi surface is rigidly occupied:
    \begin{equation}
    \text{For all } k \in \Lambda^* \text{ with } k \le N - K, \quad k \in S.
    \end{equation}
    \item \textbf{High Momentum Rigidity:} Every momentum mode high above the Fermi surface is rigidly empty:
    \begin{equation}
    \text{For all } k \in \Lambda^* \text{ with } k > N + K, \quad k \notin S.
    \end{equation}
\end{enumerate}
\end{theorem}

\begin{proof}
By Theorem~\ref{thm:rank_formula}, $e(S) = \sum_{i=0}^{h+N-1} d_i \le K$. Since $d_i \ge 0$ for all $i$, each individual displacement satisfies $d_i \le K$.
\begin{enumerate}
    \item If a mode $k\le N-K$ were unoccupied, the resulting hole would force total displacement at least $N-k+1\ge K+1$, contradicting $\sum_i d_i=e(S)\le K$. Thus every such mode is occupied.
    \item If an occupied mode $k>N+K$ existed, the highest occupied mode would satisfy $s_{h+N-1}\ge k$ while its ground position is $g_{h+N-1}=N$. Hence $d_{h+N-1}\ge k-N>K$, contradicting $\sum_i d_i\le K$. Thus every such mode is empty.
\end{enumerate}
\end{proof}

\section{Energy Budget Subspaces and Filtered Operators}

\subsection{Coordinate Subspaces and Projector Algebra}
For $U\subseteq\mathcal P(\Lambda^*)$, define the coordinate subspace and its orthogonal projector by
\begin{equation}
\operatorname{coord}(U)=\operatorname{span}_{\mathbb C}\{|S\rangle:S\in U\},\qquad
P_U|S\rangle=\begin{cases}|S\rangle,&S\in U,\\0,&S\notin U,\end{cases}
\quad P_U\in\operatorname{End}_{\mathbb C}(\mathcal F(L)).
\end{equation}
Their algebra is
\begin{align}
P_U^2&=P_U, & P_U^\dagger&=P_U, & \operatorname{Ran}(P_U)&=\operatorname{coord}(U),\\
P_UP_V&=P_{U\cap V}, & P_{U\cup V}&=P_U+P_V-P_{U\cap V}.
\end{align}
Thus disjoint coordinate sets give orthogonal subspaces and $P_UP_V=0$. These identities follow by evaluating both sides on each occupation basis ket.


\begin{definition}[Energy Budget Subspaces]
\label{def:budget_subspace}
For sector $N \in \{-h, \dots, h\}$ and cutoff $K \in \mathbb{N}$:
\begin{enumerate}
    \item The \emph{fixed-energy subspace} $\mathcal{H}(N, E)$ is the span of basis states with charge $N$ and energy $E$:
    \begin{equation}
    \mathcal{H}(N, E) \coloneqq \mathrm{span}_{\mathbb{C}} \{ |S\rangle \mid N(S) = N, \; e(S) = E \}.
    \end{equation}
    \item The \emph{fixed-charge energy budget subspace} $B(N, K)$ is the internal orthogonal direct sum:
    \begin{equation}
    B(N, K) \coloneqq \bigoplus_{E=0}^K \mathcal{H}(N, E) = \mathrm{span}_{\mathbb{C}} \{ |S\rangle \mid N(S) = N, \; e(S) \le K \}.
    \end{equation}
    \item The \emph{charge-box budget subspace} $\mathcal{B}_{K, N_{\max}}$ for $N_{\max} \ge 0$ is:
    \begin{equation}
    \mathcal{B}_{K, N_{\max}} \coloneqq \bigoplus_{|N| \le N_{\max}} B(N, K) = \mathrm{span}_{\mathbb{C}} \{ |S\rangle \mid e(S) \le K, \; |N(S)| \le N_{\max} \}.
    \end{equation}
\end{enumerate}
The orthogonal coordinate projectors onto these subspaces are denoted $P_{N, K}$ and $P_{K, N_{\max}}$, respectively.
\end{definition}

\begin{definition}[Filtered Operator Parameters: Charge Shift and Energy Increase]
Let $A \in \mathrm{End}_{\mathbb{C}}(\mathcal{F}(L))$ be a linear operator.
\begin{enumerate}
    \item $A$ carries a \emph{charge shift} $q \in \mathbb{Z}$ if for every basis state $|S\rangle$, every output state $|S'\rangle$ in the support of $A |S\rangle$ has charge:
    \begin{equation}
    N(S') = N(S) + q.
    \end{equation}
    \item $A$ carries an \emph{energy increase bound} $d \in \mathbb{N}$ if for every basis state $|S\rangle$, every output state $|S'\rangle$ in the support of $A |S\rangle$ has excitation energy:
    \begin{equation}
    e(S') \le e(S) + d.
    \end{equation}
    \item The map is an \emph{exact charge-energy shift} $(q,d)\in\mathbb Z^2$ if every output configuration in $A|S\rangle$ satisfies both $N(S')=N(S)+q$ and $e(S')=e(S)+d$.
\end{enumerate}
\end{definition}

\begin{theorem}[Filtered Operator Action]
\label{thm:filtered_action}
Let $A$ be a linear operator with charge shift $q \in \mathbb{Z}$ and energy increase bound $d \in \mathbb{N}$.
Then $A$ maps the budget subspace $B(N, K)$ into $B(N+q, K+d)$:
\begin{equation}
A\big( B(N, K) \big) \subseteq B(N+q, K+d).
\end{equation}
\end{theorem}

\begin{proof}
Every vector in $B(N, K)$ is a linear combination of basis vectors $|S\rangle$ satisfying $N(S) = N$ and $e(S) \le K$. Applying $A$, each output basis vector $|S'\rangle$ satisfies $N(S') = N + q$ and $e(S') \le e(S) + d \le K + d$, which lies in $B(N+q, K+d)$.
\end{proof}

\begin{theorem}[Unified Margin Inequality]
\label{thm:unified_margin}
Let density operators with maximum mode cutoff $M \ge 1$ act on states in the charge-box budget $\mathcal{B}_{K, N_{\max}}$. If intermediate states undergo cumulative upward energy excursions of at most $K_{\mathrm{excursion}}$, the condition ensuring that all active modes stay strictly within the linear interior of the band without hitting the Nyquist boundaries is:
\begin{equation}
2M + K + K_{\mathrm{excursion}} + N_{\max} \le h.
\end{equation}
Under this condition, no density shift triggers an Umklapp boundary wrap, so the specified mode transfers retain their unwrapped integer action on the stated budget.
\end{theorem}

\begin{proof}
For every intermediate state, any highest occupied mode is at most $N+K+K_{\mathrm{excursion}}$, while any lowest unoccupied mode is at least $N-K-K_{\mathrm{excursion}}+1$. Using $|N|\le N_{\max}$ and allowing transfers up to $M$ with the corresponding edge buffer, the unified condition $2M+K+K_{\mathrm{excursion}}+N_{\max}\le h$ keeps both extremal modes in the linear interior and prevents Umklapp wrapping.
\end{proof}

\begin{remark}[Finite Algebraic Foundations]
The rank formula, coordinate projectors, filtered maps, and word excursion bounds follow from finite sums on the occupation basis. The vacuum projectors and displacement description make the energy-budget regularization explicit without invoking an abstract dimension count or an infinite-volume limit.
\end{remark}

\section{Word Excursion Tracking and Compression Remainders}

For a product of operators $W=A_1A_2\cdots A_s$ with exact energy shifts $\Delta e_j$ and charge shifts $q_j$, the rightmost factor acts first. Define its maximum upward energy excursion by
\begin{equation}
X(W)=\max\left(0,\max_{1\le r\le s}\sum_{j=1}^{r}\Delta e_{s-j+1}\right).
\end{equation}

\begin{theorem}[Word Evaluation Safety]
\label{thm:word_evaluation_safety}
Suppose $\psi\in B(N,K)$ and $K+X(W)\le K_{\max}$. After applying the first $r$ factors from the right, every intermediate vector lies in
\begin{equation}
B\left(N+\sum_{j=1}^{r}q_{s-j+1},K_{\max}\right),
\end{equation}
so every intermediate state remains within the energy budget $K_{\max}$.
\end{theorem}

\begin{proof}
After $r$ right-to-left applications the cumulative energy shift is $\sum_{j=1}^{r}\Delta e_{s-j+1}\le X(W)$. Starting from energy at most $K$, the intermediate energy is therefore at most $K+X(W)\le K_{\max}$; the charge shift is the sum of the corresponding $r$ charge shifts.
\end{proof}

\begin{proposition}[Compression Remainder Identity]
\label{prop:compression_remainder}
For a coordinate projector $P=P_U$ and operators $A,B\in\operatorname{End}_{\mathbb C}(\mathcal F(L))$, the product of compressed operators differs from the compressed product by the exact intermediate leakage term
\begin{equation}
PABP-(PAP)(PBP)=PA(I-P)BP.
\end{equation}
\end{proposition}

\begin{proof}
Expanding $PABP-PAPBP$ and factoring on the left and right gives $PA(I-P)BP$.
\end{proof}

\section{Chapter Summary and Literature Comparison}

In this chapter, we developed the Energy Budget regularization scheme:
\begin{itemize}
    \item Admissible charge sectors $N \in \{-h, \dots, h\}$ possess unique ground configurations $\mathcal{S}_N = \{-h+1, \dots, N\}$.
    \item The sector ground reference offset $E_0(N) = N(N+1)/2$ evaluates the macroscopic charging energy in exact integer arithmetic.
    \item Defined particle displacements $d_i = s_i - g_i$ and proved the Rank Formula $e(S) = \sum d_i \ge 0$, establishing ground state uniqueness $e(S) = 0 \iff S = \mathcal{S}_N$.
    \item Connected ground state non-degeneracy to the gapless $\mathcal{O}(1/L)$ critical scaling of the 1D Fermi gas.
    \item Proved the Frozen Margins Theorem: states with $e(S) \le K$ have rigidly filled deep seas ($k \le N - K$) and rigidly empty high modes ($k > N + K$).
    \item Defined budget subspaces $B(N, K) = \bigoplus_{E=0}^K \mathcal{H}(N, E)$, specified filtered operator parameters $(q, d)$, and proved the Unified Margin Inequality $2M + K + K_{\mathrm{excursion}} + N_{\max} \le h$.
    \item Established the coordinate-projector algebra, tracked cumulative energy excursions of operator words, and identified the exact leakage term produced by compressing an operator product.
\end{itemize}

\begin{quote}
\item[]\paragraph{Comparison with Standard Continuum Literature.}
In standard continuum bosonization (Giamarchi~\cite{giamarchi2004}, Gogolin et al.~\cite{gogolin1998}, and von Delft and Schoeller~\cite{vondelft1998}), many treatments use smooth regulators, such as exponential damping factors $e^{-\alpha |k|}$ or Gaussian filters, to control momentum sums. The present construction instead specifies finite-dimensional orthogonal subspaces directly through integer charge and excitation-energy conditions.

The energy budget $B(N,K)$ is an exact finite-dimensional orthogonal subspace. Theorem~\ref{thm:unified_margin} gives a sufficient margin under which the specified finite-lattice mode transfers avoid boundary wrapping.
\end{quote}

\chapter{The Umbral Boson Fock Layer}
\label{ch:boson_fock}

\section{Introduction and Purely Algebraic Bosons}
Traditional bosonic Fock spaces require infinite-dimensional Hilbert spaces of square-integrable wavefunctions, leading to unbounded creation/annihilation operators and domain subtleties. 
In the Umbral algebraic bosonization framework, the bosonic carrier is realized purely algebraically as the ring of multivariate polynomials:
\begin{equation}
\mathcal{F}_b \coloneqq \mathbb{C}[X_1, X_2, \dots].
\end{equation}

\section{Raw Currents and the Haldane Inner Product}
For mode integer $m \ge 1$:
\begin{itemize}
    \item \textbf{Creation Operator (Current):} $C_m \coloneqq X_m$ (multiplication by indeterminate $X_m$).
    \item \textbf{Annihilation Operator:} $A_m \coloneqq m \frac{\partial}{\partial X_m}$.
\end{itemize}
These operators satisfy the exact Canonical Commutation Relations (CCR):
\begin{equation}
[A_m, C_n] = m\, \delta_{m, n}\, I.
\end{equation}
The Haldane inner product is the unique positive-definite Hermitian form making $A_m^\dagger = C_m$. Monomials $\prod_{i} X_{\lambda_i}$ are mutually orthogonal with norm squared:
\begin{equation}
\left\langle \prod X_{\lambda_i}, \prod X_{\mu_j} \right\rangle = z_\lambda\, \delta_{\lambda, \mu},
\end{equation}
where $z_\lambda = \prod_m m^{a_m} a_m!$ is the standard partition symmetry factor.

\section{Key Results and Formal Discoveries}
\begin{itemize}
    \item \textbf{Nilpotent Exponentials:} On any weight-truncated subspace $\mathcal{F}_{\le K}$, the operator $A_m$ is nilpotent. Consequently, the formal exponential series $\exp(t A_m)$ terminates after finitely many terms as an exact polynomial operator without analytic limits.
    \item \textbf{Exact Baker--Campbell--Hausdorff (BCH) Formula:} For commute-to-scalar operators, $e^{A} e^{B} = e^{A+B} e^{\frac{1}{2}[A, B]}$ is an exact finite algebraic identity on truncated polynomial spaces.
\end{itemize}

\begin{remark}[Formalization Note]
In the verified Lean~4 core (Chapter 8, Boson Fock Layer), this entire bosonic structure is proved constructively over polynomials with rational and complex coefficients without analytic convergence theorems.
\end{remark}


\part{The Density Sector and Completeness}

\chapter{Density Modes and Kinematics}
\label{ch:density_modes}

\section{Introduction and the Bridge to Bosons}
The bridge between fermionic Fock space and bosonic current algebra is formed by the fermionic density modes $\rho_m$. In the continuum, $\rho(x) = \psi^\dagger(x) \psi(x)$ is singular and requires subtractive renormalization. On the discrete band $\Lambda^*$, density modes are defined as exact truncated momentum-shift operators.

\section{Truncated and Normal-Ordered Densities}
For mode integer $m \in \mathbb{Z} \setminus \{0\}$, the raw density operator on the momentum CAR algebra is:
\begin{equation}
\rho_m^{\mathrm{raw}} \coloneqq \sum_{k \in \Lambda^*} c_{k \ominus m}^\dagger c_k.
\end{equation}
The normal-ordered density mode with respect to the Dirac sea $S_N$ is defined by:
\begin{equation}
:\!\rho_m\!: \coloneqq \sum_{k \in \Lambda^*} : \! c_{k \ominus m}^\dagger c_k \!:_{S_N},
\end{equation}
where $: \! c_p^\dagger c_k \!:_{S_N} = c_p^\dagger c_k - \langle N|_0 c_p^\dagger c_k |N\rangle_0 I$.

\section{Key Results and Formal Discoveries}
\begin{itemize}
    \item \textbf{Energy Grading:} For $m > 0$, the normal-ordered density $:\!\rho_m\!:$ acts as an annihilation-type operator that lowers excitation energy by $m$, while $:\!\rho_{-m}\!:$ raises excitation energy by $m$ within the admissible margin.
    \item \textbf{Charge Conservation:} The density modes commute with the total fermion number:
    \[
    [\hat{N}, :\!\rho_m\!:] = 0.
    \]
    Density excitations conserve net fermion charge.
    \item \textbf{Action on the Vacuum:} The ground state $|N\rangle_0$ is annihilated by positive density modes:
    \[
    :\!\rho_m\!: |N\rangle_0 = 0 \quad \text{for all } m > 0.
    \]
\end{itemize}

\begin{remark}[Formalization Note]
In the verified Lean~4 core (Chapter 9 and Appendix A04), the non-wrapping kinematics and margin inequalities are formalized constructively, ensuring that density shifts do not cause boundary wraps when applied to states within the budget.
\end{remark}

\chapter{Heisenberg Algebra and the Exact Schwinger Term}
\label{ch:heisenberg_algebra}

\section{Introduction and the Algebraic Origin of the Schwinger Anomaly}

The central mathematical miracle of quantum bosonization is that bilinear products of fermionic field operators satisfy bosonic canonical commutation relations. In continuous quantum field theory, this phenomenon is expressed by the affine Lie algebra $\widehat{\mathfrak{u}}(1)$ (the $U(1)$ Kac--Moody current algebra):
\begin{equation}
[J_m, J_n] = m\, \delta_{m+n, 0}\, I.
\end{equation}
The central term $m\, \delta_{m+n, 0}$, known as the \emph{Schwinger term} or the chiral anomaly, is traditionally derived via heuristic point-splitting, formal normal ordering, or infinite Dirac sea subtractions.

On our discrete lattice $\Lambda = \mathbb{Z}/L\mathbb{Z}$ ($L = 2h$, $h \ge 1$), the situation is fundamentally clarified. 
As proven by the No-Go Theorem~\ref{thm:no_go}, no finite-dimensional matrix algebra can admit operators satisfying $[A, B] = I$ because the trace of any commutator must vanish identically. How, then, can fermionic density operators on the finite Fock space $\mathcal{F}(L)$ satisfy bosonic commutation relations?

\begin{remark}[Energy Budget Regularization]
The answer lies in the \emph{Energy Budget regularization}:
\begin{enumerate}
    \item On the full Fock space $\mathcal{F}(L)$, the commutator $[\rho_{-m}, \rho_m]$ is not a scalar, but a non-trivial dynamical operator given by the difference between the fermion occupations at the bottom and top edges of the band.
    \item When evaluated on vectors in the energy budget subspace $B(N, K)$ satisfying explicit \emph{margin inequalities}, the deep Fermi sea is frozen full and the high momentum band is frozen empty.
    \item Consequently, the edge operators evaluate to the exact positive integer $m$ on budget states, realizing the exact Kac--Moody algebra without any infinite-dimensional or topological approximations.
\end{enumerate}
\end{remark}

\section{The Exact Diagonal Edge Formula on Fock Space}

We first evaluate the commutator of opposing density modes on the entire Fock space $\mathcal{F}(L)$ without any subspace restriction.

Recall that $\rho_m = d\Gamma(T_m)$ is the second quantization of the one-particle partial shift matrix $T_m$ on $\ell^2(\Lambda^*)$ (Chapter~\ref{ch:density_modes}). By Lemma~\ref{lem:dgamma_comm}, second quantization preserves Lie brackets:
\begin{equation}
[\rho_{-m}, \rho_m] = d\Gamma([T_{-m}, T_m]).
\end{equation}

\begin{theorem}[Exact Diagonal Edge Formula]
\label{thm:diagonal_edge_formula}
Let $L = 2h$ with $h \ge 1$. For any integer $1 \le m \le h$, the algebraic commutator $[\rho_{-m}, \rho_m]$ on the entire Fock space $\mathcal{F}(L)$ satisfies the exact operator identity:
\begin{equation}
[\rho_{-m}, \rho_m] = \sum_{q = -h+1}^{-h+m} n_q - \sum_{q = h-m+1}^{h} n_q,
\end{equation}
where $n_q \coloneqq c_q^\dagger c_q$ is the momentum occupation number operator at mode $q$.
\end{theorem}

\begin{proof}
The matrices $T_{-m}T_m$ and $T_mT_{-m}$ are diagonal: the first is the identity on modes $q\le h-m$, while the second is the identity on modes $q\ge-h+1+m$. Their diagonal entries cancel on the bulk interval $-h+1+m\le q\le h-m$; the difference is $+1$ on the bottom edge and $-1$ on the top edge, with all off-diagonal entries zero. Thus
\begin{equation}
[T_{-m},T_m]=\sum_{q=-h+1}^{-h+m}E_{q,q}-\sum_{q=h-m+1}^{h}E_{q,q}.
\end{equation}
Applying $d\Gamma(E_{q,q})=c_q^\dagger c_q=n_q$ gives the asserted identity.
\end{proof}

\begin{definition}[Edge Domains]
For any mode $1 \le m \le h$, the \emph{Bottom Edge Domain} $\mathcal{E}_{\mathrm{bot}}(m)$ and \emph{Top Edge Domain} $\mathcal{E}_{\mathrm{top}}(m)$ are the subsets of $\Lambda^*$:
\begin{align}
\mathcal{E}_{\mathrm{bot}}(m) &\coloneqq \{ q \in \Lambda^* \mid q \le -h + m \} = \{-h+1, -h+2, \dots, -h+m\}, \\
\mathcal{E}_{\mathrm{top}}(m) &\coloneqq \{ q \in \Lambda^* \mid q > h - m \} = \{h - m + 1, \dots, h\}.
\end{align}
Both sets contain exactly $m$ momentum modes: $\#\mathcal{E}_{\mathrm{bot}}(m) = \#\mathcal{E}_{\mathrm{top}}(m) = m$.
\end{definition}

\section{Regime M1: The Diagonal Buffer Zone and the Schwinger Term}

We now evaluate the edge formula on states belonging to the energy budget subspace $B(N, K)$.

\begin{theorem}[Regime M1: Exact Diagonal Schwinger Term]
\label{thm:regime_m1}
Let $\psi \in B(N, K)$ be any state in the energy budget with charge $N \in \{-h, \dots, h\}$ and excitation energy at most $K \ge 0$.
If the mode $m \ge 1$ satisfies the \textbf{M1 Margin Condition}:
\begin{equation}
m + K + |N| \le h,
\end{equation}
then the commutator $[\rho_{-m}, \rho_m]$ evaluates on $\psi$ to the exact scalar Schwinger term:
\begin{equation}
[\rho_{-m}, \rho_m] \psi = m \psi.
\end{equation}
\end{theorem}

\begin{proof}
It suffices by linearity to take a basis vector $|S\rangle$ in the budget. The edge formula gives
\begin{equation}
[\rho_{-m},\rho_m]|S\rangle=\left(\sum_{q\in\mathcal{E}_{\mathrm{bot}}(m)}n_q-\sum_{q\in\mathcal{E}_{\mathrm{top}}(m)}n_q\right)|S\rangle.
\end{equation}
The M1 condition places every bottom edge mode at or below the deep-sea rigidity threshold $N-K$, so all $m$ such modes are occupied; it places every top edge mode strictly above the high-momentum threshold $N+K$, so those $m$ modes are empty. The displayed edge sum is therefore $(m-0)|S\rangle=m|S\rangle$. Linearity gives the result for every $\psi\in B(N,K)$.
\end{proof}

\section{Regime M2: Off-Diagonal Suppression and Pauli Blocking}

We now examine off-diagonal commutators $[\rho_m, \rho_n]$ where $m + n \neq 0$.

\begin{theorem}[Regime M2: Off-Diagonal Suppression]
\label{thm:regime_m2}
Let $\psi \in B(N, K)$. For any mode transfers $m, n \in \mathbb{Z}$ with $m + n \neq 0$, if the parameters satisfy the \textbf{M2 Margin Condition}:
\begin{equation}
|m| + |n| + K + |N| \le h,
\end{equation}
then the density modes commute identically on $\psi$:
\begin{equation}
[\rho_m, \rho_n] \psi = 0.
\end{equation}
\end{theorem}

\begin{proof}
If $m,n$ have the same sign, the global same-sign commutativity theorem applies. For opposite signs, assume without loss of generality $m>0$ and $n=-k<0$, with $m\ne k$. The one-particle commutator consists of hops across the two band edges. Under the M2 margin, each such hop either has its source in the rigidly empty high-momentum region, or its destination in the rigidly occupied deep sea. In the first case annihilation kills the source particle; in the second Pauli exclusion blocks creation at the destination. Thus every edge hop annihilates each budget basis vector, and hence $[\rho_m,\rho_n]\psi=0$ throughout $B(N,K)$.
\end{proof}

\section[The Exact U(1) Kac--Moody Algebra on the Budget Subspace]{The Exact $U(1)$ Kac--Moody Algebra on the Budget Subspace}

Combining Regimes M1 and M2 yields the central theorem of constructive bosonization.

\begin{theorem}[$U(1)$ Kac--Moody Current Algebra on $B(N, K)$]
\label{thm:kac_moody}
Let $M \ge 1$ be a mode cutoff. Under the uniform margin condition:
\begin{equation}
2M + K + |N| \le h,
\end{equation}
the normal-ordered density modes $:\!\rho_m\!:$ satisfy the exact Kac--Moody algebra for all $|m|, |n| \le M$ on the budget subspace $B(N, K)$:
\begin{equation}
\left[ :\!\rho_m\!:, :\!\rho_n\!: \right] \psi = -m\, \delta_{m+n, 0}\, \psi \quad \text{for all } \psi \in B(N, K).
\end{equation}
Equivalently, setting $n = -m$ with $m \ge 1$:
\begin{equation}
\left[ :\!\rho_{-m}\!:, :\!\rho_m\!: \right] \psi = m \psi.
\end{equation}
\end{theorem}

\begin{proof}
If either mode is zero, use $:\!\rho_0\!:=N$ and charge covariance to obtain a zero commutator. When $m+n\ne0$, the uniform margin implies the M2 condition, so Theorem~\ref{thm:regime_m2} gives zero. When $m+n=0$, M1 applies to the positive mode $k=|m|$: if $m<0$, it gives the stated commutator directly, while if $m>0$, antisymmetry reverses the order and gives the stated negative sign. These cases establish the formula.
\end{proof}

\begin{definition}[Harmonic Oscillator Bosonic Modes]
For each mode $m \ge 1$, we define the normalized bosonic creation and annihilation operators on $B(N, K)$ by:
\begin{equation}
b_m \coloneqq \frac{1}{\sqrt{m}} :\!\rho_{-m}\!:, \qquad b_m^\dagger \coloneqq \frac{1}{\sqrt{m}} :\!\rho_m\!: = (b_m)^\dagger.
\end{equation}
Under the margin condition of Theorem~\ref{thm:kac_moody}, these operators satisfy the standard harmonic oscillator algebra on $B(N, K)$:
\begin{equation}
[b_m, b_n^\dagger] \psi = \delta_{m, n} \psi, \qquad [b_m, b_n] \psi = 0, \qquad [b_m^\dagger, b_n^\dagger] \psi = 0.
\end{equation}
\end{definition}

\section{Word-Excursion Tracking in Composed Products}

When evaluating products of density operators $\rho_{a_1} \rho_{a_2} \cdots \rho_{a_s}$ (such as in vertex operators or multi-mode partition states), intermediate states undergo transient energy fluctuations.

\begin{remark}[Interpreting Energy Excursions]
For example, the written word $W=\rho_{-2}\rho_3\rho_{-1}$ acts first by $\rho_{-1}$, then $\rho_3$, then $\rho_{-2}$. The successive energy changes are $-1,3,-2$, with partial sums $-1,2,0$, so the largest upward excursion is $2$. This maximum tells us how much extra energy budget must be available for every intermediate state, even though the final net change is zero.
\end{remark}

\begin{definition}[Application Order and Cumulative Energy Excursion]
For a written operator product $W = \rho_{a_1} \rho_{a_2} \cdots \rho_{a_s}$, the operators act from right to left in the sequence of signed transfers $(a_s, a_{s-1}, \dots, a_1)$.
\begin{enumerate}
    \item The partial energy shift after applying the first $r$ operators ($0 \le r \le s$) is:
    \begin{equation}
    \Delta E_r \coloneqq \sum_{j=1}^r a_{s - j + 1}.
    \end{equation}
    \item The \emph{maximum upward energy excursion} $X(W)$ is:
    \begin{equation}
    X(W) \coloneqq \max\left(0, \; \max_{0 \le r \le s} \Delta E_r \right).
    \end{equation}
\end{enumerate}
\end{definition}

\begin{theorem}[Valid Margin Excursion Accounting]
\label{thm:excursion_accounting}
Let $W = \rho_{a_1} \dots \rho_{a_s}$ be an operator word with mode cutoff $|a_j| \le M$. If the input state $\psi \in B(N, K)$ satisfies:
\begin{equation}
2M + K + X(W) + |N| \le h,
\end{equation}
then every intermediate state generated during the evaluation of $W \psi$ belongs to $B(N, K + X(W))$, and every commutator $[:\!\rho_m\!:, :\!\rho_n\!:]$ encountered along the evaluation acts as the exact scalar $-m \delta_{m+n, 0}$.
\end{theorem}

\begin{proof}
By Theorem~\ref{thm:budget_action_annihilation}, applying $\rho_{a_s}$ to $\psi \in B(N, K)$ produces a state in $B(N, K + a_s)$. By induction, after $r$ steps the state belongs to $B(N, K + \Delta E_r) \subseteq B(N, K + X(W))$.
Since $2M + (K + X(W)) + |N| \le h$, each intermediate state satisfies the uniform margin condition of Theorem~\ref{thm:kac_moody}. Therefore, all intermediate commutators evaluate to exact scalars.
\end{proof}

\begin{remark}[Budget-Local Resolution of the Trace Obstruction]
On the full finite Fock space, $[\rho_{-m},\rho_m]$ is the difference of the bottom- and top-edge occupation sums, so its trace vanishes and it is not the scalar $mI$. On a budget subspace with the M1 margins, the bottom edge is fully occupied and the top edge empty; the same operator therefore evaluates as $mI$ there. The central term is thus a local identity on the physically selected low-energy states, consistent with the finite-dimensional trace constraint.
\end{remark}

\section{Chapter Summary}
The central result is a budget-local Heisenberg relation for the finite
density modes. It resolves the apparent conflict with the trace identity:
the commutator is traceless on the full finite Fock space, but evaluates to
the central scalar on the low-energy states protected by the margin
condition. Tracking intermediate energy excursions explains exactly when
products and reordered words remain inside that valid region. This gives a
finite, domain-conscious counterpart to the continuum current algebra and
sets up the partition-counting proof of Haldane completeness.

\chapter{Haldane Completeness}
\label{ch:haldane_completeness}

\section{Introduction and the Completeness Problem}
While Chapter~\ref{ch:heisenberg_algebra} establishes that fermionic density operators satisfy the bosonic Heisenberg algebra on budget states, a crucial question remains: \emph{do the bosonic density states span the entire fermionic low-energy Hilbert space?}
Haldane completeness answers this affirmatively: the low-energy fermionic Fock space is isomorphic to the bosonic Fock space, with basis states in exact one-to-one correspondence with integer partitions.

\section{Bosonic Partition States}
An integer partition $\lambda \vdash E$ is a non-increasing sequence of positive integers $\lambda = (\lambda_1 \ge \lambda_2 \ge \dots \ge \lambda_\ell > 0)$ such that $\sum_{i=1}^\ell \lambda_i = E$.
For each partition $\lambda \vdash E$, the corresponding bosonic partition state in sector $N$ is defined by:
\begin{equation}
|\lambda; N\rangle \coloneqq :\! \rho_{-\lambda_1} \!: :\! \rho_{-\lambda_2} \!: \dots :\! \rho_{-\lambda_\ell} \!: |N\rangle_0.
\end{equation}

\section{Key Results and Proved Theorems}
\begin{itemize}
    \item \textbf{Regime R1 (Dimension Equality):} The dimension of the fermionic fixed-energy subspace $\mathcal{H}(N, E)$ equals the number of integer partitions $p(E)$ when $E$ satisfies the margin bound $E \le h - |N|$:
    \[
    \dim \mathcal{H}(N, E) = p(E).
    \]
    This is established via an explicit, constructive bijection between fermionic particle-hole excitation configurations and integer partitions.
    \item \textbf{Regime R2 (Hall--Gram Invertibility):} The Gram matrix of partition states:
    \[
    G_{\lambda, \mu} \coloneqq \langle \lambda; N | \mu; N \rangle
    \]
    is strictly positive-definite on the budget subspace under the margin hypotheses.
    \item \textbf{Whole Budget Basis:} The set of partition states:
    \[
    \mathcal{B}_{\mathrm{basis}}(N, K) \coloneqq \bigcup_{E=0}^K \{|\lambda; N\rangle \mid \lambda \vdash E\}
    \]
    forms a complete, linearly independent basis of the whole budget subspace $B(N, K)$.
\end{itemize}

\begin{remark}[Formalization Note]
In the verified Lean~4 core (Chapter 11 and Appendix A04), 61 proved theorems establish the occupation--partition bijection, the strict dimension count, the non-singularity of the Hall--Gram matrix, and the explicit basis expansion without non-constructive principles.
\end{remark}


\part{Outlook: The Bosonization Dictionary and Interactions}

\chapter*{The Forward Research Roadmap}
\addcontentsline{toc}{chapter}{The Forward Research Roadmap}
\label{ch:roadmap}

Chapters~\ref{ch:lattice_band} through~\ref{ch:haldane_completeness} establish the foundational discrete kinematics of our models: from spatial lattices and Fourier transforms to finite-dimensional CAR Fock spaces, the Heisenberg algebra of density modes, and Haldane completeness on the finite energy budget $B(N, K)$.

The subsequent chapters outline the steps to extend this foundation toward interacting models. As noted in the Preface, Chapters~12--18 are not included in this edition. While their implemented Lean developments are already promoted, their full textbook derivations remain unpublished here. We map their intended coverage here to identify the missing prerequisites for the later text:
\begin{itemize}
    \item \textbf{Chapter 12 (Sugawara Construction):} Compares the fermionic kinetic energy operator with the bosonic Sugawara Hamiltonian $H_{\mathrm{sug}}^{(M)} = \sum_{m=1}^M \rho_{-m} \rho_m$ on the finite budget $B(N, K)$.
    \item \textbf{Chapter 13 (Klein Factors and Multi-Species Fermions):} Introduces multi-species branch fermions with explicit ground-ket sign factors to map charge and chirality across commuting budget squares.
    \item \textbf{Chapter 14 (The Projected Mattis--Mandelstam Formula):} Explores finite vertex operators using nilpotent polynomial exponentials and discusses the conditional cyclic-span criterion.
    \item \textbf{Chapters 15--16 (Dual Fields and Lattice Hamiltonians):} Examines Hermitian phase and density fields $\phi(x), \theta(x)$ and computes the spatial lattice gradients and dispersion errors.
    \item \textbf{Chapters 17--18 (Interactions and Bogoliubov Transformations):} Analyzes forward scattering models and the exact canonical Bogoliubov transformation on budget vectors.
    \item \textbf{Chapters 19--21 (Correlations, Duality, and Spinful Systems):} Develops correlation functions, explicitly specified algebraic model duality, and spinful finite constructions. CH19 Phase B is accepted but not Core promoted; CH20--21 remain manuscript drafts and do not yet have implemented Lean modules.
\end{itemize}

% Manuscript drafts imported from individually supplied standalone sections.
% CH12--18 private drafts remain outside the published master pending review.
\setcounter{chapter}{18}
\chapter{States and Correlations}
\label{ch:correlations}
\section{Scope, finite expectations, and abstract state functionals}
\label{sec:ch19-scope}

A correlation function is an expectation of an \emph{ordered product} of
observables.  In elementary quantum mechanics this usually means inserting
operators between a normalized vector and its adjoint, or taking their trace
against a density matrix.  For the bosonization problem developed here, that
familiar operation must be distinguished from the evaluation of words in an
\emph{abstract} canonical-commutation-relation (CCR) algebra.  The distinction
is essential: the finite fermionic Hilbert space has genuine normalized states,
but it cannot carry nonzero scalar CCR as identities of all its operators.
Conversely, the uncompressed oscillator algebra admits an exactly positive,
Gaussian vacuum functional, but the construction of this functional does not
identify it with any finite-band fermionic ground state.

The prerequisites are the finite CAR representation and occupation-state
inner product of Chapter~\ref{ch:car_fock}, the finite energy budgets of
Chapter~\ref{ch:vacuum_budget}, and the polynomial oscillator construction of
Chapter~\ref{ch:boson_fock}.  We use the dressed-mode transformation from
Chapter~18, with real parameters obeying $c^2-s^2=1$.  No infinite-volume
limit, Hilbert-space completion of the polynomial oscillator carrier, or
identification of finite fermionic and abstract bosonic expectations is assumed.
The aim of this section is to construct and distinguish the two kinds of
state evaluation.  We explicitly distinguish the finite CAR carrier (which models actual fermionic lattice states), the ambient polynomial oscillator carrier (unbounded degree, before completion), the protected budget (a subspace with valid commutation properties), the algebraic state (the abstract expectation functional $\omega_P$), and the coefficientwise formal-series evaluation used in later models. Ordered mode contractions, their spatial Fourier sums,
and quantitative comparisons between the models are treated subsequently.
For the usual physical context of bosonized correlators, we refer the reader to standard analytical treatments such as
\cite{giamarchi2004,vondelft1998}. The exact finite hypotheses used below
are those of the present construction rather than consequences of continuum
bosonization prescriptions.

\subsection{Expectations in a finite quantum system}

Let $\mathcal H_{\mathrm{fin}}$ be a nonzero, finite-dimensional complex
Hilbert space.  We use the convention that $\langle\phi,\psi\rangle$ is
conjugate-linear in its first argument and linear in its second.  Its
observable algebra is $\mathcal B(\mathcal H_{\mathrm{fin}})$, the algebra of
all linear operators on this space, with the Hilbert adjoint $X^\dagger$.
A \emph{normalized vector state} is specified by
$\psi\in\mathcal H_{\mathrm{fin}}$ with $\langle\psi,\psi\rangle=1$:
\begin{equation}
  \langle X\rangle_{\psi,\mathrm{fin}}
   :=\langle\psi,X\psi\rangle,
  \qquad X\in\mathcal B(\mathcal H_{\mathrm{fin}}).
  \label{eq:ch19-vector-expectation}
\end{equation}
More generally, a \emph{density matrix} is an operator $\Gamma$ satisfying
$\Gamma=\Gamma^\dagger$, $\langle v,\Gamma v\rangle\geq0$ for all $v$,
and $\operatorname{Tr}\Gamma=1$.  Its expectation is
\begin{equation}
  \langle X\rangle_{\Gamma,\mathrm{fin}}
     :=\operatorname{Tr}(\Gamma X).
  \label{eq:ch19-density-expectation}
\end{equation}
The pure state in \eqref{eq:ch19-vector-expectation} is recovered by taking
$\Gamma=|\psi\rangle\langle\psi|$.  Both functionals are complex-linear,
normalized, and positive: for example,
\begin{align}
  \langle I\rangle_{\Gamma,\mathrm{fin}}&=1,\\
  \langle X^\dagger X\rangle_{\psi,\mathrm{fin}}
    &=\|X\psi\|^2\geq0,\\
  \langle X^\dagger X\rangle_{\Gamma,\mathrm{fin}}
    &=\operatorname{Tr}\bigl((X\Gamma^{1/2})^\dagger
                                  X\Gamma^{1/2}\bigr)\geq0.
  \label{eq:ch19-finite-positivity}
\end{align}
In particular, a state is a linear functional, \emph{not} generally a
multiplicative map: the number $\langle XY\rangle$ cannot be replaced by
$\langle X\rangle\langle Y\rangle$ without an additional hypothesis.
Moreover, the order of $X$ and $Y$ matters if they do not commute.

A self-adjoint finite Hamiltonian $H_{\mathrm{fin}}$ has a smallest
eigenvalue $E_0$ and at least one normalized eigenvector $\psi_0$ with
$H_{\mathrm{fin}}\psi_0=E_0\psi_0$.  Its ground expectation is therefore
well-defined by \eqref{eq:ch19-vector-expectation}; any positive trace-one
matrix supported on the ground eigenspace also defines a ground-state
expectation.  This elementary existence result does \emph{not} say that the
same vector is killed by operators obtained by compressing a Bogoliubov
transformation.  The latter is an independent, potentially inconsistent set
of equations.

In the finite CAR setting of the earlier chapters,
$\mathcal H_{\mathrm{fin}}$ may be the occupation Fock space of a finite
momentum band, or an energy-budget subspace with compressed operators.
Whenever we write a finite expectation we specify its vector or density
matrix and the operators acting on the corresponding space.  In particular,
these brackets will never denote the abstract oscillator functional
$\omega_P$ defined below.

\subsection{Why an exact scalar CCR cannot hold on a finite space}

\begin{proposition}[Trace obstruction to finite-dimensional scalar CCR]
\label{prop:ch19-trace-obstruction}
Let $\mathcal H$ be a nonzero finite-dimensional complex vector space,
$A,C\in\operatorname{End}_{\mathbb C}(\mathcal H)$, and
$\lambda\in\mathbb C$.  If
\begin{equation}
  [A,C]:=AC-CA=\lambda I_{\mathcal H},
  \label{eq:ch19-claimed-finite-ccr}
\end{equation}
then $\lambda=0$.  In particular, $[A,C]=mI$ is impossible for a
positive integer $m$.
\end{proposition}
\begin{proof}
Writing $d=\dim_{\mathbb C}\mathcal H>0$ and using cyclicity of the
ordinary matrix trace, we obtain
\[
 \operatorname{Tr}[A,C]
   =\operatorname{Tr}(AC)-\operatorname{Tr}(CA)=0.
\]
On the other hand, \eqref{eq:ch19-claimed-finite-ccr} gives
$\operatorname{Tr}[A,C]=\lambda\operatorname{Tr}I=\lambda d$.
Since the coefficient field has characteristic zero and $d>0$,
$\lambda=0$.
\end{proof}

It is important to state exactly what the obstruction excludes.  Finite
fermionic density operators can have precisely computed commutators on
selected states, sectors, or energy budgets, and these calculations may
reproduce the expected central term on those states.  What cannot hold is
a nonzero \emph{scalar operator identity on the entire finite space}.
Compression need not preserve a CCR: if $P$ is a finite-rank orthogonal
projection in an ambient representation, $A_P=PAP|_{P\mathcal H}$, and
$C_P=PCP|_{P\mathcal H}$, then on a common domain on which all products
make sense,
\begin{equation}
 [A_P,C_P]
 =\left(P[A,C]P-PA(I-P)CP+PC(I-P)AP\right)|_{P\mathcal H}.
 \label{eq:ch19-compression-remainder}
\end{equation}
The last two terms are boundary contributions: they retain the excursions
outside the projected space that are lost when operators are multiplied
after compression.  Even if $[A,C]=mI$ in the ambient algebra, the
compressed commutator cannot equal $mI_{P\mathcal H}$ for $m>0$ and
$0<\operatorname{rank}P<\infty$.  For the finite CAR densities considered
earlier, the analogous residuals must be estimated under their actual
band, margin, and excursion hypotheses; they are not removed by simply
calling the finite currents ``bosons.''

\subsection{An uncompressed CCR algebra and its polynomial vacuum}

\paragraph{Why an uncompressed description is necessary.}
In the harmonic oscillator of undergraduate quantum mechanics, the raising
operator can be applied repeatedly:
$|0\rangle,|1\rangle,|2\rangle,\ldots$.  Even one oscillator has no
largest occupation number.  The algebraic model below retains precisely this
feature without first introducing an infinite series or a Hilbert-space
completion.  We shall work with a \emph{finite set of oscillator modes}, but
allow polynomials of \emph{arbitrarily large finite degree} in their creation
operators.  ``Uncompressed'' therefore refers to the absence of a degree or
energy cutoff, \emph{not} to infinitely many momentum labels, infinite
lattice size, or an implicit thermodynamic limit.

This distinction matters because a finite-dimensional truncation can retain
many low-occupation matrix elements of an oscillator and still fail its
operator commutation relations.  The trace obstruction just proved applies
to the truncated matrices, whereas the uncompressed differential operators
below act on an infinite-dimensional vector space.  Our goal is to establish
an exact algebraic setting for the CCR before assigning expectations to
operator products.

\paragraph{Three objects that should not be confused.}
We shall use three elementary notions from algebra.  An \emph{algebra} over
$\mathbb C$ is a complex vector space with an associative, bilinear product;
we require a multiplicative identity $I$.  For example, matrices form an
algebra under matrix multiplication, while polynomials form an algebra under
ordinary multiplication.  An element of an algebra need not yet be a
concrete linear operator.  A \emph{representation} supplies that extra step:
it is a map taking abstract algebra elements to operators while respecting
sums, scalar multiplication, products, and the identity.  In particular,
representation of a product is composition, so the rightmost operator acts
first.

The \emph{carrier} of a mathematical structure means the underlying set or
vector space on which its operations are defined.  Here there are two
relevant carriers: the vector space of formal operator expressions forming
the abstract algebra, and the polynomial vector space on which these
expressions will act.  We will call the latter the \emph{representation
carrier} $\mathcal P_M$.  Neither word denotes a spatial lattice, a
single-particle momentum band, or a finite-particle Fock space.  The term
\emph{ambient} likewise means ``the larger space or algebra before a
specified restriction.''  In this construction the ambient word algebra
precedes imposition of the CCR, while the uncompressed polynomial space is
ambient to its finite-degree subspaces.  The two uses will always be
identified by their symbols.

\paragraph{Free words and the relations we wish to impose.}
Fix an integer $M\geq1$ and the finite set of mode labels
\[
  \mathcal I_M=\{(\nu,m):\nu\in\{R,L\},\;1\leq m\leq M\}.
\]
For each $j=(\nu,m)$ introduce two \emph{formal symbols}: $A_j$, intended
for annihilation, and $C_j$, intended for creation.  Before any relations
have been imposed, an expression such as $A_j C_k A_\ell$ is simply an
ordered \emph{word}; it is not yet a matrix or differential operator.
Consider all finite complex linear combinations of such words, including
the empty word $I$.  Their multiplication is concatenation followed by
linear extension.  This is the \emph{free unital algebra}, which we denote
by $\mathfrak F_M$.  In $\mathfrak F_M$ the words $A_jC_k$ and $C_kA_j$
are distinct, just as two products of noncommuting quantum observables
must be kept in their stated order.

The word algebra also has an \emph{involution}, denoted by $*$: it exchanges
$A_j$ with $C_j$, conjugates complex coefficients, and reverses every
word.  More explicitly,
\begin{equation}
 (A_j)^*=C_j,\qquad (C_j)^*=A_j,\qquad
 (\lambda u+\mu v)^*=\overline\lambda u^*+
       \overline\mu v^*,\qquad (uv)^*=v^*u^*.
 \label{eq:ch19-free-star}
\end{equation}
This is the algebraic counterpart of the Hilbert-space adjoint.  It becomes
an actual adjoint relation only after we construct a suitable inner product
and representation.

The canonical commutation relations should hold \emph{as algebraic
equalities}, independently of any particular vector.  For clarity define
the desired relation expressions
\begin{align}
 r^{AC}_{jk}&:=A_jC_k-C_kA_j-\delta_{jk}m_j I,\nonumber\\
 r^{AA}_{jk}&:=A_jA_k-A_kA_j,\qquad
 r^{CC}_{jk}:=C_jC_k-C_kC_j,
 \label{eq:ch19-relation-expressions}
\end{align}
where $j=(\nu,m_j)$ and $\delta_{jk}=1$ exactly when both the branch and
the positive mode index agree.  The expressions $r^{AC}_{jk}$,
$r^{AA}_{jk}$, and $r^{CC}_{jk}$ are generally \emph{not zero in the
free algebra}.  They are the relations that we now impose.

\paragraph{What an ideal and a quotient do.}
An \emph{ideal} $\mathfrak J$ of an algebra is a linear subspace that is
closed under multiplication by arbitrary algebra elements on the left
\emph{and} on the right: if $r\in\mathfrak J$ and $u,v\in\mathfrak F_M$,
then $urv\in\mathfrak J$.  This two-sided requirement is crucial for
operator algebra.  If $A_jC_k-C_kA_j=\delta_{jk}m_j I$, the same equality
must remain valid when inserted inside a longer word, for example after
multiplying by $C_\ell$ on the left and $A_p$ on the right.

Let $\mathfrak J_{\mathrm{CCR}}$ be the smallest two-sided ideal
containing all the expressions in
\eqref{eq:ch19-relation-expressions} and their involutions.  Thus its
elements are finite linear combinations of expressions $urv$ built from
these relations (including their stars).  The inclusion of involutions
makes the ideal \emph{$*$-stable}: $r\in\mathfrak J_{\mathrm{CCR}}$
implies $r^*\in\mathfrak J_{\mathrm{CCR}}$.  This ensures that taking
adjoints is compatible with the equations being imposed.

The \emph{quotient algebra}
\begin{equation}
 \mathfrak A_{\mathrm{CCR}}(M)
   :=\mathfrak F_M/\mathfrak J_{\mathrm{CCR}}
 \label{eq:ch19-ccr-quotient}
\end{equation}
identifies two word expressions $u$ and $v$ when
$u-v\in\mathfrak J_{\mathrm{CCR}}$.  We are not setting every word to
zero; only expressions implied by the selected relations become zero.
For instance, in the quotient $A_jC_j=C_jA_j+m_j I$, but neither
$A_jC_j$ nor $C_jA_j$ is required to vanish.  This is a systematic way
to turn formal commutation rules into exact identities, while preserving
the order of all remaining products.

The resulting unital $*$-algebra is characterized by
\begin{align}
 [A_{\nu m},C_{\eta n}]
   &=\delta_{\nu\eta}\delta_{mn}\,mI,\nonumber\\
 [A_{\nu m},A_{\eta n}]&=0,
 \qquad [C_{\nu m},C_{\eta n}]=0.
 \label{eq:ch19-abstract-ccr}
\end{align}
These equalities hold in $\mathfrak A_{\mathrm{CCR}}(M)$ for \emph{all}
finite words, not merely after taking a vacuum expectation.
One must still check that the quotient has not collapsed to the zero
algebra (equivalently, that $I$ has not been identified with zero), and
that it admits a \emph{positive} evaluation.  Both will follow from a
concrete polynomial representation.

\paragraph{From abstract words to an explicit representation carrier.}
The variables $X_j$ introduced next are \emph{commuting polynomial
variables}, unlike the noncommuting $A_j,C_j$ in $\mathfrak F_M$.
The purpose of the representation is not to equate these two algebras:
it sends the abstract generators to \emph{operators on polynomials}.
Creation becomes multiplication by a variable, whereas annihilation
becomes differentiation.  Even though the polynomial variables commute,
multiplication and differentiation generally do not, which is precisely
how the CCR arise.  A representation satisfying the relations annihilates
every element of $\mathfrak J_{\mathrm{CCR}}$ and therefore gives a
well-defined action of the quotient on its carrier.

Let $\mathcal P_M=\mathbb C[X_j:j\in\mathcal I_M]$, the space of
polynomials in commuting indeterminates.  For a multi-index
$\alpha=(\alpha_j)_{j\in\mathcal I_M}$ with $\alpha_j\in\mathbb N$,
write $X^\alpha=\prod_jX_j^{\alpha_j}$.  Give this vector space the
Hermitian form
\begin{equation}
 \langle X^\alpha,X^\beta\rangle_{\mathcal P}
 =\delta_{\alpha\beta}
     \prod_{(\nu,m)\in\mathcal I_M}\alpha_{\nu m}!\,
                                           m^{\alpha_{\nu m}}.
 \label{eq:ch19-polynomial-gram}
\end{equation}
Every weight in the product is strictly positive.  Thus if
$p=\sum_\alpha p_\alpha X^\alpha$ is nonzero, then
\begin{equation}
 \|p\|_{\mathcal P}^2
 =\sum_\alpha |p_\alpha|^2
    \prod_{(\nu,m)}\alpha_{\nu m}!\,m^{\alpha_{\nu m}}>0.
 \label{eq:ch19-polynomial-norm}
\end{equation}
Only finitely many terms occur in each polynomial; no completion or
convergence of infinite sums is involved.  Define linear operators on
$\mathcal P_M$ by
\begin{equation}
  \pi(C_{\nu m})p=X_{\nu m}p,
  \qquad
  \pi(A_{\nu m})p=m\frac{\partial p}{\partial X_{\nu m}}.
  \label{eq:ch19-polynomial-representation}
\end{equation}
They are adjoint on this polynomial pre-Hilbert space.  To verify this
without appealing to an infinite-dimensional operator theorem, fix
$j=(\nu,m)$ and monomials $X^\alpha,X^\beta$.  The only potentially
nonzero matrix element of $\pi(A_j)$ from $X^\alpha$ to $X^\beta$ occurs
when $\alpha=\beta+e_j$.  In that case,
\[
 \langle X^\beta,\pi(A_j)X^\alpha\rangle_{\mathcal P}
   =m\alpha_j\,\|X^\beta\|_{\mathcal P}^2
   =\|X^\alpha\|_{\mathcal P}^2
   =\langle\pi(C_j)X^\beta,X^\alpha\rangle_{\mathcal P}.
\]
Thus $\pi(A_j)^\dagger=\pi(C_j)$ in the algebraic sense that these
pairings agree on every pair of polynomials.  The elementary identity
$\partial_{X_j}(X_kp)-X_k\partial_{X_j}p=\delta_{jk}p$ gives
\[
 [\pi(A_{\nu m}),\pi(C_{\eta n})]
     =\delta_{\nu\eta}\delta_{mn}\,mI_{\mathcal P_M},
\]
and the two homogeneous commutators vanish.  Because the CCR hold on every polynomial, the representation of the free
word algebra sends each relation expression, including its left and
right multiples, to zero.  It therefore factors through the quotient,
with the unit represented by the nonzero identity on $\mathcal P_M$.
In particular, the quotient is not the zero algebra.  Faithfulness of
the representation is not needed to establish this conclusion.

\paragraph{One-mode calculation: where the missing top level goes.}
For a single fixed label $j=(\nu,m)$, suppress the subscripts and consider
$A=m\,d/dX$ and $C$ acting by multiplication by $X$.  On every monomial
$X^n$, $n\geq0$,
\begin{align}
 AC\,X^n&=m(n+1)X^n, & CA\,X^n&=mnX^n,\\
 [A,C]X^n&=mX^n.&&
 \label{eq:ch19-one-mode-ccr}
\end{align}
Thus the commutator is exactly $mI$ on the entire polynomial carrier,
including the vacuum $1=X^0$.  Repeated creation yields
$C^n1=X^n$ and
\begin{equation}
 \|C^n1\|_{\mathcal P}^2=n!\,m^n,
 \qquad n=0,1,2,\ldots .
 \label{eq:ch19-one-mode-norms}
\end{equation}
No polynomial has infinite degree, but for \emph{every} $n$ there is a
nonzero vector of degree $n$.  This is the elementary reason that the
representation space has infinite dimension even when only one mode is
considered.

For comparison, keep only $1,X,\ldots,X^N$ and \emph{compress} both
operators to that $(N+1)$-dimensional subspace.  The derivative still
lowers degree, but the compressed creator kills $X^N$ because $X^{N+1}$
has been discarded.  Consequently the compressed commutator has action
\begin{equation}
 [A_N,C_N]X^n=
 \begin{cases}
 mX^n,&0\leq n<N,\\
 -mN X^N,&n=N.
 \end{cases}
 \label{eq:ch19-top-degree-defect}
\end{equation}
The failure of $mI$ is concentrated at the top degree, and the trace is
$mN-mN=0$, exactly as required by
Proposition~\ref{prop:ch19-trace-obstruction}.  This example exhibits
both the usefulness and the limitation of an energy cutoff: low-degree
vectors may obey the expected commutator while the global finite matrix
identity is impossible.  The trace obstruction does not apply to
$\mathcal P_M$ itself, since it is infinite-dimensional; moreover, the
algebraic multiplication and differentiation maps are not being asserted
to be bounded operators on a Hilbert completion.

\begin{definition}[Polynomial vacuum functional]
\label{def:ch19-polynomial-vacuum}
For $a\in\mathfrak A_{\mathrm{CCR}}(M)$ set
\begin{equation}
   \omega_0(a):=\langle1,\pi(a)1\rangle_{\mathcal P}.
   \label{eq:ch19-omega-zero}
\end{equation}
The vector $1$ is the constant polynomial, of norm one.
\end{definition}

\begin{proposition}[Normalization, positivity, and vacuum orientation]
\label{prop:ch19-polynomial-state}
The linear functional $\omega_0$ satisfies
\begin{align}
 \omega_0(I)&=1,\qquad
 \omega_0(a^*a)=\|\pi(a)1\|_{\mathcal P}^2\geq0,
 \label{eq:ch19-omega0-positive}\\
 \omega_0(aA_{\nu m})&=0,\qquad
 \omega_0(C_{\nu m}a)=0
 \quad\text{for every }a\in\mathfrak A_{\mathrm{CCR}}(M).
 \label{eq:ch19-omega0-annihilation}
\end{align}
In particular, the state is not a multiplicative algebra homomorphism.
\end{proposition}
\begin{proof}
Linearity and normalization follow from the polynomial representation
and $\|1\|_{\mathcal P}=1$.  Algebraic adjointness gives
$\langle1,\pi(a^*a)1\rangle_{\mathcal P}
 =\langle\pi(a)1,\pi(a)1\rangle_{\mathcal P}$,
which proves positivity.  For the first vanishing equation,
$\pi(A_{\nu m})1=0$; the corresponding operator occurs on the
\emph{right} of the product, so it acts first.  For the second,
adjointness gives
\[
 \langle1,\pi(C_{\nu m})\pi(a)1\rangle_{\mathcal P}
 =\langle\pi(A_{\nu m})1,\pi(a)1\rangle_{\mathcal P}=0.
\]
Finally, for any $m\geq1$,
\[
 \omega_0(A_{\nu m})=\omega_0(C_{\nu m})=0,
 \quad\omega_0(A_{\nu m}C_{\nu m})=m,
 \quad\omega_0(C_{\nu m}A_{\nu m})=0.
\]
The nonzero two-point expectation is inconsistent with
multiplicativity of a map having zero one-point values.
\end{proof}

The polynomial realization also explains the term \emph{Gaussian} or
\emph{quasi-free}.  Expectations of any finite word reduce by successively
commuting annihilators past creators, until a rightmost annihilator or a
leftmost creator kills the vacuum expectation.  The recursion expresses
even moments as sums of pairings of generators, while odd moments vanish.
Equivalently, all higher truncated moments beyond degree two vanish in this
algebraic oscillator state.  This is a statement about finite polynomial
words; no analytic generating-function convergence has been assumed.

\subsection{Dressing the vacuum by a Bogoliubov automorphism}

Let $P$ stand for real transformation parameters $(c,s)$ satisfying
\begin{equation}
 c,s\in\mathbb R,\qquad c^2-s^2=1.
 \label{eq:ch19-hyperbolic}
\end{equation}
The larger Luttinger-parameter record used in the preceding chapter can
provide such a pair, but no separate value of a Luttinger parameter is
needed for the present construction.  Write $\bar R=L$, $\bar L=R$.
Define on the generators
\begin{align}
 \beta_P(C_{\nu m})&=cC_{\nu m}+sA_{\bar\nu m},
 \label{eq:ch19-bogoliubov-creator}\\
 \beta_P(A_{\nu m})&=cA_{\nu m}+sC_{\bar\nu m}.
 \label{eq:ch19-bogoliubov-annihilator}
\end{align}
The parameters are real, so these formulas preserve the involution.
Their mixed commutator is multiplied by $c^2-s^2=1$; their
same-type commutators still vanish (the cross-orientation terms cancel).
Thus $\beta_P$ preserves the defining relations
\eqref{eq:ch19-abstract-ccr} and extends to a star-algebra
automorphism of $\mathfrak A_{\mathrm{CCR}}(M)$.  Direct substitution
shows that
\begin{equation}
 \beta_P^{-1}=\beta_{(c,-s)}.
 \label{eq:ch19-inverse-beta}
\end{equation}
We call $\widetilde A_{\nu m}:=\beta_P(A_{\nu m})$ and
$\widetilde C_{\nu m}:=\beta_P(C_{\nu m})$ the \emph{dressed}
generators.  Notice that a Bogoliubov transformation of abstract
CCR generators is not, by itself, a unitary transformation on the
finite CAR Hilbert space.

\begin{definition}[Dressed abstract state]
\label{def:ch19-dressed-functional}
The dressed oscillator state is the linear functional
\begin{equation}
  \boxed{\displaystyle
  \omega_P:=\omega_0\circ\beta_P^{-1}}
  \quad\text{on }\mathfrak A_{\mathrm{CCR}}(M).
  \label{eq:ch19-dressed-state}
\end{equation}
\end{definition}

\begin{proposition}[Positive dressed vacuum]
\label{prop:ch19-dressed-positive}
Under \eqref{eq:ch19-hyperbolic}, $\omega_P$ is normalized and
positive, and for every $a\in\mathfrak A_{\mathrm{CCR}}(M)$,
\begin{equation}
 \omega_P(a\widetilde A_{\nu m})=0,
 \qquad
 \omega_P(\widetilde C_{\nu m}a)=0.
 \label{eq:ch19-dressed-annihilation}
\end{equation}
\end{proposition}
\begin{proof}
Because $\beta_P^{-1}$ is a unital star-algebra automorphism,
\[
 \omega_P(I)=\omega_0(I)=1,\qquad
 \omega_P(a^*a)
 =\omega_0\!\left((\beta_P^{-1}a)^*\beta_P^{-1}a\right)\geq0.
\]
For the first annihilation identity, multiplicativity of the
\emph{automorphism}, not of the state, gives
\begin{align*}
 \omega_P(a\widetilde A_{\nu m})
 &=\omega_0\!\left(\beta_P^{-1}(a)
                      \beta_P^{-1}(\beta_P(A_{\nu m}))\right)\\
 &=\omega_0(\beta_P^{-1}(a)A_{\nu m})=0.
\end{align*}
Similarly,
$\omega_P(\widetilde C_{\nu m}a)
 =\omega_0(C_{\nu m}\beta_P^{-1}(a))=0$.
The direction of the pullback is therefore necessary:
using $\omega_0\circ\beta_P$ instead would not, in general, yield the
vacuum annihilated by the dressed lowering generators just defined.
\end{proof}

The familiar undressed case is $c=1$, $s=0$, for which
$\beta_P=\mathrm{id}$ and $\omega_P=\omega_0$.  A nonzero value of $s$
changes the covariance of the bare modes, not the abstract CCR
itself.  The dressed functional is again quasi-free because substituting
linear combinations of $A$ and $C$ into any finite word preserves the
finite-word Gaussian reduction described above.  Nevertheless, it is not
an ordinary matrix trace on $\mathcal H_{\mathrm{fin}}$, and it must not
be denoted by $\langle\cdot\rangle_{\psi,\mathrm{fin}}$.

\paragraph{Relation to the accepted Lean construction.}
The implementation in \texttt{A08Oscillator}\allowbreak\texttt{States.lean}
represents the \emph{free} word algebra by
\texttt{OscAlgebra M}, the polynomials by \texttt{OscPoly M},
and the creators and annihilators by \texttt{oscCreation} and
\texttt{osc}\allowbreak\texttt{Annihilation}.  The functionals \texttt{omega0} and
\texttt{omegaP} are complex-linear maps on that free-word carrier,
with $\texttt{omegaP}=\texttt{omega0}\circ
\texttt{bogoliubovMap}(c,-s)$.  Theorems such as
\texttt{omega0\_}\allowbreak\texttt{gramSum}, \texttt{omegaP\_}\allowbreak\texttt{gramSum},
\texttt{omegaP\_ccr\_}\allowbreak\texttt{descends} and its homogeneous-relation variants
establish positivity of finite word combinations and invariance under
insertion of the CCR relations; \texttt{omegaP\_comp\_}\allowbreak\texttt{bogoliubovMap}
and the dressed annihilation theorems fix the orientation of the
pullback.  The quotient $\mathfrak A_{\mathrm{CCR}}(M)$ above is a
mathematical presentation of the relation algebra, \emph{not} a
separately instantiated Lean quotient type in A08.  In that
implementation, the ambient \texttt{FreeAlgebra} remains the carrier of
\texttt{omega0} and \texttt{omegaP}.  The explicit theorems showing
that their values are unchanged when any defining relation is inserted
between arbitrary words justify descending these functionals to the
corresponding relation quotient.  This distinction between a
\emph{presented quotient} and an \emph{evaluation proved to respect its
relations} is a genuine interface distinction, not merely a matter of
notation.
Likewise, the implementation's \texttt{B\_complete} status records
acceptance of that development; it is not promotion to the frozen Core.

\subsection{Worked finite example: compression can destroy the dressed vacuum}

We now return to an ordinary three-dimensional Hilbert space.  The
example makes the separation of \eqref{eq:ch19-density-expectation}
and \eqref{eq:ch19-dressed-state} concrete.  Retain the two modes
$(R,1)$ and $(L,1)$, but project the polynomial carrier to total
degree at most one:
\[
 \mathcal H_{\leq1}
  :=\operatorname{span}_{\mathbb C}\{|0\rangle,|R\rangle,|L\rangle\}
  =\operatorname{span}_{\mathbb C}\{1,X_R,X_L\}.
\]
For $m=1$ the Gram weights in \eqref{eq:ch19-polynomial-gram}
make the displayed ordered basis orthonormal.  Let $D_R,D_L$ be the
compressed derivatives and $\widehat X_R,\widehat X_L$ the compressed
multiplication operators.  Consider the compressed dressed annihilators
\begin{equation}
 Q_R=cD_R+s\widehat X_L,
 \qquad Q_L=cD_L+s\widehat X_R.
 \label{eq:ch19-finite-Q-def}
\end{equation}
Their matrices, with columns representing the action on basis vectors,
are
\begin{equation}
 Q_R=\begin{pmatrix}
 0&c&0\\ 0&0&0\\ s&0&0
 \end{pmatrix},\qquad
 Q_L=\begin{pmatrix}
 0&0&c\\ s&0&0\\ 0&0&0
 \end{pmatrix}.
 \label{eq:ch19-finite-Q-matrices}
\end{equation}
For real $c,s$, direct multiplication gives
\begin{equation}
 H_Q:=Q_R^\dagger Q_R+Q_L^\dagger Q_L
   =\operatorname{diag}(2s^2,c^2,c^2).
 \label{eq:ch19-compressed-gram}
\end{equation}
Choose the admissible hyperbolic parameters $c=5/3$, $s=4/3$.
Then $c^2-s^2=1$ but
\begin{equation}
 H_Q=\operatorname{diag}\!\left(\frac{32}{9},
                                    \frac{25}{9},\frac{25}{9}\right)
 \succeq\frac{25}{9}I.
 \label{eq:ch19-compressed-witness}
\end{equation}
For every $v\in\mathcal H_{\leq1}$,
\begin{equation}
 \|Q_Rv\|^2+\|Q_Lv\|^2
 =\langle v,H_Qv\rangle
 \geq\frac{25}{9}\|v\|^2.
 \label{eq:ch19-compressed-coercivity}
\end{equation}
Therefore $\ker Q_R\cap\ker Q_L=\{0\}$.  There is no normalized
vector $\psi$ satisfying both $Q_R\psi=Q_L\psi=0$.

The density-matrix formulation gives the same obstruction without
assuming that the state is pure.  If $\Gamma\geq0$ has trace one, then
\begin{equation}
 \operatorname{Tr}(\Gamma H_Q)
 \geq\frac{25}{9}\operatorname{Tr}\Gamma
 =\frac{25}{9}>0.
 \label{eq:ch19-no-compressed-vacuum}
\end{equation}
In fact
$\operatorname{Tr}(\Gamma Q_\nu^\dagger Q_\nu)
 =\|Q_\nu\Gamma^{1/2}\|_{\mathrm{HS}}^2$;
vanishing for both $\nu=R,L$ would force the support of $\Gamma$
into their common kernel, hence $\Gamma=0$, a contradiction.  The
accepted A08 result \texttt{compressed\_no\_vacuum\_density}
proves the stronger algebraic impossibility of a trace-one matrix
satisfying $Q_R\Gamma=Q_L\Gamma=0$, without even requiring positivity.

Yet $H_Q$ is a perfectly valid finite self-adjoint Hamiltonian.  Its
ground energy is $25/9$, with ground eigenspace
$\operatorname{span}\{|R\rangle,|L\rangle\}$.  The density matrix
\begin{equation}
 \Gamma_{\mathrm{gs}}
  =\frac12\bigl(|R\rangle\langle R|
                   +|L\rangle\langle L|\bigr)
  =\operatorname{diag}\!\left(0,\frac12,\frac12\right)
 \label{eq:ch19-finite-ground-density}
\end{equation}
is positive and normalized, and satisfies
$\langle H_Q\rangle_{\Gamma_{\mathrm{gs}},\mathrm{fin}}=25/9$.
It is a genuine ground state but is \emph{not} a joint annihilation
vacuum.  By contrast, when $s=0$ the vector $|0\rangle$ belongs to
$\ker Q_R\cap\ker Q_L$.  Thus the obstruction in this example is
specific to the nonzero-mixing case; it does not say that every finite
Hamiltonian lacks a ground vector, nor that every compressed annihilator
family has trivial common kernel.

The same conclusion arises in a less artificial setting from the
finite-Fock analysis of the accepted Chapter~19 implementation.
There, under the explicitly required positivity of the band parameter,
admissibility of the chosen charge sector, positive mode indices inside
the cutoff, and both branch-margin hypotheses, the finite undressed
ground ket need not be the interacting dressed vacuum when $s\ne0$.
The theorem \texttt{ground\_not\_dressed\_vacuum} supplies a concrete
witness.  It does not construct a different finite interacting ground
state; that would require a specified finite Hamiltonian and its own
spectral analysis.  The three-dimensional model above is a diagnostic
counterexample, not a substitute for that analysis.

\subsection{What is established, and what remains to be compared}

There are therefore two well-defined, but mathematically different,
evaluation procedures:
\begin{equation}
 \underbrace{\langle X\rangle_{\psi,\mathrm{fin}}
 \quad\text{or}\quad
 \langle X\rangle_{\Gamma,\mathrm{fin}}}_{
 \substack{\text{finite CAR or finite}\text{ compressed Hamiltonian model}}}
 \qquad\text{versus}\qquad
 \underbrace{\omega_P(a)}_{
 \substack{\text{uncompressed abstract}\text{ CCR polynomial model}}}.
 \label{eq:ch19-two-state-layers}
\end{equation}
The first evaluates operators on a specified finite Hilbert space.
The second evaluates elements of an uncompressed relation algebra and
is positive by its polynomial construction and inverse Bogoliubov
pullback.  Both are normalized positive \emph{linear} functionals;
their domains, operator relations, and vacuum conditions differ.
There is no canonical equality between them.

An actual comparison would require, at minimum, a specified map taking
chosen abstract observables into finite-CAR observables (or a common
comparison algebra), a chosen finite state, the precise mode and energy
budgets, and an estimate of the errors caused by projection and finite
band edges.  Pointwise identities on a controlled excitation sector do
not automatically extend to all operator products; any estimate must
keep track of the excursions of intermediate states.  Subsequent
sections first compute ordered contractions using $\omega_P$, then
construct their finite spatial kernels, and finally address which
finite-Fock comparisons the available hypotheses justify.  No continuum
power-law or physical finite-ground-state equality is claimed here.

\section{Mode Contractions in the Abstract State}
\label{sec:ch19-mode-correlators}

The preceding section constructed a normalized positive functional on an
\emph{uncompressed} oscillator algebra. We now ask a more concrete question:
what does that state predict for the product of two density modes? A
Bogoliubov transformation mixes an oscillator's creation operator on one
branch with an annihilation operator on the opposite branch. Consequently,
even if the transformed (dressed) modes define a vacuum, the original
(\emph{bare}) modes need not annihilate it. The outcome is a nonzero
correlation between pairs of bare operators. This is the elementary
algebraic mechanism behind the mode contractions derived below.

There is a crucial choice before computing anything: does the observable
contain $AB$ or $BA$? These are different measurements whenever $A$ and
$B$ do not commute, and a state is a \emph{linear}, not a multiplicative,
evaluator. In particular, exchanging the order of a bosonic creation and
annihilation mode changes its vacuum expectation by their scalar
commutator. We therefore compute \emph{ordered} contractions first; the
symmetrized combinations will be identified only afterward. This section
works solely in the abstract oscillator functional of
Section~\ref{sec:ch19-scope}. It does not evaluate a ground-state vector
or density matrix in the finite fermionic energy-budget space.

\subsection{Mode labels, hypotheses, and the vacuum rule}

Fix a nonnegative integer $M$. The set of positive mode labels is
\[
  \mathcal{I}_M=\{1,2,\ldots,M\},
  \qquad \nu\in\{R,L\}.
\]
When $M=0$ the statements about individual modes are empty; whenever we
choose a particular mode below we assume $M\geq1$. The formalized index
$i\in\operatorname{Fin}(M)$ corresponds to $m=i+1$, so every oscillator
weight $m$ is strictly positive. The lattice length $L$ is reserved for
the spatial Fourier sums of the next section (in the supplied finite-lattice
implementation $L=2h$); it is \emph{not} the dimension of the oscillator
polynomial carrier. The mode algebra considered here has $2M$ oscillator
species but admits arbitrarily high occupation in each species.

Write $C_{\nu m}$ for the bare creation generator and
$A_{\nu m}:=C_{\nu m}^{*}$ for its adjoint, the annihilation generator.
We use the density-mode correspondence of the source convention,
\[
  \rho_{m,\nu}=C_{\nu m},\qquad
  \rho_{-m,\nu}=A_{\nu m}\quad(m\in\mathcal I_M).
\]
The exact uncompressed canonical commutation relations (CCR) are
\begin{align}
 [A_{\nu m},C_{\mu n}]&=m\,\delta_{\nu\mu}\delta_{mn}\,I,
   \label{eq:ch19-modes-ccr}\\
 [A_{\nu m},A_{\mu n}]&=0,\qquad
 [C_{\nu m},C_{\mu n}]=0.
\end{align}
Here $\delta_{mn}$ is one if $m=n$ and zero otherwise. Operators on
different branches commute, but, as we shall see, they can still have
nonzero expectations of their products in a correlated state.

Let $\omega_0$ denote the polynomial vacuum functional defined in
Section~\ref{sec:ch19-scope}. In its explicit polynomial representation,
$C_{\nu m}$ multiplies by the variable $X_{\nu m}$, while
$A_{\nu m}=m\,\partial/\partial X_{\nu m}$. The vacuum is the
constant polynomial $1$. Therefore
\begin{align}
 A_{\nu m}1&=0,
 &C_{\nu m}1&=X_{\nu m},
 &A_{\nu m}X_{\mu n}
   &=m\delta_{\nu\mu}\delta_{mn}\,1.
\end{align}
Applying the operators from right to left, then extracting the vacuum
coefficient, yields
\begin{align}
 \omega_0(A_{\nu m}C_{\mu n})
   &=m\delta_{\nu\mu}\delta_{mn},
   &\omega_0(C_{\nu m}A_{\mu n})&=0,
     \label{eq:ch19-free-two-orders}\\
 \omega_0(A_{\nu m}A_{\mu n})&=0,
   &\omega_0(C_{\nu m}C_{\mu n})&=0.
     \label{eq:ch19-free-like-pairs}
\end{align}
For example, $A_{\nu m}C_{\nu m}1=m1$, whereas
$C_{\nu m}A_{\nu m}1=0$. This two-line calculation is the source of
the distinction between the $c^2$ and $s^2$ contractions. The identities
are exact on the uncompressed polynomial carrier: no top-occupation
boundary has been discarded.

Let $c,s\in\mathbb R$ satisfy
\begin{equation}
  c^2-s^2=1.
  \label{eq:ch19-hyperbolic-condition}
\end{equation}
This is the \emph{real hyperbolic-pair} hypothesis of Chapter~18; it
includes the choice $c=1,s=0$. No sign condition on $s$ is necessary.
The Bogoliubov generator map $\beta=\beta_{c,s}$ acts by
\begin{align}
 \beta(C_{Rm})&=cC_{Rm}+sA_{Lm},
 &\beta(A_{Rm})&=cA_{Rm}+sC_{Lm},\label{eq:ch19-beta-right}\\
 \beta(C_{Lm})&=cC_{Lm}+sA_{Rm},
 &\beta(A_{Lm})&=cA_{Lm}+sC_{Rm}.
 \label{eq:ch19-beta-left}
\end{align}
Because the coefficients are real, the second column of each line is
the adjoint of the first. Because $c^2-s^2=1$, applying the same map
with $s$ replaced by $-s$ gives its inverse. For example,
\[
 \beta^{-1}(C_{Rm})=cC_{Rm}-sA_{Lm},\qquad
 \beta^{-1}(A_{Rm})=cA_{Rm}-sC_{Lm}.
\]
The other branch follows by interchanging $R$ and $L$. These identities
are a special case of the invertible generator transformation established
in Chapter~18 and in the supplied A08 development.

The \emph{dressed} modes are $\widetilde C_{\nu m}=\beta(C_{\nu m})$
and $\widetilde A_{\nu m}=\beta(A_{\nu m})$. As established in the
preceding section, the dressed state is the \emph{inverse pullback}
\begin{equation}
  \omega_P(B)=\omega_0\bigl(\beta^{-1}(B)\bigr),
  \qquad B\in\mathfrak A_{\mathrm{CCR}}.
  \label{eq:ch19-pullback}
\end{equation}
This direction is essential. In fact,
$\omega_P(B\widetilde A_{\nu m})=0$ and
$\omega_P(\widetilde C_{\nu m}B)=0$ for every polynomial observable
$B$, since applying $\beta^{-1}$ reduces these expressions to the
undressed vacuum rules. The shorthand $\omega_{\widetilde\Omega}$ in
the original chapter notes refers to this same abstract functional;
we use $\omega_P$ consistently here. In particular, the equations below
are not assertions about finite-CAR expectations
$\langle\Psi,B\Psi\rangle$ or $\operatorname{Tr}(\varrho B)$.

\subsection{Same-branch contractions: why both orders survive}

\begin{theorem}[Both ordered right- and left-branch contractions]
\label{thm:ch19-mode-same-branch}
For any $M\in\mathbb N$, real $c,s$ satisfying
\eqref{eq:ch19-hyperbolic-condition}, $m,n\in\mathcal I_M$, and
$\nu\in\{R,L\}$, the functional~\eqref{eq:ch19-pullback} obeys
\begin{align}
 \omega_P(C_{\nu m}A_{\nu n})
   &=m\delta_{mn}s^2,
   \label{eq:ch19-mode-CA}\\
 \omega_P(A_{\nu m}C_{\nu n})
   &=m\delta_{mn}c^2,
   \label{eq:ch19-mode-AC}\\
 \omega_P(C_{\nu m}C_{\nu n})
   &=0,\qquad
 \omega_P(A_{\nu m}A_{\nu n})=0.
 \label{eq:ch19-mode-like}
\end{align}
All products are written in their actual evaluation order.
\end{theorem}

\begin{proof}
It suffices to treat the right branch; the opposite-branch exchange
$R\leftrightarrow L$ preserves the CCR and the transformation.
Multiplicativity of the \emph{algebra map} $\beta^{-1}$ (not of the
state $\omega_0$) gives
\begin{align*}
 \omega_P(C_{Rm}A_{Rn})
 &=\omega_0\bigl[(cC_{Rm}-sA_{Lm})
                     (cA_{Rn}-sC_{Ln})\bigr]\\
 &=c^2\omega_0(C_{Rm}A_{Rn})
   -cs\omega_0(C_{Rm}C_{Ln})\\
 &\quad-cs\omega_0(A_{Lm}A_{Rn})
   +s^2\omega_0(A_{Lm}C_{Ln})\\
 &=s^2m\delta_{mn}.
\end{align*}
Every vanishing term has been removed using
\eqref{eq:ch19-free-two-orders}--\eqref{eq:ch19-free-like-pairs}.
For the reversed product the expansion is different:
\begin{align*}
 \omega_P(A_{Rm}C_{Rn})
 &=\omega_0\bigl[(cA_{Rm}-sC_{Lm})
                     (cC_{Rn}-sA_{Ln})\bigr]\\
 &=c^2\omega_0(A_{Rm}C_{Rn})
   -cs\omega_0(A_{Rm}A_{Ln})\\
 &\quad-cs\omega_0(C_{Lm}C_{Rn})
   +s^2\omega_0(C_{Lm}A_{Ln})\\
 &=c^2m\delta_{mn}.
\end{align*}
The $CC$ product is likewise mapped to
\begin{align*}
 \omega_0[(cC_{Rm}-sA_{Lm})(cC_{Rn}-sA_{Ln})]
 &=c^2\omega_0(C_{Rm}C_{Rn})\\
 &\quad-cs\omega_0(C_{Rm}A_{Ln})
   -cs\omega_0(A_{Lm}C_{Rn})\\
 &\quad+s^2\omega_0(A_{Lm}A_{Ln})=0,
\end{align*}
since in the potentially surviving $AC$ term the branches differ.
The $AA$ identity follows by adjunction or by the same four-term
expansion. Nothing in the proof requires a finite-dimensional
representation of the CCR.
\end{proof}

Subtracting the first contraction from the second gives the exact
commutator check
\begin{equation}
 \omega_P(A_{\nu m}C_{\nu n})
 -\omega_P(C_{\nu n}A_{\nu m})
 =m\delta_{mn}(c^2-s^2)=m\delta_{mn}.
 \label{eq:ch19-mode-commutator-check}
\end{equation}
The right-hand side is fixed by the algebra, independently of the
chosen state. Thus the coefficient $s^2$ cannot simply replace
$c^2$ in the reversed contraction. More physically, even in a free
vacuum an annihilator followed by a creator can produce a nonzero
fluctuation: the creator first excites the vacuum, then the annihilator
returns it to the vacuum.

\subsection{Cross-branch contractions and the sign of the mixing}

The transformation also couples opposite branches. Correlation does not
require a nonzero cross-branch commutator: here the $R$ and $L$
oscillators commute, but $\omega_P$ is not a product state of the
\emph{bare} right- and left-branch vacuum functionals. The distinction
between algebraic commutativity and statistical independence is
important in interpreting a Bogoliubov-transformed ground state.

\begin{theorem}[Opposite-branch anomalous contractions]
\label{thm:ch19-mode-cross-branch}
Under the hypotheses of Theorem~\ref{thm:ch19-mode-same-branch},
\begin{align}
 \omega_P(C_{Rm}C_{Ln})
   &=-m\delta_{mn}cs,
   \label{eq:ch19-cross-CC}\\
 \omega_P(A_{Rm}A_{Ln})
   &=-m\delta_{mn}cs,
   \label{eq:ch19-cross-AA}\\
 \omega_P(C_{Rm}A_{Ln})
   &=0,\qquad
 \omega_P(A_{Rm}C_{Ln})=0.
   \label{eq:ch19-cross-mixed}
\end{align}
The same results hold after exchanging $R$ and $L$.
\end{theorem}

\begin{proof}
Consider first the product of two creators. Applying the inverse
transformation and expanding without commuting factors gives
\begin{align*}
 \omega_P(C_{Rm}C_{Ln})
 &=\omega_0\bigl[(cC_{Rm}-sA_{Lm})
                     (cC_{Ln}-sA_{Rn})\bigr]\\
 &=c^2\omega_0(C_{Rm}C_{Ln})
   -cs\omega_0(C_{Rm}A_{Rn})\\
 &\quad-cs\omega_0(A_{Lm}C_{Ln})
   +s^2\omega_0(A_{Lm}A_{Rn})\\
 &=-cs\,m\delta_{mn}.
\end{align*}
The single surviving term has precisely the minus sign introduced
by $\beta^{-1}$; using the forward map instead would reverse this
sign and would describe a different state.
For two annihilators the analogous computation is
\begin{align*}
 \omega_P(A_{Rm}A_{Ln})
 &=\omega_0\bigl[(cA_{Rm}-sC_{Lm})
                     (cA_{Ln}-sC_{Rn})\bigr]\\
 &=c^2\omega_0(A_{Rm}A_{Ln})
   -cs\omega_0(A_{Rm}C_{Rn})\\
 &\quad-cs\omega_0(C_{Lm}A_{Ln})
   +s^2\omega_0(C_{Lm}C_{Rn})\\
 &=-cs\,m\delta_{mn}.
\end{align*}
Finally, the mixed-branch product becomes
\begin{align*}
 \omega_P(C_{Rm}A_{Ln})
 &=c^2\omega_0(C_{Rm}A_{Ln})
   -cs\omega_0(C_{Rm}C_{Rn})\\
 &\quad-cs\omega_0(A_{Lm}A_{Ln})
   +s^2\omega_0(A_{Lm}C_{Rn})=0.
\end{align*}
The other mixed-branch contraction vanishes by the same vacuum
identities. Exchanges of branches are permitted because they commute
in the CCR algebra. This completes the proof.
\end{proof}

These cross-branch quantities are often called \emph{anomalous}
contractions: they pair two creators or two annihilators, which would
have zero expectation in the undressed polynomial vacuum. Their
nonzero value in $\omega_P$ encodes the pair mixing; the sign records
our choice of the inverse-pullback convention and the sign of $s$.

\subsection{Ordered versus symmetrized expectations}

For any two algebra elements $B,D$, introduce their symmetrized
product $\{B,D\}:=BD+DB$. The ordered and symmetrized functionals
measure different algebra elements:
\[
 \omega_P(BD)
 =\frac12\omega_P(\{B,D\})
  +\frac12\omega_P([B,D]).
\]
For one right- or left-branch creator/annihilator pair, the previous
theorem implies
\begin{equation}
 \frac12\omega_P(\{A_{\nu m},C_{\nu n}\})
   =\frac{m\delta_{mn}}2(c^2+s^2).
 \label{eq:ch19-mode-anticommutator}
\end{equation}
The antisymmetric half is
$\frac12\omega_P([A_{\nu m},C_{\nu n}])=
\frac12m\delta_{mn}$. Thus symmetrization keeps the sum of the
fluctuations but suppresses the information about which ordered
product carried the vacuum term. The anomalous $CC$ and $AA$
contractions do not have this ordering ambiguity across branches:
the corresponding generators commute.

It is useful to express these coefficients in the parameter convention
introduced in Chapter~18. Put
\begin{equation}
 g=(c-s)^2>0,\qquad g^{-1}=(c+s)^2.
 \label{eq:ch19-mode-g}
\end{equation}
The inverse equality follows from $(c-s)(c+s)=1$. Hence
\begin{align}
 c^2+s^2&=\frac{g+g^{-1}}2,
 & c^2&=\frac{g+g^{-1}+2}{4},
 & s^2&=\frac{g+g^{-1}-2}{4},\\
 cs&=\frac{g^{-1}-g}{4}.
 \label{eq:ch19-mode-g-coefficients}
\end{align}
In particular,
\begin{align*}
 \frac12\omega_P(\{A_{\nu m},C_{\nu n}\})
   &=\frac{m\delta_{mn}}4(g+g^{-1}),\\
 \omega_P(C_{Rm}C_{Ln})
   &=-\frac{m\delta_{mn}}4(g^{-1}-g).
\end{align*}
The coupling $g$ here is a \emph{hyperbolic/Luttinger parameter}; it
is not the Umklapp coupling $g_U$ used in a later chapter.
These finite algebraic identities do not invoke a continuum limit.

\subsection{Worked example: two modes and a squeezed combination}

Choose $M=2$ and the real hyperbolic pair
\[
 c=\frac53,\qquad s=\frac43,\qquad
 c^2-s^2=\frac{25-16}{9}=1,
 \qquad g=(c-s)^2=\frac19.
\]
At $m=n=1$ the right-branch ordered contractions are
\[
 \omega_P(C_{R1}A_{R1})=\frac{16}{9},\qquad
 \omega_P(A_{R1}C_{R1})=\frac{25}{9}.
\]
Their difference is exactly $1$, as required by the uncompressed
commutator. The cross-branch anomalous contractions are
\[
 \omega_P(C_{R1}C_{L1})
 =\omega_P(A_{R1}A_{L1})=-\frac{20}{9}.
\]
For comparison, mode orthogonality yields
$\omega_P(C_{R1}A_{R2})=\omega_P(C_{R1}C_{L2})=0$.
At the second mode the nonzero contractions simply acquire the
weight $m=2$:
\[
 \omega_P(C_{R2}A_{R2})=\frac{32}{9},\quad
 \omega_P(A_{R2}C_{R2})=\frac{50}{9},\quad
 \omega_P(C_{R2}C_{L2})=-\frac{40}{9}.
\]

There is also an instructive positive quadratic observable. Fix one
mode $m$ and define $Y_m=C_{Rm}+A_{Lm}$. Its adjoint is
$Y_m^*=A_{Rm}+C_{Lm}$. Expanding $Y_m^*Y_m$ in precisely this order,
\begin{align*}
 \omega_P(Y_m^*Y_m)
 &=\omega_P(A_{Rm}C_{Rm})
   +\omega_P(A_{Rm}A_{Lm})\\
 &\quad+\omega_P(C_{Lm}C_{Rm})
   +\omega_P(C_{Lm}A_{Lm})\\
 &=m(c^2-2cs+s^2)=m(c-s)^2=mg.
\end{align*}
In our numerical example, $\omega_P(Y_1^*Y_1)=1/9>0$.
Although individual contributions include negative cross-branch
correlations, the positive functional still gives a nonnegative
expectation to an element of the form $Y^*Y$. This calculation makes
the underlying tradeoff visible: the state can encode correlations
that are simple in the full oscillator representation, while a
finite energy-budget compression cannot be assumed to preserve these
CCR-based manipulations without additional boundary analysis.

\subsection{Free-pair limit and exact scope of the result}

The supplied implementation chooses the free hyperbolic pair
$(c,s)=(1,0)$, so $\beta$ and $\beta^{-1}$ are the identity and
$\omega_P=\omega_0$. The complete nonzero two-point data then reduce
to
\begin{equation}
 \omega_0(A_{\nu m}C_{\mu n})
 =m\delta_{\nu\mu}\delta_{mn},
 \label{eq:ch19-mode-free-limit}
\end{equation}
while $\omega_0(C_{\nu m}A_{\mu n})=0$ and all $CC$, $AA$
and opposite-branch contractions vanish. Thus the normal-ordered
free-vacuum occupation is zero, but the reverse-ordered vacuum
fluctuation is \emph{not} zero. Omitting it would make the subsequent
spatial density correlator incorrect already in the free theory.

All results in this section are \emph{exact algebraic evaluations} for
finitely many positive mode labels in the positive abstract state
$\omega_P$. They follow from the polynomial vacuum rules, the real
hyperbolic condition, and the inverse Bogoliubov pullback. In the
supplied Lean sources these evaluations correspond to the same-branch
and cross-branch contraction theorems in
	exttt{BosonizeStubs/Ch19Correlations.lean}, together with its separate
vanishing lemmas. Their acceptance in the Chapter~19/A08
Phase~B development does not constitute promotion to the frozen Core.

The next section will combine the ordered contractions with a
\emph{specified} finite Fourier orientation and normalization to
construct spatial density kernels. That step must retain both the
$c^2$ and $s^2$ terms; replacing the ordered product by a cosine-only
symmetrization changes the observable. A further, independent task is
to compare these abstract values with expectations in a constructed
finite fermionic ground state, including any projection or boundary
residuals. No such equality is assumed here.

\section{Finite Spatial Correlation Kernels}
\label{sec:ch19-spatial-correlators}

A density mode is a useful algebraic coordinate, but a physical density is
usually discussed at a \emph{site}.  To pass from the ordered mode
contractions of Section~\ref{sec:ch19-mode-correlators} to an observable
at two positions, we must reconstruct the spatial field by a finite Fourier
sum.  This operation appears innocuous: multiply each mode by a phase and
sum.  Yet the order of the two resulting fields matters, because each field
contains both creation and annihilation generators.  Reversing their product
changes the answer by an exact canonical commutator.

The question of this section is therefore deliberately finite and precise:
\emph{given the positive abstract CCR functional constructed in
Section~\ref{sec:ch19-scope}, what complex number does it assign to the
ordered product of two spatial density operators?}  We will first derive
that number, then separate its symmetric (fluctuation) and antisymmetric
(commutator) contributions.  Only afterward will we interpret positivity
and evaluate a concrete lattice example.  Neither this reconstruction nor
the finite Fourier calculation identifies the abstract functional with a
ground-state expectation on a compressed fermionic energy-budget space.

\subsection{A finite lattice and its Fourier phases}

In the source implementation, the lattice length is \emph{even}:
\begin{equation}
 L=2h,\qquad h\in\mathbb N,\quad h>0,\qquad
 \Lambda_L=\mathbb Z/L\mathbb Z.
 \label{eq:ch19-spatial-lattice}
\end{equation}
A point $x\in\Lambda_L$ is a residue class, represented when necessary
by its unique integer $x_0\in\{0,\ldots,L-1\}$.  Choose the canonical
primitive $L$th root
\begin{equation}
 \zeta=\exp(2\pi i/L),\qquad \zeta^L=1,
 \quad |\zeta|=1.
 \label{eq:ch19-spatial-zeta}
\end{equation}
More generally, every result written in terms of residue-class
characters holds when $\zeta$ is \emph{any} primitive $L$th root.
The word \emph{primitive} means that $\zeta^k\ne1$ for
$1\leq k<L$: the powers of $\zeta$ distinguish the $L$ sites.
The source also proves some intermediate algebraic identities for a
merely nonzero complex number $\zeta$.  That weaker hypothesis suffices
for powers involving the chosen representatives $x_0,y_0$, but the
periodic character interpretation and the Hermiticity arguments below
use the root-of-unity condition.

For $m\in\mathbb Z$ define the character
\begin{equation}
 \chi_m(x):=\zeta^{m x_0}.
 \label{eq:ch19-spatial-character}
\end{equation}
This is well defined on residue classes because replacing $x_0$ by
$x_0+kL$ leaves its value unchanged.  Moreover,
\begin{equation}
 \chi_m(x+y)=\chi_m(x)\chi_m(y),\qquad
 \chi_{-m}(x)=\overline{\chi_m(x)}=\chi_m(x)^{-1},\qquad
 \chi_m(x-y)=\chi_m(x)\chi_{-m}(y).
 \label{eq:ch19-spatial-character-laws}
\end{equation}
The last identity is the step that turns a double Fourier sum into a
function of the \emph{separation} $x-y$, rather than of the two
positions independently.  Notice that the separation belongs to
$\Lambda_L$; no choice of an infinite-lattice coordinate is needed.

Fix a positive-mode cutoff $M\in\mathbb N$ and retain the indices
$m\in\{1,\ldots,M\}$.  The source algebraic theorem allows $M=0$
(the sum is then empty), and does not assume $M<L/2$.  If one wants
these labels to represent distinct nonzero lattice frequencies away
from the Nyquist mode, the extra condition $M<h$ is natural, but it is
\emph{not a hypothesis} of the abstract kernel identity.  For $M\geq L$
some characters repeat even though the abstract oscillator generators
remain distinct.  The number of modes $M$ is also not an occupation
cutoff: each mode still acts on the uncompressed polynomial carrier.

\subsection{The spatial density as a Hermitian algebra element}

Let $C_{Rm}$ and $A_{Rm}=C_{Rm}^{*}$ denote the right-branch bare
oscillator generators, with exact uncompressed CCR
\begin{equation}
 [A_{Rm},C_{Rn}]=m\delta_{mn}I,\qquad
 [A_{Rm},A_{Rn}]=[C_{Rm},C_{Rn}]=0.
 \label{eq:ch19-spatial-ccr}
\end{equation}
We identify the positive- and negative-index density modes as
$\rho_{m,R}=C_{Rm}$ and $\rho_{-m,R}=A_{Rm}$, respectively.
Define the right-branch \emph{oscillatory density} by
\begin{equation}
 \rho_R(x):=\frac1L\sum_{m=1}^{M}
 \bigl(\chi_{-m}(x)C_{Rm}+\chi_m(x)A_{Rm}\bigr).
 \label{eq:ch19-spatial-rho}
\end{equation}
It contains no charge zero mode.  This convention uses $1/L$ rather
than $1/\sqrt L$; the normalizations of the Fourier transform and of
the oscillator weights should not be silently combined.  In
particular, the cutoff $M$ labels a selected collection of oscillatory
modes and does not turn \eqref{eq:ch19-spatial-rho} into a claim about
the full microscopic lattice density.

Why is this definition an observable rather than an arbitrary complex
combination?  The conjugate character multiplies the adjoint operator:
\begin{align*}
 \rho_R(x)^*
 &=\frac1L\sum_{m=1}^{M}
   \bigl(\overline{\chi_{-m}(x)}A_{Rm}
        +\overline{\chi_m(x)}C_{Rm}\bigr)\\
 &=\frac1L\sum_{m=1}^{M}
   \bigl(\chi_m(x)A_{Rm}+\chi_{-m}(x)C_{Rm}\bigr)
 =\rho_R(x).
\end{align*}
The reality of $1/L$ and the unit modulus of $\zeta$ were used here.
This simple check will be essential when we deduce positive
semidefiniteness of the correlation matrix; a generic nonzero
complex $\zeta$ would not justify the same argument.

Let $P=(c,s)$ be a real hyperbolic pair,
\begin{equation}
 c,s\in\mathbb R,\qquad c^2-s^2=1.
 \label{eq:ch19-spatial-hyperbolic}
\end{equation}
We evaluate in the normalized positive \emph{abstract} functional
$\omega_P=\omega_0\circ\beta_{c,s}^{-1}$ of the preceding sections.
Its mode contractions, in the exact source ordering, are
\begin{align}
 \omega_P(C_{Rm}A_{Rn})&=m\delta_{mn}s^2,
 &\omega_P(A_{Rm}C_{Rn})&=m\delta_{mn}c^2,
 \label{eq:ch19-spatial-input-contractions}\\
 \omega_P(C_{Rm}C_{Rn})&=0,
 &\omega_P(A_{Rm}A_{Rn})&=0.
 \label{eq:ch19-spatial-input-zeros}
\end{align}
These are recalled as \emph{inputs}; their Bogoliubov derivation is not
repeated here.  The notation $\omega_P$ is deliberately distinct from
a finite-Fock vector expectation $\langle\Psi,\cdot\,\Psi\rangle$
or a density-matrix expectation $\operatorname{Tr}(\varrho\,\cdot)$.

\subsection{The exact ordered two-point kernel}

\begin{theorem}[Ordered spatial density kernel]
\label{thm:ch19-spatial-ordered}
Assume \eqref{eq:ch19-spatial-lattice}--\eqref{eq:ch19-spatial-hyperbolic}
and a nonnegative integer cutoff $M$.  For every $x,y\in\Lambda_L$,
the ordered correlation
$D_P(x,y):=\omega_P\bigl(\rho_R(x)\rho_R(y)\bigr)$ equals
\begin{equation}
 \boxed{\displaystyle
 D_P(x,y)=\frac1{L^2}\sum_{m=1}^{M}m
 \bigl(s^2\chi_{-m}(x-y)+c^2\chi_m(x-y)\bigr).}
 \label{eq:ch19-spatial-ordered}
\end{equation}
For the canonical root $\zeta=e^{2\pi i/L}$, the characters
in this formula are the exponentials
$\chi_{\pm m}(x-y)=e^{\pm2\pi i m(x_0-y_0)/L}$.
The value is independent of the chosen representatives.
\end{theorem}

\begin{proof}
Write $u_m(x)=\chi_{-m}(x)$ and $v_m(x)=\chi_m(x)$.
Expanding the product of the two finite sums \emph{before}
applying the functional gives
\begin{align*}
 L^2\rho_R(x)\rho_R(y)
 &=\sum_{m=1}^{M}\sum_{n=1}^{M}
  \bigl[u_m(x)u_n(y)C_{Rm}C_{Rn}\\
 &\qquad+u_m(x)v_n(y)C_{Rm}A_{Rn}\\
 &\qquad+v_m(x)u_n(y)A_{Rm}C_{Rn}\\
 &\qquad+v_m(x)v_n(y)A_{Rm}A_{Rn}\bigr].
\end{align*}
Because the sum is finite and $\omega_P$ is complex linear,
we may evaluate its four kinds of terms individually.  The $CC$
and $AA$ contributions vanish by
\eqref{eq:ch19-spatial-input-zeros}; the remaining two retain their
original order:
\begin{align*}
 L^2D_P(x,y)
 &=\sum_{m,n=1}^{M}u_m(x)v_n(y)
       \underbrace{\omega_P(C_{Rm}A_{Rn})}_{m\delta_{mn}s^2}\\
 &\qquad+\sum_{m,n=1}^{M}v_m(x)u_n(y)
       \underbrace{\omega_P(A_{Rm}C_{Rn})}_{m\delta_{mn}c^2}\\
 &=\sum_{m=1}^{M}m
  \bigl(s^2\chi_{-m}(x)\chi_m(y)
       +c^2\chi_m(x)\chi_{-m}(y)\bigr).
\end{align*}
The Kronecker delta removed one summation index.  Finally, the
character multiplication law
\eqref{eq:ch19-spatial-character-laws} replaces each product of two
phases by the corresponding character of $x-y$.  Dividing by $L^2$
proves the claim.  No infinite sum, reordering limit, or continuum
approximation enters the calculation.
\end{proof}

The notation $D_P(x,y)$ emphasizes the parameter dependence.  The
original chapter notes also call this quantity $D_c(x,y)$; that
subscript should not be misread as dependence on $c$ alone, since
both $c^2$ and $s^2$ occur.  The \emph{ordered} kernel is generally
complex for $x\neq y$.  Replacing the two exponentials immediately by
a cosine would discard its antisymmetric part and calculate a
different observable.

\subsection{The state-independent commutator and the symmetrized kernel}

For two algebra elements, write
$[B,C]=BC-CB$ and $\{B,C\}=BC+CB$.  By linearity,
\begin{equation}
 D_P(x,y)=\frac12\omega_P(\{\rho_R(x),\rho_R(y)\})
          +\frac12\omega_P([\rho_R(x),\rho_R(y)]).
 \label{eq:ch19-spatial-decomposition}
\end{equation}
The first term describes symmetric two-point fluctuations, while the
second remembers the operator ordering.  They are not interchangeable.

\begin{proposition}[Exact density commutator]
\label{prop:ch19-spatial-commutator}
Under the hypotheses of Theorem~\ref{thm:ch19-spatial-ordered},
\begin{align}
 [\rho_R(x),\rho_R(y)]
 &=\frac1{L^2}\sum_{m=1}^{M}m
   \bigl(\chi_m(x-y)-\chi_{-m}(x-y)\bigr)I,
 \label{eq:ch19-spatial-commutator}\\
 D_P(x,y)-D_P(y,x)
 &=\frac{2i}{L^2}\sum_{m=1}^{M}m
   \sin\!\left(\frac{2\pi m(x_0-y_0)}L\right)
 \label{eq:ch19-spatial-commutator-sine}
\end{align}
for the canonical root in the second line.  In particular, this
commutator expectation is independent of $c,s$ and of the chosen
normalized state on the exact CCR algebra.
\end{proposition}

\begin{proof}
The equal-sign generator commutators vanish.  In the expansion of
the commutator of the two sums, only the mixed terms survive:
\begin{align*}
 L^2[\rho_R(x),\rho_R(y)]
 &=\sum_{m,n=1}^{M}\bigl[
  \chi_{-m}(x)\chi_n(y)[C_{Rm},A_{Rn}]\\
 &\hspace{49mm}+
  \chi_m(x)\chi_{-n}(y)[A_{Rm},C_{Rn}]\bigr].
\end{align*}
By \eqref{eq:ch19-spatial-ccr},
$[C_{Rm},A_{Rn}]=-m\delta_{mn}I$ and
$[A_{Rm},C_{Rn}]=m\delta_{mn}I$.  Collapsing the deltas and using
\eqref{eq:ch19-spatial-character-laws} gives the first equation.
Applying $\omega_P(I)=1$ yields the second after
$e^{i\theta}-e^{-i\theta}=2i\sin\theta$.
Alternatively, subtracting \eqref{eq:ch19-spatial-ordered} with
$x,y$ interchanged produces the same expression because
$c^2-s^2=1$.  The direct CCR calculation demonstrates the stronger
fact: the commutator is already a scalar multiple of the unit
\emph{before} a state is chosen.
\end{proof}

\begin{proposition}[Symmetrized density correlation]
\label{prop:ch19-spatial-symmetrized}
Let
$S_P(x,y):=\tfrac12\omega_P(\{\rho_R(x),\rho_R(y)\})$.
Then for the canonical root
\begin{equation}
 \boxed{\displaystyle
 S_P(x,y)=\frac{c^2+s^2}{L^2}
 \sum_{m=1}^{M}m\cos\!\left(\frac{2\pi m(x_0-y_0)}L\right).}
 \label{eq:ch19-spatial-symmetrized}
\end{equation}
Set $g=(c-s)^2>0$, so $g^{-1}=(c+s)^2$.  Equivalently,
\begin{equation}
 c^2+s^2=\frac{g+g^{-1}}2,
 \qquad
 S_P(x,y)=\frac{g+g^{-1}}{2L^2}
 \sum_{m=1}^{M}m\cos\!\left(\frac{2\pi m(x_0-y_0)}L\right).
 \label{eq:ch19-spatial-symmetrized-g}
\end{equation}
\end{proposition}

\begin{proof}
Write $d=x_0-y_0$ and $\theta_m=2\pi md/L$.
The ordered formula takes the form
\begin{align*}
 D_P(x,y)
 &=\frac1{L^2}\sum_{m=1}^{M}m
   \left[s^2(\cos\theta_m-i\sin\theta_m)
        +c^2(\cos\theta_m+i\sin\theta_m)\right]\\
 &=\frac1{L^2}\sum_{m=1}^{M}m
   \left[(c^2+s^2)\cos\theta_m+i(c^2-s^2)\sin\theta_m\right]\\
 &=S_P(x,y)+\frac{i}{L^2}\sum_{m=1}^{M}m\sin\theta_m.
\end{align*}
The last term changes sign when $x$ and $y$ are exchanged,
whereas the first does not.  Their half-sum proves
\eqref{eq:ch19-spatial-symmetrized}, and the half-difference
agrees with Proposition~\ref{prop:ch19-spatial-commutator}.
Finally, $(c-s)(c+s)=1$ gives
$g^{-1}=(c+s)^2$ and hence
$g+g^{-1}=2(c^2+s^2)$.
\end{proof}

The decomposition is particularly transparent:
\begin{equation}
 \boxed{\displaystyle
 D_P(x,y)=S_P(x,y)
  +\frac{i}{L^2}\sum_{m=1}^{M}m
    \sin\!\left(\frac{2\pi m(x_0-y_0)}L\right).}
 \label{eq:ch19-spatial-ordered-realimag}
\end{equation}
The imaginary, antisymmetric contribution is fixed by the CCR; the
interaction-dependent parameter $g$ affects the real symmetric part.
This is an \emph{exact finite-mode} separation, not an assertion that
the imaginary part becomes negligible at large distance or large $L$.

\subsection{Coincident points, positivity, and what positivity means}

\begin{theorem}[Diagonal evaluation and positivity]
\label{thm:ch19-spatial-diagonal}
Under the hypotheses of Theorem~\ref{thm:ch19-spatial-ordered},
for every $x\in\Lambda_L$,
\begin{equation}
 \boxed{\displaystyle
 D_P(x,x)=\frac{c^2+s^2}{L^2}\frac{M(M+1)}2
           =\frac{g+g^{-1}}{2L^2}\frac{M(M+1)}2.}
 \label{eq:ch19-spatial-diagonal}
\end{equation}
This value is real and nonnegative.  It is strictly positive when
$M\geq1$ (and $h>0$, as assumed), whereas $M=0$ gives zero.
\end{theorem}

\begin{proof}
At $x=y$, every character in \eqref{eq:ch19-spatial-ordered}
is $\chi_{\pm m}(0)=1$, so
\[
 D_P(x,x)=\frac{c^2+s^2}{L^2}\sum_{m=1}^{M}m.
\]
The elementary identity $\sum_{m=1}^{M}m=M(M+1)/2$ proves the
formula, including the empty-sum case.  Since $c^2-s^2=1$,
\[
 c^2+s^2=1+2s^2\geq1.
\]
The denominator $L^2$ is positive for $h>0$.  Therefore the right
side is strictly positive precisely when the cutoff contains at least
one positive mode.  No restriction such as $M<h$ is used.
\end{proof}

This diagonal is the second moment of the Hermitian oscillatory
density in the abstract state:
$D_P(x,x)=\omega_P(\rho_R(x)^*\rho_R(x))$.
Its finiteness results from the \emph{mode cutoff} $M$, not from a
finite-dimensional oscillator representation.  Increasing $M$ at
fixed $L$ makes the diagonal grow quadratically; thus its finite
value does not, by itself, define a cutoff-independent continuum
coincident-point observable.

There is a stronger, more useful positivity statement than
$D_P(x,x)\geq0$.  A complex two-point function is a
\emph{positive-semidefinite kernel} if for any finite collection of
sites $x_1,\ldots,x_N$ and complex numbers $z_1,\ldots,z_N$,
\begin{equation}
 \sum_{a,b=1}^{N}\overline{z_a}\,D_P(x_a,x_b)\,z_b\geq0.
 \label{eq:ch19-spatial-psd}
\end{equation}
This inequality is an inequality of \emph{real} numbers: its left
side is automatically real.

\begin{proposition}[Positive-semidefinite ordered kernel]
\label{prop:ch19-spatial-psd}
The kernel $D_P$ of Theorem~\ref{thm:ch19-spatial-ordered} satisfies
\eqref{eq:ch19-spatial-psd}; it is also Hermitian,
$D_P(y,x)=\overline{D_P(x,y)}$.
\end{proposition}

\begin{proof}
For the chosen sites and coefficients, set
$B=\sum_{b=1}^{N}z_b\rho_R(x_b)$.
Each density is self-adjoint, so
$B^*=\sum_{a=1}^{N}\overline{z_a}\rho_R(x_a)$.
Consequently,
\begin{align*}
 \sum_{a,b}\overline{z_a}D_P(x_a,x_b)z_b
 &=\omega_P\!\left(
    \sum_{a,b}\overline{z_a}z_b\rho_R(x_a)\rho_R(x_b)
    \right)\\
 &=\omega_P(B^*B)\geq0,
\end{align*}
by positivity of the abstract functional.  Hermiticity follows
from either $\omega_P(B^*)=\overline{\omega_P(B)}$ or the explicit
formula \eqref{eq:ch19-spatial-ordered} and
$\chi_{-m}(d)=\overline{\chi_m(d)}$.
\end{proof}

Equivalently, \eqref{eq:ch19-spatial-ordered} expresses the matrix
$D_P(x_a,x_b)$ as a sum of rank-one Hermitian matrices with
nonnegative weights $ms^2/L^2$ and $mc^2/L^2$.  This gives a second
way to see positivity and shows its rank is at most $2M$, possibly
less because lattice Fourier characters can coincide.  Positive
semidefiniteness does \emph{not} mean that an off-diagonal entry is
real or positive: in general $D_P(x,y)$ is complex, and even its
symmetrized value $S_P(x,y)$ may be negative.

\subsection{Worked finite example: eight sites and two positive modes}

Take an eight-site periodic lattice and two abstract positive modes,
\begin{equation}
 L=8,\quad h=4,\quad M=2,\quad
 \zeta=e^{i\pi/4},\qquad c=\frac53,\quad s=\frac43.
 \label{eq:ch19-spatial-example-data}
\end{equation}
These choices satisfy $c^2-s^2=1$ and
$g=(c-s)^2=1/9$.  The two retained frequencies are distinct and
below the Nyquist mode $m=h=4$.  The normalized density at $x=0$
is explicitly
\[
 \rho_R(0)=\frac18(C_{R1}+A_{R1}+C_{R2}+A_{R2}),
\]
while at $x=1$ the corresponding four generators are weighted by
$\zeta^{-1},\zeta,\zeta^{-2},\zeta^2$, respectively.  Only the
pairings of opposite adjoint type and equal mode index survive when
the two densities are multiplied and evaluated.

For the separation $d=x_0-y_0$, the ordered kernel becomes
\begin{align}
 D_P(d)
 &=\frac1{64}\sum_{m=1}^{2}m
   \left(\frac{16}{9}e^{-i m\pi d/4}
        +\frac{25}{9}e^{i m\pi d/4}\right)\notag\\
 &=\frac1{576}\left[
  41\left(\cos\frac{\pi d}{4}
            +2\cos\frac{\pi d}{2}\right)
  +9i\left(\sin\frac{\pi d}{4}
            +2\sin\frac{\pi d}{2}\right)\right].
 \label{eq:ch19-spatial-example-finite}
\end{align}
In particular,
\begin{align}
 D_P(0)&=\frac{41}{192},
 \label{eq:ch19-spatial-example-diagonal}\\
 D_P(1)&=\frac{41\sqrt2}{1152}
          +i\,\frac{\sqrt2+4}{128},
 \label{eq:ch19-spatial-example-d1}\\
 D_P(-1)&=\frac{41\sqrt2}{1152}
          -i\,\frac{\sqrt2+4}{128}.
 \label{eq:ch19-spatial-example-dminus1}
\end{align}
Thus the finite example checks all three operations independently:
\begin{align*}
 S_P(1)&=\frac{D_P(1)+D_P(-1)}2
       =\frac{41\sqrt2}{1152},\\
 \omega_P([\rho_R(1),\rho_R(0)])
 &=D_P(1)-D_P(-1)
   =i\,\frac{\sqrt2+4}{64}.
\end{align*}
The commutator is independent of $c,s$: evaluating the same expression
with the free pair $(c,s)=(1,0)$ gives the \emph{same} imaginary
difference, although its real symmetric part changes.
For example, in that free pair,
\begin{equation}
 D_0(1)=\frac{\sqrt2}{128}
          +i\,\frac{\sqrt2+4}{128},
 \qquad D_0(0)=\frac3{64}.
 \label{eq:ch19-spatial-example-free}
\end{equation}
Both states obey the same exact CCR; only their fluctuation
covariances differ.

For completeness, the $2\times2$ matrix associated with the sites
$0$ and $1$ is
\[
 \begin{pmatrix}
  41/192&D_P(-1)\\
  D_P(1)&41/192
 \end{pmatrix}.
\]
It is Hermitian and positive semidefinite by
Proposition~\ref{prop:ch19-spatial-psd}, even though the entries
$D_P(\pm1)$ are not real.  This is the correct interpretation of
positive spatial correlations in a noncommutative theory.

\subsection{What has, and has not, been established}

The passage from Section~\ref{sec:ch19-mode-correlators} to this
section is an exact finite Fourier reconstruction within the
\emph{abstract, uncompressed} CCR algebra.  The main formula is
valid for every integer $M\geq0$ and every pair of lattice sites,
under the stated root-of-unity and hyperbolic-pair conditions.
The commutator kernel depends only on the CCR and the chosen Fourier
normalization.  The symmetric kernel depends on the squeezing
parameter through $c^2+s^2=(g+g^{-1})/2$.  The diagonal is finite at
each cutoff and positive when $M\geq1$, and the full ordered kernel
is positive semidefinite in the sense of quadratic forms.

Three distinct extrapolations require additional work.  First,
a finite periodic character sum is not a proof of an infinite-volume
power law: any such statement needs a specified limit of $L$, $M$,
separations, and operator normalization.  Second, because the
oscillator carrier is uncompressed, the kernel is not automatically
a correlator of compressed finite-CAR density matrices; the next
comparison section must explicitly track maps and boundary
remainders.  Third, the oscillatory density used here omits a charge
zero mode, so the result should not be presented as the complete
microscopic density--density function without supplying that missing
sector and its cross terms.  These are limitations of the identified
observable, not failures of the finite algebraic calculation.

The corresponding statements in the supplied
\texttt{BosonizeStubs/Ch19Correlations.lean} snapshot are:
\begin{itemize}
 \item \texttt{ordered\_density\_kernel} and
       \texttt{ordered\_density\_kernel\_lattice};
 \item \texttt{density\_commutator\_expectation};
 \item \texttt{symmetrized\_density\_kernel};
 \item \texttt{density\_kernel\_diagonal\_pos}.
\end{itemize}
Their Phase~B status is not Core promotion, and standalone LaTeX
compilation does not replace independent mathematical review or
whole-book integration.

\section{Finite-Fock Expectations and Comparison Limits}
\label{sec:ch19-finite-comparison}

\subsection{The comparison problem: the same formula is not the same state}

Sections~\ref{sec:ch19-scope}--\ref{sec:ch19-spatial-correlators}
used an abstract positive functional on the uncompressed oscillator algebra.
One might now try to interpret its Gaussian contractions as expectation values
in the ground state of the original finite fermionic model.  That identification
is tempting: both constructions have current operators with the same labels,
and the protected low-energy currents satisfy a Schwinger relation on suitable
vectors.  Yet the two constructions have different mathematical carriers and
different hypotheses.  An equality of formal current commutators on a
restricted vector is not an equality of states on the whole operator algebra.

Here is the actual question we can answer.  \emph{Starting from a concrete,
normalized fermionic sea vector, which ordered current products can be
evaluated exactly?  Do these evaluations coincide with the abstract CCR
functional for any choice of the Bogoliubov parameters?}  The first answer
is affirmative in a bounded, explicitly protected regime.  The second is
affirmative for the \emph{free} pair, and cannot be extended by simply
replacing the free pair by an interacting one.  A short calculation below
exhibits the obstruction without taking a large-lattice limit.

To keep the distinction visible, write
\begin{equation}
 \langle X\rangle^{\mathrm F}_{N}:=
 \langle\Omega_N,X\Omega_N\rangle_{\mathcal H_{h}},
 \qquad \omega_P(a):=\omega_0\!\left(\beta_P^{-1}(a)\right).
 \label{eq:ch19-comp-two-evaluators}
\end{equation}
The first expression acts on \emph{finite CAR/Fock operators} $X$;
the second acts on \emph{abstract CCR words} $a$.  Even when the same
letter $\rho$ is used for corresponding currents, these are not, without
an additional map, elements of the same algebra.  A density-matrix state
on the finite carrier is equally possible:
$\langle X\rangle^{\mathrm F}_{\Gamma}
=\operatorname{Tr}_{\mathcal H_h}(\Gamma X)$ for
$\Gamma\geq0$ and $\operatorname{Tr}\Gamma=1$.  The proof below uses the
pure state $\Gamma=|\Omega_N\rangle\langle\Omega_N|$.

\subsection{The finite objects and the protected-current hypotheses}

Fix an integer $h\geq1$ and put $L=2h$.  The finite one-particle band has
$L$ momentum labels per species, and the fermionic Fock space
$\mathcal H_h$ built from this band is finite dimensional.  Let $\sigma$
run through the finite, ordered set of species and fix a charge assignment
$N:\sigma\to\mathbb Z$ admissible for the sea construction in
Chapters~\ref{ch:fermions} and~\ref{ch:vacuum_budget}.  Denote its
normalized sea vector by $\Omega_N$.  In the accepted implementation this
is \texttt{A05.groundKet h Ns}; its normalization is an exact CAR/Fock
statement, not an oscillator assumption:
\begin{equation}
 \langle\Omega_N,\Omega_N\rangle=1,
 \qquad \langle I\rangle_N^{\mathrm F}=1.
 \label{eq:ch19-comp-normalized}
\end{equation}
The term \emph{ground} here refers to the constructed sea or sector ground
for the finite bare setting.  It does not identify this vector as the ground
state of an arbitrary interacting Hamiltonian.

Choose the right species $\nu_R$ in the two-branch record, a positive
mode cutoff $M\in\mathbb N$, and current operators
\begin{equation}
 R_m:=\rho_{m,\nu_R}^{\mathrm F},\qquad
 R_{-m}:=\rho_{-m,\nu_R}^{\mathrm F},\qquad
 R_m^*=R_{-m},\qquad 1\leq m\leq M.
 \label{eq:ch19-comp-current-def}
\end{equation}
The current label is \emph{signed}: $R_m$ creates an excitation of
positive energy grade, whereas $R_{-m}$ lowers it.  It is important not to
replace these signed indices by their absolute values inside an operator
product.  Products act from right to left.

The finite current formulae below assume all of the following:
$h>0$; $N$ admissible; $1\leq m,n\leq M$; and the
\emph{right-branch zero-budget margin condition}
\begin{equation}
 \mathrm{marginAt}(h,M,0,N,\nu_R).
 \label{eq:ch19-comp-margin}
\end{equation}
This is the frozen-margin predicate of the finite implementation, not a
claim about arbitrary lattice parameters.  Its use here is precise: it
ensures that the signed-current commutator theorem applies on $\Omega_N$
for every pair under consideration; in particular its required excursion
inequality
\begin{equation}
 m+n+|N(\nu_R)|\leq h
 \label{eq:ch19-comp-excursion}
\end{equation}
holds for $1\leq m,n\leq M$.  The same construction supplies the
exact lowering property
\begin{equation}
 R_{-m}\Omega_N=0\qquad(1\leq m\leq M).
 \label{eq:ch19-comp-lowering}
\end{equation}
These are \emph{local, vector-dependent facts}.  They do not assert the
operator identity $[R_{-m},R_n]=m\delta_{mn}I$ on all of
$\mathcal H_h$, which would contradict the finite-dimensional trace
obstruction discussed in Section~\ref{sec:ch19-scope}.

Under the margin hypothesis the signed-current theorem gives, \emph{on
the constructed ground vector},
\begin{equation}
 [R_{-m},R_n]\Omega_N=m\delta_{mn}\Omega_N.
 \label{eq:ch19-comp-ground-commutator}
\end{equation}
This is the only commutator relation needed for the protected quadratic
expectations.  No representation of the full uncompressed CCR algebra by
finite matrices is being assumed.

\subsection{Exact ordered contractions in the bare finite ground}

\begin{proposition}[Finite ground-state current contractions]
\label{prop:ch19-comp-ground-contractions}
Under the hypotheses of~\eqref{eq:ch19-comp-margin}, for
$1\leq m,n\leq M$ one has
\begin{align}
 \langle R_{-m}R_n\rangle_N^{\mathrm F}
     &=m\delta_{mn},
 \label{eq:ch19-comp-lr}\\
 \langle R_mR_{-n}\rangle_N^{\mathrm F}
     &=0,
 \label{eq:ch19-comp-rl}\\
 \langle R_mR_n\rangle_N^{\mathrm F}
     &=0,\qquad
 \langle R_{-m}R_{-n}\rangle_N^{\mathrm F}=0.
 \label{eq:ch19-comp-same-sign}
\end{align}
In particular the finite excitation $R_m\Omega_N$ has
$\|R_m\Omega_N\|^2=m$.
\end{proposition}

\begin{proof}
Begin with the first product.  The definition of a commutator,
Equation~\eqref{eq:ch19-comp-ground-commutator}, and the rightmost
lowering relation give
\begin{align*}
 R_{-m}R_n\Omega_N
 &=\bigl([R_{-m},R_n]+R_nR_{-m}\bigr)\Omega_N\\
 &=m\delta_{mn}\Omega_N+R_n(0)
 =m\delta_{mn}\Omega_N.
\end{align*}
Taking the inner product with $\Omega_N$ and using
\eqref{eq:ch19-comp-normalized} proves
\eqref{eq:ch19-comp-lr}.  For the reverse order the lowering mode acts
first, so
$R_mR_{-n}\Omega_N=R_m(0)=0$ and
\eqref{eq:ch19-comp-rl} follows.

Similarly, $R_{-m}R_{-n}\Omega_N=0$.  For the two raising
operators, use the adjoint relation and~\eqref{eq:ch19-comp-lowering}:
\begin{align*}
 \langle\Omega_N,R_mR_n\Omega_N\rangle
 &=\langle R_m^*\Omega_N,R_n\Omega_N\rangle\\
 &=\langle R_{-m}\Omega_N,R_n\Omega_N\rangle=0.
\end{align*}
Finally, adjointness and the already-proved contraction imply
\begin{align*}
 \|R_m\Omega_N\|^2
 &=\langle R_m\Omega_N,R_m\Omega_N\rangle
 =\langle\Omega_N,R_{-m}R_m\Omega_N\rangle=m.
\end{align*}
The argument is finite-dimensional throughout; the nonzero norm is a
certificate for a concrete excited ket in the protected regime.
\end{proof}

Two points are worth separating.  First, the $m$ on the right-hand side
is not an occupation expectation: it is the coefficient of the Schwinger
commutator on the sea vector.  Second, there is no contradiction between
$\langle R_{-m}R_m\rangle=m$ and
$\langle R_mR_{-m}\rangle=0$.  The difference is exactly the
commutator expectation, and its order cannot be changed casually.

\subsection{Finite spatial density and the free-state comparison}

For spatial evaluation let $\Lambda_L=\mathbb Z/L\mathbb Z$ and choose a
primitive $L$th root $\zeta$ (the canonical choice is
$e^{2\pi i/L}$).  With the character
$\chi_m(x)=\zeta^{m x_0}$, where $x_0$ is any integer representative,
introduce the finite right-branch oscillatory density
\begin{equation}
 \rho_R^{\mathrm F}(x):=\frac1L\sum_{m=1}^{M}
 \bigl(\chi_{-m}(x)R_m+\chi_m(x)R_{-m}\bigr).
 \label{eq:ch19-comp-fock-density}
\end{equation}
This operator uses exactly the same selected mode set and Fourier
normalization as the abstract density in
Section~\ref{sec:ch19-spatial-correlators}.  It remains a finite CAR
operator; the equality of its \emph{formula} with a symbolic CCR
expression is not an identification of their representations.
For this character choice the finite density is self-adjoint.

\begin{theorem}[Protected finite-density kernel and free comparison]
\label{thm:ch19-comp-free-kernel}
Assume $h>0$, admissible $N$, the margin condition
\eqref{eq:ch19-comp-margin}, $M\geq0$, and a nonzero root $\zeta$ as
above.  Then for all $x,y\in\Lambda_L$,
\begin{equation}
 \langle \rho_R^{\mathrm F}(x)\rho_R^{\mathrm F}(y)
 \rangle_N^{\mathrm F}
 =\frac{1}{L^2}\sum_{m=1}^{M}m\chi_m(x-y).
 \label{eq:ch19-comp-fock-kernel}
\end{equation}
Moreover, this equals the \emph{abstract} ordered density kernel
evaluated in the free Bogoliubov pair $(c,s)=(1,0)$:
\begin{equation}
 \langle \rho_R^{\mathrm F}(x)\rho_R^{\mathrm F}(y)
 \rangle_N^{\mathrm F}
 =\omega_{\mathrm{free}}\bigl(\rho_R(x)\rho_R(y)\bigr).
 \label{eq:ch19-comp-free-agreement}
\end{equation}
The assertion concerns these quadratic observables and the protected
bare ground, not equality of states on their entire algebras.
\end{theorem}

\begin{proof}
Expand the product in the indicated order:
\begin{align*}
 &\rho_R^{\mathrm F}(x)\rho_R^{\mathrm F}(y)
 =\frac{1}{L^2}\sum_{m,n=1}^{M}\bigl(
 \chi_{-m}(x)\chi_{-n}(y)R_mR_n\\
 &\hspace{15mm}+\chi_{-m}(x)\chi_n(y)R_mR_{-n}
 +\chi_m(x)\chi_{-n}(y)R_{-m}R_n\\
 &\hspace{15mm}+\chi_m(x)\chi_n(y)R_{-m}R_{-n}\bigr).
\end{align*}
Linearity of the finite vector expectation and
Proposition~\ref{prop:ch19-comp-ground-contractions} kill three of
these four products.  The surviving term is the \emph{third}, not the
second:
\begin{align*}
 \langle\rho_R^{\mathrm F}(x)\rho_R^{\mathrm F}(y)
 \rangle_N^{\mathrm F}
 &=\frac{1}{L^2}\sum_{m,n=1}^{M}
 \chi_m(x)\chi_{-n}(y)n\delta_{mn}\\
 &=\frac{1}{L^2}\sum_{m=1}^{M}
 m\chi_m(x)\chi_{-m}(y)\\
 &=\frac{1}{L^2}\sum_{m=1}^{M}m\chi_m(x-y).
\end{align*}
For a merely nonzero complex $\zeta$ the corresponding source identity
is formulated in terms of chosen representatives and integer powers;
primitivity and unit modulus supply the periodic, Hermitian spatial
interpretation here.

The abstract ordered kernel of Section~\ref{sec:ch19-spatial-correlators}
is $L^{-2}\sum_m m\bigl(s^2\chi_{-m}(x-y)
+c^2\chi_m(x-y)\bigr)$.  Substituting $s=0$, $c=1$ leaves
exactly~\eqref{eq:ch19-comp-fock-kernel}.  This proves
\eqref{eq:ch19-comp-free-agreement} by comparing explicitly derived
numbers; no intertwining theorem between the finite CAR algebra and the
abstract CCR algebra has been invoked.
\end{proof}

At coincidence $x=y$ the free kernel gives the exact finite variance
\begin{equation}
 \langle(\rho_R^{\mathrm F}(x))^2\rangle_N^{\mathrm F}
 =\frac{1}{L^2}\sum_{m=1}^{M}m
 =\frac{M(M+1)}{2L^2}.
 \label{eq:ch19-comp-fock-diagonal}
\end{equation}
It is positive for $M\geq1$, and zero for the empty mode set $M=0$.
The evaluation is a genuine nonnegative vector norm since
$\rho_R^{\mathrm F}(x)$ is self-adjoint.  It is \emph{not} an
infinite-volume decay law; both $L$ and $M$ remain finite and explicit.

\subsection{Why the bare ground is not a dressed vacuum}

It remains possible for the \emph{finite} model to have dressed current
combinations, but the existence of those operators does not supply a
vector annihilated by all of them.  Let $\eta\in\{R,L\}$ and
$\bar\eta$ denote the opposite branch.  For a real hyperbolic pair
\begin{equation}
 c,s\in\mathbb R,\qquad c^2-s^2=1,
 \label{eq:ch19-comp-hyperbolic}
\end{equation}
use the finite Bogoliubov operator convention carried by the accepted
implementation:
\begin{align}
 \widetilde A_{\eta m}^{\mathrm F}
 &=c\,\rho_{-m,\eta}^{\mathrm F}
   +s\,\rho_{m,\bar\eta}^{\mathrm F},
 \label{eq:ch19-comp-dressed-ann}\\
 \widetilde C_{\eta m}^{\mathrm F}
 &=c\,\rho_{m,\eta}^{\mathrm F}
   +s\,\rho_{-m,\bar\eta}^{\mathrm F}
 =\bigl(\widetilde A_{\eta m}^{\mathrm F}\bigr)^*.
 \label{eq:ch19-comp-dressed-cre}
\end{align}
Here the signs and the \emph{opposite-branch} mixing are conventions,
not freely adjustable by a generic squeezing formula.  To control the
both-branch matrix elements assume the stronger predicate
$\mathrm{branchMargins}(h,M,0,N,\mathrm{branches})$, so the
protected relations hold for both chosen species.

For a fixed $1\leq m\leq M$ the bare lowering action gives
\begin{equation}
 \widetilde A_{\eta m}^{\mathrm F}\Omega_N
 =s\,\rho_{m,\bar\eta}^{\mathrm F}\Omega_N.
 \label{eq:ch19-comp-ann-action}
\end{equation}
Under the opposite-branch margin its squared norm is $m s^2$.
Thus the obstruction to the dressed vacuum is particularly transparent:
when $s\ne0$, the dressed \emph{annihilator itself creates a nonzero
bare excitation}.  This statement is stronger than merely noticing
that one numerical contraction differs from the abstract result.

\begin{proposition}[Dressed-number witness on the finite bare ground]
\label{prop:ch19-comp-dressed-number}
Under the admissibility, positivity, mode-range, and both-branch margin
hypotheses above,
\begin{align}
 \bigl\langle\widetilde C_{\eta m}^{\mathrm F}
 \widetilde A_{\eta m}^{\mathrm F}\bigr\rangle_N^{\mathrm F}
 &=m s^2,
 \label{eq:ch19-comp-dressed-ca}\\
 \bigl\langle\widetilde A_{\eta m}^{\mathrm F}
 \widetilde C_{\eta m}^{\mathrm F}\bigr\rangle_N^{\mathrm F}
 &=m c^2.
 \label{eq:ch19-comp-dressed-ac}
\end{align}
If $s\ne0$, $\Omega_N$ is not annihilated by the dressed lowering
operator $\widetilde A_{\eta m}^{\mathrm F}$.
\end{proposition}

\begin{proof}
Adjointness, followed by~\eqref{eq:ch19-comp-ann-action}, yields
\begin{align*}
 \bigl\langle\widetilde C_{\eta m}^{\mathrm F}
 \widetilde A_{\eta m}^{\mathrm F}\bigr\rangle_N^{\mathrm F}
 &=\bigl\|\widetilde A_{\eta m}^{\mathrm F}\Omega_N\bigr\|^2\\
 &=s^2\bigl\|\rho_{m,\bar\eta}^{\mathrm F}\Omega_N\bigr\|^2
 =m s^2.
\end{align*}
For the reverse order,
$\widetilde C_{\eta m}^{\mathrm F}\Omega_N
=c\rho_{m,\eta}^{\mathrm F}\Omega_N$, since the opposite-branch
lowering operator kills the bare ground.  Therefore
\begin{align*}
 \bigl\langle\widetilde A_{\eta m}^{\mathrm F}
 \widetilde C_{\eta m}^{\mathrm F}\bigr\rangle_N^{\mathrm F}
 &=\bigl\|\widetilde C_{\eta m}^{\mathrm F}\Omega_N\bigr\|^2
 =c^2\bigl\|\rho_{m,\eta}^{\mathrm F}\Omega_N\bigr\|^2
 =m c^2.
\end{align*}
Since $m>0$, the first value is strictly positive whenever
$s\ne0$, proving that the dressed annihilator cannot kill $\Omega_N$.
Neither step uses a finite global CCR identity.
\end{proof}

The apparent numerical parallel to the abstract contractions requires
care.  In the \emph{abstract interacting} state $\omega_P$, the dressed
annihilators annihilate the state in the algebraic sense, whereas the
bare contractions obey $\omega_P(C_{Rm}A_{Rm})=ms^2$ and
$\omega_P(A_{Rm}C_{Rm})=mc^2$.  Above we obtained the same two
numbers for \emph{dressed} finite operators evaluated in the
\emph{bare} finite ground.  These are different ordered products in
different algebras and different states.  Equality of the displayed
coefficients does not license swapping their roles.  In particular,
$\omega_P(\widetilde C_{Rm}\widetilde A_{Rm})=0$, whereas
\eqref{eq:ch19-comp-dressed-ca} is nonzero for $s\ne0$.

\subsection{A complete finite witness: two modes on sixteen sites}

\begin{example}[Source-verified finite/abstract discrepancy]
\label{ex:ch19-comp-16-sites}
Take $h=8$, so $L=16$, positive-mode cutoff $M=2$, and the explicit
admissible charge/branch data from the CH18 concrete witness.  Their
both-branch margin predicates hold.  The hyperbolic coefficients used
in that witness are
\begin{equation}
 c=\frac54,\qquad s=\frac34,
 \qquad c^2-s^2=\frac{25-9}{16}=1.
 \label{eq:ch19-comp-example-pair}
\end{equation}
For $m=2$, the protected current contraction and excited-vector norm
are
\begin{equation}
 \langle R_{-2}R_2\rangle_N^{\mathrm F}=2,
 \qquad\|R_2\Omega_N\|^2=2,
 \qquad\langle R_2R_{-2}\rangle_N^{\mathrm F}=0.
 \label{eq:ch19-comp-example-bare}
\end{equation}
The dressed-number test gives
\begin{equation}
 \bigl\langle\widetilde C_{R2}^{\mathrm F}
 \widetilde A_{R2}^{\mathrm F}\bigr\rangle_N^{\mathrm F}
 =2\left(\frac34\right)^2=\frac98>0.
 \label{eq:ch19-comp-example-dressed}
\end{equation}
Consequently the \emph{actual finite bare ground} is not an
annihilation vacuum for these nontrivial dressed combinations.

At the same site the finite bare density has
\begin{equation}
 \langle\rho_R^{\mathrm F}(0)\rho_R^{\mathrm F}(0)
 \rangle_N^{\mathrm F}
 =\frac{1+2}{16^2}=\frac3{256}.
 \label{eq:ch19-comp-example-free}
\end{equation}
In contrast, in the \emph{abstract interacting} state the two-mode
ordered kernel is
\begin{align}
 \omega_P\bigl(\rho_R(0)\rho_R(0)\bigr)
 &=\frac{1+2}{16^2}(c^2+s^2)\notag\\
 &=\frac3{256}\frac{34}{16}
 =\frac{51}{2048}.
 \label{eq:ch19-comp-example-abstract}
\end{align}
The two finite numbers are unequal:
$3/256=24/2048\ne51/2048$.  This comparison does \emph{not}
show that no interacting finite ground state exists.  It shows exactly
that the constructed \emph{bare} finite ground and the interacting
abstract quasi-free functional cannot be identified on this chosen
density product.  The first value is certified by the finite CAR
margin calculation; the second by the abstract CCR construction.
\end{example}

Notice that this example uses $c=5/4$ and $s=3/4$.  The distinct
$5/3,4/3$ pair used for the separate compressed-three-dimensional
common-kernel counterexample in Section~\ref{sec:ch19-scope} must not be substituted
here.  These are different mathematical witnesses with different
operators and hypotheses.

\subsection{What an interacting finite comparison would actually require}

The preceding computations settle a useful \emph{free bare-ground
comparison}, but they leave a genuine interacting-state problem.  It
is worth stating the missing pieces because the later order-parameter
chapters will otherwise invite a false shortcut.

\paragraph{A comparison requires an operator map and a specified state.}
Let $\mathfrak A_{\mathrm{CCR}}$ be the algebra used in
Sections~\ref{sec:ch19-scope}--\ref{sec:ch19-spatial-correlators}.
A proposed comparison would have to specify, at least on a chosen
set of words $\mathcal W$, concrete finite-Fock operators
$\Phi_h(a)$ and a finite state $\Gamma_h$.  One could then define
an actual defect, for each \emph{specified} word $a\in\mathcal W$,
\begin{equation}
 \varepsilon_h(a)
 :=\operatorname{Tr}\!\bigl(\Gamma_h\Phi_h(a)\bigr)-\omega_P(a).
 \label{eq:ch19-comp-comparison-defect}
\end{equation}
No theorem in this section establishes that such a map is a unital
$*$-homomorphism on the full CCR algebra; the finite-dimensional trace
obstruction rules out a globally exact nonzero scalar-CCR
representation.  Neither is a bound $\varepsilon_h(a)\to0$ claimed.
For the density products at the free pair we have an explicitly
constructed \emph{wordwise numerical equality}; for arbitrary
interacting pairs no analogous theorem follows from this packet.

\paragraph{Compression introduces a boundary correction.}
Even if one begins with an ambient space $\mathcal H$ and a
finite-rank orthogonal projection $P$ onto an energy-budget subspace,
compression does not commute with operator multiplication.  Put
$Q=I-P$.  The exact algebraic identity on $P\mathcal H$ is
\begin{equation}
 [PXP,PYP]
 =P[X,Y]P-PXQYP+PYQXP.
 \label{eq:ch19-comp-compression-identity}
\end{equation}
It follows by inserting $I=P+Q$ between the factors of
$P[X,Y]P$.  The two last terms measure excursions out of the budget
space and back.  Thus even when an ambient commutator formula is known,
the \emph{compressed} commutator is controlled only after its boundary
remainder has been estimated or shown to vanish on the chosen states.
The protected relation~\eqref{eq:ch19-comp-ground-commutator} is a
local certificate, not a license to discard these terms for every word.

\paragraph{Dressed-vacuum conditions and Hamiltonian ground states
are independent obligations.}
The positive finite operator
$\widetilde C\widetilde A=\widetilde A^*\widetilde A$
is a norm-square test.  The nonzero result
\eqref{eq:ch19-comp-dressed-ca} disproves an annihilation-vacuum
claim for $\Omega_N$; it does not rule out another ground vector of
an interacting finite Hamiltonian $H_{h,\mathrm{int}}$.
To compare that ground to $\omega_P$, one must construct its
normalized vector or density matrix, prove the relevant spectral
statement, and evaluate the same observables with a controlled
projection or comparison mechanism.

\paragraph{Prerequisites for later order parameters.}
Physical CDW, pairing and spinful correlation functions require more
than oscillator two-point contractions: their concrete finite-CAR
operator definitions, charge-sector shifts and Klein-factor signs,
ordering and adjoints, as well as any projection residuals, must be
fixed before assigning a finite expectation.  Formal vertex
exponentials in a power-series algebra are a separate coefficientwise
construction; their identities alone neither define convergent
operators at a numerical parameter nor compare them with finite CAR
observables.  Any claimed continuum power law additionally needs a
specified lattice/cutoff scaling and uniform error estimates.
The CH20 and CH21 developments are not established by the CH19
Phase-B results.

\subsection{Conclusion: exactly what has been compared}

We have constructed the finite expectation
$\langle\cdot\rangle_N^{\mathrm F}$ on a normalized fermionic sea,
used the frozen-margin and lowering hypotheses to prove its ordered
current contractions, and derived an exact finite spatial kernel.
Its agreement with the abstract CCR kernel at $(c,s)=(1,0)$ is a
\emph{quadratic free-state comparison}.  For $s\ne0$ the same bare
ground exhibits a strictly nonzero dressed-number expectation,
providing a direct witness against treating it as the interacting
algebraic vacuum.  The next task is therefore not to repeat the
Gaussian calculation, but to specify and verify a finite interacting
state together with its observable comparison map and boundary
estimates.  The accepted CH19/A08 Phase-B Lean results support the
claims made above; they have not thereby been promoted to the frozen
Core, nor does successful compilation of this prose supply an
independent mathematical review.


\chapter{Projected Order Parameters, Formal Vertices, and Model Duality}
\label{ch:order-parameters}
\section{Exact Finite CAR Products: Density Waves and Pairing}
\label{sec:ch20-projected-products}

A single fermion field changes particle number by one.  A charge-density-wave
(CDW) operator, by contrast, removes a left-moving fermion and creates a
right-moving fermion; a superconducting (SC) pair-annihilation operator removes
one fermion of each chirality.  Both are products of two fields.  This simple
observation raises a question which does not arise for a single field:
\emph{what happens between the first operation and the second when our
calculation is confined to an energy-budget subspace?}

The full finite fermionic Fock space has an unambiguous answer: multiply the
operators in their stated order.  Budget restrictions introduce a different
operation, however.  The product of \emph{individually projected} fields
passes through a selected intermediate subspace, whereas a projection of the
\emph{full product} allows intermediate states to leave that subspace and
return before the final projection.  Those discarded paths form an explicit
remainder.  Keeping the remainder is the purpose of this section.  The
construction also provides a convenient setting in which to check chirality,
charge labels, Fourier phases, and fermionic signs before introducing the
formal vertex operators of Section~\ref{sec:ch20-formal-vertices}.

Throughout this section, all identities concerning finite fermion operators
are elementary statements on a finite-dimensional Hilbert space; their
application to the chapter's proposed bosonization dictionary remains a
separate mathematical verification task.  In particular, no new Chapter~20
Lean theorem is being claimed here.

\subsection{The ambient finite-CAR system}
\label{subsec:ch20-ambient}

Fix an integer \(h\geq 1\) and put \(L=2h\).  The spatial lattice is
\(\mathbb Z/L\mathbb Z\), represented by integers \(x\) modulo \(L\); the
centered momentum band is
\begin{equation}
  \mathcal B_L=\{-h+1,-h+2,\ldots,h\},\qquad
  \zeta=\exp(2\pi i/L).
  \label{eq:ch20-band}
\end{equation}
There are two distinct species, called right and left and denoted
\(R\) and \(L\) in subscripts.  The letter \(L\) without a subscript
continues to mean lattice length.  We order the single-particle orbitals
\emph{first by species} \(R<L\), then by increasing momentum within each
species.  This order fixes the occupation-basis signs and is not a harmless
convention if intermediate formulas are compared term by term.

Let
\(\mathcal F=\bigwedge(\mathbb C^{\{R,L\}\times\mathcal B_L})\)
be the fermionic Fock Hilbert space.  Its mode annihilators and creators
\(a_{\nu k}\) and \(a_{\nu k}^{\dagger}\) obey the canonical
anticommutation relations (CAR)
\begin{align}
 \{a_{\nu k},a_{\eta q}\}&=0,
 &\{a_{\nu k}^{\dagger},a_{\eta q}^{\dagger}\}&=0,
 &\{a_{\nu k},a_{\eta q}^{\dagger}\}&=
             \delta_{\nu\eta}\delta_{kq}I.
 \label{eq:ch20-car}
\end{align}
Here \(\{A,B\}=AB+BA\), and the rightmost factor acts first.  The
periodic-twist fields used in the calculations below are
\begin{equation}
 c_{\nu,x}=\frac1{\sqrt L}\sum_{k\in\mathcal B_L}
       \zeta^{kx}a_{\nu k},\qquad
 c_{\nu,x}^{\dagger}=\frac1{\sqrt L}\sum_{k\in\mathcal B_L}
       \zeta^{-kx}a_{\nu k}^{\dagger}.
 \label{eq:ch20-fields}
\end{equation}
Thus the annihilation field carries a \emph{positive} Fourier exponent,
whereas the creator carries its complex conjugate.  Both species use this
same computational band; the left-moving physical branch is encoded by the
branch identification rather than by reversing the exponent in
\eqref{eq:ch20-fields}.  In particular, a continuum oscillation
\(e^{\pm2ik_Fx}\) has \emph{not} been inserted by hand.  Formula
\eqref{eq:ch20-fields} is the periodic-twist specialization of the
finite fields from Chapter~14; a general boundary twist contributes
additional factors and is not covered by the phase formulas derived below.

For later comparison with positive-frequency currents, we reserve
\(M\in\mathbb N\) for an optional mode cutoff with labels
\(m=1,\ldots,M\) (an empty set when \(M=0\)).  The finite-CAR fields
\eqref{eq:ch20-fields} sum over the \emph{whole} band \(\mathcal B_L\);
there is no cutoff \(M\) in their definition or in the exact product
identities of this section.  Restrictions on \(M\) needed for a particular
non-wrapping current representation must be imposed in that representation.

\begin{definition}[Finite CDW and pair-annihilation words]
\label{def:ch20-order-words}
At a lattice site \(x\), define bounded endomorphisms of \(\mathcal F\) by
\begin{equation}
  O_{\mathrm{CDW}}(x)=c_{R,x}^{\dagger}c_{L,x},\qquad
  O_{\mathrm{SC}}(x)=c_{R,x}c_{L,x}.
  \label{eq:ch20-order-words}
\end{equation}
The order of the factors in each definition is part of the definition.
\(O_{\mathrm{SC}}\) annihilates a pair; its adjoint
\(O_{\mathrm{SC}}^{\dagger}=c_{L,x}^{\dagger}c_{R,x}^{\dagger}\)
creates a pair in the reversed order dictated by the adjoint.
\end{definition}

Both operators are \emph{even CAR words}: a fermionic parity transformation
changes the sign of each field but not of their product.  Evenness does not
mean charge neutrality.  The CDW operator is neutral only with respect to
the \emph{total} charge, whereas the SC operator lowers it by two.

To state this exactly, let
\begin{equation}
  \widehat n_{\nu}=\sum_{k\in\mathcal B_L}
     a_{\nu k}^{\dagger}a_{\nu k},\qquad
  Q_{\nu}=\widehat n_{\nu}-hI,
 \label{eq:ch20-charges}
\end{equation}
so that the reference half-filled sea carries zero relative charge
\(Q_R=Q_L=0\).  For a charge label
\(N=(N_R,N_L)\in\mathbb Z^2\), let
\(e_R=(1,0)\) and \(e_L=(0,1)\).

\begin{proposition}[Exact charge shifts]
\label{prop:ch20-charge-shifts}
For either \(\nu\in\{R,L\}\), the following operator identities hold on
\emph{all} of \(\mathcal F\):
\begin{align}
 [Q_{\nu},O_{\mathrm{CDW}}(x)]
   &=\bigl(\delta_{\nu R}-\delta_{\nu L}\bigr)
                       O_{\mathrm{CDW}}(x),
 \label{eq:ch20-cdw-charge}\\
 [Q_{\nu},O_{\mathrm{SC}}(x)]
   &=-\bigl(\delta_{\nu R}+\delta_{\nu L}\bigr)
                       O_{\mathrm{SC}}(x).
 \label{eq:ch20-sc-charge}
\end{align}
Consequently, with \(Q_{\mathrm{tot}}=Q_R+Q_L\),
\([Q_{\mathrm{tot}},O_{\mathrm{CDW}}]=0\) and
\([Q_{\mathrm{tot}},O_{\mathrm{SC}}]=-2O_{\mathrm{SC}}\).
\end{proposition}
\begin{proof}
For a single orbital \(j=(\eta,q)\), the CAR imply
\([a_j^{\dagger}a_j,a_i]=-\delta_{ji}a_i\) and
\([a_j^{\dagger}a_j,a_i^{\dagger}]=\delta_{ji}a_i^{\dagger}\).
For example,
\(a_j^{\dagger}a_ja_i-a_i a_j^{\dagger}a_j
 =-a_j^{\dagger}a_i a_j-(\delta_{ij}-a_j^{\dagger}a_i)a_j
 =-\delta_{ij}a_j\), which equals \(-\delta_{ij}a_i\).
Summing over \(q\) and using the linear Fourier transforms gives
\begin{equation}
 [Q_\nu,c_{\eta,x}]=-\delta_{\nu\eta}c_{\eta,x},
 \qquad
 [Q_\nu,c_{\eta,x}^{\dagger}]
       =\delta_{\nu\eta}c_{\eta,x}^{\dagger}.
 \label{eq:ch20-single-charge}
\end{equation}
The scalar \(-hI\) commutes with every operator.  Finally apply the
commutator product rule \([Q,AB]=[Q,A]B+A[Q,B]\) to each ordered
product in \eqref{eq:ch20-order-words}; summing the two charge identities
proves the final assertions.
\end{proof}

In particular, the source, intermediate and target charges for our two
products are forced by the operators, not chosen independently:
\begin{equation}
\begin{aligned}
 N_s=N\quad&\xrightarrow{\ c_{L,x}\ }
       N_m=N-e_L\\[-2pt]
       &\xrightarrow{\ c_{R,x}^{\dagger}\ }
       N_t^{\mathrm{CDW}}=N+e_R-e_L,\\[3pt]
 N_s=N\quad&\xrightarrow{\ c_{L,x}\ }
       N_m=N-e_L\\[-2pt]
       &\xrightarrow{\ c_{R,x}\ }
       N_t^{\mathrm{SC}}=N-e_R-e_L.
\end{aligned}
\label{eq:ch20-charge-arrows}
\end{equation}
The arrows are drawn in \emph{time-of-application order}, from right to
left in the written CAR word.  Because \(c_R\) and \(c_L\)
anticommute, reversing their positions introduces a minus sign even
though the charge shift is unchanged.

\subsection{Three budget carriers, not one}
\label{subsec:ch20-budget}

For each admissible charge \(N_\nu\in[-h,h]\cap\mathbb Z\), the
\emph{sector ground occupation} fills momenta
\(\{-h+1,\ldots,N_\nu\}\) for species \(\nu\); if
\(N_\nu=-h\) the set is empty.  Write \(|\Omega_N\rangle\) for the
normalized occupation ket obtained by creating all these particles in
the fixed \(R<L\), increasing-momentum order.  Thus
\begin{equation}
  Q_\nu|\Omega_N\rangle=N_\nu|\Omega_N\rangle.
 \label{eq:ch20-ground-charge}
\end{equation}
An occupation configuration \(S=(S_R,S_L)\) in this sector has a
nonnegative integer excitation grade
\begin{equation}
 E_N(S)=\sum_{\nu=R,L}\left(
      \sum_{k\in S_\nu}k
      -\sum_{k=-h+1}^{N_\nu}k\right).
 \label{eq:ch20-excitation}
\end{equation}
Indeed, among subsets of \(\mathcal B_L\) of cardinality \(h+N_\nu\),
the lowest \(h+N_\nu\) momenta minimize the sum.  Equivalently,
\(E_N(S)\) counts the net upward displacement of particles relative to
the sector ground.  The same grading underlies the budget spaces of
Chapter~\ref{ch:vacuum_budget}; in the present finite-dimensional carrier
it can be read directly from occupation kets.

For \(K\in\mathbb N\), define the cutoff carrier
\begin{equation}
 V_{N,K}=\operatorname{span}\{\,|S\rangle:
       Q_\nu|S\rangle=N_\nu|S\rangle\ (\nu=R,L),
       \ E_N(S)\leq K\,\}\ \subseteq\mathcal F.
 \label{eq:ch20-budget-carrier}
\end{equation}
For an inadmissible charge we take \(V_{N,K}=\{0\}\).  Each such space
has the inner product inherited from \(\mathcal F\).  Denote its
isometric inclusion by \(i_{N,K}:V_{N,K}\hookrightarrow\mathcal F\)
and its orthogonal projection by
\(p_{N,K}:\mathcal F\to V_{N,K}\).  The associated ambient
projection is
\begin{equation}
 P_{N,K}=i_{N,K}p_{N,K},\quad
 P_{N,K}^{2}=P_{N,K}=P_{N,K}^{\dagger},\quad
 p_{N,K}i_{N,K}=I_{V_{N,K}}.
 \label{eq:ch20-typed-projections}
\end{equation}
The notation distinguishes an endomorphism of the ambient Fock space
from a map between two \emph{different} budget carriers.  These are not
interchangeable objects, even if they happen to have the same dimension.

Choose three cutoff integers \(K_s,K_m,K_t\geq0\).  In the formulas
below, the abbreviations
\begin{equation}
\begin{gathered}
 V_s=V_{N,K_s},\qquad V_m=V_{N-e_L,K_m},\\
 V_t^{J}=V_{N_t^{J},K_t},\qquad
 i_s=i_{N,K_s},\quad i_m=i_{N-e_L,K_m},\quad
 p_m=p_{N-e_L,K_m},\quad P_m=i_mp_m
\end{gathered}
\label{eq:ch20-three-carriers}
\end{equation}
are \emph{typed} by the charge labels in
\eqref{eq:ch20-charge-arrows}; the target projection is written
\(p_t^J=p_{N_t^J,K_t}\).  No condition such as
\(K_s=K_m=K_t\) is imposed.  The intermediate cutoff is a separate
choice and must not be silently identified with the target cutoff.

\begin{definition}[Restricted one-field arrows]
\label{def:ch20-field-arrows}
Set
\begin{equation}
 B_L=p_m c_{L,x}i_s:V_s\longrightarrow V_m,
 \label{eq:ch20-BL}
\end{equation}
and define the second arrows, with the indicated \emph{different targets},
by
\begin{align}
 B_R^{\mathrm{CDW},\dagger}
   &=p_t^{\mathrm{CDW}}c_{R,x}^{\dagger}i_m:
         V_m\longrightarrow V_t^{\mathrm{CDW}},
 \label{eq:ch20-BR-CDW}\\
 B_R^{\mathrm{SC}}
   &=p_t^{\mathrm{SC}}c_{R,x}i_m:
         V_m\longrightarrow V_t^{\mathrm{SC}}.
 \label{eq:ch20-BR-SC}
\end{align}
The dagger in \(B_R^{\mathrm{CDW},\dagger}\) refers to the underlying
ambient creator.  It does \emph{not} purport to take the adjoint of
\(B_L\), nor of a field with unrelated source or target data.
\end{definition}

There is a genuine Hilbert-adjoint statement behind the creator notation.
Define the oppositely directed map
\(D_R=p_m c_{R,x}i_{N_t^{\mathrm{CDW}},K_t}:
V_t^{\mathrm{CDW}}\to V_m\).  Then
\begin{equation}
 D_R^{\dagger}
   =p_t^{\mathrm{CDW}}c_{R,x}^{\dagger}i_m
   =B_R^{\mathrm{CDW},\dagger}.
 \label{eq:ch20-reversed-adjoint}
\end{equation}
To verify this rather than infer it from notation, take
\(u\in V_m\) and \(v\in V_t^{\mathrm{CDW}}\):
\begin{align*}
 \langle D_Rv,u\rangle_{V_m}
 &=\langle c_{R,x}i_tv,i_mu\rangle_{\mathcal F}\\
 &=\langle i_tv,c_{R,x}^{\dagger}i_mu\rangle_{\mathcal F}
  =\langle v,p_t^{\mathrm{CDW}}c_{R,x}^{\dagger}i_mu\rangle_{V_t^{\mathrm{CDW}}}.
\end{align*}
Here \(i_t=i_{N_t^{\mathrm{CDW}},K_t}\) in this display.
Thus the adjoint reverses \emph{both} the charge data and the cutoffs.
This is the precise finite Hilbert-space meaning of Chapter~14's
budget-adjoint construction.

\begin{theorem}[Two-field projection insertion and retained residual]
\label{thm:ch20-projection-residual}
For every admissible or inadmissible \(N\), every choice of
\(K_s,K_m,K_t\geq0\), and every \(x\in\mathbb Z\), the following
identities of linear maps hold with the indicated source and target:
\begin{align}
 p_t^{\mathrm{CDW}}O_{\mathrm{CDW}}(x)i_s
 &=B_R^{\mathrm{CDW},\dagger}\circ B_L
   +\mathcal R_{\mathrm{CDW}}^{smt}(x),
 &V_s&\longrightarrow V_t^{\mathrm{CDW}},
 \label{eq:ch20-cdw-residual}\\
 p_t^{\mathrm{SC}}O_{\mathrm{SC}}(x)i_s
 &=B_R^{\mathrm{SC}}\circ B_L
   +\mathcal R_{\mathrm{SC}}^{smt}(x),
 &V_s&\longrightarrow V_t^{\mathrm{SC}},
 \label{eq:ch20-sc-residual}
\end{align}
where the \emph{uncompressed intermediate excursions} are
\begin{align}
 \mathcal R_{\mathrm{CDW}}^{smt}(x)
   &=p_t^{\mathrm{CDW}}c_{R,x}^{\dagger}(I_{\mathcal F}-P_m)
                                   c_{L,x}i_s,
 \label{eq:ch20-cdw-R}\\
 \mathcal R_{\mathrm{SC}}^{smt}(x)
   &=p_t^{\mathrm{SC}}c_{R,x}(I_{\mathcal F}-P_m)
                                   c_{L,x}i_s.
 \label{eq:ch20-sc-R}
\end{align}
In particular, the factorization into the two projected single fields
is exact \emph{if and only if} the corresponding displayed residual map
vanishes.  No such vanishing is assumed in the identities.
\end{theorem}

\begin{proof}
The ambient projection satisfies the tautological decomposition
\(I_{\mathcal F}=P_m+(I_{\mathcal F}-P_m)\).
Insert it between the two factors of the CDW word:
\begin{align*}
 p_t^{\mathrm{CDW}}c_{R,x}^{\dagger}c_{L,x}i_s
 & =p_t^{\mathrm{CDW}}c_{R,x}^{\dagger}P_m c_{L,x}i_s
   +p_t^{\mathrm{CDW}}c_{R,x}^{\dagger}
                        (I_{\mathcal F}-P_m)c_{L,x}i_s\\
 & =\bigl(p_t^{\mathrm{CDW}}c_{R,x}^{\dagger}i_m\bigr)
    \circ\bigl(p_m c_{L,x}i_s\bigr)
   +\mathcal R_{\mathrm{CDW}}^{smt}(x),
\end{align*}
where the second equality uses \(P_m=i_mp_m\).  This is
\eqref{eq:ch20-cdw-residual}.  Replace the first creator by
\(c_{R,x}\) to obtain \eqref{eq:ch20-sc-residual}.  All expressions
are maps between the spaces explicitly shown in
\eqref{eq:ch20-three-carriers}; the proof makes no assumption about
how many intermediate excitations a field can generate.
Finally, subtract the factored term from either equality: it vanishes
exactly when its residual does.
\end{proof}

An immediate \emph{sufficient}, though not necessary, vanishing
criterion is
\begin{equation}
 (I_{\mathcal F}-P_m)c_{L,x}i_s=0,
 \quad\text{equivalently}\quad
 c_{L,x}(V_s)\subseteq i_m(V_m).
 \label{eq:ch20-safe-intermediate}
\end{equation}
This is the strongest, easiest-to-check form of an intermediate
\emph{range condition}: every state produced by the first field already
lies in the proposed intermediate budget.  A weaker possibility is that
the second field and final projection kill all excluded states:
\(p_t^J c_{R,x}^{\sharp}(I-P_m)c_{L,x}i_s=0\), where
\(c_{R,x}^{\sharp}=c_{R,x}^{\dagger}\) for CDW and
\(c_{R,x}^{\sharp}=c_{R,x}\) for SC.  Such a cancellation is about the
\emph{complete} product, and cannot be inferred by checking each factor
in isolation.

\begin{remark}[Why an intermediate projection changes multiplication]
\label{rem:ch20-not-multiplicative}
Compression \(A\mapsto p A i\) is linear, but generally not
multiplicative.  An intermediate state can lie outside \(V_m\) and still
be mapped into \(V_t^J\) by the second field.  Inserting \(P_m\)
deletes that route.  Even a correct identity for each \emph{one-field}
projected map therefore does not license multiplying those identities
and discarding \(\mathcal R_J^{smt}\).  In particular, Chapter~14's
finite Taylor candidates, if separately shown to equal the appropriate
individual projected fields under its conditional cyclic criterion,
still require the present residual and the actual transition cutoffs when
composed.  No factorization, commutation, or Baker--Campbell--Hausdorff
identity for \emph{compressed exponentials} follows from
\eqref{eq:ch20-cdw-residual}--\eqref{eq:ch20-sc-residual}.
\end{remark}

\subsection{Charge phases and the ordered Klein signs}
\label{subsec:ch20-klein}

The operator identities above require no vacuum or special reference
state.  It is nevertheless useful to compute the matrix elements between
\emph{sector ground kets}, both as a stringent sign check and as the
zero-excitation portion of a possible order-parameter dictionary.  We now
assume the periodic twist in \eqref{eq:ch20-fields} and the fixed
occupation ordering \(R<L\).

The normalized occupation ket at charge \(N\) contains exactly
\(h+N_\nu\) particles of species \(\nu\).  In the occupation basis,
annihilation of an occupied orbital contributes \((-1)^n\), with \(n\)
the number of occupied orbitals preceding it; creation of an unoccupied
orbital obeys the same sign rule.  Define the signed Klein coefficient for
removal of the \emph{highest occupied} orbital of a species by
\begin{equation}
 \begin{aligned}
 P_R(N)&=(-1)^{h+N_R-1},\\
 P_L(N)&=(-1)^{(h+N_R)+(h+N_L-1)}.
 \end{aligned}
 \label{eq:ch20-klein-phases}
\end{equation}
Equivalently,
\begin{equation}
 P_\nu(N)=(-1)^{\displaystyle
     \sum_{\eta<\nu}(h+N_\eta)+h+N_\nu-1},
 \label{eq:ch20-klein-general}
\end{equation}
where \(\eta<\nu\) refers to the fixed species ordering and
\((-1)^n\) is defined for every integer \(n\) by its parity.
The formula is the charge-sector sign convention of Chapter~13,
not an independently adjustable choice of relative Klein phases.

For removal of the last occupied momentum \(k=N_\nu\), the annihilation
field supplies the Fourier coefficient \(L^{-1/2}\zeta^{N_\nu x}\).
For addition of the first unoccupied momentum \(k=N_R+1\), the creator
supplies \(L^{-1/2}\zeta^{-(N_R+1)x}\).  The creation sign equals
\(P_R(N+e_R)\): it counts \(h+N_R\) occupied right orbitals
\emph{before} the new particle is added.  This shifted argument matters.
Likewise, after the left annihilation in an SC product the right removal
sign is \(P_R(N-e_L)\); no factor may be moved past the other without
recomputing its source label.

\begin{proposition}[Exact sector-ground matrix elements]
\label{prop:ch20-ground-coeff}
Let \(-h\leq N_R,N_L\leq h\), and let
\(\langle\Omega_{N'}|A|\Omega_N\rangle\) denote the finite-Fock
Hilbert-space matrix element.  If \(N_L>-h\) and \(N_R<h\), then
\begin{align}
 &\langle\Omega_{N+e_R-e_L}|
                O_{\mathrm{CDW}}(x)|\Omega_N\rangle\\
 &\hspace{1cm}=\frac1L
     \zeta^{(N_L-N_R-1)x}
     P_L(N)P_R(N-e_L+e_R).
 \label{eq:ch20-cdw-ground}
\end{align}
If \(N_L>-h\) and \(N_R>-h\), then
\begin{align}
 &\langle\Omega_{N-e_R-e_L}|
                O_{\mathrm{SC}}(x)|\Omega_N\rangle\\
 &\hspace{1cm}=\frac1L
     \zeta^{(N_R+N_L)x}
     P_L(N)P_R(N-e_L).
 \label{eq:ch20-sc-ground}
\end{align}
If a required removal or creation is Pauli-forbidden at the corresponding
empty or full boundary, the respective ground-to-ground matrix element is
zero.  The nonzero coefficients above have modulus \(1/L\).
\end{proposition}
\begin{proof}
Expand the product using \eqref{eq:ch20-fields}:
\begin{align*}
 O_{\mathrm{CDW}}(x)
  &=\frac1L\sum_{p,q\in\mathcal B_L}
      \zeta^{(-p+q)x}a_{Rp}^{\dagger}a_{Lq},\\
 O_{\mathrm{SC}}(x)
  &=\frac1L\sum_{p,q\in\mathcal B_L}
      \zeta^{(p+q)x}a_{Rp}a_{Lq}.
\end{align*}
In the CDW case the target ket contains precisely the occupied left
orbitals \(-h+1,\ldots,N_L-1\); hence the \emph{only} term that
can connect the two specified ground configurations has \(q=N_L\).
Similarly, the target right orbitals are \(-h+1,\ldots,N_R+1\),
which forces \(p=N_R+1\).  The left field acts first, giving
\(P_L(N)\).  The right creator then adds the first empty right
orbital to the intermediate state \(N-e_L\), with sign
\((-1)^{h+N_R}=P_R(N-e_L+e_R)\).  The two Fourier characters
multiply to \(\zeta^{(N_L-N_R-1)x}\).  All other terms are either
zero by Pauli exclusion or produce an occupation ket orthogonal to
the stated target ground.

For SC the target has one fewer particle of each species.  Its ground
occupation forces \(q=N_L\) and \(p=N_R\).  Acting first with the
left field yields \(P_L(N)\); the subsequent right removal sees the
charge \(N-e_L\), and gives \(P_R(N-e_L)\).  Its Fourier factor is
\(\zeta^{(N_L+N_R)x}\).  The asserted boundary cases follow because
a missing highest occupied momentum cannot be annihilated, and a full
right band offers no unoccupied orbital for CDW creation.  Because
\(|\zeta|=|P_\nu|=1\), the modulus of each allowed matrix element
is \(1/L\).
\end{proof}

It is helpful to separate the phase coming from the \emph{source
occupation} from the signed shift.  For these ground matrix elements
write
\begin{align}
 \chi_{\mathrm{CDW}}(N,x)&=\zeta^{(N_L-N_R-1)x},
 &\kappa_{\mathrm{CDW}}(N)
   &=P_L(N)P_R(N-e_L+e_R),
 \label{eq:ch20-cdw-chi-kappa}\\
 \chi_{\mathrm{SC}}(N,x)&=\zeta^{(N_R+N_L)x},
 &\kappa_{\mathrm{SC}}(N)
   &=P_L(N)P_R(N-e_L).
 \label{eq:ch20-sc-chi-kappa}
\end{align}
The scalar \(\chi_J(N,x)\) acts on the \emph{source} charge label
before either shift is performed.  The left removal acts first in
both products.  Its Klein sign is evaluated at \(N\); the next
Klein sign uses the appropriate shifted charge.  These functions are
\emph{matrix-element coefficients}; they are not an assertion that
\(O_J\) equals a charge shift multiplied by a scalar on arbitrary
excited states.  Excited states require the full CAR word and, when
budget-restricted, its residual.

\subsection{A four-site calculation: the residual is genuinely present}
\label{subsec:ch20-four-site}

\begin{example}[Both order parameters and a missed intermediate path]
\label{ex:ch20-four-site}
Take \(h=2\), hence \(L=4\),
\(\mathcal B_4=\{-1,0,1,2\}\),
\(\zeta=i\), and \(x=1\).  For compactness write
\(|R\{\ldots\};L\{\ldots\}\rangle\) for the normalized occupation
ket with the indicated occupied momenta, in the \(R<L\) ordering.

\emph{First: the CDW ground coefficient.}
The source at \(N=(0,0)\) is
\[
 |\Omega_{(0,0)}\rangle
  =|R\{-1,0\};L\{-1,0\}\rangle.
\]
After a left removal and right creation the target sector is
\(N_t=(1,-1)\), with ground
\[
 |\Omega_{(1,-1)}\rangle
  =|R\{-1,0,1\};L\{-1\}\rangle.
\]
Only \(a_{R,1}^{\dagger}a_{L,0}\) connects these two ground kets.
The Fourier product is \(\zeta^{-1}= -i\), the left removal sign is
\(P_L(0,0)=(-1)^3=-1\), and the right creation sign is
\(P_R(1,-1)=(-1)^2=+1\).  Therefore
\begin{equation}
 \langle\Omega_{(1,-1)}|
    O_{\mathrm{CDW}}(1)|\Omega_{(0,0)}\rangle
 =\frac14(-i)(-1)(+1)=\frac{i}{4}.
 \label{eq:ch20-four-cdw}
\end{equation}

\emph{Second: the SC ground coefficient.}
Choose the source \(N=(1,0)\):
\[
 |\Omega_{(1,0)}\rangle
   =|R\{-1,0,1\};L\{-1,0\}\rangle.
\]
The target charge is \((0,-1)\), with ground
\( |\Omega_{(0,-1)}\rangle
   =|R\{-1,0\};L\{-1\}\rangle\).
Only \(a_{R,1}a_{L,0}\) connects these states.  Its Fourier
character is \(\zeta^{1+0}=i\), and the two successive signs are
\(P_L(1,0)=+1\) and \(P_R(1,-1)=+1\).  It follows that
\begin{equation}
 \langle\Omega_{(0,-1)}|
    O_{\mathrm{SC}}(1)|\Omega_{(1,0)}\rangle
 =\frac{i}{4}.
 \label{eq:ch20-four-sc}
\end{equation}
These two coefficients are exact finite-Fock numbers, not
large-system approximations or formal vertex-state evaluations.

\emph{Third: a discarded intermediate excitation survives the
second field.}
Return to the CDW source \(N=(0,0)\) and select
\(K_s=K_m=0\), \(K_t=1\).  The intermediate charge is
\(N_m=(0,-1)\) and its zero-budget state has occupied orbitals
\(R\{-1,0\},L\{-1\}\).  Consider instead the term in which
\(c_{L,1}\) removes the \emph{left} particle at \(q=-1\):
\[
 |R\{-1,0\};L\{-1,0\}\rangle
 \xrightarrow{\ a_{L,-1}\ }
 |R\{-1,0\};L\{0\}\rangle.
\]
The latter configuration has the correct intermediate \emph{charge}
\((0,-1)\), but its left particle occupies \(0\) rather than \(-1\).
By \eqref{eq:ch20-excitation} it has excitation \(E_{N_m}=1\), and
is therefore killed by \(P_m=P_{(0,-1),0}\).  The right creator can
nevertheless add \(p=1\):
\[
 |R\{-1,0\};L\{0\}\rangle
 \xrightarrow{\ a_{R,1}^{\dagger}\ }
 |R\{-1,0,1\};L\{0\}\rangle.
\]
This final state belongs to \(V_{(1,-1),1}\), having left excitation
one and right excitation zero.  With the stated ordering, both mode
operations have occupation sign \(+1\); their Fourier factors are
\(\zeta^{-1}=-i\) for the left annihilator and
\(\zeta^{-1}=-i\) for the right creator.  No other excluded
intermediate path contributes to this final occupation ket.  Hence
\begin{equation}
 \begin{aligned}
 &\mathcal R_{\mathrm{CDW}}^{smt}(1)
                |\Omega_{(0,0)}\rangle\\
 &\qquad=-\frac14
      |R\{-1,0,1\};L\{0\}\rangle\ne0
       \quad (K_s=K_m=0,\ K_t=1).
 \end{aligned}
 \label{eq:ch20-four-residual}
\end{equation}
In particular, equating the projected product with
\(B_R^{\mathrm{CDW},\dagger}B_L\) in this example would
be \emph{false}, despite the fact that both constituent fields and
all three projections are correctly defined.

If we change only the \emph{target} cutoff to \(K_t=0\), the state
on the right of \eqref{eq:ch20-four-residual} is also discarded.
There is then no surviving excluded path from this source ground to
the target ground, and
\begin{equation}
 \mathcal R_{\mathrm{CDW}}^{smt}(1)
      |\Omega_{(0,0)}\rangle=0
 \qquad(K_s=K_m=K_t=0).
 \label{eq:ch20-four-residual-zero}
\end{equation}
The restricted, ground-to-ground equality here is a property of the
\emph{final projection} and this source ket; it does not turn
\(\mathcal R_{\mathrm{CDW}}^{smt}\) into the zero operator for
arbitrary larger budgets.
\end{example}

The example illustrates two conceptually different uses of a cutoff.
The intermediate budget \(K_m=0\) declares an excited left-hole state
inadmissible \emph{before} the second field acts.  The target budget
\(K_t=1\) still accepts its eventual image.  This mismatch creates the
nonzero remainder; it is not an error in the CAR algebra or the Fourier
normalization.  Conversely, if one is interested only in a particular
ground-to-ground coefficient, orthogonality of the untouched left
occupation can eliminate the excluded paths \emph{after} the full
product.  Neither conclusion justifies a blanket rule that energy-budget
projection commutes with multiplication.

\subsection{What these exact results do and do not provide}
\label{subsec:ch20-scope}

The finite CAR construction now supplies three distinct pieces of data:
a globally defined even word \(O_J(x)\); a collection of typed maps
\(V_s\to V_t^J\) with explicit intermediate remainder
\(\mathcal R_J^{smt}\); and a charge-ground scalar coefficient with
specified Fourier and Klein ordering.  These statements require no
choice of finite interacting vacuum and no canonical-commutation
relations for compressed density modes.

For a normalized \(\psi\in\mathcal F\), an observable expectation such
as \(\langle\psi,O_J(x)^{\dagger}O_J(y)\psi\rangle\) is an ordinary
finite Hilbert-space number.  A density matrix \(\rho\geq0\),
\(\operatorname{Tr}\rho=1\), similarly defines
\(\operatorname{Tr}(\rho O_J(x)^{\dagger}O_J(y))\).  These should
not be identified without a comparison theorem with the abstract
positive functional on the \emph{uncompressed} oscillator algebra
discussed in Chapter~\ref{ch:correlations} and the construction
underlying Section~\ref{sec:ch20-formal-vertices}.  In particular,
a formal vertex correlator evaluated coefficientwise in a power-series
parameter is not automatically the finite-CAR expectation just defined.

A proposed Mattis--Mandelstam description may eventually replace a
single projected field by a finite polynomial or oscillator expression,
but Chapter~14's conditional criterion must first be checked for the
\emph{specified source and target cutoffs}.  The exact product theorem
then retains \(\mathcal R_J^{smt}\) unless an explicit range or
final-projection annihilation argument removes it.  This distinction is
especially important for later pairing, duality and spinful/Umklapp
constructions: a scalar agreement on sector grounds does not certify an
operator identity on all excited budget states.

Finally, the sign and phase conventions fixed here are provisional
cross-chapter conventions: \(R<L\), common positive-exponent finite
annihilation Fourier transform, left field acting first, source charge
phases, and separate cutoffs for each arrow.  Their compatibility with
future chapter integrations must be reviewed against the actual committed
interfaces.  Neither a thermodynamic limit nor continuum scaling,
spectral equivalence or a finite interacting-state identification has
been established in this section.

\section{Constructed Formal Vertices and Coefficientwise Correlators}
\label{sec:ch20-formal-vertices}

Section~\ref{sec:ch20-projected-products} treated density-wave and pairing
operators as \emph{literal products of finite-lattice fermion fields}.  We
could evaluate these products on the full CAR Fock space, and we saw why an
intermediate energy-budget projection introduces an explicit residual.  We
now ask a related but mathematically different question: can we construct a
model in which the oscillator expressions suggested by bosonization are
well-defined, possess an actual positive vacuum functional, and admit
exactly computable two-point functions?

A tempting shortcut would be to replace each fermion field by its
Mattis--Mandelstam expression, multiply exponentials, and invoke Gaussian
averaging.  This is not justified by the preceding finite-CAR calculation:
the projected product may have a nonzero residual, and compression does not
preserve every exponential identity.  Instead we shall \emph{construct a
second carrier}, with uncompressed canonical commutation relations (CCR),
independent charge shifts, and a formal-power-series parameter.  It makes
the algebraic steps behind the familiar vertex correlation formula exact,
while keeping the comparison with a finite interacting fermion state open.
We develop the uncompressed oscillator and charge construction here, giving
each definition and the arguments needed at the point where they first enter
the calculation.

\subsection{An uncompressed oscillator algebra and its positive vacuum}
\label{subsec:ch20b-oscillator}

Fix an integer \(h\geq 1\) and set \(L=2h\).  Spatial positions
\(x,y\) belong to \(\mathbb Z/L\mathbb Z\), and we use the canonical
Fourier character
\begin{equation}
 \zeta=\exp(2\pi i/L), \qquad
 \zeta^{m(x+L)}=\zeta^{mx}.
 \label{eq:ch20b-character}
\end{equation}
Let \(M\geq1\) be a \emph{finite} positive-mode cutoff and write
\(\mathcal I_M=\{(\nu,m):\nu\in\{R,L\},\ 1\leq m\leq M\}\).
The formal construction works for any finite \(M\); its intended
low-momentum comparison with the lattice usually has \(M\leq h\).
We use \(L\) for lattice length even when it appears as a subscript
for the left-moving species, where it is distinguishable by context.

Recall that an algebra is a vector space equipped with an associative
multiplication and a unit.  A \emph{star algebra} additionally has an
antilinear involution \(a\mapsto a^*\) satisfying
\((ab)^*=b^*a^*\).  The \emph{free} star algebra on symbols \(A_{\nu m}\)
and \(C_{\nu m}\) contains all finite linear combinations of their
noncommuting words.  Imposing an equation in this algebra means quotienting
by the two-sided, star-stable ideal generated by its difference: an
\emph{ideal} is a subspace closed under multiplication on either side by
arbitrary ambient words.  In the resulting quotient the imposed
relations hold for every word, not only when applied to a particular state.

\begin{definition}[The uncompressed oscillator algebra]
\label{def:ch20b-ccr}
Let \(\mathfrak A_{\mathrm{osc}}\) be the unital complex star-algebra
of finite words in \(A_{\nu m},C_{\nu m}\), modulo the relations
\begin{align}
 [A_{\nu m},C_{\eta n}]&=m\delta_{\nu\eta}\delta_{mn}I,
 & [A_{\nu m},A_{\eta n}]&=0,
 & [C_{\nu m},C_{\eta n}]&=0,
 \label{eq:ch20b-ccr}\\
 A_{\nu m}^*&=C_{\nu m},& C_{\nu m}^*&=A_{\nu m}.
 \nonumber
\end{align}
Here \([a,b]=ab-ba\).  There is no restriction on word degree or on
oscillator occupation number.
\end{definition}

The word \emph{uncompressed} is essential.  There are only \(2M\)
oscillator labels, but each creation generator may be applied any number
of times, so the representation space below is infinite-dimensional.
The trace of a commutator vanishes on a finite-dimensional space, whereas
the trace of \(mI\) does not: a faithful finite-dimensional realization of
\eqref{eq:ch20b-ccr} is therefore impossible.  This is the familiar
finite-dimensional obstruction from Chapter~\ref{ch:boson_fock}, not an
obstruction to the following polynomial representation.

\begin{proposition}[Concrete positive oscillator vacuum]
\label{prop:ch20b-vacuum}
There is a nonzero representation \(\pi_0\) of
\(\mathfrak A_{\mathrm{osc}}\) on the polynomial space
\(\mathcal P=\mathbb C[X_{\nu m}:(\nu,m)\in\mathcal I_M]\) and a
normalized, complex-linear, positive functional
\(\omega_0:\mathfrak A_{\mathrm{osc}}\to\mathbb C\), defined by
\begin{align}
 \pi_0(C_{\nu m})f&=X_{\nu m}f,
 &\pi_0(A_{\nu m})f&=m\,\frac{\partial f}{\partial X_{\nu m}},
 \label{eq:ch20b-polynomial-action}\\
 \langle X^\alpha,X^\beta\rangle_0
 &=\delta_{\alpha\beta}\prod_{(\nu,m)\in\mathcal I_M}
       \alpha_{\nu m}!\,m^{\alpha_{\nu m}},
 &\omega_0(a)&=\langle1,\pi_0(a)1\rangle_0.
 \label{eq:ch20b-gram}
\end{align}
The inner product is conjugate-linear in its first argument.  In
particular,
\(\omega_0(I)=1\) and \(\omega_0(a^*a)\geq0\) for every
\(a\in\mathfrak A_{\mathrm{osc}}\).
\end{proposition}
\begin{proof}
A multi-index \(\alpha\) has finitely many nonnegative integer
components.  Every weight in \eqref{eq:ch20b-gram} is strictly positive,
so the monomials form an orthogonal basis of a positive-definite
pre-Hilbert space; no completion is required.  Multiplication and
weighted differentiation satisfy
\([m\partial_{X_{\nu m}},X_{\eta n}]
 =m\delta_{\nu\eta}\delta_{mn}I\), and distinct multiplications or
differentiations commute.  Thus the defining commutator relations
annihilate every polynomial and the free-word action descends to the
quotient.  Moreover, if \(\beta=\alpha+e_{\nu m}\),
\begin{align*}
 \langle X_{\nu m}X^\alpha,X^\beta\rangle_0
 &=\prod_j\beta_j!\,m_j^{\beta_j}\\
 &=m(\alpha_{\nu m}+1)
         \prod_j\alpha_j!\,m_j^{\alpha_j}
 =\langle X^\alpha,m\partial_{X_{\nu m}}X^\beta\rangle_0.
\end{align*}
Both sides vanish if \(\beta\neq\alpha+e_{\nu m}\).  Hence creation
and annihilation are mutually adjoint on polynomials.  By reversing
words, \(\pi_0(a^*)\) is the algebraic adjoint of \(\pi_0(a)\).
Consequently
\[
 \omega_0(a^*a)=\langle\pi_0(a)1,\pi_0(a)1\rangle_0\geq0,
 \qquad\omega_0(I)=\langle1,1\rangle_0=1.
\]
In particular the quotient is not the zero algebra.  The vacuum
conditions \(\pi_0(A_{\nu m})1=0\) and
\(\langle1,\pi_0(C_{\nu m})f\rangle_0=0\) will determine its ordered
contractions.
\end{proof}

\begin{remark}[What positivity does and does not assert]
A normalized positive functional is \emph{linear}, not multiplicative.
Indeed a multiplicative scalar-valued map would send every commutator to
zero, contradicting \([A_{\nu m},C_{\nu m}]=mI\) for \(m\geq1\).
This polynomial vacuum is an algebraic state on an unbounded-operator
star algebra; we have not constructed a Weyl \(C^*\)-algebra or a
finite-dimensional fermionic density matrix.
\end{remark}

\subsection{Bogoliubov dressing and the charge-sector carrier}
\label{subsec:ch20b-dressed}

Let \(c,s\in\mathbb R\) satisfy \(c^2-s^2=1\).  When the remaining
coupling data matter, we denote their record by \(P\); in the present
section only \(c,s\) enter the oscillator formula.  Write
\begin{equation}
 g=(c-s)^2>0,\qquad g^{-1}=(c+s)^2,\qquad
 (c-s)(c+s)=1.
 \label{eq:ch20b-g}
\end{equation}
This Luttinger parameter \(g\) is not the Umklapp coupling \(g_U\) of
Chapter~21.  For each fixed positive mode, set
\begin{align}
 \widetilde C_{Rm}&=cC_{Rm}+sA_{Lm},
 &\widetilde C_{Lm}&=cC_{Lm}+sA_{Rm},
 \label{eq:ch20b-bog-creator}\\
 \widetilde A_{Rm}&=cA_{Rm}+sC_{Lm},
 &\widetilde A_{Lm}&=cA_{Lm}+sC_{Rm}.
 \label{eq:ch20b-bog-ann}
\end{align}
All four formulas retain the operator order inherited from the word
algebra.  The tilde refers to a change of generators, not to a
projection onto a smaller finite-energy space.

\begin{proposition}[Dressed CCR and positive state]
\label{prop:ch20b-dressing}
The substitution \(\beta_P(C_{\nu m})=\widetilde C_{\nu m}\),
\(\beta_P(A_{\nu m})=\widetilde A_{\nu m}\) extends to a star
algebra automorphism of \(\mathfrak A_{\mathrm{osc}}\).  Its inverse is
the same construction with \(s\) replaced by \(-s\).  The functional
\begin{equation}
 \omega_P=\omega_0\circ\beta_P^{-1}
 \label{eq:ch20b-omegaP}
\end{equation}
is normalized and positive.  Its dressed vacuum contractions are
\begin{equation}
 \omega_P(\widetilde A_{\nu m}\widetilde C_{\eta n})
 =m\delta_{\nu\eta}\delta_{mn},\quad
 \omega_P(\widetilde C_{\nu m}\widetilde A_{\eta n})=0,
 \quad\omega_P(\widetilde A\widetilde A)
 =\omega_P(\widetilde C\widetilde C)=0.
 \label{eq:ch20b-dressed-contractions}
\end{equation}
The last identities apply to every pair of indicated dressed modes.
\end{proposition}
\begin{proof}
As a representative calculation,
\[
 [\widetilde A_{Rm},\widetilde C_{Rn}]
 =c^2[A_{Rm},C_{Rn}]+s^2[C_{Lm},A_{Ln}]
 =(c^2-s^2)m\delta_{mn}I=m\delta_{mn}I.
\]
For the potentially nonzero cross-chirality annihilator commutator,
\[
 [\widetilde A_{Rm},\widetilde A_{Ln}]
 =cs\bigl([A_{Rm},C_{Rn}]+[C_{Lm},A_{Ln}]\bigr)=0.
\]
The other relations follow by exchanging \(R,L\), taking adjoints,
and using the bare relations.  The substitutions therefore respect the
relation ideal and star operation.  Applying the transformation with
\(-s\) composes the pairs
\((C_R,A_L)\) and \((C_L,A_R)\) by
\(\left(\begin{smallmatrix}c&-s\\-s&c\end{smallmatrix}\right)
  \left(\begin{smallmatrix}c&s\\s&c\end{smallmatrix}\right)=I\).
Thus \(\beta_P\) is invertible.  For every \(a\),
\[
 \omega_P(a^*a)
 =\omega_0\bigl((\beta_P^{-1}a)^*(\beta_P^{-1}a)\bigr)\geq0,
 \qquad\omega_P(I)=1.
\]
Since \(\beta_P^{-1}(\widetilde A)=A\) and
\(\beta_P^{-1}(\widetilde C)=C\), the contractions reduce to
\(\omega_0(AC)=m\) and
\(\omega_0(CA)=\omega_0(AA)=\omega_0(CC)=0\), mode by mode.
\end{proof}

The vertices must also change \emph{charge}.  For this reason the
oscillator algebra alone does not provide their carrier.  Let
\(\mathcal C=\mathbb C^{(\mathbb Z^2)}\) be the vector space of finite
linear combinations of kets \(\delta_N\), with
\(N=(N_R,N_L)\in\mathbb Z^2\) and
\(\langle\delta_N,\delta_{N'}\rangle=\delta_{NN'}\).
It is an algebraic direct sum, not the finite-band charge space of
Section~20.1.  The \(N_\nu\) act diagonally; hence all integer charge
labels are present.  Let \(e_R=(1,0)\) and \(e_L=(0,1)\).  Fix the
species ordering \(R<L\) and the reference half-band offset \(h\).
Define the signed bilateral shifts
\begin{align}
 F_\nu\delta_N&=P_\nu(N)\delta_{N-e_\nu},
 &P_\nu(N)&=(-1)^{\sum_{\eta<\nu}(h+N_\eta)+h+N_\nu-1},
 \label{eq:ch20b-signed-shift}\\
 F_\nu^*\delta_N&=P_\nu(N+e_\nu)\delta_{N+e_\nu},
 &Z_\nu(x)\delta_N&=\zeta^{N_\nu x}\delta_N.
 \label{eq:ch20b-phase-action}
\end{align}
The exponents of \((-1)\) are integers, even for negative charges.
Each \(F_\nu\) is unitary on this algebraic orthonormal basis:
its inverse and adjoint both have the displayed raising action.
Crucially, the left sign depends on the right charge, which encodes
the interspecies fermionic ordering.

\begin{theorem}[Anticommutation of Auxiliary Charge Shifts]
\label{thm:ch20b-shift-anticommute}
On the unbounded auxiliary module $\mathcal C = \mathbb C^{(\mathbb Z^2)}$, the charge shifts
satisfy exact mixed anticommutation and number shift relations:
\begin{equation}
 \{F_R, F_L\} = 0, \qquad
 \{F_R, F_L^*\} = 0, \qquad
 [N_\nu, F_\eta] = -\delta_{\nu\eta} F_\eta.
 \label{eq:ch20b-shift-car}
\end{equation}
\end{theorem}
\begin{proof}
Evaluate both operator orders on a basis vector $\delta_N$:
\begin{align*}
 F_R F_L \delta_N
  &= F_R \bigl(P_L(N) \delta_{N-e_L}\bigr)
   = P_L(N) P_R(N-e_L) \delta_{N-e_R-e_L}, \\
 F_L F_R \delta_N
  &= F_L \bigl(P_R(N) \delta_{N-e_R}\bigr)
   = P_R(N) P_L(N-e_R) \delta_{N-e_R-e_L}.
\end{align*}
By \eqref{eq:ch20b-signed-shift}, $P_R(N) = (-1)^{h+N_R-1}$, which is completely independent of $N_L$;
hence $P_R(N-e_L) = P_R(N)$.
On the other hand, $P_L(N) = (-1)^{(h+N_R) + (h+N_L-1)}$, which depends directly on $N_R$.
Shifting $N_R \to N_R - 1$ flips the parity:
\[
 P_L(N-e_R) = (-1)^{(h+N_R-1) + (h+N_L-1)} = - P_L(N).
\]
Therefore,
\[
 P_R(N) P_L(N-e_R) = - P_R(N) P_L(N) = - P_L(N) P_R(N-e_L),
\]
which proves $F_L F_R \delta_N = - F_R F_L \delta_N$ for all $N$, so $\{F_R, F_L\} = 0$.
The adjoint anticommutator $\{F_R, F_L^*\} = 0$ follows by an identical sign calculation.
Finally, since $F_\eta$ maps $\delta_N$ to $\delta_{N-e_\eta}$,
$[N_\nu, F_\eta]\delta_N = (N_\nu - \delta_{\nu\eta})F_\eta\delta_N - N_\nu F_\eta\delta_N = -\delta_{\nu\eta}F_\eta\delta_N$.
\end{proof}

\begin{theorem}[General Group 2-Cocycle Identity]
\label{thm:ch20b-cocycle-identity}
For arbitrary integer charge vectors $p, q \in \mathbb Z^d$, define the bilinear sign functional
\begin{equation}
 \varepsilon(p, q) := (-1)^{\sum_{1 \le i < j \le d} p_j q_i}.
 \label{eq:ch20b-cocycle-general}
\end{equation}
Then $\varepsilon$ satisfies the 2-cocycle condition:
\begin{equation}
 \varepsilon(p, q) \varepsilon(p+q, t) = \varepsilon(q, t) \varepsilon(p, q+t)
 \qquad (\forall p, q, t \in \mathbb Z^d).
 \label{eq:ch20b-cocycle-identity}
\end{equation}
\end{theorem}
\begin{proof}
Evaluating the exponent of the left-hand side modulo $2$ gives
$\sum_{i<j} [ p_j q_i + (p_j + q_j) t_i ] = \sum_{i<j} [ p_j q_i + p_j t_i + q_j t_i ]$.
The exponent of the right-hand side modulo $2$ is
$\sum_{i<j} [ q_j t_i + p_j (q_i + t_i) ] = \sum_{i<j} [ q_j t_i + p_j q_i + p_j t_i ]$.
These two sums are identical in $\mathbb Z$, proving the cocycle identity in $\mathbb C$.
\end{proof}

Let \(\mathfrak A_{\mathrm{ch}}\) be the \emph{generated} star algebra of
these operators and their finite sums and products, and set
\(\mathfrak A=\mathfrak A_{\mathrm{ch}}\otimes_{\mathbb C}
\mathfrak A_{\mathrm{osc}}\).  Operators from distinct tensor factors
commute.  On the algebraic tensor carrier
\(\mathcal C\otimes\mathcal P\), write
\begin{equation}
 \omega_{P,N}(a)=
 \big\langle\delta_N\otimes1,\,
      (\pi_{\mathrm{ch}}\otimes(\pi_0\circ\beta_P^{-1}))(a)
      (\delta_N\otimes1)\big\rangle.
 \label{eq:ch20b-tensor-state}
\end{equation}
A tensor product of these represented polynomial and finite-support
charge spaces has a positive-definite inner product.  The argument of
Proposition~\ref{prop:ch20b-vacuum} therefore gives
\(\omega_{P,N}(a^*a)\geq0\) and \(\omega_{P,N}(I)=1\).
The state depends on the \emph{source charge} \(N\).  It is not a
finite-CAR ground-state expectation: to make such a claim would require
a separately constructed comparison of states and operators.

\subsection{Hermitian phases and their commutators}
\label{subsec:ch20b-phases}

The opposite Fourier orientations of the two chiralities are important.
At position \(x\in\mathbb Z/L\mathbb Z\) define
\begin{align}
 Q_R(x)&=\sum_{m=1}^M\frac{1}{m}
 \bigl(\zeta^{-mx}\widetilde C_{Rm}
       +\zeta^{mx}\widetilde A_{Rm}\bigr),
 \label{eq:ch20b-QR}\\
 Q_L(x)&=\sum_{m=1}^M\frac{1}{m}
 \bigl(\zeta^{mx}\widetilde C_{Lm}
       +\zeta^{-mx}\widetilde A_{Lm}\bigr),
 \label{eq:ch20b-QL}\\
 X_{\mathrm{CDW}}(P,x)&=(c-s)\bigl(Q_R(x)+Q_L(x)\bigr),
 \label{eq:ch20b-Xcdw}\\
 X_{\mathrm{SC}}(P,x)&=(c+s)\bigl(Q_R(x)-Q_L(x)\bigr).
 \label{eq:ch20b-Xsc}
\end{align}
The factor \(1/m\) is fixed by our normalization
\([\widetilde A_{\nu m},\widetilde C_{\nu m}]=mI\).
The phases are oscillator operators; none of these symbols denotes a
finite-CAR field or its compressed image.

\begin{lemma}[Adjoints and cancellation of equal-channel commutators]
\label{lem:ch20b-hermitian-commuting}
For every \(x,y\),
\begin{equation}
 Q_\nu(x)^*=Q_\nu(x),\qquad
 X_J(P,x)^*=X_J(P,x),\qquad
 [X_J(P,x),X_J(P,y)]=0
 \quad(J\in\{\mathrm{CDW},\mathrm{SC}\}).
 \label{eq:ch20b-xcommute}
\end{equation}
\end{lemma}
\begin{proof}
Since \(\zeta\) has unit modulus,
\((\zeta^{-mx}\widetilde C_{Rm})^*
=\zeta^{mx}\widetilde A_{Rm}\); the other term is its adjoint.
The same applies to the left chirality.  Because \(c\pm s\) are real,
all the \(Q_\nu\) and \(X_J\) are Hermitian.

The CCR reduce the right-chiral commutator to
\begin{align*}
 [Q_R(x),Q_R(y)]
 &=\sum_{m=1}^M\frac1m
   \bigl(\zeta^{m(x-y)}-\zeta^{-m(x-y)}\bigr)I\\
 &=2i\sum_{m=1}^M\frac{\sin(2\pi m(x-y)/L)}m\,I.
\end{align*}
The reversed spatial orientation gives exactly the negative expression
for \([Q_L(x),Q_L(y)]\).  Mixed commutators vanish because distinct
dressed chiralities obey independent CCR.  Expanding a commutator of
either \(Q_R+Q_L\) or \(Q_R-Q_L\) with itself at two positions leaves
the \emph{sum} of the two same-chirality commutators, so the signs cancel.
Multiplication by \((c-s)^2\) or \((c+s)^2\) proves
\eqref{eq:ch20b-xcommute}.  Notice that a single \(Q_R\) does
\emph{not} generally commute with itself at different positions.
\end{proof}

We shall need a finite kernel rather than a limiting logarithm:
\begin{definition}[Finite logarithmic kernel]
\label{def:ch20b-D1}
\begin{equation}
 D_1(x,y)=\sum_{m=1}^{M}\frac{1-\cos(2\pi m(x-y)/L)}{m}.
 \label{eq:ch20b-D1}
\end{equation}
\end{definition}
Every sum here is finite.  Thus \(D_1(x,x)=0\),
\(D_1(x,y)=D_1(y,x)\), and \(D_1(x,y)\geq0\).
Since \(M\geq1\), the \(m=1\) term implies strict positivity if
\(x\not\equiv y\pmod L\).  No continuum limit is implicit in the
term ``logarithmic.''

\begin{proposition}[Exact difference variances]
\label{prop:ch20b-variance}
Put \(\Delta X_J=X_J(P,y)-X_J(P,x)\).  In the positive state
\(\omega_{P,N}\), which factorizes on oscillator and charge words,
\begin{equation}
 \omega_{P,N}(\Delta X_{\mathrm{CDW}})=0,
 \qquad\omega_{P,N}(\Delta X_{\mathrm{CDW}}^2)=4gD_1(x,y),
 \label{eq:ch20b-variance-cdw}
\end{equation}
\begin{equation}
 \omega_{P,N}(\Delta X_{\mathrm{SC}})=0,
 \qquad\omega_{P,N}(\Delta X_{\mathrm{SC}}^2)
       =4g^{-1}D_1(x,y).
 \label{eq:ch20b-variance-sc}
\end{equation}
\end{proposition}
\begin{proof}
Each \(Q_\nu\) is a finite linear combination of single dressed
creation and annihilation operators, so its expectation is zero by
\eqref{eq:ch20b-dressed-contractions}.  For the right chirality the
nonzero ordered contraction is
\[
 \omega_P(Q_R(x)Q_R(y))
 =\sum_{m,n=1}^M\frac{\zeta^{mx}\zeta^{-ny}}{mn}
   \omega_P(\widetilde A_{Rm}\widetilde C_{Rn})
 =\sum_{m=1}^M\frac{\zeta^{m(x-y)}}{m}.
\]
For the left chirality it is the conjugate expression,
\(\sum_m\zeta^{-m(x-y)}/m\).  In either chirality
\begin{align*}
 \omega_P\bigl((Q_\nu(y)-Q_\nu(x))^2\bigr)
 &=\omega_P(Q_\nu(y)^2)+\omega_P(Q_\nu(x)^2)\\
 &\quad-\omega_P(Q_\nu(y)Q_\nu(x))
       -\omega_P(Q_\nu(x)Q_\nu(y))\\
 &=2\sum_{m=1}^M\frac{1-\cos(2\pi m(x-y)/L)}m
 =2D_1(x,y).
\end{align*}
Mixed \(R,L\) two-point contractions vanish by the dressed vacuum
formula.  The variance of either \(\Delta Q_R+\Delta Q_L\) or
\(\Delta Q_R-\Delta Q_L\) is therefore \(4D_1\).
Multiply respectively by \((c-s)^2=g\) and
\((c+s)^2=g^{-1}\).  The charge tensor factor acts as the identity
throughout this calculation.
\end{proof}

\subsection{Charge words, formal series, and Gaussian coefficients}
\label{subsec:ch20b-formal}

It is instructive to separate the pieces of a vertex which change
charge from those which create oscillator excitations.  In the charge
algebra introduce the \emph{ordered} Klein--phase words
\begin{align}
 K_{\mathrm{CDW}}(x)
 &=F_R^*F_L\,\zeta^{(N_L-N_R-1)x},
 \label{eq:ch20b-Kcdw}\\
 K_{\mathrm{SC}}(x)
 &=F_RF_L\,\zeta^{(N_R+N_L)x}.
 \label{eq:ch20b-Ksc}
\end{align}
The rightmost diagonal phase acts on the source charge before either
shift.  Replacing the product by one with the phase on the left changes
its coefficient: these factors need not commute.  The star in
\(F_R^*\) is the adjoint shift, not complex conjugation of a scalar.

\begin{lemma}[Klein products at fixed source charge]
\label{lem:ch20b-klein}
For every \(N=(N_R,N_L)\in\mathbb Z^2\) and \(x,y\in\mathbb Z/L\mathbb Z\),
\begin{align}
 K_{\mathrm{CDW}}(x)^*K_{\mathrm{CDW}}(y)\delta_N
 &=\zeta^{(N_L-N_R-1)(y-x)}\delta_N,
 \label{eq:ch20b-klein-cdw}\\
 K_{\mathrm{SC}}(x)^*K_{\mathrm{SC}}(y)\delta_N
 &=\zeta^{(N_R+N_L)(y-x)}\delta_N.
 \label{eq:ch20b-klein-sc}
\end{align}
In particular \(K_J(x)^*K_J(x)=I\).
\end{lemma}
\begin{proof}
Write \(\chi_{\mathrm{CDW}}(N,x)
=\zeta^{(N_L-N_R-1)x}\) and
\(\chi_{\mathrm{SC}}(N,x)=\zeta^{(N_R+N_L)x}\).
The rightmost factor in \eqref{eq:ch20b-Kcdw} produces
\(\chi_{\mathrm{CDW}}(N,x)\).  Next, \(F_L\) lowers the left
charge and contributes \(P_L(N)\), after which \(F_R^*\) raises the
right charge and contributes \(P_R(N-e_L+e_R)\).  Thus
\[
 K_{\mathrm{CDW}}(x)\delta_N
 =\chi_{\mathrm{CDW}}(N,x)
   P_L(N)P_R(N-e_L+e_R)
   \delta_{N+e_R-e_L}.
\]
For the pairing word, both shifts lower charge, in the given order:
\[
 K_{\mathrm{SC}}(x)\delta_N
 =\chi_{\mathrm{SC}}(N,x)
   P_L(N)P_R(N-e_L)
   \delta_{N-e_R-e_L}.
\]
In each expression the product of the two signed Klein coefficients
is a real number of absolute value one, independent of \(x\).
The shift target is unique, so the adjoint at \(x\) takes that target
back to \(\delta_N\) and multiplies by the complex conjugate of its
coefficient at \(x\).  The two signs square to one, leaving
\(\overline{\chi_J(N,x)}\chi_J(N,y)\), which is precisely the
stated phase.  This computation explicitly preserves the source,
intermediate, and target charge labels.
\end{proof}

A formal exponential requires no analytic convergence.  Introduce a
\emph{new central indeterminate} \(t\) with \(t^*=t\).
The ring \(\mathfrak A[[t]]\) consists of all formal sequences
\(\sum_{n\geq0}a_nt^n\), \(a_n\in\mathfrak A\), with multiplication
by the finite convolution
\begin{equation}
 [t^n]\biggl(\sum_{j\geq0}a_jt^j\biggr)
        \biggl(\sum_{k\geq0}b_kt^k\biggr)
 =\sum_{j=0}^n a_jb_{n-j}.
 \label{eq:ch20b-convolution}
\end{equation}
There is no requirement that the coefficient algebra be commutative.
The notation \([t^n]\) means the coefficient of \(t^n\), not an
operator commutator.  For \(a\in\mathfrak A\) define
\begin{equation}
 \operatorname{Exp}_{\mathrm f}(ta)
 =\sum_{n=0}^\infty\frac{t^na^n}{n!}\in\mathfrak A[[t]],
 \qquad a^0=I.
 \label{eq:ch20b-formal-exp}
\end{equation}
Factorials are invertible \emph{complex scalar} coefficients.
More generally, if \(u\in t\mathfrak A[[t]]\) has zero constant
coefficient, then
\(\operatorname{Exp}_{\mathrm f}(u)=\sum_{k\geq0}u^k/k!\)
is a well-defined formal series: for each degree \(n\), only
\(k\leq n\) can contribute.  In particular,
\(\operatorname{Exp}_{\mathrm f}(\lambda t^2)\) for
\(\lambda\in\mathbb C\) denotes
\(\sum_{k\geq0}\lambda^kt^{2k}/k!\).
The positive functional \(\omega_{P,N}\) extends coefficientwise to a
\(\mathbb C[[t]]\)-linear evaluator
\begin{equation}
 \widehat\omega_{P,N}\biggl(\sum_{n\geq0}a_nt^n\biggr)
 =\sum_{n\geq0}\omega_{P,N}(a_n)t^n.
 \label{eq:ch20b-omega-hat}
\end{equation}
Although each original \(\omega_{P,N}\) is positive, this map is not
an ordinary scalar-valued state on arbitrary formal series, for which
there is no assumed coefficientwise nonnegative order.  Above all,
\(t\) is \emph{formal}: setting \(t=1\) is not defined by
\eqref{eq:ch20b-convolution}--\eqref{eq:ch20b-formal-exp}.

\begin{lemma}[Formal exponential product, with its actual hypothesis]
\label{lem:ch20b-exp-product}
If \(a,b\) belong to a unital complex algebra and \(ab=ba\), then
in its noncommutative-coefficient series ring
\begin{equation}
 \operatorname{Exp}_{\mathrm f}(ta)
 \operatorname{Exp}_{\mathrm f}(tb)
 =\operatorname{Exp}_{\mathrm f}(t(a+b)).
 \label{eq:ch20b-formal-binomial}
\end{equation}
In particular
\(\operatorname{Exp}_{\mathrm f}(-ta)
 \operatorname{Exp}_{\mathrm f}(ta)=I\).
No such factorization is claimed for noncommuting \(a,b\).
\end{lemma}
\begin{proof}
At degree \(n\), the product on the left has coefficient
\[
 \sum_{k=0}^n\frac{a^kb^{n-k}}{k!(n-k)!}
 =\frac1{n!}\sum_{k=0}^n\binom nk a^kb^{n-k}
 =\frac{(a+b)^n}{n!}.
\]
The last equality is the binomial formula, valid because \(ab=ba\).
Every coefficient is a finite sum, so no limit or rearrangement of
infinite analytic sums is involved.  Taking \(b=-a\) gives the inverse.
\end{proof}

The final ingredient is a Gaussian moment identity.  We prove it
rather than invoking a probability distribution on the infinite
polynomial space.
\begin{lemma}[Vacuum Wick recurrence and formal Gaussian identity]
\label{lem:ch20b-wick}
Let \(Y\) be a Hermitian \emph{linear combination} of dressed
creation and annihilation generators, with no constant term.  Put
\(V=\omega_P(Y^2)\).  Then \(V\geq0\),
\begin{align}
 \omega_P(Y^{2k+1})&=0,
 &\omega_P(Y^{2k})&=\frac{(2k)!}{2^k k!}\,V^k,
 \qquad k\geq0,
 \label{eq:ch20b-gauss-moments}\\
 \widehat\omega_P\bigl(\operatorname{Exp}_{\mathrm f}(itY)\bigr)
 &=\operatorname{Exp}_{\mathrm f}(-Vt^2/2)
 \quad\text{in }\mathbb C[[t]].
 \label{eq:ch20b-gauss-exp}
\end{align}
For \(k=0\), the even moment is understood to equal \(\omega_P(I)=1\).
\end{lemma}
\begin{proof}
Pull back by \(\beta_P^{-1}\) and denote the resulting operator
again by \(Y\).  Hermiticity allows the decomposition
\(Y=C_f+A_{\overline f}\), where
\(C_f=\sum_j f_jC_j\) and
\(A_{\overline f}=\sum_j\overline{f_j}A_j\), for finitely many
oscillator labels \(j=(\nu,m)\).  Consequently
\([A_{\overline f},Y]=V I\), where
\begin{equation}
 V=\sum_j m_j|f_j|^2
   =\omega_0(Y^2)\geq0.
 \label{eq:ch20b-GaussianV}
\end{equation}
The equality follows because only the ordered \(AC\) term survives
between vacuum vectors.  In the polynomial representation,
\(A_{\overline f}1=0\) and
\(\langle1,C_fq\rangle_0=0\) for all polynomials \(q\).
Since the commutator \([A_{\overline f},Y]\) is central, induction on
\(n\) gives \([A_{\overline f},Y^n]=nVY^{n-1}\).  Hence for
\(n\geq1\)
\begin{align*}
 \omega_0(Y^{n+1})
 &=\langle1,(C_f+A_{\overline f})Y^n1\rangle_0
 =\langle1,A_{\overline f}Y^n1\rangle_0\\
 &=\langle1,[A_{\overline f},Y^n]1\rangle_0
 =nV\,\omega_0(Y^{n-1}).
\end{align*}
Together with \(\omega_0(Y)=0\) and \(\omega_0(I)=1\), this
recurrence yields the claimed odd and even moments.  The pullback
\(\omega_P=\omega_0\circ\beta_P^{-1}\) restores the dressed formulation.
Finally, the coefficient of \(t^{2k}\) on the left of
\eqref{eq:ch20b-gauss-exp} is
\[
 \frac{i^{2k}}{(2k)!}\omega_P(Y^{2k})
 =\frac{(-V/2)^k}{k!},
\]
while all odd coefficients vanish.  These are exactly the coefficients
of the right-hand scalar formal exponential.
\end{proof}

\subsection{The formal vertices and their two-point functions}
\label{subsec:ch20b-correlators}

\begin{definition}[Weyl-type formal CDW and pairing vertices]
\label{def:ch20b-vertices}
For \(J\in\{\mathrm{CDW},\mathrm{SC}\}\), with \(K_J\) given in
\eqref{eq:ch20b-Kcdw}--\eqref{eq:ch20b-Ksc} and \(X_J\) given in
\eqref{eq:ch20b-Xcdw}--\eqref{eq:ch20b-Xsc}, define the formal series
\begin{align}
 W_{\mathrm{CDW}}(P,t,x)
 &=\frac1L\,F_R^*F_L\zeta^{(N_L-N_R-1)x}
    \operatorname{Exp}_{\mathrm f}(itX_{\mathrm{CDW}}(P,x)),
 \label{eq:ch20b-Wcdw}\\
 W_{\mathrm{SC}}(P,t,x)
 &=\frac1L\,F_RF_L\zeta^{(N_R+N_L)x}
    \operatorname{Exp}_{\mathrm f}(itX_{\mathrm{SC}}(P,x)).
 \label{eq:ch20b-Wsc}
\end{align}
These are \emph{Weyl-type formal vertices}, not normal-ordered
Mattis--Mandelstam fields.  The oscillator series commute with the
charge-word factors because they live in independent tensor factors.
\end{definition}

The choice not to normal-order is substantive.  Normal ordering of
creation and annihilation parts changes vacuum contractions and the
prefactor needed to identify a particular physical observable.  Neither
normal ordering nor exponentiation of \emph{compressed} field operators
can be inserted here without a new theorem.

\begin{theorem}[Coefficientwise formal vertex correlators; draft]
\label{thm:ch20b-correlators}
Assume \(L=2h\) with \(h\geq1\), \(M\geq1\) finite,
\(N\in\mathbb Z^2\), and real \(c,s\) satisfying \(c^2-s^2=1\).
On the constructed algebraic oscillator/charge carrier, the
coefficientwise evaluator satisfies
\begin{align}
 &\widehat\omega_{P,N}\!\left(
    W_{\mathrm{CDW}}(P,t,x)^*W_{\mathrm{CDW}}(P,t,y)\right)
 \nonumber\\[-2pt]
 &\hspace{1cm}=
 \frac{\zeta^{(N_L-N_R-1)(y-x)}}{L^2}
 \operatorname{Exp}_{\mathrm f}\!\left(-2gD_1(x,y)t^2\right),
 \label{eq:ch20b-cdw-corr}\\[4pt]
 &\widehat\omega_{P,N}\!\left(
    W_{\mathrm{SC}}(P,t,x)^*W_{\mathrm{SC}}(P,t,y)\right)
 \nonumber\\[-2pt]
 &\hspace{1cm}=
 \frac{\zeta^{(N_R+N_L)(y-x)}}{L^2}
 \operatorname{Exp}_{\mathrm f}\!\left(-2g^{-1}D_1(x,y)t^2\right).
 \label{eq:ch20b-sc-corr}
\end{align}
The exponentials on the right are scalar \emph{formal} series in
\(\mathbb C[[t]]\).  Equality means equality of every coefficient;
it does not involve an analytic value at \(t=1\).
\end{theorem}
\begin{proof}
We give each algebraic step, as the ordering is important.  The
complex conjugate of \(i\) is \(-i\), \(t^*=t\), and the \(X_J\)
are Hermitian.  Therefore
\begin{equation}
 \bigl(\operatorname{Exp}_{\mathrm f}(itX_J(P,x))\bigr)^*
 =\operatorname{Exp}_{\mathrm f}(-itX_J(P,x))
 \label{eq:ch20b-star-exp}
\end{equation}
coefficient by coefficient.  Adjunction reverses product order.  Since
every charge factor commutes with every oscillator factor,
\begin{align*}
 W_J(P,t,x)^*W_J(P,t,y)
 &=L^{-2}\operatorname{Exp}_{\mathrm f}(-itX_J(P,x))
      K_J(x)^*K_J(y)
      \operatorname{Exp}_{\mathrm f}(itX_J(P,y))\\
 &=L^{-2}K_J(x)^*K_J(y)
      \operatorname{Exp}_{\mathrm f}(-itX_J(P,x))
      \operatorname{Exp}_{\mathrm f}(itX_J(P,y)).
\end{align*}
Here we have \emph{not} commuted two arbitrary charge factors past
one another.  We moved only a tensor-factor charge product past an
oscillator product.

Lemma~\ref{lem:ch20b-hermitian-commuting} establishes
\([X_J(P,x),X_J(P,y)]=0\), so the
coefficientwise binomial identity
\eqref{eq:ch20b-formal-binomial} gives the \emph{exact} series identity
\begin{equation}
 W_J(P,t,x)^*W_J(P,t,y)
 =L^{-2}K_J(x)^*K_J(y)
   \operatorname{Exp}_{\mathrm f}(it\Delta X_J).
 \label{eq:ch20b-product-series}
\end{equation}
Apply \(\widehat\omega_{P,N}\) coefficient by coefficient.
Because the state is a tensor-product vector functional and the
charge operator is independent of the oscillator operators, every
coefficient factorizes into a charge expectation times an oscillator
expectation.  By Lemma~\ref{lem:ch20b-klein}, the charge factor equals
\(\zeta^{\alpha_J(N)(y-x)}\), with
\[
 \alpha_{\mathrm{CDW}}(N)=N_L-N_R-1,
 \qquad\alpha_{\mathrm{SC}}(N)=N_R+N_L.
\]
The Gaussian lemma applies to the Hermitian \emph{linear} oscillator
combination \(\Delta X_J\).  By
Proposition~\ref{prop:ch20b-variance}, its variance is respectively
\(V_{\mathrm{CDW}}=4gD_1(x,y)\) or
\(V_{\mathrm{SC}}=4g^{-1}D_1(x,y)\).  Thus the oscillator factor is
\(\operatorname{Exp}_{\mathrm f}(-V_Jt^2/2)\).  Substitution of the
two values gives \eqref{eq:ch20b-cdw-corr} and
\eqref{eq:ch20b-sc-corr}, including the factor \(L^{-2}\).
Every identity was established within a specified algebra and every
infinite-looking equality was checked coefficientwise.
\end{proof}

Two immediate consequences make the normalization explicit.
At \(x=y\), \(D_1(x,x)=0\) and the charge phase is one.  Thus
\begin{equation}
 \widehat\omega_{P,N}(W_J(P,t,x)^*W_J(P,t,x))=L^{-2}
 \quad\text{as a constant formal series.}
 \label{eq:ch20b-coincident}
\end{equation}
In particular, if one writes
\(\widehat\omega(W_J(x)^*W_J(y))=
  L^{-2}\chi_J(N;y-x)\sum_{k\geq0}C_k^Jt^{2k}\),
then \(C_0^{\mathrm{CDW}}=C_0^{\mathrm{SC}}=1\).  For any \(x,y\)
the remaining coefficients can be read directly:
\begin{equation}
 C_k^{\mathrm{CDW}}=\frac{(-2gD_1(x,y))^k}{k!},\qquad
 C_k^{\mathrm{SC}}=\frac{(-2g^{-1}D_1(x,y))^k}{k!}.
 \label{eq:ch20b-coefficients}
\end{equation}
They are complex-scalar coefficients of a formal series, not
measurements of a lattice decay exponent.

\subsection{Worked four-site example and a finite-CAR counterexample}
\label{subsec:ch20b-example}

Take \(L=4\), \(h=2\), \(M=1\), and
\(\zeta=i\), with positions \(x=0\), \(y=1\).
In the free Bogoliubov case \(c=1,s=0\), one has \(g=g^{-1}=1\).
For the source charge \(N=(0,0)\),
\[
 D_1(0,1)=1-\cos(\pi/2)=1,
 \quad\alpha_{\mathrm{CDW}}(0,0)=-1,
 \quad\alpha_{\mathrm{SC}}(0,0)=0.
\]
The theorem therefore specializes to
\begin{align}
 \widehat\omega_{P,(0,0)}(W_{\mathrm{CDW}}(t,0)^*
                                  W_{\mathrm{CDW}}(t,1))
 &=-\frac{i}{16}\operatorname{Exp}_{\mathrm f}(-2t^2),
 \label{eq:ch20b-example-cdw}\\
 \widehat\omega_{P,(0,0)}(W_{\mathrm{SC}}(t,0)^*
                                  W_{\mathrm{SC}}(t,1))
 &=\frac1{16}\operatorname{Exp}_{\mathrm f}(-2t^2).
 \label{eq:ch20b-example-sc}
\end{align}
To see the first nonconstant coefficient without trusting the closed
formula, compute \(V_J=4\) directly from the two chiral contributions.
The \(t^2\) coefficient of the oscillator expectation is
\[
 \frac{i^2}{2!}\,\omega_P((\Delta X_J)^2)
 =-\frac42=-2,
\]
as predicted.  Higher even coefficients are
\(2,-4/3,\ldots\), since
\begin{equation}
 \operatorname{Exp}_{\mathrm f}(-2t^2)
 =1-2t^2+2t^4-\frac43t^6+\cdots.
 \label{eq:ch20b-example-series}
\end{equation}
The only distinction between the two answers in this example is the
fixed, ordered Klein/charge phase.  The oscillator kernel is the same
because \(g=1\).

\begin{example}[Equal-point finite-CAR normalization is different]
\label{ex:ch20b-finite-counterexample}
On the \(L=4\) CAR Fock space of Section~20.1, choose a normalized
Slater sea \(\lvert\Omega\rangle\) with exactly two occupied momentum
orbitals in each of the independent \(R\) and \(L\) species.  The
normalized Fourier fields obey
\(\{c_{\nu,x},c_{\eta,y}^\dagger\}=
 \delta_{\nu\eta}\delta_{xy}\).
Translation invariance of a filled set of plane-wave orbitals gives
\[
 \langle\Omega\lvert n_{R,x}\rvert\Omega\rangle
 =\langle\Omega\lvert n_{L,x}\rvert\Omega\rangle=\frac24=\frac12,
 \qquad n_{\nu,x}=c_{\nu,x}^\dagger c_{\nu,x}.
\]
The even number operators on distinct species factorize in this
product Slater state.  Using the exact CAR words
\(O_{\mathrm{CDW}}=c_{R,x}^\dagger c_{L,x}\) and
\(O_{\mathrm{SC}}=c_{R,x}c_{L,x}\), we obtain
\begin{align*}
 \langle\Omega\lvert O_{\mathrm{CDW}}(x)^\dagger
               O_{\mathrm{CDW}}(x)\rvert\Omega\rangle
 &=\langle n_{L,x}(1-n_{R,x})\rangle_\Omega=\tfrac14,\\
 \langle\Omega\lvert O_{\mathrm{SC}}(x)^\dagger
               O_{\mathrm{SC}}(x)\rvert\Omega\rangle
 &=\langle n_{R,x}n_{L,x}\rangle_\Omega=\tfrac14.
\end{align*}
But \eqref{eq:ch20b-coincident} gives \(L^{-2}=1/16\) for either
formal Weyl-type vertex.  Therefore \emph{even at coincident points and
in this elementary free example}, these formally defined vertices are
not equal, as normalized observables and states, to the finite-CAR
bilinears.  The discrepancy is neither a numerical error nor a failure
of the theorem: the two sides were defined on different carriers and
with different operator normalization prescriptions.  The value
\(1/4\) is an example-specific finite-state calculation, not a
universal CAR identity.
\end{example}

\subsection{Vertex exchange, zero modes, and finite-cutoff approximations}
\label{subsec:ch20b-vertex-exchange}

We now examine how fermionic statistics are reconstructed from bosonic vertices.
Consider bare normal-ordered oscillator vertices $V_R(x) := :\!e^{i X_R(x)}\!:$ and
$V_L(y) := :\!e^{i X_L(y)}\!:$.
Because $[b_{R,n}, b_{L,m}^\dagger] = 0$, the two oscillator families commute.
By the normal-ordering multiplication law \eqref{eq:ch19-compression-remainder},
\[
 V_R(x) V_L(y) = :\!e^{i X_R(x) + i X_L(y)}\!: = V_L(y) V_R(x).
\]
Therefore, naive oscillator vertices are strictly \emph{commuting}:
\begin{equation}
 [V_R(x), V_L(y)] = 0, \qquad \{V_R(x), V_L(y)\} = 2 V_R(x) V_L(y) \ne 0.
 \label{eq:ch20b-naive-commuting}
\end{equation}
Bogoliubov dressing does not alter this: by \eqref{eq:ch19-cross-chiral-comm-zero},
$[A_R(x), A_L(y)] = 0$ as well.

\begin{theorem}[Mixed-Species Fermionic Anticommutation]
\label{thm:ch20b-mixed-car}
Let $\Psi_r(x) := F_r V_r(x)$ be the Klein-dressed vertex operators on
$\mathcal C \otimes \mathcal P$.  Then the mixed-species fields satisfy exact CAR:
\begin{equation}
 \{\Psi_R(x), \Psi_L(y)\} = 0, \qquad \{\Psi_R(x), \Psi_L^*(y)\} = 0.
 \label{eq:ch20b-mixed-car}
\end{equation}
\end{theorem}
\begin{proof}
Because oscillator modes and charge shifts live in independent tensor factors,
$F_R$ commutes with $V_L(y)$, and $F_L$ commutes with $V_R(x)$.  Using
$[V_R(x), V_L(y)] = 0$ and Theorem~\ref{thm:ch20b-shift-anticommute}:
\begin{align*}
 \Psi_R(x) \Psi_L(y) + \Psi_L(y) \Psi_R(x)
  &= F_R V_R(x) F_L V_L(y) + F_L V_L(y) F_R V_R(x) \\
  &= (F_R F_L + F_L F_R) V_R(x) V_L(y) = 0.
\end{align*}
The adjoint relation $\{\Psi_R(x), \Psi_L^*(y)\} = 0$ follows identically from $\{F_R, F_L^*\} = 0$.
\end{proof}

Next, consider \textbf{same-chirality exchange}.
For $x \ne y$, let $t = 2\pi(x-y)/L$ with $0 < t < 2\pi$.
Using the central commutator \eqref{eq:ch19-field-commutator}, Baker--Campbell--Hausdorff
multiplication yields
\begin{equation}
 V_R(x) V_R(y) = \exp\left( -2i \sum_{n=1}^M \frac{e^{-a q_n}}{n} \sin(q_n(x-y)) \right) V_R(y) V_R(x).
 \label{eq:ch20b-finite-osc-phase}
\end{equation}
In the continuum limit ($M \to \infty, a \to 0^+$), we use the Fourier sawtooth series:
\begin{equation}
 \sum_{n=1}^\infty \frac{\sin(n t)}{n} = \frac{\pi - t}{2} \qquad (0 < t < 2\pi).
 \label{eq:ch20b-sawtooth}
\end{equation}
The continuum oscillator exchange phase is therefore
\[
 \exp\left( -2i \frac{\pi - t}{2} \right) = e^{-i(\pi - t)} = -e^{it}.
\]
Notice that the oscillator phase $-e^{it}$ alone is \emph{not} equal to $-1$: it contains
an extra coordinate-dependent factor $e^{it}$.

The remaining factor is canceled by the zero-mode charge shift.
On the charge lattice $\mathcal C$, define the position-dependent zero-mode operator
\begin{equation}
 Z_R(x) := F_R \exp\left( i \frac{2\pi}{L}(N_R + \nu_R) x \right), \qquad \nu_R \in \{0, \tfrac{1}{2}\}.
 \label{eq:ch20b-zero-mode-field}
\end{equation}
Because $[N_R, F_R] = -F_R$, moving $F_R$ past $\exp(i \frac{2\pi}{L} N_R y)$ shifts $N_R \to N_R - 1$, yielding
\begin{equation}
 Z_R(x) Z_R(y) = e^{-it} Z_R(y) Z_R(x), \qquad t = \frac{2\pi(x-y)}{L}.
 \label{eq:ch20b-zero-mode-exchange}
\end{equation}
Combining the zero-mode and continuum oscillator factors for the full field $\Psi_R(x) = Z_R(x) V_R(x)$:
\begin{equation}
 \Psi_R(x) \Psi_R(y) = (-e^{it})(e^{-it}) \Psi_R(y) \Psi_R(x) = -\Psi_R(y) \Psi_R(x) \qquad (x \ne y \pmod L).
 \label{eq:ch20b-separated-fermion-exchange}
\end{equation}
The spin structure offset $\nu_R$ selects periodic ($\nu_R = 1/2$) or antiperiodic ($\nu_R = 0$) boundary conditions.

\begin{note}[Finite Mode Cutoff Warning]
\label{warn:ch20b-finite-cutoff-exchange}
Identity \eqref{eq:ch20b-separated-fermion-exchange} holds \textbf{only in the continuum limit}
$M \to \infty, a \to 0^+$.  At finite mode cutoff $M < \infty$, the finite Fourier sum
$\sum_{n=1}^M \frac{\sin(nt)}{n}$ differs from the sawtooth line $(\pi-t)/2$ by the Gibbs
oscillation residual:
\[
 \Delta_M(t) := \frac{\pi - t}{2} - \sum_{n=1}^M \frac{\sin(nt)}{n} \ne 0.
\]
Consequently, at finite $M$, the exchange of two right-moving vertex operators acquires
the non-fermionic phase factor $-e^{2i\Delta_M(t)} \ne -1$.
A finite-mode oscillator vertex is therefore \emph{an approximation to a fermion field},
not an exact fermionic operator.  Exact CAR on the finite lattice is provided exclusively
by the discrete Fourier CAR of Theorem~\ref{thm:ch20-discrete-fourier-car} and the
Jordan--Wigner occupation cocycle of Theorem~\ref{thm:ch20-constructive-car}.
\end{note}

\subsection{Scope and remaining comparison problem}
\label{subsec:ch20b-scope}

The results of this section are \emph{exact mathematical identities in
the constructed algebraic model}, subject to the displayed finite-mode,
real-parameter, charge-ordering, and state hypotheses.  Positivity of
\(\omega_{P,N}\) is a property of the uncompressed polynomial/charge
representation.  The series evaluator \(\widehat\omega_{P,N}\) is an
algebraic coefficientwise extension; it was never asserted to be a
convergent function of \(t\).  In particular there is no evaluation at
\(t=1\), no thermodynamic limit, and no power-law asymptotics in these
proofs.

Equally important, these formulas do not remove the projection
residuals of Section~20.1 or identify a compressed interacting finite-
CAR ground state with \(\omega_{P,N}\).  Such an identification would
require a precise operator dictionary, a comparison of the actual
states, control of cutoff and projection errors, and, if one wants an
analytic vertex, a topology and a convergence theorem.  The same
qualification applies to the duality construction proposed in
Section~\ref{sec:ch20-model-duality}: a star-algebra substitution may
transport this \emph{family} of formal states without being a unitary
symmetry of every finite-charge fermionic sector.

Finally, the relative normalization of raw charge/spin currents in a
later multicomponent theory cannot be changed silently by inserting a
\(1/\sqrt2\); and any use of the Fourier or Klein convention must
retain the stated orientation and source action.  Both conventions are
explicit here to make a future integration audit possible.  The formal-vertex
derivations in this section are included manuscript drafts awaiting independent
review and potential Lean formalization.  Section~20.1 has a separately accepted
finite-CAR theorem package in frozen Core; that acceptance does not extend to
the formal vertices or their comparison with an interacting finite-CAR state.

\section{Explicit Algebraic Model Duality}
\label{sec:ch20-model-duality}

\subsection{Why the order-parameter exchange needs a map of algebras}
\label{subsec:ch20c-motivation}

In Section~\ref{sec:ch20-formal-vertices}, we defined the charge-density-wave (CDW) and superconducting (SC) vertices as formal series within an uncompressed oscillator--charge algebra. Their evaluated oscillator correlators depend on the reciprocal interaction parameters $g$ and $g^{-1}$, suggesting a duality between the density-wave and pairing channels. However, the similarity between two resulting formulas does not constitute a proof of duality. To establish duality, we construct a single, well-defined algebraic map that consistently transforms all oscillator generators, charge labels, fermionic Klein signs, and ordered position-dependent phases. This explicit map allows us to transport Hamiltonians and states between the two models.

The charge sectors reveal a structural hurdle before analyzing the oscillators. The CDW Klein word shifts the source charge by $e_R-e_L=(1,-1)$. In contrast, our SC word is an annihilation pair, shifting the charge by $-e_R-e_L=(-1,-1)$. A simple reflection of the left charge maps $(1,-1)$ to $(1,1)$, which is the charge shift of the adjoint operator $W_{\mathrm{SC}}^*$, not $W_{\mathrm{SC}}$ itself. Consequently, any algebraic exchange must incorporate an adjoint operation. We will also see that resolving the remaining source-phase mismatch requires a factor of $\zeta^x$. This complexity shows why simply swapping $g$ and $g^{-1}$ is insufficient.

We consider an even finite lattice length $L=2h$ ($h\geq1$), an integer position $x$ modulo $L$, the canonical phase $\zeta=e^{2\pi i/L}$, and a finite oscillator-mode cutoff $1\leq m\leq M$ ($M\geq1$). The left and right chiralities remain ordered as $R<L$. We briefly recall the definitions required to make this duality explicit, relying on the positive polynomial vacuum and formal vertices from Section~\ref{sec:ch20-formal-vertices}.

An \emph{algebra automorphism} is a bijective linear map preserving
products and the unit.  A \emph{star automorphism} also satisfies
$D(a^*)=D(a)^*$.  In particular it does not reverse products:
$D(ab)=D(a)D(b)$.  We shall construct a map of this kind rather
than infer one from its proposed action on a few expectation values.

\subsection{The parameter and oscillator involutions}
\label{subsec:ch20c-oscillators}

Write the real interaction parameters as
\begin{equation}
 P=(u,c,s,g,g^{-1}),\qquad
 c^2-s^2=1,\qquad g=(c-s)^2>0,\qquad
 g^{-1}=(c+s)^2,
 \label{eq:ch20c-params}
\end{equation}
with $u>0$.  In this notation $g^{-1}$ is the reciprocal of $g$,
not an independent coefficient.  If the Hamiltonian is parametrized
by $v_1,v_2$, we use
\begin{equation}
 v_1=u(c^2+s^2),\qquad v_2=2ucs.
 \label{eq:ch20c-velocities}
\end{equation}
The symbol $g$ should not be confused with the Umklapp coefficient
$g_U$ of Chapter~21 or with an energy-budget bound.

\begin{definition}[Interaction-family involution]
\label{def:ch20c-param}
Define
\begin{equation}
 P^*=(u,c,-s,g^{-1},g).
 \label{eq:ch20c-Pstar}
\end{equation}
The star on $P^*$ designates a \emph{parameter involution}; it is
not complex conjugation or an adjoint of a Hamiltonian.
\end{definition}

\begin{lemma}[Parameter exchange]
\label{lem:ch20c-param}
The map $P\mapsto P^*$ preserves $u$ and $v_1$, reverses $v_2$,
exchanges $g$ with $g^{-1}$, and satisfies $(P^*)^*=P$.
\end{lemma}
\begin{proof}
Since $c^2-(-s)^2=1$, the transformed record is admissible.
Furthermore
\[
 (c-(-s))^2=(c+s)^2=g^{-1},\qquad
 (c+(-s))^2=(c-s)^2=g.
\]
In~\eqref{eq:ch20c-velocities}, $s$ enters $v_1$ only through
$s^2$ and enters $v_2$ linearly.  Changing its sign therefore gives
$v_1^*=v_1$ and $v_2^*=-v_2$.  Changing the sign twice returns the
initial record.  The reciprocal identity also follows directly from
$(c-s)(c+s)=c^2-s^2=1$.
\end{proof}

Let $\mathfrak A_{\mathrm{osc}}$ be the unital complex star-algebra
generated by $A_{\nu m},C_{\nu m}$ ($\nu\in\{R,L\}$,
$1\leq m\leq M$), with
\begin{align}
 [A_{\nu m},C_{\eta n}]&=m\delta_{\nu\eta}\delta_{mn}I,
 & [A_{\nu m},A_{\eta n}]&=0,
 & [C_{\nu m},C_{\eta n}]&=0,\nonumber\\
 A_{\nu m}^*&=C_{\nu m},& C_{\nu m}^*&=A_{\nu m}.
 \label{eq:ch20c-ccr}
\end{align}
It is a quotient of a free word algebra by the two-sided star-stable
ideal generated by these relations, not a finite-dimensional matrix
algebra.  Its familiar polynomial representation has $C_{\nu m}$
acting by multiplication by $X_{\nu m}$ and $A_{\nu m}$ by
$m\partial/\partial X_{\nu m}$.

\begin{proposition}[Oscillator sign automorphism]
\label{prop:ch20c-Dosc}
There is a unique complex-linear star automorphism
$D_{\mathrm{osc}}$ of $\mathfrak A_{\mathrm{osc}}$ specified by
\begin{align}
 D_{\mathrm{osc}}(A_{Rm})&=-A_{Rm},&
 D_{\mathrm{osc}}(C_{Rm})&=-C_{Rm},\nonumber\\
 D_{\mathrm{osc}}(A_{Lm})&= A_{Lm},&
 D_{\mathrm{osc}}(C_{Lm})&= C_{Lm}.
 \label{eq:ch20c-Dosc}
\end{align}
It is involutive: $D_{\mathrm{osc}}^2=\mathrm{id}$.
\end{proposition}
\begin{proof}
Define the substitution on arbitrary finite words multiplicatively
and on finite linear combinations linearly.  For the $R$ generators,
\[
 [-A_{Rm},-C_{Rn}]=[A_{Rm},C_{Rn}]
 =m\delta_{mn}I,
\]
and for $L$ nothing changes.  Mixed-chirality commutators are zero
before and after the substitution; the vanishing $AA$ and $CC$
commutators are likewise preserved.  Because both members of each
adjoint pair acquire the same real sign,
$D_{\mathrm{osc}}(A_{\nu m}^*)
 =D_{\mathrm{osc}}(A_{\nu m})^*$.
Thus every defining relation maps into the relation ideal and the
substitution descends to the quotient.  Repeating it changes each
$R$ generator by two minus signs and fixes every $L$ generator,
so the induced map is its own inverse.
\end{proof}

In the notation of Section~\ref{sec:ch20-formal-vertices}, define the
Bogoliubov automorphism $\beta_P$ by
\begin{align}
 \widetilde C_{Rm}(P)&=cC_{Rm}+sA_{Lm},&
 \widetilde A_{Rm}(P)&=cA_{Rm}+sC_{Lm},\nonumber\\
 \widetilde C_{Lm}(P)&=cC_{Lm}+sA_{Rm},&
 \widetilde A_{Lm}(P)&=cA_{Lm}+sC_{Rm}.
 \label{eq:ch20c-dressed}
\end{align}
All these equalities are inside the uncompressed algebra.  They
must not be read as identities of compressed finite-CAR matrices.

\begin{lemma}[Intertwining with Bogoliubov dressing]
\label{lem:ch20c-intertwine}
On the entire oscillator star algebra,
\begin{equation}
 D_{\mathrm{osc}}\circ\beta_P
 =\beta_{P^*}\circ D_{\mathrm{osc}}.
 \label{eq:ch20c-intertwine}
\end{equation}
In particular,
\begin{align}
 D_{\mathrm{osc}}\widetilde C_{Rm}(P)
   &=-\widetilde C_{Rm}(P^*),&
 D_{\mathrm{osc}}\widetilde A_{Rm}(P)
   &=-\widetilde A_{Rm}(P^*),\nonumber\\
 D_{\mathrm{osc}}\widetilde C_{Lm}(P)
   &=\widetilde C_{Lm}(P^*),&
 D_{\mathrm{osc}}\widetilde A_{Lm}(P)
   &=\widetilde A_{Lm}(P^*).
 \label{eq:ch20c-dressed-transform}
\end{align}
\end{lemma}
\begin{proof}
For the right creation generator,
\[
 D_{\mathrm{osc}}(cC_{Rm}+sA_{Lm})
   =-cC_{Rm}+sA_{Lm}
   =-(cC_{Rm}-sA_{Lm}),
\]
which is $\beta_{P^*}(-C_{Rm})$.  For the left creation generator,
\[
 D_{\mathrm{osc}}(cC_{Lm}+sA_{Rm})
   =cC_{Lm}-sA_{Rm}=\beta_{P^*}(C_{Lm}).
\]
The annihilation formulas follow by taking adjoints.  Since the two
sides of~\eqref{eq:ch20c-intertwine} are algebra homomorphisms,
agreement on generators proves agreement on every finite word and
hence on the quotient.
\end{proof}

Define the Hermitian dressed chiral fields with their opposite
Fourier orientations,
\begin{align}
 Q_R(P,x)&=\sum_{m=1}^{M}\frac{
  \zeta^{-mx}\widetilde C_{Rm}(P)
 +\zeta^{mx}\widetilde A_{Rm}(P)}{m},\nonumber\\
 Q_L(P,x)&=\sum_{m=1}^{M}\frac{
  \zeta^{mx}\widetilde C_{Lm}(P)
 +\zeta^{-mx}\widetilde A_{Lm}(P)}{m},
 \label{eq:ch20c-Q}\end{align}
\begin{equation}
 X_{\mathrm{CDW}}(P,x)=(c-s)(Q_R(P,x)+Q_L(P,x)),\quad
 X_{\mathrm{SC}}(P,x)=(c+s)(Q_R(P,x)-Q_L(P,x)).
 \label{eq:ch20c-X}
\end{equation}
These are well-defined finite sums of oscillator words, even though
their formal exponentials have infinitely many coefficients.
\begin{corollary}[Exchange of oscillator phases]
\label{cor:ch20c-X}
For all integer positions $x$,
\begin{align}
 D_{\mathrm{osc}}X_{\mathrm{CDW}}(P,x)
  &=-X_{\mathrm{SC}}(P^*,x),\nonumber\\
 D_{\mathrm{osc}}X_{\mathrm{SC}}(P,x)
  &=-X_{\mathrm{CDW}}(P^*,x).
 \label{eq:ch20c-Xtransform}
\end{align}
\end{corollary}
\begin{proof}
By~\eqref{eq:ch20c-dressed-transform},
$D_{\mathrm{osc}}Q_R(P,x)=-Q_R(P^*,x)$ and
$D_{\mathrm{osc}}Q_L(P,x)=Q_L(P^*,x)$.
For example,
\[
 D_{\mathrm{osc}}X_{\mathrm{CDW}}(P,x)
 =(c-s)(-Q_R(P^*,x)+Q_L(P^*,x)).
\]
The SC coefficient at $P^*$ is $c+(-s)=c-s$, giving the first
identity.  The second uses $c-(-s)=c+s$.  No sign is supplied by
complex conjugation: these are linear algebra substitutions.
\end{proof}

\subsection{Reflecting charge and preserving the Klein signs}
\label{subsec:ch20c-charge}

The charge representation occupies the algebraic direct sum
$\mathcal C=\mathbb C^{(\mathbb Z^2)}$ of finite-support vectors.
Its distinguished orthonormal kets are $\delta_N$ with
$N=(N_R,N_L)$.  Fix the same integer half-band offset $h$ as in
Section~\ref{sec:ch20-formal-vertices}.  Write $\mathsf N_\nu$ for the diagonal charge operator and
$N_\nu$ for the integer coordinate of a source ket.  The charge
operators and root-of-unity phases act by
\begin{equation}
 \mathsf N_\nu\delta_N=N_\nu\delta_N,
 \qquad Z_\nu(x)\delta_N=\zeta^{N_\nu x}\delta_N.
 \label{eq:ch20c-NZ}
\end{equation}
The signed lowering operators and their adjoints are
\begin{align}
 F_\nu\delta_N&=P_\nu(N)\delta_{N-e_\nu},\nonumber\\
 F_\nu^*\delta_N&=P_\nu(N+e_\nu)\delta_{N+e_\nu},\nonumber\\
 P_R(N)&=(-1)^{h+N_R-1},\qquad
 P_L(N)=(-1)^{N_R+N_L-1}.
 \label{eq:ch20c-Klein}
\end{align}
The left sign follows from the ordered-phase convention:
\[
 P_L(N)=(-1)^{(h+N_R)+(h+N_L-1)}.
\]
The integer $2h$ drops out of parity.  Thus $F_L$ contains a sign
depending on the \emph{right} charge; this convention cannot be
suppressed without changing the product coefficients.

\begin{definition}[Charge reflection and induced map]
\label{def:ch20c-charge}
Set
\begin{equation}
 f(N_R,N_L)=(N_R,-N_L),\qquad
 U_f\delta_N=\delta_{f(N)}.
 \label{eq:ch20c-f}
\end{equation}
The permutation $U_f$ is an involutive isometry of the algebraic
inner-product space $\mathcal C$.  On the \emph{generated adjointable}
star algebra $\mathfrak A_{\mathrm{ch}}$ of the operators in
\eqref{eq:ch20c-NZ}--\eqref{eq:ch20c-Klein}, define
\begin{equation}
 D_{\mathrm{ch}}(B)=U_fBU_f^{-1}=U_fBU_f.
 \label{eq:ch20c-Dcharge}
\end{equation}
The map is not defined here as a star automorphism of the entire
space of arbitrary endomorphisms of $\mathcal C$.
\end{definition}

\begin{lemma}[Exact action on the charge generators]
\label{lem:ch20c-charge}
The map $D_{\mathrm{ch}}$ preserves the generated star algebra and
satisfies
\begin{align}
 D_{\mathrm{ch}}(\mathsf N_R)&=\mathsf N_R,&
 D_{\mathrm{ch}}(\mathsf N_L)&=-\mathsf N_L,\nonumber\\
 D_{\mathrm{ch}}(Z_R(x))&=Z_R(x),&
 D_{\mathrm{ch}}(Z_L(x))&=Z_L(x)^*,\nonumber\\
 D_{\mathrm{ch}}(F_R)&=F_R,&
 D_{\mathrm{ch}}(F_L)&=-F_L^*.
 \label{eq:ch20c-charge-images}
\end{align}
It is a complex-linear, involutive star automorphism.
\end{lemma}
\begin{proof}
For each diagonal operator the claim follows by applying
$U_fBU_f$ to $\delta_N$: its eigenvalue is evaluated at $f(N)$.
In particular $\zeta^{-N_Lx}=\overline{\zeta^{N_Lx}}$, proving
the left-phase identity.  For the right signed shift,
\[
 U_fF_RU_f\delta_N=P_R(fN)\delta_{N-e_R}
 =P_R(N)\delta_{N-e_R}=F_R\delta_N.
\]
For the left shift, its reflected action is
\[
 U_fF_LU_f\delta_N
  =P_L(fN)\delta_{N+e_L},
\]
while the adjoint acts with coefficient $P_L(N+e_L)$.  Their ratio is
\[
 \frac{P_L(fN)}{P_L(N+e_L)}
 =\frac{(-1)^{N_R-N_L-1}}{(-1)^{N_R+N_L}}
 =(-1)^{-2N_L-1}=-1.
\]
Thus $D_{\mathrm{ch}}(F_L)=-F_L^*$, including the minus sign.
Every displayed image is again in the generated star algebra,
so conjugation restricts to an endomorphism of it.  As $U_f^2=I$,
this endomorphism is invertible and involutive.  Finally,
$U_f$ preserves the algebraic inner product, hence
$(U_fBU_f)^*=U_fB^*U_f$ on finite-support kets.  This proves
star preservation without assuming that the whole endomorphism
algebra carries an appropriate adjoint operation.
\end{proof}

\subsection{Ordered zero-mode words and exact vertex exchange}
\label{subsec:ch20c-vertices}

The CDW and pairing Klein--phase factors, in precisely the ordering
used for the finite-CAR coefficients and the formal vertices, are
\begin{align}
 K_{\mathrm{CDW}}(x)
 &=F_R^*F_L\,\zeta^{(\mathsf N_L-\mathsf N_R-1)x},
 \label{eq:ch20c-Kcdw}\\
 K_{\mathrm{SC}}(x)
 &=F_RF_L\,\zeta^{(\mathsf N_R+\mathsf N_L)x}.
 \label{eq:ch20c-Ksc}
\end{align}
The rightmost character acts on the \emph{source} ket before the
first shift.  This is why the power of $\zeta$ is not obtained by
simply replacing $N_L$ with $-N_L$ in a final answer.

\begin{proposition}[Klein--phase transformation with source shifts]
\label{prop:ch20c-Ktransform}
With $R<L$, the canonical root $\zeta$, and the signs
\eqref{eq:ch20c-Klein}, one has the exact operator identities
\begin{align}
 D_{\mathrm{ch}}K_{\mathrm{CDW}}(x)
 &=\zeta^x K_{\mathrm{SC}}(x)^*,\nonumber\\
 D_{\mathrm{ch}}K_{\mathrm{SC}}(x)
 &=\zeta^x K_{\mathrm{CDW}}(x)^*.
 \label{eq:ch20c-Kdual}
\end{align}
The scalar $\zeta^x$ and both adjoints are indispensable.
\end{proposition}
\begin{proof}
First record the action of the original words on any source
$N=(N_R,N_L)$:
\begin{align}
 K_{\mathrm{CDW}}(x)\delta_N
 &=a_C(N)\,\zeta^{(N_L-N_R-1)x}
       \delta_{N+e_R-e_L},\nonumber\\
 a_C(N)&=P_L(N)P_R(N-e_L+e_R),
 \label{eq:ch20c-actionC}\\
 K_{\mathrm{SC}}(x)\delta_N
 &=a_S(N)\,\zeta^{(N_R+N_L)x}
       \delta_{N-e_R-e_L},\nonumber\\
 a_S(N)&=P_L(N)P_R(N-e_L).
 \label{eq:ch20c-actionS}
\end{align}
The coefficients $a_C,a_S$ are real signs.  Therefore their
adjoint actions are obtained by reversing the charge shifts and
conjugating the \emph{phase evaluated at the original source}.

Apply the reflected CDW word to $\delta_N$.  Its amplitude is
$a_C(fN)\zeta^{(-N_L-N_R-1)x}$, and its target is
$N+e_R+e_L$.  On the other hand, $K_{\mathrm{SC}}(x)^*$ has
this same target, because its inverse source is
$N+e_R+e_L$.  The corresponding signed coefficient is
$a_S(N+e_R+e_L)$ and its character is
$\zeta^{-(N_R+N_L+2)x}$.  Direct use of
\eqref{eq:ch20c-Klein} gives
\begin{align*}
 a_C(fN)
 &=(-1)^{N_R-N_L-1}(-1)^{h+N_R}
 =(-1)^{h+N_L+1},\\
 a_S(N+e_R+e_L)
 &=(-1)^{N_R+N_L+1}(-1)^{h+N_R}
 =(-1)^{h+N_L+1}.
\end{align*}
The signs match, and their characters differ by precisely
$\zeta^x$:
\[
 \zeta^{(-N_L-N_R-1)x}
 =\zeta^x\zeta^{-(N_R+N_L+2)x}.
\]
Hence $D_{\mathrm{ch}}K_{\mathrm{CDW}}(x)
=\zeta^xK_{\mathrm{SC}}(x)^*$ on every basis vector.

For the reflected SC word, the target is $N-e_R+e_L$.
The inverse source for the CDW adjoint is then $N-e_R+e_L$.
The signs are
\begin{align*}
 a_S(fN)&=(-1)^{N_R-N_L-1}(-1)^{h+N_R-1}
 =(-1)^{h+N_L},\\
 a_C(N-e_R+e_L)
 &=(-1)^{N_R+N_L-1}(-1)^{h+N_R-1}
 =(-1)^{h+N_L},
\end{align*}
while the phases obey
\[
 \zeta^{(N_R-N_L)x}
 =\zeta^x\zeta^{-(N_L-N_R+1)x}.
\]
This proves the second operator identity on every $\delta_N$.
Because these vectors span $\mathcal C$, basis extensionality
finishes the proof.  Reordering the Klein operators or evaluating
the diagonal phase after a shift would spoil this calculation.
\end{proof}

Set $\mathfrak A=\mathfrak A_{\mathrm{ch}}\otimes_{\mathbb C}
\mathfrak A_{\mathrm{osc}}$.  Since charge and oscillator generators
act on different tensor factors, they commute in $\mathfrak A$.
Define
\begin{equation}
 D=D_{\mathrm{ch}}\otimes D_{\mathrm{osc}}:
  \mathfrak A\longrightarrow\mathfrak A.
 \label{eq:ch20c-Dtensor}
\end{equation}
On pure tensors $D(B\otimes a)=D_{\mathrm{ch}}(B)\otimes
D_{\mathrm{osc}}(a)$; bilinearity ensures that this is
well-defined on the algebraic tensor product.  It is a unital,
involutive star automorphism by the two preceding constructions.

For the indeterminate $t$ assume $t$ is central and $t^*=t$.
Extend $D$ coefficientwise to the noncommutative formal-series algebra
$\mathfrak A[[t]]$ by
\[
 D\left(\sum_{n\geq0}a_nt^n\right)
  =\sum_{n\geq0}D(a_n)t^n.
\]
This preserves multiplication: the coefficient of $t^n$ in a
product is the \emph{finite} sum
$\sum_{j=0}^{n}a_jb_{n-j}$; applying $D$ distributes over the
sum and preserves the order of each product.  It preserves star,
and its square is the identity.  No convergence or substitution
$t=1$ is used.

Recall that the Weyl-type vertices are, with the formal exponential
$\operatorname{Exp}_{\mathrm f}(tY)
=\sum_{n\geq0}Y^nt^n/n!$,
\begin{equation}
 W_J(P,t,x)=\frac1L K_J(x)
   \operatorname{Exp}_{\mathrm f}(itX_J(P,x)),
 \qquad J\in\{\mathrm{CDW},\mathrm{SC}\}.
 \label{eq:ch20c-W}
\end{equation}
These are not normal-ordered operators and their formal parameter
is not an analytic coupling.

\begin{theorem}[Exact coefficientwise CDW--pairing duality]
\label{thm:ch20c-Wdual}
For every admissible $P$, finite $M$, integer $x$, and formal $t$,
\begin{align}
 D W_{\mathrm{CDW}}(P,t,x)
 &=\zeta^xW_{\mathrm{SC}}(P^*,t,x)^*,\nonumber\\
 D W_{\mathrm{SC}}(P,t,x)
 &=\zeta^xW_{\mathrm{CDW}}(P^*,t,x)^*.
 \label{eq:ch20c-vertex-duality}
\end{align}
Both equalities hold in $\mathfrak A[[t]]$ coefficient by
coefficient.  In particular they assert no value at $t=1$.
\end{theorem}
\begin{proof}
For a Hermitian phase $X$ and $t^*=t$,
\[
 \operatorname{Exp}_{\mathrm f}(itX)^*
 =\operatorname{Exp}_{\mathrm f}(-itX)
\]
coefficient by coefficient: taking adjoints of $X^n$ reverses the
word order but leaves it equal to $X^n$, and conjugates $i^n$.
By Corollary~\ref{cor:ch20c-X},
\[
 D\operatorname{Exp}_{\mathrm f}
       (itX_{\mathrm{CDW}}(P,x))
 =\operatorname{Exp}_{\mathrm f}
       (-itX_{\mathrm{SC}}(P^*,x)).
\]
By Proposition~\ref{prop:ch20c-Ktransform}, the full CDW image is
\[
 \frac{\zeta^x}{L} K_{\mathrm{SC}}(x)^*
 \operatorname{Exp}_{\mathrm f}
       (-itX_{\mathrm{SC}}(P^*,x)).
\]
The charge word commutes with the oscillator series.  Thus this is
$\zeta^xW_{\mathrm{SC}}(P^*,t,x)^*$, with the \emph{reversed}
order required by taking the adjoint of $W_{\mathrm{SC}}$.
The second equality is identical with CDW and SC exchanged.
\end{proof}

\subsection{Hamiltonian covariance and transported states}
\label{subsec:ch20c-Hstate}

An identity of interaction families becomes stronger when the same
map intertwines their Hamiltonians and prescribed states.  Both
statements require an explicit choice of Hamiltonian and of the
source charge: they are not assertions about every finite lattice
Hamiltonian or every fixed-charge state.

\begin{theorem}[Symmetric Hamiltonian covariance]
\label{thm:ch20c-H}
Assume the symmetric, chemically adjusted zero-mode convention and
real $v_1,v_2$ as in~\eqref{eq:ch20c-velocities}.  In the
uncompressed oscillator--charge algebra define
\begin{align}
 H(P)={}&\frac{2\pi}{L}\sum_{m=1}^{M}\Bigl[
 v_1\bigl(C_{Rm}A_{Rm}+C_{Lm}A_{Lm}\bigr)
 +v_2\bigl(C_{Rm}C_{Lm}+A_{Rm}A_{Lm}\bigr)\Bigr]
 \nonumber\\
 &+\frac{\pi}{L}\Bigl[
 v_1(\mathsf N_R^2+\mathsf N_L^2)
 +2v_2\mathsf N_R\mathsf N_L\Bigr].
 \label{eq:ch20c-H}
\end{align}
Then, with $v_1(P^*)=v_1(P)$ and $v_2(P^*)=-v_2(P)$,
\begin{equation}
 D H(P)=H(P^*).
 \label{eq:ch20c-Hdual}
\end{equation}
\end{theorem}
\begin{proof}
For each $m$, both same-chirality number words are unchanged:
\[
 D(C_{Rm}A_{Rm})=(-C_{Rm})(-A_{Rm})=C_{Rm}A_{Rm},
\]
and the $L$ word is fixed.  Every mixed-chirality pairing word
contains exactly one $R$ factor, hence
\[
 D(C_{Rm}C_{Lm})=-C_{Rm}C_{Lm},\qquad
 D(A_{Rm}A_{Lm})=-A_{Rm}A_{Lm}.
\]
Thus the oscillator quadratic expression retains its form with
$v_2$ replaced by $-v_2$.  For the zero-mode polynomial,
\[
 D(\mathsf N_R^2+\mathsf N_L^2)
 =\mathsf N_R^2+\mathsf N_L^2,\quad
 D(\mathsf N_R\mathsf N_L)=-\mathsf N_R\mathsf N_L.
\]
Applying these identities term by term to the finite sum gives
\eqref{eq:ch20c-Hdual}.  No commutation of a creator past an
annihilator was needed, so the given normal and pairing order was
preserved.
\end{proof}

The corresponding states must be transported along the charge
permutation.  Recall from Section~\ref{sec:ch20-formal-vertices}
that the bare oscillator vacuum $\omega_0$ is represented by the
constant polynomial $1$ and that
$\omega_P=\omega_0\circ\beta_P^{-1}$.  For a source label $N$, the
full state $\omega_{P,N}$ is the vector functional at
$\delta_N\otimes1$ in the tensor product representation with the
dressed oscillator vacuum.  This is a normalized positive
\emph{complex-valued} functional on $\mathfrak A$; its formal-series
extension will be denoted $\widehat\omega_{P,N}$.

\begin{theorem}[Charge-indexed state transport]
\label{thm:ch20c-state}
For every $a\in\mathfrak A$ and every $N\in\mathbb Z^2$,
\begin{equation}
 \omega_{P^*,f(N)}(D a)=\omega_{P,N}(a).
 \label{eq:ch20c-statedual}
\end{equation}
The identity extends coefficientwise to all
$A(t)\in\mathfrak A[[t]]$:
\begin{equation}
 \widehat\omega_{P^*,f(N)}(D A(t))
   =\widehat\omega_{P,N}(A(t)).
 \label{eq:ch20c-formal-statedual}
\end{equation}
It is a \emph{family} identity; it does not claim that an arbitrary
single fixed-charge state is invariant under $D$.
\end{theorem}
\begin{proof}
First show that the polynomial vacuum is invariant under the bare
oscillator sign change.  In its polynomial representation define
$V_R$ by
\[
 (V_Rp)(X_{Rm},X_{Lm})
   =p(-X_{Rm},X_{Lm})\qquad(1\leq m\leq M),
\]
where the notation abbreviates the simultaneous change of sign of
all right-mode variables.  Every monomial is multiplied by $\pm1$;
therefore $V_R$ is an involutive isometry for the positive weighted
monomial Gram form, and $V_R1=1$.  Directly on polynomials,
\[
 V_R\pi_0(a)V_R^{-1}=\pi_0(D_{\mathrm{osc}}a)
\]
for the multiplication and differentiation generators, hence for
every oscillator word.  Consequently
\[
 \omega_0(D_{\mathrm{osc}}a)
 =\langle1,V_R\pi_0(a)V_R1\rangle
 =\omega_0(a).
\]
By Lemma~\ref{lem:ch20c-intertwine} and invertibility of the
Bogoliubov maps,
\[
 \beta_{P^*}^{-1}D_{\mathrm{osc}}
  =D_{\mathrm{osc}}\beta_P^{-1}.
\]
Thus
\begin{align*}
 \omega_{P^*}(D_{\mathrm{osc}}a)
 &=\omega_0\bigl(\beta_{P^*}^{-1}
                     D_{\mathrm{osc}}a\bigr)\\
 &=\omega_0\bigl(D_{\mathrm{osc}}\beta_P^{-1}a\bigr)
 =\omega_P(a).
\end{align*}

On the charge side the permutation itself implements
$D_{\mathrm{ch}}$, so for each charge word $B$,
\begin{align*}
 \langle\delta_{f(N)},D_{\mathrm{ch}}(B)
              \delta_{f(N)}\rangle
 &=\langle U_f\delta_N,U_fBU_f\,U_f\delta_N\rangle\\
 &=\langle\delta_N,B\delta_N\rangle.
\end{align*}
For a pure tensor $B\otimes a$ the full expectation is the
product of these two factor expectations.  Combining the identities
proves~\eqref{eq:ch20c-statedual} on pure tensors and, by
linearity, on finite sums of pure tensors, which exhaust the
algebraic tensor product.  Applying the established equality to
one coefficient at a time yields
\eqref{eq:ch20c-formal-statedual}; this uses no notion of positivity
for a general formal power series.
\end{proof}

The Hamiltonian and state laws together may be summarized as
\[
 H(P)\ \xrightarrow{\ D\ }\ H(P^*),\qquad
 (P,N)\ \xrightarrow{\ D\ }\ (P^*,fN).
\]
Notice that the second arrow refers to a \emph{pair of labels}, not
to an invariant state at $N$.  In particular, for $N_L\neq0$,
$fN\neq N$.  The map on operators is an automorphism of the
uncompressed algebra; its existence alone does not construct a
finite-energy-budget isomorphism of the Fock subspaces from
Section~\ref{sec:ch20-projected-products}.

\begin{remark}[Chemical potentials must be transformed or adjusted]
\label{rem:ch20c-mu}
An unadjusted zero mode can contain linear terms.  For example, if
\[
 H_{\mu}(P)=H(P)-\mu(\mathsf N_R+\mathsf N_L),
\]
then
\[
 D H_{\mu}(P)
 =H(P^*)-\mu(\mathsf N_R-\mathsf N_L)
 =H_{\mu}(P^*)+2\mu\mathsf N_L.
\]
The simple covariance law fails for nonzero $\mu$ on the full
charge algebra.  One may instead use the chemically adjusted
symmetric zero-mode convention above, restrict to a sector where
the residual vanishes, or explicitly reflect the coefficient of the
left linear term in a larger Hamiltonian family.  Such adjustments
must be stated; one cannot silently discard a chemical-potential
contribution.  Additive scalar offsets likewise require consistent
parameter transformation if they are included.
\end{remark}

\subsection{A four-site calculation}
\label{subsec:ch20c-example}

\begin{example}[Duality with signs, phases, and a nonzero coupling]
\label{ex:ch20c-L4}
Take $L=4$, $h=2$, $M=1$, $x=1$, $\zeta=i$, and the hyperbolic
parameters
\[
 u=8,\qquad c=\frac54,\qquad s=\frac34.
\]
Indeed $c^2-s^2=(25-9)/16=1$, while
\[
 (v_1,v_2)=(17,15),\quad
 (g,g^{-1})=\left(\frac14,4\right),\quad
 (v_1^*,v_2^*)=(17,-15).
\]
Thus this example is interacting at the quadratic level: $v_2$
does not vanish.  Let the charge source be $N=(0,0)$.
The ordered Klein signs are
\[
 P_L(0,0)=-1,\qquad P_R(1,-1)=+1,
\]
and the CDW source phase is $\zeta^{-1}=-i$.  Accordingly
\[
 K_{\mathrm{CDW}}(1)\delta_{(0,0)}
 =(-i)(-1)(+1)\delta_{(1,-1)}
 =i\delta_{(1,-1)}.
\]
Applying the charge reflection to its initial and final kets gives
\[
 D_{\mathrm{ch}}K_{\mathrm{CDW}}(1)\delta_{(0,0)}
  =i\delta_{(1,1)}.
\]
To check the right-hand side of the duality law independently,
$K_{\mathrm{SC}}(1)$ sends $\delta_{(1,1)}$ to
$\delta_{(0,0)}$ with coefficient $+1$: its Klein signs multiply
to $-1$ and its source phase is $i^{2}=-1$.  Hence
\[
 iK_{\mathrm{SC}}(1)^*\delta_{(0,0)}
 =i\delta_{(1,1)},
\]
in exact agreement.  Omitting the factor $\zeta^x=i$ would
produce an incorrect complex coefficient, and omitting the adjoint
would produce the opposite charge shift.

The oscillator transformation is equally explicit:
\[
 D_{\mathrm{osc}}\widetilde C_{R1}(P)
 =-\frac54 C_{R1}+\frac34 A_{L1}
 =-\widetilde C_{R1}(P^*),
\]
with the left dressed generator mapped without an overall minus
sign.  Therefore
\(D X_{\mathrm{CDW}}(P,1)=-X_{\mathrm{SC}}(P^*,1)\),
so the entire formal vertex identity, not only its constant
coefficient, follows coefficient by coefficient.

For a distinct illustration of the Hamiltonian statement, choose
$N=(1,2)$.  Its zero-mode energy is
\[
 E_0(P;1,2)
 =\frac\pi4[17(1^2+2^2)+30(1)(2)]
 =\frac{145\pi}{4}.
\]
The reflected charge is $(1,-2)$.  With $v_2^*=-15$,
\[
 E_0(P^*;1,-2)
 =\frac\pi4[17(1^2+(-2)^2)+2(-15)(1)(-2)]
 =\frac{145\pi}{4}.
\]
These are equal because the charge \emph{and} the coupling are
reflected.  They are not a demonstration of spectral equivalence of
finite-band fermionic Hamiltonians.
\end{example}

\subsection{Logical limits of the duality}
\label{subsec:ch20c-limits}

The preceding proofs establish exact algebraic identities at any
finite $M$ on the \emph{uncompressed} oscillator--charge model,
and coefficientwise identities for its formal vertex series.  They
do not establish the following further statements.

\begin{enumerate}
\item \emph{Finite-CAR equivalence.}  The finite fermionic Hilbert
space has Pauli constraints and the energy-budget projections of
Section~\ref{sec:ch20-projected-products}.  An isomorphism between
those budget carriers, compatible with their source, intermediate,
and target sectors, requires its own construction.  The intermediate
projection residuals derived there remain present.  In particular,
$D$ is not being asserted as a unitary conjugation on the finite-CAR
Fock space.

\item \emph{Finite-budget spectral equivalence.}  Even if two
quadratic expressions have the same transformation law in the
uncompressed algebra, compressing them with projections that are not
transported by an established intertwiner can change their
spectra.  A statement of unitary equivalence of compressed
Hamiltonians requires a specified map between budget subspaces and
a proved intertwining equality on those spaces.

\item \emph{Continuum or compactification duality.}  The algebraic
exchange $g\leftrightarrow g^{-1}$ echoes the familiar
$\phi\leftrightarrow\theta$ interchange in Luttinger-liquid
language (see, for context, \cite{giamarchi2004}).  With the bare
field convention $\phi=Q_R+Q_L$ and
$\theta=Q_R-Q_L$, the sign change of the right component produces
$\phi\mapsto-\theta$ and $\theta\mapsto-\phi$ (with the corresponding
dressed-family comparison).  A continuum $T$-duality, however,
needs a constructed continuum theory, compactification radius and
charge lattice, and an appropriate treatment of vertex sectors and
boundary conditions; none follows from the finite-mode star map.

\item \emph{Analytic vertices or physical phases.}  The maps on
$\mathfrak A[[t]]$ are coefficientwise.  Neither setting $t=1$ nor
claiming power-law asymptotics, a phase transition, or a gap is
justified without convergence estimates or a separate limiting and
spectral argument.  Nor does the positive state
$\omega_{P,N}$ automatically identify with a ground-state
expectation of an interacting finite-CAR Hamiltonian.
\end{enumerate}

In summary, the duality involution has been explicitly verified on the formal algebra constructed here. We have mapped the bare generators, their dressed families, the signed charge shifts, the ordered characters, the finite polynomial Hamiltonians, the charge-indexed positive states, and the coefficientwise Weyl-type vertices. Crucially, this algebraic model duality is not a statement of full interacting-CAR equivalence. Extending this formal mapping into a proven unitary symmetry of the finite-lattice CAR Hamiltonian, or evaluating it as a continuum dynamical equivalence, remains an open problem. As throughout Chapter 20, these derivations are unverified mathematical drafts awaiting formalization in Lean.

The models discussed in this chapter incorporate only two species of fermions (right and left movers). Physical models often include spin, which doubles the mode content to four species. In Chapter~21, we will motivate this four-species extension, establishing the species ordering, charge shifts, and parity conventions needed to define spinful interacting systems without ambiguity.


\chapter{Spinful Systems, Sector Leakage, and Effective Hamiltonians}
\label{ch:spinful}
\section{Spinful Species, Raw Currents, and the Parity Sublattice}
\label{sec:ch21-spin-charge}

\subsection*{From two chiral branches to four fermion species}

For spinless fermions, the two chiral branches carry separate density
fluctuations.  Adding an electron spin label produces four branches, but it
also raises a question which is invisible if one starts immediately with
continuum ``charge'' and ``spin'' fields: \emph{which pairs of charge and spin
quantum numbers actually describe finite fermionic states?}  A second question
concerns the familiar claim of spin--charge separation.  The operator algebra
will separate into simple linear combinations, whereas commutation on a
finite energy budget and factorization of the state space require distinct
arguments.  We first formulate these questions entirely in the finite CAR
(carrier) representation, before using any oscillator interpretation.

Fix an integer $h\geq 1$, and write
\begin{equation}
  L=2h,\qquad
  \mathcal B_L=(-h,h]\cap\mathbb Z=\{1-h,\ldots,h\},\qquad
  \Lambda_L=\mathbb Z/L\mathbb Z.
  \label{eq:ch21-finite-band}
\end{equation}
Thus $L$ is the number of lattice sites and also the number of allowed
momenta per species.  We reserve $M\geq 1$ for a \emph{positive-mode cutoff}:
when a collection of modes is considered, $1\leq m\leq M$; an individual
signed density shift $m$ may instead be any integer.  The Fourier root and
orientation are fixed by
\begin{equation}
  \zeta=\exp(2\pi i/L),\qquad
  c_{\alpha,x}=\frac{1}{\sqrt L}
     \sum_{k\in\mathcal B_L}\zeta^{kx}c_{\alpha,k},
  \qquad x\in\Lambda_L.
  \label{eq:ch21-fourier-orientation}
\end{equation}
We use periodic boundary conditions.  The inverse transform carries
$\zeta^{-kx}$.  The finite \emph{non-wrapping} density modes below must not be
silently identified with the complete spatial Fourier coefficients of
$c_{\alpha,x}^{\dagger}c_{\alpha,x}$: the latter also contain momenta
wrapped across the ends of $\mathcal B_L$ (see Chapter~\ref{ch:density_modes}).

\begin{definition}[Ordered spinful CAR carrier]
\label{def:ch21-species}
Let
\begin{equation}
  \mathcal C=\{L,R\}\times\{\uparrow,\downarrow\}.
  \label{eq:ch21-species}
\end{equation}
The chiral label is $\eta\in\{L,R\}$ and the spin label is
$\sigma\in\{\uparrow,\downarrow\}$.  For explicit sign computations we
choose once and for all the total order
\begin{equation}
  (L,\uparrow)<(L,\downarrow)<(R,\uparrow)<(R,\downarrow),
  \label{eq:ch21-species-order}
\end{equation}
followed by increasing momentum within each species.  The one-particle mode
set is $\mathcal I=\mathcal C\times\mathcal B_L$.  The fermionic Hilbert
space has the orthonormal occupation basis
\begin{equation}
  \mathcal H_L=\ell^2\!\bigl(\mathcal P(\mathcal I)\bigr),
  \quad |S\rangle,\ S\subseteq\mathcal I,
  \qquad \dim\mathcal H_L=2^{4L}.
\end{equation}
For $i\in\mathcal I$ write $c_i$ and $c_i^\dagger$ for the
annihilation and creation maps.  With $\{A,B\}=AB+BA$ they obey
\begin{equation}
 \{c_i,c_j\}=\{c_i^\dagger,c_j^\dagger\}=0,
 \qquad \{c_i,c_j^\dagger\}=\delta_{ij}I.
 \label{eq:ch21-four-species-car}
\end{equation}
\end{definition}

The order in \eqref{eq:ch21-species-order} is not a spin anisotropy.  It is a
convention needed to assign signs to occupation kets, CAR words and, later,
Klein operators.  To see the mechanism, put
$\epsilon(i,S)=(-1)^{\#\{j\in S:j<i\}}$.  Then
\begin{align}
 c_i|S\rangle
   &=\begin{cases}\epsilon(i,S)|S\setminus\{i\}\rangle,&i\in S,\\
                     0,&i\notin S,\end{cases}\label{eq:ch21-ann-action}\\
 c_i^\dagger|S\rangle
   &=\begin{cases}\epsilon(i,S)|S\cup\{i\}\rangle,&i\notin S,\\
                     0,&i\in S.\end{cases}\label{eq:ch21-cre-action}
\end{align}
These formulas reproduce the CAR: for $i\ne j$ the two possible orders
of inserting or deleting distinct occupied modes have opposite signs,
whereas for $i=j$ exactly one of $c_ic_i^\dagger$ and
$c_i^\dagger c_i$ acts as the identity on each basis ket.
The construction and its general CAR proof were established in
Chapter~\ref{ch:car_fock}; its multispecies version and the canonical
occupation tensor coordinates are supplied by the Core A05 module.
This matters because fermion operators of different species
\emph{anticommute}, even though the corresponding even density operators
will commute.

Let $n_{\alpha,k}=c_{\alpha,k}^\dagger c_{\alpha,k}$ and define the
sea-subtracted species-charge observable
\begin{equation}
  \widehat N_\alpha=\sum_{k\in\mathcal B_L}n_{\alpha,k}-hI,
  \qquad \alpha\in\mathcal C.
  \label{eq:ch21-relative-species-charge}
\end{equation}
Here the reference sea occupies the $h$ momenta $k\leq 0$ in every
species.  Its relative charges vanish; an occupation configuration has
$N_\alpha\in[-h,h]\cap\mathbb Z$.  This subtraction is essential when
talking about zero modes: $\sum_k n_{\alpha,k}$ counts \emph{absolute}
particles, while $\widehat N_\alpha$ counts the integer charge relative
to the chosen sea.

For an admissible vector of species charges
$\mathbf N=(N_\alpha)_{\alpha\in\mathcal C}$ and an integer $K\geq0$, let
\begin{equation}
 \mathcal B_{\mathbf N,K}
  =\operatorname{span}\Bigl\{|S\rangle:
     N_\alpha(S)=N_\alpha\ \text{for all }\alpha,
     \ \mathcal E(S)=\sum_{\alpha\in\mathcal C}\mathcal E_\alpha(S)
          \leq K\Bigr\}.
 \label{eq:ch21-joint-budget}
\end{equation}
Each $\mathcal E_\alpha$ is the nonnegative, integer excitation grading
\emph{above the ground configuration of its own charge sector}, constructed
in Chapter~\ref{ch:vacuum_budget}.  Equation
\eqref{eq:ch21-joint-budget} defines a coordinate subspace, not an
operator-order inequality of the form $H\psi\leq K\psi$.
We write $P_{\mathbf N,K}$ for its orthogonal projection.  The budget
constrains the \emph{sum} of the four excitation energies; it does not
place a separate independent cutoff $K$ on every species.

\subsection*{Raw currents and their exact inverse}

For $m\in\mathbb Z$ and $\alpha\in\mathcal C$, define the ordered
non-wrapping density operator
\begin{equation}
  \rho_{\alpha,m}
   =\sum_{\substack{k\in\mathcal B_L\\k+m\in\mathcal B_L}}
        c_{\alpha,k+m}^{\dagger}c_{\alpha,k}.
  \label{eq:ch21-species-density}
\end{equation}
The sign convention is deliberate: $m>0$ raises the momentum and the
excitation grading, and
$\rho_{\alpha,m}^{\dagger}=\rho_{\alpha,-m}$.  Since
$\rho_{\alpha,m}$ is even in fermion operators, the global CAR imply
\begin{equation}
 [\rho_{\alpha,m},\rho_{\beta,n}]=0\quad(\alpha\ne\beta).
 \label{eq:ch21-cross-species-density}
\end{equation}
For example, moving a two-fermion word for $\rho_{\alpha,m}$ past
one for $\rho_{\beta,n}$ entails four minus signs.  No energy budget is
needed for \eqref{eq:ch21-cross-species-density}.

\begin{definition}[Unnormalized charge and spin currents]
\label{def:ch21-raw-currents}
At each chirality $\eta$ put
\begin{align}
 R^c_{\eta,m}&=\rho_{(\eta,\uparrow),m}
                 +\rho_{(\eta,\downarrow),m},
     \label{eq:ch21-raw-charge}\\
 R^s_{\eta,m}&=\rho_{(\eta,\uparrow),m}
                 -\rho_{(\eta,\downarrow),m}.
     \label{eq:ch21-raw-spin}
\end{align}
Here $R^{c,s}$ denotes a \emph{raw current}, not the right-moving
chirality $R$.  Neither current contains a factor $1/\sqrt2$.
\end{definition}

The change of variables is algebraically invertible over $\mathbb C$:
\begin{equation}
 \begin{pmatrix}R^c_{\eta,m}\\R^s_{\eta,m}\end{pmatrix}
 =\begin{pmatrix}1&1\\1&-1\end{pmatrix}
   \begin{pmatrix}\rho_{(\eta,\uparrow),m}\\
                   \rho_{(\eta,\downarrow),m}\end{pmatrix},
 \qquad
 \begin{pmatrix}\rho_{(\eta,\uparrow),m}\\
                   \rho_{(\eta,\downarrow),m}\end{pmatrix}
 =\frac12\begin{pmatrix}1&1\\1&-1\end{pmatrix}
   \begin{pmatrix}R^c_{\eta,m}\\R^s_{\eta,m}\end{pmatrix}.
 \label{eq:ch21-raw-inverse}
\end{equation}
Indeed, the square of the displayed integer matrix is $2I$.
Taking adjoints and using that its entries are real proves
$(R^{c,s}_{\eta,m})^\dagger=R^{c,s}_{\eta,-m}$.  The factors
$1/2$ in the inverse are harmless for operators over $\mathbb C$;
their meaning for \emph{integer charge eigenvalues} will be quite
different.

\subsection*{The edge obstruction and the finite common-margin theorem}

A tempting argument says that up- and down-spin densities commute and
therefore so do the charge and spin currents.  The conclusion does
\emph{not} follow.  The two same-spin commutators still remain, and a
finite band has edge-dependent corrections to the usual current algebra.
Define
\begin{equation}
 E_{\eta,\sigma}(m,n)
  =[\rho_{(\eta,\sigma),m},\rho_{(\eta,\sigma),n}].
 \label{eq:ch21-edge-operator}
\end{equation}
These are genuine finite-dimensional operators, not a priori scalar
multiples of the identity.

\begin{proposition}[Global mixed-current identity]
\label{prop:ch21-global-edge}
For every $m,n\in\mathbb Z$ and either chirality $\eta$,
\begin{align}
 [R^c_{\eta,m},R^s_{\eta,n}]
   &=E_{\eta,\uparrow}(m,n)-E_{\eta,\downarrow}(m,n),
       \label{eq:ch21-mixed-edge}\\
 [R^c_{\eta,m},R^c_{\eta,n}]
   &=[R^s_{\eta,m},R^s_{\eta,n}]
     =E_{\eta,\uparrow}(m,n)+E_{\eta,\downarrow}(m,n).
       \label{eq:ch21-self-edge}
\end{align}
If $\eta\ne\eta'$, every commutator between
$R^a_{\eta,m}$ and $R^b_{\eta',n}$ vanishes for
$a,b\in\{c,s\}$, as an identity on all of $\mathcal H_L$.
\end{proposition}
\begin{proof}
Use bilinearity of $[A,B]=AB-BA$.  For instance, suppressing
$\eta,m,n$ when unambiguous,
\begin{align*}
 [R^c_m,R^s_n]
  &=[\rho_{\uparrow,m}+\rho_{\downarrow,m},
      \rho_{\uparrow,n}-\rho_{\downarrow,n}]\\
  &=[\rho_{\uparrow,m},\rho_{\uparrow,n}]
    -[\rho_{\uparrow,m},\rho_{\downarrow,n}]
    +[\rho_{\downarrow,m},\rho_{\uparrow,n}]
    -[\rho_{\downarrow,m},\rho_{\downarrow,n}].
\end{align*}
The middle two terms vanish by
\eqref{eq:ch21-cross-species-density}, giving
\eqref{eq:ch21-mixed-edge}.  With two plus signs, or two minus signs,
the same expansion gives \eqref{eq:ch21-self-edge}.  For different
chiralities all four species in the expansion are distinct, so again
every constituent commutator vanishes.  No rearrangement of fermionic
\emph{odd} factors was used.
\end{proof}

The single-species Schwinger coefficient in the convention of this book is
\begin{equation}
 \kappa(m,n)=
 \begin{cases}-m,&m+n=0,\\0,&m+n\ne0.\end{cases}
 \label{eq:ch21-schwinger-coefficient}
\end{equation}
It is the coefficient of a commutator acting on an appropriate budget,
not a globally valid nonzero scalar CCR on the finite Fock space.  The
origin of this distinction is the exact diagonal edge formula and the
frozen-margin theorem of Chapters~\ref{ch:heisenberg_algebra} and
\ref{ch:vacuum_budget}; their multispecies current version is provided
in the Core A06 support module.

\begin{theorem}[Raw charge/spin current algebra on a common budget]
\label{thm:ch21-budget-ccr}
Assume $h>0$, $K\geq0$, and an admissible charge vector
$\mathbf N$.  Fix $\eta\in\{L,R\}$ and $m,n\in\mathbb Z$.
Suppose \emph{both} spin species at that chirality satisfy
\begin{equation}
 \begin{split}
 |m|+|n|+K+|N_{(\eta,\uparrow)}|&\leq h,\\
 |m|+|n|+K+|N_{(\eta,\downarrow)}|&\leq h.
 \end{split}
 \label{eq:ch21-common-margin}
\end{equation}
For every $v\in\mathcal B_{\mathbf N,K}$, with the commutators
formed in the \emph{ambient} finite CAR representation,
\begin{align}
 [R^c_{\eta,m},R^s_{\eta,n}]v&=0,
   \label{eq:ch21-cross-ccr-budget}\\
 [R^c_{\eta,m},R^c_{\eta,n}]v
   &=[R^s_{\eta,m},R^s_{\eta,n}]v
     =2\kappa(m,n)v.
   \label{eq:ch21-self-ccr-budget}
\end{align}
In particular, if $m\geq1$ and
$2m+K+\max_{\sigma}|N_{(\eta,\sigma)}|\leq h$, then
\begin{equation}
 [R^c_{\eta,-m},R^c_{\eta,m}]v
   =[R^s_{\eta,-m},R^s_{\eta,m}]v=2m\,v,
 \qquad [R^c_{\eta,-m},R^s_{\eta,m}]v=0.
 \label{eq:ch21-two-m}
\end{equation}
\end{theorem}
\begin{proof}
The established single-species signed-current theorem states that,
under the corresponding margin in \eqref{eq:ch21-common-margin},
\begin{equation*}
 E_{\eta,\sigma}(m,n)v=\kappa(m,n)v
 \qquad(\sigma=\uparrow,\downarrow).
\end{equation*}
The same vector $v$ belongs to the joint budget, so each of these
identities applies \emph{to that vector}; the two separate conditions
are needed because the relative species charges can differ.  Substitute
them into \eqref{eq:ch21-mixed-edge} to obtain
$(\kappa-\kappa)v=0$.  Substituting into
\eqref{eq:ch21-self-edge} gives $(\kappa+\kappa)v=2\kappa v$.
Finally, taking the ordered pair of signed shifts $(-m,m)$ gives
$\kappa(-m,m)=m$ and $|{-m}|+|m|=2m$.
This proves all statements without postulating a finite-dimensional
Heisenberg representation.
\end{proof}

For all $1\leq m,n\leq M$, the simple \emph{sufficient} uniform
condition
\begin{equation}
 2M+K+\max_{\substack{\eta\in\{L,R\}\\
                         \sigma\in\{\uparrow,\downarrow\}}}
                   |N_{(\eta,\sigma)}|\leq h
 \label{eq:ch21-uniform-margin}
\end{equation}
allows the theorem to be applied to every pair of shifts among
$\{-M,\ldots,M\}$ at both chiralities.  It is not claimed to be
necessary.  In particular, a proof involving a \emph{product of several
currents} must additionally account for the intermediate excursions;
checking a one-commutator bound does not certify a whole word.

\begin{example}[The first nontrivial charge/spin oscillator]
\label{ex:ch21-first-mode}
Take $h=2$, so $L=4$ and
$\mathcal B_L=\{-1,0,1,2\}$, with all species charges zero and
$K=0$.  The ground ket $\Omega$ occupies $-1$ and $0$ in each
species.  For a fixed chirality and spin, the only unblocked
$m=1$ hop from this ground configuration is $0\to1$; write
$|e_\sigma\rangle=\rho_{(\eta,\sigma),1}\Omega$.
The CAR sign formula gives
\begin{equation*}
 \rho_{(\eta,\sigma),-1}\Omega=0,
 \qquad \rho_{(\eta,\sigma),-1}|e_\sigma\rangle=\Omega,
 \qquad \langle e_\uparrow,e_\downarrow\rangle=0.
\end{equation*}
Thus $R^c_{\eta,1}\Omega=|e_\uparrow\rangle+|e_\downarrow\rangle$
whereas $R^s_{\eta,1}\Omega=|e_\uparrow\rangle-|e_\downarrow\rangle$.
Their squared norms are both $2$, and direct application gives
\begin{equation}
 [R^c_{\eta,-1},R^c_{\eta,1}]\Omega=2\Omega,
 \quad [R^s_{\eta,-1},R^s_{\eta,1}]\Omega=2\Omega,
 \quad [R^c_{\eta,-1},R^s_{\eta,1}]\Omega=0.
 \label{eq:ch21-example-two}
\end{equation}
The bounds hold because $2(1)+0+0\leq h$.  Notice that the
\emph{positive} shift creates an excitation in either spin species:
the relative plus or minus sign distinguishes charge and spin
combinations, while both retain the raw normalization $2m$.
\end{example}

The budget hypothesis is substantive.  For example, keep one chirality,
put its up-spin species into the \emph{completely filled} band, its
down-spin species into the sea ground, and leave the other two species
in their sea grounds.  Every nonzero-momentum density hop annihilates
the full up-spin band, whereas the down-spin sector obeys
$[\rho_{\downarrow,-1},\rho_{\downarrow,1}]\Omega_\downarrow
 =\Omega_\downarrow$ for $h\geq2$.  Consequently on the resulting
nonzero product ket $v$,
\begin{equation}
  [R^c_{\eta,-1},R^s_{\eta,1}]v=-v\ne0.
  \label{eq:ch21-off-margin-example}
\end{equation}
The up-spin relative charge is $h$, so the common bound fails.
This is a direct counterexample to unrestricted global spin/charge
commutation on the finite CAR carrier.

There is another important boundary effect even \emph{inside} the
regime of Theorem~\ref{thm:ch21-budget-ccr}.  Set
$P=P_{\mathbf N,K}$, $Q=I-P$, and let $A,B$ be any ambient maps.
Insertion of $I=P+Q$ gives the exact identity
\begin{equation}
 [PAP,PBP]
   =P[A,B]P-PAQBP+PBQAP.
 \label{eq:ch21-compression-residual}
\end{equation}
Here the left side contains \emph{compressed} operators, unlike the
ambient commutators of the theorem.  In the preceding $K=0$ example,
$P$ has rank one, $PR^{c}_{\eta,\pm1}P=0$ and therefore the
compressed commutator is zero, although
$P[R^c_{\eta,-1},R^c_{\eta,1}]P=2P$.  The difference is precisely
supplied by the excursions through $Q\mathcal H_L$.  This example
makes concrete why no finite-dimensional compression can realize a
nonzero scalar CCR as an identity on its entire range: taking the
trace of $[A,B]=\lambda I$ in finite dimension would give
$0=\lambda\dim\mathcal H$.

\subsection*{Integer zero modes and the parity sublattice}

The same linear transformation must now be applied to the integer
\emph{zero-mode} charges, not to oscillator operators.  For each
chirality, define the commuting self-adjoint observables
\begin{equation}
 \widehat Q_{c,\eta}=\widehat N_{(\eta,\uparrow)}
                         +\widehat N_{(\eta,\downarrow)},
 \qquad
 \widehat Q_{s,\eta}=\widehat N_{(\eta,\uparrow)}
                         -\widehat N_{(\eta,\downarrow)}.
 \label{eq:ch21-zero-modes}
\end{equation}
They are called charge and spin zero modes in \emph{raw integer units};
$Q_s$ is twice the $S^z$ polarization in the usual spin-$\tfrac12$
normalization.  On $\mathcal B_{\mathbf N,K}$ their eigenvalues are
\begin{equation}
 Q_{c,\eta}=N_{(\eta,\uparrow)}+N_{(\eta,\downarrow)},
 \qquad
 Q_{s,\eta}=N_{(\eta,\uparrow)}-N_{(\eta,\downarrow)}.
 \label{eq:ch21-zero-eigenvalues}
\end{equation}
This follows immediately from the diagonal species-charge observable
and its established action on a fixed sector.

\begin{proposition}[Parity sublattice and integral inverse]
\label{prop:ch21-parity-lattice}
Consider the map of integer charge lattices
\begin{equation}
 T:\mathbb Z^2\longrightarrow\mathbb Z^2,
 \qquad T(a,b)=(a+b,a-b).
 \label{eq:ch21-parity-map}
\end{equation}
Its image is exactly the index-two subgroup
\begin{equation}
 \Lambda_{c,s}
 =\{(q_c,q_s)\in\mathbb Z^2:q_c\equiv q_s\pmod2\}.
 \label{eq:ch21-parity-sublattice}
\end{equation}
On that subgroup $T$ is a bijection, with inverse
\begin{equation}
 T^{-1}(q_c,q_s)
 =\left(\frac{q_c+q_s}{2},\frac{q_c-q_s}{2}\right).
 \label{eq:ch21-parity-inverse}
\end{equation}
The inclusion $\Lambda_{c,s}\subsetneq\mathbb Z^2$ is proper;
for example, $(1,0)\notin\Lambda_{c,s}$.
\end{proposition}
\begin{proof}
For $a,b\in\mathbb Z$,
$T(a,b)_c-T(a,b)_s=(a+b)-(a-b)=2b$;
hence both coordinates have the same parity.  Conversely, assume
$q_c\equiv q_s\pmod2$.  Then $(q_c-q_s)/2$ is an integer, and
$q_c+q_s=(q_c-q_s)+2q_s$ is also even.  Therefore the two coordinates
in \eqref{eq:ch21-parity-inverse} are integers.  Applying $T$
to that pair yields
\begin{equation*}
  \left(\frac{q_c+q_s}{2}+\frac{q_c-q_s}{2},
        \frac{q_c+q_s}{2}-\frac{q_c-q_s}{2}\right)=(q_c,q_s).
\end{equation*}
Applying the proposed inverse to $T(a,b)$ returns $(a,b)$ in the
same way.  This proves injectivity, surjectivity onto the displayed
image, and the inverse formula.  Finally $1-0$ is odd, so $(1,0)$
is excluded.  Equivalently, the integer matrix for $T$ has determinant
$-2$; its image has index two and is not all of $\mathbb Z^2$.
\end{proof}

For the full four-species problem, the proposition applies
\emph{separately} at $\eta=L$ and $\eta=R$.  In other words, the
physical integer labels belong to
$\Lambda_{c,s}^{(L)}\times\Lambda_{c,s}^{(R)}$, with additional
finite-band charge bounds inherited from $-h\leq N_\alpha\leq h$.
A single parity statement for the \emph{sum} of the two chiralities
would be weaker and would not recover both independent species-charge
pairs.  The parity rule does not project out fermionic states: it is an
arithmetic description of the image of the existing species labels
under a non-unimodular change of coordinates.

\begin{example}[An allowed zero mode and a forbidden formal label]
Take $h=2$ and one chirality with
$N_\uparrow=2$ and $N_\downarrow=-1$.  The raw labels are
$(Q_c,Q_s)=(1,3)$.  They have equal parity, and
\eqref{eq:ch21-parity-inverse} recovers
$(N_\uparrow,N_\downarrow)=(2,-1)$.  By contrast, a hypothetical
$(Q_c,Q_s)=(1,0)$ would give
$(N_\uparrow,N_\downarrow)=(1/2,1/2)$; it is a perfectly reasonable
point of $\mathbb R^2$ but not a finite fermionic charge sector.
This is why replacing integer charges by two unconstrained
independent bosonic zero modes changes the state-label set.
\end{example}

\subsection*{What factorizes, and what does not}

Because the one-particle labels of different species are disjoint, the
occupation basis gives a unitary \emph{ordered} identification of the
unrestricted finite Hilbert space with the tensor product of the four
single-species occupation spaces.  For odd CAR operators the ordering
requires fermionic parity strings; for even density operators their
action can be taken on the corresponding tensor factor.  This
occupation-space factorization is a theorem of the finite CAR
construction, but it is \emph{not} a factorization into independent
charge and spin zero-mode lattices.

To make the energy restriction explicit, denote by
$\mathcal H_{\alpha,N_\alpha}^{(e)}$ the subspace of a single
species with relative charge $N_\alpha$ and excitation energy exactly
$e\in\mathbb N$.  Regrouping the occupation basis of
\eqref{eq:ch21-joint-budget} yields the canonical graded-vector-space
identification
\begin{equation}
 \mathcal B_{\mathbf N,K}\ \simeq\!
 \bigoplus_{\substack{(e_\alpha)_{\alpha\in\mathcal C}\in\mathbb N^4\\
                    \sum_\alpha e_\alpha\leq K}}
 \bigotimes_{\alpha\in\mathcal C}
       \mathcal H_{\alpha,N_\alpha}^{(e_\alpha)}.
 \label{eq:ch21-budget-energy-sum}
\end{equation}
The direct sum is orthogonal, and any zero-dimensional energy-grade
factor contributes nothing.  To prove the identity, map an occupation
ket to the ordered tensor product of its four restrictions to species.
This is a bijection of orthonormal bases; by definition of the total
excitation grading, the ket belongs to the budget precisely when the
four resulting grades have sum at most $K$.  Linearity proves the
claimed unitary identification.  No statement about charge/spin
\emph{oscillator} tensor factors is required for this proof.

In general,
\begin{equation}
  \mathcal B_{\mathbf N,K}
  \not\simeq_{\mathrm{natural}}
  \bigotimes_{\alpha\in\mathcal C}
      \Bigl(\bigoplus_{e=0}^{K}
             \mathcal H_{\alpha,N_\alpha}^{(e)}\Bigr),
 \label{eq:ch21-not-product-cutoff}
\end{equation}
where the notation emphasizes that the \emph{natural inclusion} of
the left-hand space into the right-hand tensor product is generally
strict; it is not a statement that abstract finite-dimensional vector
spaces could never have accidentally equal dimensions.  For example,
at $K=1$ a state with one energy-one excitation in each of two
species lies in the product of separate $K=1$ cuts, but has total
energy two and is absent from the joint budget.  In particular, a
putative $\mathcal B_{c,K}\otimes\mathcal B_{s,K}$ cannot be inferred
from a change of current variables, even before addressing the parity
sublattice.

\subsection*{Exchange-symmetric quadratic couplings}

We now isolate the precise sense in which spin and charge \emph{do}
separate algebraically.  Let $X_\uparrow,X_\downarrow$ and
$Y_\uparrow,Y_\downarrow$ be operators on a common vector space.
They may be density modes at different chiralities or momenta, and
we do \emph{not} assume that $X$ and $Y$ commute.  Introduce the raw
combinations $X_c=X_\uparrow+X_\downarrow$,
$X_s=X_\uparrow-X_\downarrow$, and similarly for $Y$.

\begin{proposition}[Ordered quadratic block identity]
\label{prop:ch21-quadratic-block}
For arbitrary $a,b\in\mathbb C$,
\begin{align}
 &a(X_\uparrow Y_\uparrow+X_\downarrow Y_\downarrow)
  +b(X_\uparrow Y_\downarrow+X_\downarrow Y_\uparrow)
 \notag\\
 &\hspace{1.5cm}=
  \frac{a+b}{2}\,X_cY_c+
  \frac{a-b}{2}\,X_sY_s.
 \label{eq:ch21-quadratic-block}
\end{align}
Thus a quadratic coupling matrix of the exchange-symmetric form
$\begin{psmallmatrix}a&b\\b&a\end{psmallmatrix}$ is block diagonal
in the raw charge/spin variables, with the coefficients displayed
above.  The statement preserves the product order throughout.
\end{proposition}
\begin{proof}
Expanding only by distributivity,
\begin{align*}
 X_cY_c&=X_\uparrow Y_\uparrow+X_\uparrow Y_\downarrow
          +X_\downarrow Y_\uparrow+X_\downarrow Y_\downarrow,\\
 X_sY_s&=X_\uparrow Y_\uparrow-X_\uparrow Y_\downarrow
          -X_\downarrow Y_\uparrow+X_\downarrow Y_\downarrow.
\end{align*}
The coefficients of the two same-spin terms become
$((a+b)+(a-b))/2=a$; those of the cross-spin terms become
$((a+b)-(a-b))/2=b$.  At no point do we commute an $X$ past a $Y$.
Equivalently, the change-of-basis matrix in
\eqref{eq:ch21-raw-inverse} diagonalizes
$\begin{psmallmatrix}a&b\\b&a\end{psmallmatrix}$ over
$\mathbb C$, but the explicit expansion shows that the matrix
argument remains valid for noncommutative operator entries.
\end{proof}

For a quadratic current Hamiltonian built as a sum of such ordered
bilinears, with an exchange-symmetric spin matrix \emph{for each}
choice of chirality and momentum labels, the proposition may be
applied term by term.  One obtains a charge-current part and a
spin-current part with no mixed bilinear.  Hermiticity of the
Hamiltonian requires the usual conjugate/reversed-index conditions
on its coefficients separately; spin exchange symmetry alone is not
an adjointness hypothesis.

The hypothesis is also necessary to inspect rather than assume.
For instance, a spin-dependent diagonal coupling with coefficients
$a_\uparrow\ne a_\downarrow$ becomes
\begin{equation}
 a_\uparrow X_\uparrow Y_\uparrow
 +a_\downarrow X_\downarrow Y_\downarrow
 =\frac{a_\uparrow+a_\downarrow}{4}(X_cY_c+X_sY_s)
  +\frac{a_\uparrow-a_\downarrow}{4}(X_cY_s+X_sY_c).
 \label{eq:ch21-asymmetry-mixes}
\end{equation}
The last two terms mix charge and spin.  Even when
\eqref{eq:ch21-quadratic-block} holds, it implies neither global
commutation of its finite-CAR charge and spin Hamiltonian parts nor a
unitary factorization of a \emph{compressed} Hamiltonian.  Global
current commutators retain the edge difference of
\eqref{eq:ch21-mixed-edge}; products require further excursion
bounds, and compressions retain the residual terms of
\eqref{eq:ch21-compression-residual}.  Nonquadratic interactions,
including the operators considered in later sections, are not
covered by this proposition.

\subsection*{Do the charge and spin modes have different velocities?}

A quadratic block decomposition permits distinct effective mode
velocities but does not \emph{force} them.  To state the condition
without presupposing a continuum limit, consider for each formal
charge/spin oscillator block $\nu\in\{c,s\}$ the stable two-chirality
pairing Hamiltonian
\begin{equation}
 H_{\nu,m}
 =m\Bigl[v_{1,\nu}
   (a_{R,\nu}^{\dagger}a_{R,\nu}
    +a_{L,\nu}^{\dagger}a_{L,\nu})
  +v_{2,\nu}(a_{R,\nu}^{\dagger}a_{L,\nu}^{\dagger}
                         +a_{R,\nu}a_{L,\nu})\Bigr],
 \quad m\geq1.
 \label{eq:ch21-pairing-hamiltonian}
\end{equation}
Here the $a$'s satisfy exact CCR on an \emph{abstract oscillator
carrier}; this is a model for the budget-stable current algebra, not
a global canonical representation by the finite-dimensional CAR
operators.  Assume real coefficients and
$v_{1,\nu}>|v_{2,\nu}|$.  Using the CCR gives
\begin{equation}
 \left[\begin{pmatrix}a_{R,\nu}\\a_{L,\nu}^{\dagger}\end{pmatrix},
                  H_{\nu,m}\right]
 =m\begin{pmatrix}
  v_{1,\nu}&v_{2,\nu}\\-v_{2,\nu}&-v_{1,\nu}
 \end{pmatrix}
 \begin{pmatrix}a_{R,\nu}\\a_{L,\nu}^{\dagger}\end{pmatrix}.
 \label{eq:ch21-bogoliubov-dynamic}
\end{equation}
The characteristic polynomial of this $2\times2$ matrix is
$\lambda^2-(v_{1,\nu}^2-v_{2,\nu}^2)$; therefore its positive
frequency, divided by $m$, is
\begin{equation}
 u_\nu=\sqrt{v_{1,\nu}^2-v_{2,\nu}^2}>0.
 \label{eq:ch21-block-speed}
\end{equation}
This is an exact calculation \emph{within the stated quadratic CCR
model}.  It follows that
\begin{equation}
 u_c\ne u_s
 \quad\Longleftrightarrow\quad
 v_{1,c}^2-v_{2,c}^2\ne v_{1,s}^2-v_{2,s}^2,
 \label{eq:ch21-speed-nondegeneracy}
\end{equation}
provided both radicands are strictly positive.  The equivalence
uses only injectivity of the positive square root.  Its reverse
implication would require re-examination without stability and
positivity.

If one parametrizes a block by real hyperbolic numbers
$c_\nu,s_\nu$ satisfying $c_\nu^2-s_\nu^2=1$, then one may use
\begin{equation}
 v_{1,\nu}=u_\nu(c_\nu^2+s_\nu^2),\qquad
 v_{2,\nu}=2u_\nu c_\nu s_\nu,\qquad
 g_\nu=(c_\nu-s_\nu)^2,\qquad
 g_\nu^{-1}=(c_\nu+s_\nu)^2.
 \label{eq:ch21-hyperbolic-parameters}
\end{equation}
The last identity follows by multiplication of the two squares and
$c_\nu^2-s_\nu^2=1$.  The symbol $g_\nu$ denotes a quadratic-model
Luttinger parameter, \emph{not} the real Umklapp coupling $g_U$
used later.  We also reserve $K$ exclusively for the finite energy
budget, not for the Luttinger parameter.

\begin{example}[Exchange symmetry without velocity separation]
\label{ex:ch21-equal-speeds}
Start with a spin-independent free linearized model of velocity
$v>0$.  Its spin-exchange-symmetric quadratic coupling has no
up/down mixing, $b=0$, and the same diagonal coefficient $a=v$
for both spins.  Equation~\eqref{eq:ch21-quadratic-block} gives
$\tfrac v2 X_cY_c+\tfrac v2 X_sY_s$; it \emph{does} have charge
and spin blocks.  After consistent canonical oscillator
normalization, both blocks have
\begin{equation*}
  v_{1,c}=v_{1,s}=v,\qquad v_{2,c}=v_{2,s}=0,
  \qquad u_c=u_s=v,
  \qquad c_c=c_s=1,\quad s_c=s_s=0.
\end{equation*}
Hence symmetry and even a quadratic charge/spin decomposition do not
imply unequal propagation speeds.  For contrast, the stable parameter
choices $(v_{1,c},v_{2,c})=(17,15)$ and
$(v_{1,s},v_{2,s})=(5,0)$ give $u_c=8$ and $u_s=5$.
The latter distinction comes from the \emph{parameters}, not merely
from the existence of the charge/spin basis.
\end{example}

\begin{remark}[Mathematical and formalization scope]
\label{rem:ch21-formalization-scope}
Three separate conclusions have been established in this section.
The four-species CAR and the global density identities concern the
\emph{entire finite Fock carrier}; the coefficient $2m$ and mixed
current cancellation concern the \emph{ambient action on specified
frozen-margin budget vectors}; and the dispersion computation
concerns a \emph{specified abstract quadratic CCR model}.  None
alone identifies its oscillator vacuum with a finite interacting
fermionic ground state, proves a thermodynamic limit, or provides a
spectral gap.  A charge/spin oscillator tensor-product description
would also have to respect the parity lattice and the total-energy
cutoff.  Such a construction is not inferred here.

The finite multispecies CAR and its support results are present in
the supplied Core modules, but the Chapter~21 spin/charge statements
as a combined Lean development are \emph{not yet implemented or
verified}.  The explicit conventions to reconcile at whole-book
integration are the momentum/Fourier orientation, sea subtraction for
zero modes, raw-versus-normalized current factors, ambient-versus-
compressed products, and the inherited sign order
\eqref{eq:ch21-species-order}.  For physical background on the
charge/spin description of one-dimensional fermions, see
\cite{giamarchi2004,gogolin1998,vondelft1998}; these references
are context rather than certificates for the finite margin theorem.
\end{remark}

\section{The Spin-Singlet Cooper Pair and Its Two Charge Channels}
\label{sec:ch21-singlet-pair}

\subsection*{Why one singlet operator has two different targets}

A spin-singlet Cooper pair is often represented by a single compact field expression in the continuum literature. However, on a finite fermionic Hilbert space, we must first ask a more elementary question: \emph{which occupation-number sectors does this operator connect?} When assembling a pair from a right mover and a left mover, the two antisymmetrized spin assignments remove particles from different species. Although they carry the same \emph{total} charge and spin, they do not possess the same four-component charge vector. Consequently, when their finite energy budgets are treated as typed linear maps, their operators have distinctly different target spaces. Furthermore, their relative fermionic signs must be explicitly computed and cannot simply be discarded under the heuristic phrase ``the Klein factors cancel.''

We continue with the ordered four-species CAR carrier of
Section~\ref{sec:ch21-spin-charge}.  Set
\begin{equation}
 L=2h,\quad h\in\mathbb N_{>0},\qquad
 \mathcal B_L=\{1-h,2-h,\ldots,h\},\qquad
 \Lambda_L=\mathbb Z/L\mathbb Z,\qquad
 \zeta=e^{2\pi i/L}.
 \label{eq:ch21-pair-setting}
\end{equation}
Every one of the four species
\(\mathcal C=\{L,R\}\times\{\uparrow,\downarrow\}\) has these \(L\)
momentum modes; the lexicographic occupation order is
\begin{equation}
 L\uparrow<L\downarrow<R\uparrow<R\downarrow,
 \quad\text{then increasing momentum within each species}.
 \label{eq:ch21-pair-order}
\end{equation}
We retain the positive-mode cutoff notation \(1\leq m\leq M\) for later
current or oscillator expansions, but the exact CAR identities in this
section require no such cutoff.  All site labels \(x\) are periodic lattice
positions, represented by integers modulo \(L\).  With periodic Fourier
orientation fixed as in Section~\ref{sec:ch21-spin-charge}, the finite site
annihilator is
\begin{equation}
 c_{\alpha,x}=\frac1{\sqrt L}\sum_{k\in\mathcal B_L}
        \zeta^{kx}c_{\alpha,k},\qquad
 c_{\alpha,x}^{\dagger}=\frac1{\sqrt L}\sum_{k\in\mathcal B_L}
        \zeta^{-kx}c_{\alpha,k}^{\dagger}.
 \label{eq:ch21-pair-site-field}
\end{equation}
These are genuine operators on the finite occupation Fock space
\(\mathcal H_L\), not formal vertex exponentials.  Orthogonality of the
finite characters and the momentum CAR give
\(\{c_{\alpha,x},c_{\beta,y}^{\dagger}\}
=\delta_{\alpha\beta}\delta_{xy}I\) and
\(\{c_{\alpha,x},c_{\beta,y}\}=0\).
Indeed, the first anticommutator is
\(L^{-1}\delta_{\alpha\beta}\sum_{k\in\mathcal B_L}
\zeta^{k(x-y)}I\), which is the stated periodic Kronecker delta.
There is no non-wrapping current approximation in this step.

\begin{definition}[Annihilating spin-singlet pair]
\label{def:ch21-singlet}
The ordered, spin-antisymmetric two-fermion operator is
\begin{equation}
 \begin{split}
 O_{SSC}(x)&=A_1(x)-A_2(x),\\
 A_1(x)&=c_{R\uparrow,x}\,c_{L\downarrow,x},\\
 A_2(x)&=c_{R\downarrow,x}\,c_{L\uparrow,x}.
 \end{split}
 \label{eq:ch21-singlet-def}
\end{equation}
The rightmost field acts first.  In particular, \(A_1\) annihilates an
\(L\downarrow\) particle before an \(R\uparrow\) particle; \(A_2\)
annihilates an \(L\uparrow\) particle before an \(R\downarrow\) particle.
\end{definition}

The minus sign in \eqref{eq:ch21-singlet-def} is the antisymmetric spin
contraction \(\sum_{\sigma,\tau}\epsilon_{\sigma\tau}
 c_{R\sigma,x}c_{L\tau,x}\), with
\(\epsilon_{\uparrow\downarrow}=1\) and
\(\epsilon_{\downarrow\uparrow}=-1\).  For \(U\in SU(2)\), the elementary
identity \(U^{T}\epsilon U=(\det U)\epsilon=\epsilon\) shows that this
contraction is a spin scalar when both spin labels transform by the same
matrix.  This representation-theoretic assertion does \emph{not} say that
the Hamiltonian possesses spin-rotation symmetry.  It also does not
eliminate the antisymmetric CAR ordering.  For example, taking the adjoint
and reordering the fermions gives
\begin{equation}
 O_{SSC}(x)^\dagger
  =-c_{R\uparrow,x}^\dagger c_{L\downarrow,x}^\dagger
    -c_{L\uparrow,x}^\dagger c_{R\downarrow,x}^\dagger.
 \label{eq:ch21-singlet-adjoint}
\end{equation}
Thus the annihilating convention here differs by an overall sign upon
adjunction from a frequently written spin-singlet \emph{creation} convention.

\subsection*{Four-component charge shifts and the common total quantum numbers}

Let \(\widehat N_\alpha=\sum_{k\in\mathcal B_L}
 c_{\alpha,k}^\dagger c_{\alpha,k}-hI\) be the relative species charge;
write \(\mathbf N\in\mathbb Z^{\mathcal C}\) for its simultaneous
eigenvalues, and \(e_\alpha\) for the unit charge vector in species
\(\alpha\).  Only admissible components \(-h\leq N_\alpha\leq h\)
occur on the finite carrier.  The \(c_{\alpha,k}\) preserve neither
\(\mathbf N\) nor the energy budget, but their \emph{charge shifts} are
exact on the entire ambient finite Fock space.

\begin{lemma}[Exact species shifts of the two singlet channels]
\label{lem:ch21-singlet-shifts}
For every \(\gamma\in\mathcal C\) and every \(x\in\Lambda_L\),
\begin{align}
 [\widehat N_\gamma,A_1(x)]
   &=-(\delta_{\gamma,R\uparrow}+\delta_{\gamma,L\downarrow})A_1(x),
   \label{eq:ch21-pair-shift-one}\\
 [\widehat N_\gamma,A_2(x)]
   &=-(\delta_{\gamma,R\downarrow}+\delta_{\gamma,L\uparrow})A_2(x).
   \label{eq:ch21-pair-shift-two}
\end{align}
Consequently, if \(\mathcal H_{\mathbf N}\subseteq\mathcal H_L\)
is the full joint eigenspace of species charges \(\mathbf N\), then
\begin{equation}
 A_j(x)\mathcal H_{\mathbf N}\subseteq\mathcal H_{\mathbf T_j},
 \qquad\begin{cases}
 \mathbf T_1=\mathbf N-e_{L\downarrow}-e_{R\uparrow},\\
 \mathbf T_2=\mathbf N-e_{L\uparrow}-e_{R\downarrow}.
 \end{cases}
 \label{eq:ch21-pair-targets}
\end{equation}
These two targets are distinct integer vectors.  Their charge eigenspaces
are orthogonal, including when one of the eigenspaces is zero.
\end{lemma}

\begin{proof}
The CAR imply
\([c_{\gamma,k}^\dagger c_{\gamma,k},c_{\beta,q}]
=-\delta_{\gamma\beta}\delta_{kq}c_{\beta,q}\).
Summing over \(k\), subtracting the scalar sea occupation \(hI\), and
then using \eqref{eq:ch21-pair-site-field} gives
\([\widehat N_\gamma,c_{\beta,x}]
=-\delta_{\gamma\beta}c_{\beta,x}\).
For any operators \(D,X,Y\), the identity
\([D,XY]=[D,X]Y+X[D,Y]\) holds without a commutativity
assumption.  Apply it to the \emph{ordered} words \(A_1,A_2\) to obtain
\eqref{eq:ch21-pair-shift-one}--\eqref{eq:ch21-pair-shift-two}.
If \(\widehat N_\gamma\psi=N_\gamma\psi\), then
\(\widehat N_\gamma A_j\psi=(N_\gamma+\Delta_{j,\gamma})A_j\psi\),
which gives \eqref{eq:ch21-pair-targets}.  Its two shift vectors are
different because, for instance, their \(L\uparrow\) components are
\(0\) and \(-1\).  Finally, distinct eigenspaces of a self-adjoint
operator are orthogonal: if \(\widehat N_\gamma u=au\) and
\(\widehat N_\gamma v=bv\), then
\(a\langle u,v\rangle=\langle\widehat N_\gamma u,v\rangle
=\langle u,\widehat N_\gamma v\rangle=b\langle u,v\rangle\).
Choose a component for which \(a\ne b\).
\end{proof}

It is instructive to express these shifts using the raw charge and spin
labels of Section~\ref{sec:ch21-spin-charge}:
\(Q_{\eta,c}=N_{\eta\uparrow}+N_{\eta\downarrow}\) and
\(Q_{\eta,s}=N_{\eta\uparrow}-N_{\eta\downarrow}\).
\begin{center}
\begin{tabular}{@{}lrrrrrr@{}}
\toprule
 & $\Delta Q_{R,c}$ & $\Delta Q_{R,s}$
 & $\Delta Q_{L,c}$ & $\Delta Q_{L,s}$
 & $\Delta Q_{\mathrm{tot},c}$ & $\Delta Q_{\mathrm{tot},s}$\\
\midrule
$A_1$ & $-1$ & $-1$ & $-1$ & $+1$ & $-2$ & $0$\\
$A_2$ & $-1$ & $+1$ & $-1$ & $-1$ & $-2$ & $0$\\
\bottomrule
\end{tabular}
\end{center}
The common total \(\Delta Q_{\mathrm{tot},c}=-2\) and
\(\Delta Q_{\mathrm{tot},s}=0\) do not identify the two targets.  In
particular, the chirality-resolved spin shifts differ, and each branch
continues to obey the integer parity constraint
\(Q_{\eta,c}\equiv Q_{\eta,s}\pmod 2\).  Forgetting chirality-resolved
charges before projecting would erase information needed to determine the
correct codomain.

\subsection*{The signed Klein factors are ordered charge-shift maps}

While a Klein factor is sometimes introduced as an abstract anticommuting symbol attached to each fermion field, on our finite occupation representation the source of its sign is entirely concrete: removing a fermion requires commuting it past all occupied modes that precede it in the fixed species--momentum order. Thus, the signed factor is a rigorously defined \emph{charge-dependent map}, rather than an arbitrary numerical phase assigned loosely to each species.

For each admissible integer vector \(\mathbf N\), write
\(\Omega_{\mathbf N}\) for the normalized charge-sector ground ket, with
species \(\alpha\) filled at
\(\{1-h,\ldots,N_\alpha\}\).  Its number of particles of species
\(\alpha\) is \(h+N_\alpha\).  Assume below that a named annihilation is
possible, namely \(N_\alpha\geq1-h\), and that all indicated lowered
charge vectors remain admissible.  In Chapter~\ref{ch:car_fock}'s
occupation action, the annihilator of the \emph{top occupied momentum}
\(k=N_\alpha\) has preceding-occupation exponent
\begin{equation}
 \begin{split}
 E_h(\mathbf N,\alpha)
   &=\sum_{\beta<\alpha}(h+N_\beta)
       +(h+N_\alpha-1),\\
 \kappa_h(\mathbf N,\alpha)
   &=(-1)^{E_h(\mathbf N,\alpha)}\in\{+1,-1\}.
 \end{split}
 \label{eq:ch21-pair-klein-phase}
\end{equation}
The first term counts occupied modes in preceding species; the second
counts those below the occupied top mode within species \(\alpha\).
The signed finite Klein kernel constructed in Core Chapter~13
reproduces precisely this phase on sector ground states, under the
charge and budget margins stated for that construction.

To keep the zero-mode notation transparent, denote by \(F_\alpha\) the
\emph{typed sector shift} whose ground-state action is
\begin{equation}
 F_\alpha\Omega_{\mathbf N}
  =\kappa_h(\mathbf N,\alpha)\,
    \Omega_{\mathbf N-e_\alpha}.
 \label{eq:ch21-pair-klein-action}
\end{equation}
This notation describes the Core Klein action on admissible ground
states; extending it to particular non-ground budgets requires the
corresponding finite kernel theorem.  It is not an assertion that one
bounded lowering operator acts isometrically on every charge sector,
including the empty boundary sector.  Define also the diagonal site phase
\(Z_{\alpha,x}\Omega_{\mathbf N}=\zeta^{xN_\alpha}
\Omega_{\mathbf N}\).  Since lowering species \(\beta\ne\alpha\)
does not change \(N_\alpha\), the ground-state actions of
\(Z_{\alpha,x}\) and \(F_\beta\) commute for distinct species.

\begin{proposition}[Two different signed zero-mode products]
\label{prop:ch21-two-klein-products}
With \(\mathbf T_1,\mathbf T_2\) as in
\eqref{eq:ch21-pair-targets}, the two ordered zero-mode combinations are
\begin{align}
 \mathcal K_1(x)&=
 F_{R\uparrow}F_{L\downarrow}
 Z_{R\uparrow,x}Z_{L\downarrow,x}:
  \mathbf N\longrightarrow\mathbf T_1,
 \label{eq:ch21-klein-one}\\
 \mathcal K_2(x)&=
 F_{R\downarrow}F_{L\uparrow}
 Z_{R\downarrow,x}Z_{L\uparrow,x}:
  \mathbf N\longrightarrow\mathbf T_2.
 \label{eq:ch21-klein-two}
\end{align}
On an admissible sector ground ket, the actions are
\begin{equation}
 \mathcal K_j(x)\Omega_{\mathbf N}
  =\zeta^{x(N_{a_j}+N_{b_j})}
    \kappa_h(\mathbf N-e_{b_j},a_j)
    \kappa_h(\mathbf N,b_j)\,\Omega_{\mathbf T_j},
 \label{eq:ch21-klein-channel-ground}
\end{equation}
where \((a_1,b_1)=(R\uparrow,L\downarrow)\) and
\((a_2,b_2)=(R\downarrow,L\uparrow)\).
The rightmost \(F_{b_j}\) acts first; in particular the second
Klein phase must be evaluated on \(\mathbf N-e_{b_j}\), not on
\(\mathbf N\).
\end{proposition}

\begin{proof}
Act with the rightmost diagonal \(Z\) operators on
\(\Omega_{\mathbf N}\); their product is
\(\zeta^{x(N_{a_j}+N_{b_j})}\).  Next \(F_{b_j}\) lowers
\(\mathbf N\) to \(\mathbf N-e_{b_j}\) with phase
\(\kappa_h(\mathbf N,b_j)\).  Finally \(F_{a_j}\) acts on that
\emph{intermediate} charge vector, giving
\(\kappa_h(\mathbf N-e_{b_j},a_j)\).
This proves \eqref{eq:ch21-klein-channel-ground}.
Notice that in our fixed ordering both \(b_1<a_1\) and
\(b_2<a_2\).  Hence
\(E_h(\mathbf N-e_b,a)=E_h(\mathbf N,a)-1\) whenever
\(b<a\) and \(a\ne b\).  Reversing the two shifts changes their
relative sign, which is the finite occupation-string origin of the
fermionic cross-species anticommutation.  No freedom to suppress that
sign has been introduced.
\end{proof}

The \(Z\) factors specify zero-mode phases, not the full oscillator
content of the site fields.  The formulas above are \emph{not} a
factorization theorem expressing \(c_{a,x}c_{b,x}\) on every budget
as a product of two compressed bosonic exponentials.  Such a statement
requires the single-field identifications, a common domain of composition,
and control of intermediate projection remainders.

\begin{lemma}[Matching the ordered CAR action on sector grounds]
\label{lem:ch21-pair-ground-action}
For distinct species \(a,b\), assume that the charge-sector ground
\(\Omega_{\mathbf N}\) and both lowerings are admissible.  Let
\(k_b=N_b\) and \(k_a=N_a\) denote the occupied top momenta in
\(\mathcal B_L\).  Then
\begin{equation}
 c_{a,k_a}c_{b,k_b}\Omega_{\mathbf N}
  =\kappa_h(\mathbf N-e_b,a)\kappa_h(\mathbf N,b)
      \Omega_{\mathbf N-e_b-e_a},
 \label{eq:ch21-pair-ordered-top}
\end{equation}
and, for the site fields,
\begin{equation}
 \big\langle\Omega_{\mathbf N-e_b-e_a},
  c_{a,x}c_{b,x}\Omega_{\mathbf N}\big\rangle
   =\frac{\zeta^{x(N_a+N_b)}}{L}
     \kappa_h(\mathbf N-e_b,a)\kappa_h(\mathbf N,b).
 \label{eq:ch21-pair-channel-coefficient}
\end{equation}
The right-hand side is \(L^{-1}\) times the corresponding
zero-mode matrix coefficient in \eqref{eq:ch21-klein-channel-ground}.
\end{lemma}

\begin{proof}
For a ground configuration, deleting its top occupied momentum
\(k_b=N_b\) leaves the ground configuration of charge
\(\mathbf N-e_b\); the occupation-basis annihilation rule gives the
phase \(\kappa_h(\mathbf N,b)\).  Deleting \(k_a=N_a\) from this
intermediate ground gives the second phase
\(\kappa_h(\mathbf N-e_b,a)\), proving
\eqref{eq:ch21-pair-ordered-top}.
To compute the matrix element for site annihilators, expand both
finite Fourier sums in \eqref{eq:ch21-pair-site-field}.
The target ket is the \emph{ground} configuration at
\(\mathbf N-e_a-e_b\).  Removing any momentum other than the
occupied top mode of species \(b\) leaves a different occupation
configuration, and removal of a non-top momentum of species \(a\)
likewise cannot produce the target ground.  Orthogonality of
occupation kets therefore selects the unique summand
\((k_a,k_b)=(N_a,N_b)\).  Its two Fourier amplitudes multiply to
\(L^{-1}\zeta^{x(N_a+N_b)}\), while its CAR sign is the product
already obtained.  No current-CCR or energy-margin hypothesis enters
this matrix coefficient; it is an exact finite-CAR statement.
\end{proof}

\subsection*{A finite sign calculation: both Klein phases matter}

\begin{example}[Four species, \(L=4\), and two opposite singlet coefficients]
\label{ex:ch21-pair-sign}
Take \(h=2\), \(L=4\), \(x=1\), \(\zeta=i\), and the zero-charge sea
\(\Omega_{\mathbf0}\).  Each of the four species occupies momenta
\(\{-1,0\}\), ordered as in \eqref{eq:ch21-pair-order}.
Both singlet channels remove the top momentum \(k=0\) from their
respective species.

For channel 1, annihilating the rightmost \(L\downarrow\) field
encounters \(2+1=3\) preceding occupied modes, hence contributes
\(-1\).  In the intermediate state, annihilating \(R\uparrow\) at
\(k=0\) encounters \(2+1+1=4\) preceding modes, hence contributes
\(+1\).  Thus
\begin{equation}
 \kappa_2(\mathbf0,L\downarrow)=-1,\quad
 \kappa_2(-e_{L\downarrow},R\uparrow)=+1,\quad
 \langle\Omega_{\mathbf T_1},A_1(1)\Omega_{\mathbf0}\rangle=-\tfrac14.
 \label{eq:ch21-pair-example-one}
\end{equation}
For channel 2, removing \(L\uparrow\) at \(k=0\) encounters one
preceding occupied mode and contributes \(-1\).  Removing
\(R\downarrow\) at \(k=0\) from the intermediate configuration
encounters \(1+2+2+1=6\) preceding modes and contributes \(+1\).
Therefore
\begin{equation}
 \kappa_2(\mathbf0,L\uparrow)=-1,\quad
 \kappa_2(-e_{L\uparrow},R\downarrow)=+1,\quad
 \langle\Omega_{\mathbf T_2},A_2(1)\Omega_{\mathbf0}\rangle=-\tfrac14.
 \label{eq:ch21-pair-example-two}
\end{equation}
The Fourier phase is \(\zeta^{(0+0)x}=1\) in both cases, and each
pair contributes the normalization \(1/L=1/4\).  In the \emph{singlet
operator} the second channel is subtracted.  Accordingly,
\begin{equation}
 \boxed{\quad
 \langle\Omega_{\mathbf T_1},O_{SSC}(1)\Omega_{\mathbf0}\rangle
   =-\tfrac14,\qquad
 \langle\Omega_{\mathbf T_2},O_{SSC}(1)\Omega_{\mathbf0}\rangle
   =+\tfrac14.\quad}
 \label{eq:ch21-pair-example-final}
\end{equation}
There is no contradiction in the opposite signs: the target kets
belong to different, mutually orthogonal charge sectors.  Their
amplitudes must not be summed as if they were two contributions to
the same basis coefficient.
\end{example}

The sign pattern has a useful exact generalization.  With
\(\mathbf N=\mathbf0\) and arbitrary \(h\geq1\), the channel-1
preceding exponents are \(2h-1\) for \(L\downarrow\) and
\(3h-2\) for \(R\uparrow\) after the first removal.  Channel 2 has
exponents \(h-1\) for \(L\uparrow\) and \(4h-2\) for
\(R\downarrow\) after the first removal.  Therefore
\begin{equation}
 \kappa_h(-e_{L\downarrow},R\uparrow)
 \kappa_h(\mathbf0,L\downarrow)
 =\kappa_h(-e_{L\uparrow},R\downarrow)
 \kappa_h(\mathbf0,L\uparrow)=(-1)^{h+1}.
 \label{eq:ch21-pair-general-sign}
\end{equation}
At \(\mathbf N=\mathbf0\) the ground-to-ground coefficients are
\(\pm(-1)^{h+1}/L\) in the two mutually orthogonal targets,
with the overall negative sign applied to channel 2.

\subsection*{Typed finite budgets, the direct-sum target, and exact residuals}

The finite-field product is an ambient endomorphism on \(\mathcal H_L\), meaning \(A_1-A_2\) is well defined there. Trouble arises only if each annihilator is first compressed to a map between \emph{different} charge-sector energy budgets and these maps are then subtracted as if their codomains agreed. The remedy is to rigorously specify each intermediate space, calculate its omitted component, and embed both final targets into the same encompassing Hilbert space.

Let \(\mathcal E(S)\) be the nonnegative integer grading relative to
the ground configuration of the charge sector of \(S\), as in
Chapter~\ref{ch:vacuum_budget}.  For admissible \(\mathbf N\) and
\(J\in\mathbb N\), define the coordinate budget
\begin{equation}
 V_{\mathbf N,J}
  =\operatorname{span}\{\,|S\rangle:\mathbf N(S)=\mathbf N,
                          \ \mathcal E(S)\leq J\,\}
  \subseteq\mathcal H_L.
 \label{eq:ch21-pair-budget}
\end{equation}
Let \(i_{\mathbf N,J}:V_{\mathbf N,J}\hookrightarrow\mathcal H_L\)
be inclusion, \(p_{\mathbf N,J}=i_{\mathbf N,J}^\dagger\) its
orthogonal component projection, and
\(P_{\mathbf N,J}=i_{\mathbf N,J}p_{\mathbf N,J}\) the ambient
orthogonal projector.  A cutoff \(J\) is an excitation-energy budget,
not a bound on the scalar product \(\langle\psi,H\psi\rangle\).

For \(j=1,2\), retain the ordered pair \((a_j,b_j)\) above and choose
\emph{independent} source, intermediate, and target cutoffs
\(J_s,J_{m,j},J_{t,j}\in\mathbb N\).  Define
\begin{equation}
 \begin{aligned}
 V_s&=V_{\mathbf N,J_s},\\
 \mathbf M_j&=\mathbf N-e_{b_j},
    &V_{m,j}&=V_{\mathbf M_j,J_{m,j}},\\
 \mathbf T_j&=\mathbf N-e_{b_j}-e_{a_j},
    &V_{t,j}&=V_{\mathbf T_j,J_{t,j}}.
 \end{aligned}
 \label{eq:ch21-pair-three-spaces}
\end{equation}
If any shifted charge vector is not admissible, its corresponding
budget is interpreted as the zero subspace of \(\mathcal H_L\).
Write \(i_s,p_{m,j},i_{m,j},p_{t,j}\) for the associated inclusion
and coordinate projections.  Finally form the \emph{orthogonal sum}
\begin{equation}
 W=V_{t,1}\oplus V_{t,2},\qquad
 j_j:V_{t,j}\longrightarrow W,\qquad
 p_W:\mathcal H_L\longrightarrow W.
 \label{eq:ch21-pair-common-target}
\end{equation}
Here \(W\) is identified isometrically with the subspace
\(V_{t,1}+V_{t,2}\subseteq\mathcal H_L\); the sum is orthogonal by
Lemma~\ref{lem:ch21-singlet-shifts}.  Thus each \(j_j\) is the
canonical inclusion of one target into the common codomain.

\begin{theorem}[Exact two-channel projected singlet identity]
\label{thm:ch21-pair-projected}
For every finite \(L=2h\), site \(x\), admissible source charge,
and chosen nonnegative cutoffs, define the typed field maps
\begin{equation}
 \begin{aligned}
 B_{b,j}&=p_{m,j}c_{b_j,x}i_s:
                  V_s\longrightarrow V_{m,j},\\
 B_{a,j}&=p_{t,j}c_{a_j,x}i_{m,j}:
                  V_{m,j}\longrightarrow V_{t,j},\\
 \mathcal R_j&=p_{t,j}c_{a_j,x}
       (I-P_{m,j})c_{b_j,x}i_s:
                  V_s\longrightarrow V_{t,j},\\
 P_{m,j}&=i_{m,j}p_{m,j}:
                  \mathcal H_L\longrightarrow\mathcal H_L.
 \end{aligned}
 \label{eq:ch21-pair-typed-maps}
\end{equation}
Then, as an equality of linear maps \(V_s\longrightarrow W\),
\begin{equation}
 \boxed{\displaystyle
 p_W O_{SSC}(x)i_s
  =j_1\!\left(B_{a,1}\circ B_{b,1}+\mathcal R_1\right)
   -j_2\!\left(B_{a,2}\circ B_{b,2}+\mathcal R_2\right).}
 \label{eq:ch21-pair-typed-identity}
\end{equation}
No intermediate or final cutoff inequality is required.  In
particular, \(\mathcal R_1\) and \(\mathcal R_2\) are separately
specified, and neither is silently set to zero.
\end{theorem}

\begin{proof}
Fix a channel \(j\).  The insertion of the ambient orthogonal
identity \(I=P_{m,j}+(I-P_{m,j})\) \emph{between the two site
annihilators} is exact:
\begin{equation}
 \begin{aligned}
 p_{t,j}c_{a_j,x}c_{b_j,x}i_s
 &=p_{t,j}c_{a_j,x}
       P_{m,j}c_{b_j,x}i_s
   +p_{t,j}c_{a_j,x}(I-P_{m,j})c_{b_j,x}i_s\\
 &=\bigl(p_{t,j}c_{a_j,x}i_{m,j}\bigr)
       \bigl(p_{m,j}c_{b_j,x}i_s\bigr)
      +\mathcal R_j\\
 &=B_{a,j}\circ B_{b,j}+\mathcal R_j.
 \end{aligned}
 \label{eq:ch21-pair-insert-projection}
\end{equation}
This equality holds inside \(\operatorname{Hom}(V_s,V_{t,j})\).
By Lemma~\ref{lem:ch21-singlet-shifts}, the image of
\(c_{a_j,x}c_{b_j,x}i_s\) lies entirely in the
\emph{full} charge sector \(\mathcal H_{\mathbf T_j}\).
Since \(V_{t,1}\) and \(V_{t,2}\) occupy different charge sectors,
projecting the channel-\(j\) image to \(W\) yields precisely
\(j_jp_{t,j}c_{a_j,x}c_{b_j,x}i_s\); the projection to the other
target is zero.  Inserting the two channel identities
\eqref{eq:ch21-pair-insert-projection} into
\(p_W(A_1-A_2)i_s\) proves
\eqref{eq:ch21-pair-typed-identity}.  The two terms are subtracted
\emph{after} composition with \(j_1\) and \(j_2\), so their
codomain is genuinely the same \(W\).  This completes the proof.
\end{proof}

The theorem is a finite-Fock identity, not an approximate bosonization
formula.  Both residuals represent concrete vectors that passed through
the \emph{complement} of the selected intermediate energy budget before
being annihilated by the second field.  In particular,
\begin{equation}
 \mathcal R_j=0
 \quad\Longleftrightarrow\quad
 p_{t,j}c_{a_j,x}(I-P_{m,j})c_{b_j,x}i_s=0.
 \label{eq:ch21-pair-residual-criterion}
\end{equation}
The stronger range condition
\((I-P_{m,j})c_{b_j,x}i_s=0\) suffices, but is not necessary:
the second field and final projection can remove all discarded
components even when the intermediate field leaves the budget.  A
single-field bosonization candidate from Core Chapter~14 may replace
\(B_{a,j}\) or \(B_{b,j}\) only after its \emph{individual}
finite-budget identity and margins have been established.  A verified product of independently
compressed candidates is not automatically the compressed physical
product: precisely the \(\mathcal R_j\) term measures the difference.

\begin{example}[A nonzero intermediate-projection residual]
\label{ex:ch21-pair-residual}
Return to \(h=2\), \(L=4\), \(x=1\), source charge
\(\mathbf N=\mathbf0\) and source budget \(J_s=0\).  In channel 1,
take \(J_{m,1}=0\) but \(J_{t,1}=1\).  Starting from the sea
\(\Omega_{\mathbf0}\), remove momentum \(k=-1\), rather than the
top momentum \(k=0\), in species \(L\downarrow\).  The resulting
species occupation is \(\{0\}\), whereas the ground occupation of
relative charge \(-1\) is \(\{-1\}\).  It therefore has one unit
of excitation energy and is annihilated by the intermediate projection
\(P_{m,1}\).

Next remove the top momentum \(k=0\) in species \(R\uparrow\).
The final occupation ket \(\Xi\) has charge
\(\mathbf T_1=-e_{L\downarrow}-e_{R\uparrow}\), one excitation
unit in \(L\downarrow\), and zero excitation in every other
species, so \(\Xi\in V_{t,1}\).  No other momentum pair
produces this exact configuration.  The preceding counts are \(2\)
for the first \(L\downarrow\) removal and \(4\) for the
\(R\uparrow\) removal; both CAR signs are positive.  The two
Fourier factors contribute
\(L^{-1}\zeta^{-1}=(-i)/4\).  Consequently,
\begin{equation}
 \left\langle\Xi,\mathcal R_1\Omega_{\mathbf0}\right\rangle
     =-\frac{i}{4}\ne0,\qquad
 \left\langle\Xi,(B_{a,1}\circ B_{b,1})
          \Omega_{\mathbf0}\right\rangle=0.
 \label{eq:ch21-pair-residual-witness}
\end{equation}
Thus even a source \emph{ground state} does not make a truncated
intermediate projection exact for an arbitrary final cutoff.  Taking
\(J_{m,1}=1\) retains the displayed intermediate ket, but that alone
is not a proof that the entire residual vanishes for all source
vectors and all final cutoffs.
\end{example}

\begin{remark}[An enlarged final codomain]
\label{rem:ch21-pair-enlarged-target}
Equation~\eqref{eq:ch21-pair-typed-identity} uses \(W\) \emph{equal to}
the indicated two-budget orthogonal sum.  If \(W'\) is any larger
closed subspace containing \(W\), with inclusion
\(\iota:W\hookrightarrow W'\), its orthogonal projection satisfies
\begin{equation}
 p_{W'}O_{SSC}(x)i_s
  =\iota\,p_WO_{SSC}(x)i_s
     +p_{W'}(I-P_W)O_{SSC}(x)i_s.
 \label{eq:ch21-pair-larger-target}
\end{equation}
The last term includes whatever part of either channel lies in the
larger target outside the two selected budgets; it must be retained.
It vanishes if \(W'=W\), or under an independently proved range
condition.  Thus merely saying that \(W'\) ``contains both target
budgets'' is insufficient to infer the same two-term formula.
\end{remark}

\subsection*{What is established, and what remains conditional}

In summary, we have derived three logically separate statements.  First,
\eqref{eq:ch21-pair-shift-one}--\eqref{eq:ch21-pair-targets} are exact
species-charge selection rules in the full finite CAR algebra.  They
show why the singlet has two, rather than one, typed charge channels.
Second, \eqref{eq:ch21-pair-klein-phase}--\eqref{eq:ch21-pair-general-sign}
compute their ordered fermionic phases on charge-sector ground kets;
these phases are the ground-state content of the finite signed Klein
construction.  Third, \eqref{eq:ch21-pair-typed-identity} is an exact
linear-map identity on arbitrary finite coordinate budgets, and
Example~\ref{ex:ch21-pair-residual} shows why its two residuals cannot
in general be omitted.

In a charge/spin collective-field description, one may substitute the
algebraic inverse currents
\(\rho_{\eta\uparrow}=(R^c_\eta+R^s_\eta)/2\) and
\(\rho_{\eta\downarrow}=(R^c_\eta-R^s_\eta)/2\) into candidate
oscillator phases.  That substitution is \emph{not} by itself a proof
that compressed exponentials factor into a single common charge factor
and a common spin factor.  The two singlet channels have different
Klein products, different charge-sector targets, and potentially
different intermediate-budget errors.  Any such rewriting must specify
the uncompressed algebra in which exponentials are defined, prove the
required commutation and ordering identities, and establish their
compatibility with projection.  The formal series of
Section~\ref{sec:ch20-formal-vertices} likewise cannot be evaluated
at a numerical parameter or identified with finite-CAR expectations
without further arguments.

No thermodynamic-limit correlation exponent, spectral gap, or
interacting-state identification follows from the operator identities
of this section.  The exact CAR results used here rest on the
previously constructed finite carrier and its supported Klein and
coordinate-projection results; the spinful Chapter~21 implementations
and the present manuscript remain to be independently formalized,
compiled against the full book, and reviewed.  In the next section,
Umklapp scattering changes four species charges at once, making the
same sector-typing and projection questions essential rather than
merely notational.

\section{Umklapp Scattering, Charge-Sector Leakage, and Finite Projection}
\label{sec:ch21-umklapp}

\subsection*{From a four-fermion process to a precise invariance question}

In the preceding sections the charge and spin currents were built from four
ordered fermion species, and the spin-singlet pair was found to have two
separate species-charge targets.  We now confront a related question for a
Hamiltonian rather than an order parameter: \emph{does a chosen finite
low-energy space remain invariant under a local interaction?}  An answer
requires more than the name ``Umklapp.''  It requires an ordered CAR word,
an explicit charge shift, and a state on which the operator acts nontrivially
\emph{outside} the chosen space.  Pauli exclusion can make a seemingly
outward process vanish; the opposite process may remain nonzero while staying
inside the permitted charge box.

The discussion is entirely on a finite lattice and uses the occupation
representation from Chapter~\ref{ch:car_fock}.  We prove exact selection rules,
self-adjointness, and a conditional failure of invariance.  These statements
do not imply a thermodynamic charge gap or establish continuum renormalization
group equations.  For physical motivation and conventional bosonized models,
see, for example, \cite{giamarchi2004,gogolin1998}; our claims below are
finite algebraic statements with their hypotheses displayed.

As in Sections~\ref{sec:ch21-spin-charge} and
\ref{sec:ch21-singlet-pair}, choose \(h\in\mathbb N_{>0}\) and set
\begin{equation}
  L=2h,\qquad \mathcal B_L=(-h,h]\cap\mathbb Z
       =\{1-h,\ldots,h\},\qquad
  \Lambda_L=\mathbb Z/L\mathbb Z,\qquad \zeta=e^{2\pi i/L}.
  \label{eq:ch21-u-band}
\end{equation}
The four species have the fixed occupation ordering
\begin{equation}
 L\uparrow<L\downarrow<R\uparrow<R\downarrow,
 \quad\text{with increasing momentum inside each species}.
 \label{eq:ch21-u-order}
\end{equation}
For \(\alpha\in\mathcal C=\{L,R\}\times\{\uparrow,\downarrow\}\),
let \(c_{\alpha,k}\) satisfy the finite canonical anticommutation relations
(CAR).  We use periodic position fields in the Fourier convention
\begin{equation}
 c_{\alpha,x}=\frac1{\sqrt L}\sum_{k\in\mathcal B_L}
                \zeta^{kx}c_{\alpha,k},\qquad
 c_{\alpha,x}^{\dagger}=\frac1{\sqrt L}\sum_{k\in\mathcal B_L}
                \zeta^{-kx}c_{\alpha,k}^{\dagger}.
 \label{eq:ch21-u-site-fields}
\end{equation}
The positive-mode cutoff, if density excitations are considered separately,
is denoted \(M\) (with \(1\leq m\leq M\)); no density-mode or
current-CCR cutoff enters the CAR calculation in this section.

\subsection*{The ordered Umklapp word and its Hermitian Hamiltonian}

\begin{definition}[Finite-site Umklapp interaction]
\label{def:ch21-u-word}
Define the ordered product
\begin{equation}
  O_U(x)=c_{L\uparrow,x}^{\dagger}
         c_{L\downarrow,x}^{\dagger}
         c_{R\downarrow,x}c_{R\uparrow,x},
       \qquad x\in\Lambda_L.
  \label{eq:ch21-u-word}
\end{equation}
For a \emph{real} coupling \(g_U\in\mathbb R\), set
\begin{equation}
 H_U=\frac{g_U}{L}\sum_{x\in\Lambda_L}
           \bigl(O_U(x)+O_U(x)^\dagger\bigr).
 \label{eq:ch21-u-hamiltonian}
\end{equation}
Products act from right to left: the right-up and right-down particles
are removed, in that order, before the left-down and left-up particles are
created.  The factor \(1/L\) in \eqref{eq:ch21-u-hamiltonian} is fixed
throughout, including the finite coefficient examples below.
\end{definition}

In the finite Fock representation these are bounded operators on the
same Hilbert space \(\mathcal H_L\), so no operator-domain problem arises.
The adjoint reverses \emph{the order} of the four factors:
\begin{equation}
 O_U(x)^\dagger
  =c_{R\uparrow,x}^\dagger c_{R\downarrow,x}^\dagger
    c_{L\downarrow,x}c_{L\uparrow,x}.
 \label{eq:ch21-u-adjoint}
\end{equation}
It is not permissible to reorder these odd operators without keeping
the resulting CAR signs.

\begin{proposition}[Self-adjointness]
\label{prop:ch21-u-selfadjoint}
For every \(h>0\) and real \(g_U\), the operator \(H_U\) in
\eqref{eq:ch21-u-hamiltonian} satisfies \(H_U^\dagger=H_U\).
\end{proposition}
\begin{proof}
Write \(A=\sum_xO_U(x)\).  Since the sum is finite,
\(A^\dagger=\sum_xO_U(x)^\dagger\); hence
\[
 H_U^\dagger
 =\overline{\frac{g_U}{L}}(A^\dagger+A)
 =\frac{g_U}{L}(A+A^\dagger)=H_U.
\]
The middle equality uses the reality of \(g_U\).  No commutation
between \(A\) and \(A^\dagger\) was assumed.
\end{proof}

An Umklapp word is a possible lattice process rather than a statement of
its nonzero action on all states.  Its momentum-space expansion follows
exactly from \eqref{eq:ch21-u-site-fields}:
\begin{align}
 O_U(x)=\frac1{L^2}\!
 \sum_{k_1,k_2,k_3,k_4\in\mathcal B_L}
 &\zeta^{(-k_1-k_2+k_3+k_4)x}\,
 c_{L\uparrow,k_1}^\dagger c_{L\downarrow,k_2}^\dagger
 c_{R\downarrow,k_3}c_{R\uparrow,k_4}.
 \label{eq:ch21-u-fourier-expansion}
\end{align}
Summation over \(x\) retains precisely those ordered words satisfying
\(k_1+k_2-k_3-k_4\equiv0\pmod L\).  Thus momentum is conserved
\emph{modulo} the reciprocal-lattice period, rather than as an
unrestricted equality between the chosen integer representatives.
This distinction will be essential for the explicit witness.

\subsection*{Charge shifts and the surviving conservation laws}

The relative species number operators are
\begin{equation}
 \widehat N_\alpha=\sum_{k\in\mathcal B_L}
     c_{\alpha,k}^\dagger c_{\alpha,k}-hI,
 \qquad \alpha\in\mathcal C.
 \label{eq:ch21-u-relative-charge}
\end{equation}
They have integer eigenvalues \(-h,\ldots,h\) on the finite Fock space.
Let \(e_\alpha\) denote a basis vector in the four-dimensional
species-charge lattice.  With the order \eqref{eq:ch21-u-order} we define
\begin{equation}
 \Delta\mathbf N
  =e_{L\uparrow}+e_{L\downarrow}
     -e_{R\uparrow}-e_{R\downarrow}
  =(+1,+1,-1,-1).
 \label{eq:ch21-u-charge-vector}
\end{equation}
Every \emph{species} shift has magnitude one, whereas each chiral
\emph{branch total} changes by two.

\begin{lemma}[One-particle charge commutators]
\label{lem:ch21-u-fundamental-commutators}
For any species \(\gamma,\beta\), momentum \(q\in\mathcal B_L\), and
site \(x\),
\begin{align}
 [\widehat N_\gamma,c_{\beta,q}]
   &=-\delta_{\gamma\beta}c_{\beta,q},&
 [\widehat N_\gamma,c_{\beta,q}^\dagger]
   &=\delta_{\gamma\beta}c_{\beta,q}^\dagger,
 \label{eq:ch21-u-elementary-commutators}\\
 [\widehat N_\gamma,c_{\beta,x}]
   &=-\delta_{\gamma\beta}c_{\beta,x},&
 [\widehat N_\gamma,c_{\beta,x}^\dagger]
   &=\delta_{\gamma\beta}c_{\beta,x}^\dagger.
 \label{eq:ch21-u-site-commutators}
\end{align}
\end{lemma}
\begin{proof}
The CAR imply, for individual modes \(i,j\),
\[
 [c_i^\dagger c_i,c_j]
   =c_i^\dagger\{c_i,c_j\}-\{c_i^\dagger,c_j\}c_i
   =-\delta_{ij}c_j,
\]
and similarly \([c_i^\dagger c_i,c_j^\dagger]=\delta_{ij}c_j^\dagger\).
Summing over the momenta of species \(\gamma\) proves
\eqref{eq:ch21-u-elementary-commutators}; the scalar sea subtraction
commutes with everything.  Fourier linearity gives
\eqref{eq:ch21-u-site-commutators}.
\end{proof}

\begin{theorem}[Exact four-species Umklapp selection rule]
\label{thm:ch21-u-charge-shift}
For every \(x\) and species \(\gamma\),
\begin{align}
 [\widehat N_\gamma,O_U(x)]
   &=\Delta N_\gamma\,O_U(x),
 \label{eq:ch21-u-commutator}\\
 [\widehat N_\gamma,O_U(x)^\dagger]
   &=-\Delta N_\gamma\,O_U(x)^\dagger.
 \label{eq:ch21-u-adjoint-commutator}
\end{align}
Consequently, if \(\mathcal H_{\mathbf N}\) is the simultaneous
relative-charge eigenspace of \(\mathbf N\in\mathbb Z^{\mathcal C}\), then
\begin{equation}
 O_U(x)\mathcal H_{\mathbf N}\subseteq
     \mathcal H_{\mathbf N+\Delta\mathbf N},\qquad
 O_U(x)^\dagger\mathcal H_{\mathbf N}\subseteq
     \mathcal H_{\mathbf N-\Delta\mathbf N}.
 \label{eq:ch21-u-charge-targets}
\end{equation}
If a listed charge is outside the allowed finite range, its eigenspace
is understood to be zero.  The total relative particle number and the
spin imbalance satisfy
\begin{equation}
 \left[\sum_{\alpha\in\mathcal C}\widehat N_\alpha,H_U\right]=0,
 \qquad
 [\widehat S,H_U]=0,
 \quad \widehat S=\sum_{\eta=L,R}
        (\widehat N_{\eta\uparrow}-\widehat N_{\eta\downarrow}).
 \label{eq:ch21-u-conserved-totals}
\end{equation}
By contrast, \(H_U\) need not conserve each species charge or the
separate left/right branch charge.
\end{theorem}
\begin{proof}
For arbitrary operators, the commutator is a derivation:
\([D,AB]=[D,A]B+A[D,B]\).  Expanding the ordered word gives
\begin{align*}
 [\widehat N_\gamma,O_U(x)]={}&
   [\widehat N_\gamma,c_{L\uparrow,x}^\dagger]
           c_{L\downarrow,x}^\dagger c_{R\downarrow,x}c_{R\uparrow,x}
 +c_{L\uparrow,x}^\dagger
   [\widehat N_\gamma,c_{L\downarrow,x}^\dagger]
           c_{R\downarrow,x}c_{R\uparrow,x}\\
 &+c_{L\uparrow,x}^\dagger c_{L\downarrow,x}^\dagger
   [\widehat N_\gamma,c_{R\downarrow,x}]c_{R\uparrow,x}
 +c_{L\uparrow,x}^\dagger c_{L\downarrow,x}^\dagger
   c_{R\downarrow,x}[\widehat N_\gamma,c_{R\uparrow,x}].
\end{align*}
Lemma~\ref{lem:ch21-u-fundamental-commutators} reduces the coefficient
to
\(\delta_{\gamma,L\uparrow}+\delta_{\gamma,L\downarrow}
-\delta_{\gamma,R\downarrow}-\delta_{\gamma,R\uparrow}\), which is
exactly \(\Delta N_\gamma\).  Taking adjoints of the operator
commutator, using self-adjointness of \(\widehat N_\gamma\), gives
\eqref{eq:ch21-u-adjoint-commutator}.  For
\(\widehat N_\gamma\psi=N_\gamma\psi\),
\[
 \widehat N_\gamma O_U(x)\psi
  =(N_\gamma+\Delta N_\gamma)O_U(x)\psi,
\]
proving the first inclusion in \eqref{eq:ch21-u-charge-targets};
the second is identical with the opposite shift.  Finally the
coefficients of the total-number commutator sum to \(1+1-1-1=0\),
while those of \(\widehat S\) sum to \((1-1)+(-1+1)=0\).
The identities persist under the finite site sum and the adjoint.
\end{proof}

Writing \(\widehat Q_c^\eta=\widehat N_{\eta\uparrow}
 +\widehat N_{\eta\downarrow}\) and
\(\widehat Q_s^\eta=\widehat N_{\eta\uparrow}
 -\widehat N_{\eta\downarrow}\) as in \S\,21.1, the same computation
immediately gives
\begin{equation}
 [\widehat Q_c^L,O_U]=2O_U,\qquad
 [\widehat Q_c^R,O_U]=-2O_U,\qquad
 [\widehat Q_s^L,O_U]=[\widehat Q_s^R,O_U]=0.
 \label{eq:ch21-u-branch-shifts}
\end{equation}
Here \(O_U\) denotes either the fixed-site word or its outward site sum.
These equations describe an interbranch transfer of two particles,
not the creation or destruction of total electric charge.

\begin{corollary}[Vanishing same-sector compression]
\label{cor:ch21-u-sector-compression}
For any admissible charge vector \(\mathbf N\) and any two nonnegative
energy cutoffs \(K,K'\),
\begin{equation}
 P_{\mathbf N,K}\,H_U\,P_{\mathbf N,K'}=0.
 \label{eq:ch21-u-zero-compression}
\end{equation}
\end{corollary}
\begin{proof}
Both projected source and projected target are contained in
\(\mathcal H_{\mathbf N}\).  Theorem~\ref{thm:ch21-u-charge-shift}
places the two summands of \(H_U\) in
\(\mathcal H_{\mathbf N+\Delta\mathbf N}\) and
\(\mathcal H_{\mathbf N-\Delta\mathbf N}\), respectively.  Neither
charge vector equals \(\mathbf N\), and joint eigenspaces with
different charge labels are orthogonal.  The left projection therefore
kills both contributions, independently of \(K,K'\).
\end{proof}

\subsection*{Charge boxes, projections, and a conditional leakage theorem}

A single-charge-sector projection is not the same as a charge
\emph{box}, which admits several sectors but limits their extent.
To define the latter without using an informal inequality between
vectors, let \(S\subseteq\mathcal C\times\mathcal B_L\) be an
occupation configuration.  Its relative charge vector is
\(\mathbf N(S)\), and its nonnegative integer excitation energy
\(\mathcal E(S)\) is measured above the ground configuration of the
\emph{same} charge vector, as in Chapter~\ref{ch:vacuum_budget}.

\begin{definition}[Per-species charge box and budget]
\label{def:ch21-u-charge-box}
Given bounds \(B_\alpha\in\mathbb N\) and an energy budget
\(K\in\mathbb N\), let
\begin{equation}
 \mathcal B_{K,\mathbf B}
   =\operatorname{span}\Bigl\{|S\rangle:\;
       |N_\alpha(S)|\le B_\alpha\ \text{for all }\alpha,
       \quad\mathcal E(S)\le K\Bigr\}
       \subseteq\mathcal H_L,
 \label{eq:ch21-u-box}
\end{equation}
and let \(P_{K,\mathbf B}\) be its orthogonal coordinate projection.
The charge bounds and the energy cutoff play logically distinct roles.
A box may be proper even when the energy cutoff is large enough to
include every state whose charge vector satisfies the bounds.
\end{definition}

\begin{theorem}[Conditional outward leakage]
\label{thm:ch21-u-leakage}
Assume \(h\ge1\), \(g_U\in\mathbb R\setminus\{0\}\), and
\(\psi\in\mathcal B_{K,\mathbf B}\) has definite charge
\(\mathbf N\), i.e.\ \(\widehat N_\alpha\psi=N_\alpha\psi\)
for all \(\alpha\).  Suppose
\begin{equation}
  \sum_{x\in\Lambda_L}O_U(x)\psi\ne0,
  \qquad
  \exists\alpha\in\mathcal C:\quad
      |N_\alpha+\Delta N_\alpha|>B_\alpha.
 \label{eq:ch21-u-leakage-hypotheses}
\end{equation}
Then \(H_U\psi\notin\mathcal B_{K,\mathbf B}\), and the ambient
Hamiltonian does not preserve this budget.  In addition,
\begin{equation}
 [P_{K,\mathbf B},H_U]\ne0.
 \label{eq:ch21-u-projector-noncommute}
\end{equation}
An analogous conclusion holds for a nonzero adjoint (inward) site sum
whose target \(\mathbf N-\Delta\mathbf N\) violates a charge bound.
\end{theorem}
\begin{proof}
For any \(\mathbf T\), denote the orthogonal full-charge-sector
projection by \(\Pi_{\mathbf T}\).  Distinct sectors are orthogonal
because the \(\widehat N_\alpha\) are commuting self-adjoint matrices.
The two shifts in \eqref{eq:ch21-u-charge-targets} differ, hence
\begin{equation}
 \Pi_{\mathbf N+\Delta\mathbf N}H_U\psi
   =\frac{g_U}{L}\sum_xO_U(x)\psi\ne0.
 \label{eq:ch21-u-projected-outward}
\end{equation}
The adjoint contribution has charge \(\mathbf N-\Delta\mathbf N\)
and cannot cancel the displayed component.  By the second hypothesis,
\(\mathcal B_{K,\mathbf B}\) has no vector in the charge sector
\(\mathbf N+\Delta\mathbf N\).  Equation
\eqref{eq:ch21-u-projected-outward} therefore exhibits a nonzero
component of \((I-P_{K,\mathbf B})H_U\psi\).
Since \(P_{K,\mathbf B}\psi=\psi\), we also have
\[
 [P_{K,\mathbf B},H_U]\psi
  =P_{K,\mathbf B}H_U\psi-H_UP_{K,\mathbf B}\psi
  =-(I-P_{K,\mathbf B})H_U\psi\ne0.
\]
The proof for the opposite shift exchanges \(O_U\) with \(O_U^\dagger\).
\end{proof}

The same elementary argument detects \emph{energy} leakage: it suffices
to exhibit an occupation ket \(|T\rangle\) excluded by the energy
cutoff for which \(\langle T,H_U\psi\rangle\ne0\).
The charge-box hypothesis in Theorem~\ref{thm:ch21-u-leakage} is a
convenient sufficient condition, not a necessary condition for
non-invariance.  Conversely, \(H_U\psi\ne0\) by itself is
\emph{not} a leakage certificate.  The latter distinction will be
visible in the second finite example.

\subsection*{An explicit outward witness: the sign and the Fourier coefficient}

The next calculation checks every part of the leakage hypothesis with
a concrete occupation configuration.  It also illustrates why a
fermionic sign cannot be inferred from charge conservation alone.

\begin{theorem}[Zero-charge sea leaves every zero-charge box]
\label{thm:ch21-u-zero-box-witness}
For every \(h\ge1\), write \(L=2h\).  In each species fill the sea
momenta \(\{1-h,\ldots,0\}\), and let \(\Omega\) be the corresponding
normalized, ordered occupation ket.  Define \(T\) by removing
momentum \(1-h\) from \(R\uparrow\) and \(R\downarrow\), and
inserting momentum \(1\) in \(L\uparrow\) and \(L\downarrow\).
Then
\begin{equation}
 \langle T,H_U\Omega\rangle=-\frac{g_U}{L^2}.
 \label{eq:ch21-u-witness-coefficient}
\end{equation}
If \(g_U\ne0\), the zero-charge sea \(\Omega\) belongs to
\(\mathcal B_{0,\mathbf0}\), but
\(H_U\Omega\notin\mathcal B_{K,\mathbf0}\) for \emph{every}
\(K\ge0\).  Thus \(H_U\) does not preserve the zero-charge budget
and does not commute with its projector.
\end{theorem}
\begin{proof}
Let \(S_0\) be the set of \(4h\) occupied species-momentum modes
constituting \(\Omega\).  The target configuration is
\[
 T=\bigl(S_0\setminus\{(R\uparrow,1-h),(R\downarrow,1-h)\}\bigr)
     \cup\{(L\uparrow,1),(L\downarrow,1)\}.
\]
For \(h\ge1\), the removed momenta lie in the sea and the inserted
momentum \(1\) lies in the band but outside the sea.  Therefore no
Pauli factor kills this particular word.  The relative-charge vector
of \(T\) is \(\Delta\mathbf N=(1,1,-1,-1)\), whereas
\(\mathbf N(\Omega)=\mathbf0\).

For an occupied mode \(i\), annihilation contributes the sign
\((-1)^{\#\{j\in S:j<i\}}\), and creation at an unoccupied mode
contributes the corresponding sign of the updated configuration.
The ordered word \eqref{eq:ch21-u-word} acts from right to left.
The number of occupied modes \emph{preceding} each operation, after
all previous operations have been made, is
\begin{center}
\begin{tabular}{@{}lll@{}}
\toprule
Operation & Preceding occupied modes & Sign \\
\midrule
\(c_{R\uparrow,1-h}\) & \(2h\) & \((-1)^{2h}\) \\
\(c_{R\downarrow,1-h}\) & \(3h-1\) & \((-1)^{3h-1}\) \\
\(c_{L\downarrow,1}^\dagger\) & \(2h\) & \((-1)^{2h}\) \\
\(c_{L\uparrow,1}^\dagger\) & \(h\) & \((-1)^h\) \\
\bottomrule
\end{tabular}
\end{center}
For example, the first removal has the two filled left species before
it; after that removal the right-down species has \(2h+(h-1)\)
occupied modes before its lowest sea momentum.  The insertion in
left-down sees all \(h\) left-up and its own \(h\) filled sea modes,
while the final insertion in left-up sees its original \(h\) sea modes.
The exponent is \(2h+(3h-1)+2h+h=8h-1\), an odd integer.
Thus the selected momentum word produces \(-|T\rangle\).

In \eqref{eq:ch21-u-fourier-expansion}, the four momenta are
\(k_1=k_2=1\) and \(k_3=k_4=1-h\).  Their Fourier phase is
\[
 \zeta^{(-1-1+(1-h)+(1-h))x}
    =\zeta^{-2hx}=\zeta^{-Lx}=1
       \quad(x\in\Lambda_L).
\]
Each site word has scalar Fourier normalization \(L^{-2}\).
The four changed occupations uniquely fix the four participating
momenta: another term in the momentum sum cannot produce \(|T\rangle\)
from \(|\Omega\rangle\).  The adjoint word changes the charge by
\(-\Delta\mathbf N\), hence also cannot contribute to this target.
The \(L\) equal site contributions multiplied by the Hamiltonian
prefactor \(g_U/L\) therefore give
\[
 \langle T,H_U\Omega\rangle
  =\frac{g_U}{L}\,L\,\frac{-1}{L^2}
  =-\frac{g_U}{L^2}.
\]
This is nonzero when \(g_U\ne0\), and \(|T\rangle\) is orthogonal
to every zero-species-charge box.  Theorem~\ref{thm:ch21-u-leakage}
then proves the final assertions.  No frozen-margin assumption for
finite density currents has been used.
\end{proof}

\begin{example}[Four-site value]
\label{ex:ch21-u-four-site}
At \(h=2\), \(L=4\), \(\zeta=i\), and \(g_U=1\),
the sea fills \(\{-1,0\}\) in every species.  The witness target
fills \(\{-1,0,1\}\) in each left species and \(\{0\}\) in
each right species.  Equation~\eqref{eq:ch21-u-witness-coefficient}
gives the exact matrix element
\[
 \langle T,H_U\Omega\rangle=-\frac1{16}.
\]
The charge vector \((1,1,-1,-1)\) already suffices to exclude this
ket from the zero-charge box, independently of the excitation budget.
\end{example}

\subsection*{Why a nonzero action need not leak: a Pauli-blocked counterexample}

\begin{example}[Only the inward channel survives]
\label{ex:ch21-u-inward-only}
Take \(h=2\), \(L=4\), and \(g_U\ne0\).  Choose a normalized occupation
ket \(|S\rangle\) with both left species \emph{fully} occupied,
\(\{-1,0,1,2\}\), and both right species occupied only at \(\{-1\}\).
Its relative charges are
\[
  \mathbf N(S)=(2,2,-1,-1).
\]
Use the proper charge box with bounds
\(\mathbf B=(2,2,1,1)\) in the species order
\((L\uparrow,L\downarrow,R\uparrow,R\downarrow)\), and the energy
cutoff \(K=16\).  The unprojected outward sum annihilates \(|S\rangle\):
\[
 O_U(x)|S\rangle=0\qquad\text{for every }x,
\]
because it would attempt to create a fermion in an already full
left species.  On the other hand, the adjoint sum is nonzero and
shifts the charges to \((1,1,0,0)\), which obey the same bounds.

To exhibit a coefficient, let \(|T'\rangle\) be the occupation ket with
left species \(\{0,1,2\}\), right-up \(\{-1,2\}\), and
right-down \(\{-1,0\}\).  The adjoint word removes \(-1\) from
each left species and inserts \(2\) in right-up and \(0\) in
right-down.  Its right-to-left CAR preceding counts are
\(0,3,8,7\), with an even sum \(18\), so the sign is positive.
The Fourier phase is \(\zeta^{(-1-1-2-0)x}=\zeta^{-4x}=1\).
Exactly as in the preceding proof,
\begin{equation}
 \langle T',H_U S\rangle=\frac{g_U}{16}\ne0.
 \label{eq:ch21-u-inward-coefficient}
\end{equation}
For this four-site band, the excitation energy above a fixed-number
sector ground is at most \(n(4-n)\le4\) for a species with \(n\)
occupied modes; summing four species gives \(\mathcal E\le16\).
Thus \emph{every} ket generated from \(|S\rangle\) by the adjoint
Umklapp term belongs to \(\mathcal B_{16,\mathbf B}\).  We have
\(H_U|S\rangle\ne0\) but \((I-P_{16,\mathbf B})H_U|S\rangle=0\).
This proves only that \emph{this state} does not leak; it does not
establish invariance of the whole box under \(H_U\).
\end{example}

\subsection*{What finite leakage does, and does not, establish}

The exact conclusions of this section can now be separated.  The CAR
calculation proves a four-species selection rule and identifies precisely
which charge sectors an Umklapp interaction couples.  Same-sector
first-order compression vanishes, but a charge-bounded or energy-bounded
subspace can fail to be invariant, as witnessed by a nonzero
\emph{outside} matrix element.  The finite compression
\(P_{K,\mathbf B}H_UP_{K,\mathbf B}\) remains a legitimate Hermitian
matrix on the retained carrier; it is simply not the restriction of an
invariant ambient dynamics when the leakage is nonzero.

The finite energy-shell machinery developed in Section~\ref{sec:ch21-finite-effective}
can subsequently be used to organize perturbative block reduction. Its
quadratic Schrieffer--Wolff/Baker--Campbell--Hausdorff identities concern
stated projectors and polynomial coefficients through a fixed order; they do
not automatically provide an analytic conjugation, error estimate, or
RG semigroup.  In particular, removing a shell of \emph{total excitation
energy} is not identical to removing a particular Fourier momentum
mode.  The shell-compression formalism must be instantiated with its
projector and block-diagonality hypotheses before being applied to
\(H_U\).

Finally, a Mott gap would require a spectral statement for a specified
many-body Hamiltonian, and a physically meaningful thermodynamic claim
would require a controlled sequence of systems as \(L\to\infty\),
with a specified density, coupling scaling, and limiting observable.
Likewise the Kosterlitz--Thouless flow equations belong to a
separately constructed continuum effective theory; the presence of a
nonzero finite matrix element \(-g_U/L^2\) neither derives their beta
functions nor proves gap formation.  Any proposed use of bosonized
charge vertices also requires care with the distinction between a
finite-CAR operator, an abstract oscillator representation, and a
coefficientwise formal series.  Chapter~21 remains unimplemented in
Lean at the time of this draft: prior Core support for CAR, charge
sectors and finite projections is not a verification of the new
Umklapp theorems.

\section{Finite Energy-Shell Projection and Effective Hamiltonians}
\label{sec:ch21-finite-effective}

The preceding sections built finite operators before asking which states they
preserve.  In particular, the Umklapp word has a definite change of four
species charges, but this selection rule is not a spectral statement, and its
nonzero action need not preserve a chosen finite-energy carrier.  We now face
a different, constructive question: \emph{if we intentionally retain only a
low-energy part of a finite Hilbert space, what operator describes its
interaction with the discarded states?}  Simply deleting the discarded
states removes the very virtual excursions that can modify low-energy
matrix elements.  A finite Schrieffer--Wolff transformation supplies a
controlled algebraic answer through second order in a coupling parameter.

This section develops the method rather than presupposing familiarity with
renormalization.  The starting objects are ordinary finite-dimensional
Hilbert spaces, orthogonal projections and matrix commutators.  We shall
separate (i) exact partition identities, (ii) a solvable non-resonant matrix
equation, (iii) an \emph{exact identity between polynomials of degree two},
and (iv) approximation of an analytic unitary rotation.  Their hypotheses
are different.  In particular, none of these steps by itself constructs a
charge gap or a continuum renormalization-group (RG) flow.

\subsection{What does an energy shell discard?}
\label{subsec:ch21-energy-shell}

Fix a finite lattice of length $L=2h$ with $h\in\mathbb N_{>0}$, a
positive-mode cutoff $M\geq1$, a set of species labels (in this chapter,
$\{R,L\}\times\{\uparrow,\downarrow\}$), and an admissible integer
species-charge vector $\boldsymbol N$.  The finite CAR Fock carrier will be
denoted $\mathcal F_L$.  We use the occupation-space excitation grading
$\mathsf E\colon\{\text{occupation configurations}\}\to\mathbb N_0$,
relative to the specified charge-sector ground configuration.  The
charge-sector \emph{budget} is the span
\begin{equation}
\mathcal B_{\boldsymbol N,K}
 :=\operatorname{span}\bigl\{\lvert A\rangle:
 \boldsymbol N(A)=\boldsymbol N,\ \mathsf E(A)\leq K\bigr\}
 \subseteq\mathcal F_L,
\label{eq:ch21-budget-definition}
\end{equation}
where $K\in\mathbb N_0$.  We emphasize that \eqref{eq:ch21-budget-definition}
is a span of occupation kets selected by a scalar \emph{grade}; it is not
a vector inequality of the form $H\psi\leq K\psi$.  The allowed grades,
charges and finite-volume margins are those of the earlier energy-budget
construction (see Chapter~\ref{ch:vacuum_budget}).

For $K\geq1$, let $\Pi_{\boldsymbol N,K}$ be the orthogonal coordinate
projection onto this budget and define the two \emph{ambient} projections
\begin{align}
P&:=\Pi_{\boldsymbol N,K-1},\\
Q&:=\Pi_{\boldsymbol N,K}-\Pi_{\boldsymbol N,K-1}
  =\operatorname{proj}_{\mathcal H_{\boldsymbol N,K}},\qquad
\mathcal H_{\boldsymbol N,K}:=
 \operatorname{span}\{\lvert A\rangle:\boldsymbol N(A)=\boldsymbol N,
                                  \mathsf E(A)=K\}.
\label{eq:ch21-shell-projections}
\end{align}
Both projections act on $\mathcal F_L$, even though their ranges are
subspaces of $\mathcal B_{\boldsymbol N,K}$.  Their defining coordinate
supports are disjoint.

\begin{proposition}[Exact finite shell partition]
\label{prop:ch21-shell-partition}
For $K\geq1$ the ambient projections satisfy
\begin{equation}
 P^2=P=P^\dagger,\quad Q^2=Q=Q^\dagger,
 \quad PQ=QP=0,\quad P+Q=\Pi_{\boldsymbol N,K}.
\label{eq:ch21-ambient-partition}
\end{equation}
Consequently, with $\mathcal H:=\mathcal B_{\boldsymbol N,K}$ as the
\emph{new carrier}, the restricted operators obey $P+Q=I_{\mathcal H}$
and give the orthogonal decomposition
$\mathcal H=\operatorname{ran}P\oplus\operatorname{ran}Q$.
\end{proposition}
\begin{proof}
In the orthonormal occupation basis, every coordinate projector acts as
identity on a selected ket and as zero on the others.  Its square therefore
equals itself, and its matrix is real diagonal, hence self-adjoint.  The
supports $\mathsf E\leq K-1$ and $\mathsf E=K$ are disjoint: applying
$PQ$ or $QP$ to each basis ket gives zero.  Their union is exactly the
support $\mathsf E\leq K$, within the fixed charge sector.  Thus their
sum equals $\Pi_{\boldsymbol N,K}$ on every occupation ket, proving
\eqref{eq:ch21-ambient-partition}.  This last projection is the identity
\emph{only after} restricting the carrier to its range.  Orthogonality and
direct-sum decomposition follow by writing every $\psi\in\mathcal H$ as
$\psi=P\psi+Q\psi$ and observing that the ranges intersect trivially.
\end{proof}

\begin{example}[Removing energy two is not removing mode two]
\label{ex:ch21-shell-versus-mode}
Consider the polynomial oscillator carrier with two positive modes,
$X_1$ and $X_2$, weighted by energies $1$ and $2$, respectively.  For
$M\geq2$, the illustrative filtered monomial basis at budget two is
\begin{equation}
 \mathcal B_2^{\mathrm{osc}}
  =\operatorname{span}\{1,X_1,X_1^2,X_2\},\qquad
 \mathcal B_1^{\mathrm{osc}}=\operatorname{span}\{1,X_1\}.
\label{eq:ch21-osc-shell-example}
\end{equation}
The shell removed by $P$ is
$\operatorname{span}\{X_1^2,X_2\}$.  Deleting the \emph{mode} $X_2$
instead would retain $X_1^2$ and yield
$\operatorname{span}\{1,X_1,X_1^2\}$.  These are distinct projections.
Their difference persists in a concrete finite-Fock realization: under
the charge-admissibility and frozen-margin hypotheses required for the
repeated positive mode-one density excitation, the vector
$\rho_{1,\nu}^{\,2}\lvert\boldsymbol N\rangle_0$ is nonzero, has
excitation grade two, is killed by the budget-one projector, and is
fixed by the exact-grade-two shell projector.  Here $\rho_{1,\nu}$
uses the convention that positive $m$ raises the excitation grade.
The needed nonvanishing is a finite CAR/current statement, not an
assumption that the polynomial oscillator is literally the full Fock
carrier.  The supporting Core A10 result gives this witness with
$h\geq2$ and its stated admissibility/margin conditions.
\end{example}

The letter $K$ will continue to denote the \emph{energy budget}, not a
Luttinger parameter.  The shell splitting does not assert that its
grade is the interacting Hamiltonian's spectrum: the grade is a basis
filtration fixed in advance.  A Hamiltonian may mix these spaces, which
is precisely why an effective Hamiltonian is needed.

\subsection{Block matrices and the denominator hypothesis}
\label{subsec:ch21-block-setup}

Work now on an arbitrary finite-dimensional complex Hilbert space
$\mathcal H=\mathcal P\oplus\mathcal Q$, with orthogonal projections
$P$ and $Q=I_{\mathcal H}-P$.  This could be the restricted budget
carrier from Proposition~\ref{prop:ch21-shell-partition}; alternatively it
could be a different explicitly selected finite spectral window.  Let
\begin{equation}
H(t)=H_0+tV,\qquad H_0^\dagger=H_0,\quad V^\dagger=V,
\quad t\in\mathbb R,\quad [P,H_0]=0.
\label{eq:ch21-H-perturbation}
\end{equation}
The last hypothesis means that $H_0$ preserves $\mathcal P$ and
$\mathcal Q$.  It is not automatic for a chosen interacting baseline.
In particular, an interbranch pairing term can mix bare Haldane
partitions: one must diagonalize the proposed $H_0$ on an appropriate
carrier or use a baseline, such as the free Sugawara Hamiltonian, for
which the requisite invariance actually holds.

Insert $I=P+Q$ to write each operator in a $2\times2$ array of blocks:
\begin{equation}
H_0=\begin{pmatrix}A&0\\0&D\end{pmatrix},\qquad
V=\begin{pmatrix}V_{PP}&B\\B^\dagger&V_{QQ}\end{pmatrix},
\quad A=PH_0P\big|_{\mathcal P},\quad D=QH_0Q\big|_{\mathcal Q},
\label{eq:ch21-block-matrices}
\end{equation}
where $B=PVQ\colon\mathcal Q\to\mathcal P$.  Define
\begin{equation}
 V_{\mathrm d}:=PVP+QVQ,\qquad
 V_{\mathrm o}:=PVQ+QVP,\qquad V=V_{\mathrm d}+V_{\mathrm o}.
\label{eq:ch21-diag-offdiag}
\end{equation}
The subscripts refer to this \emph{chosen decomposition}, not to
momentum conservation or locality.  A block-off-diagonal perturbation
moves a state between retained and discarded sectors.  Its second-order
virtual excursion may start and end in $\mathcal P$ even when the first
order transition does not stay there.

Because $H_0$ is self-adjoint and preserves the two summands, we may
choose an orthonormal eigenbasis
$\{\lvert\lambda\rangle\}_{\lambda\in I_P}$ of $A$ and
$\{\lvert\mu\rangle\}_{\mu\in I_Q}$ of $D$,
with \emph{real} eigenvalues $E_\lambda$ and $E_\mu$.  No relation between
the two eigenvalue sets has yet been assumed.  A virtual transition
requires an energy denominator; it is therefore necessary to state the
non-resonance condition rather than dividing by a possibly zero number.

\begin{definition}[Coupling-wise non-resonance]
\label{def:ch21-coupling-nonresonance}
The pair $(H_0,V)$ is \emph{non-resonant across $P/Q$ on its nonzero
couplings} if
\begin{equation}
 \langle\lambda\lvert V\rvert\mu\rangle\ne0
 \quad\Longrightarrow\quad E_\lambda\ne E_\mu,
 \qquad \lambda\in I_P,\ \mu\in I_Q.
\label{eq:ch21-nonresonance}
\end{equation}
Degenerate pairs with zero cross-block matrix element are allowed.
\end{definition}

The stronger assumption
$\operatorname{spec}(A)\cap\operatorname{spec}(D)=\varnothing$
guarantees invertibility of the Sylvester operator introduced below.
The weaker, coupling-wise assumption is sufficient for \emph{existence}
of the particular solution used in this section, but not generally for
uniqueness.  Since there are finitely many coupled pairs, its nonzero
energy differences have a strictly positive minimum when any such
pairs exist; that finite number need not remain bounded away from zero
as $L$ or the budget varies.

\subsection{Solving the Sylvester equation}
\label{subsec:ch21-sylvester}

We seek a first-order change of basis generated by an anti-Hermitian,
block-off-diagonal operator $S_1$:
\begin{equation}
 S_1^\dagger=-S_1,\qquad PS_1P=QS_1Q=0,
 \qquad [S_1,H_0]=-V_{\mathrm o}.
\label{eq:ch21-generator-equation}
\end{equation}
Anti-Hermiticity ensures $e^{tS_1}$ is unitary for real $t$, while the
commutator condition cancels the first-order coupling between
$\mathcal P$ and $\mathcal Q$.  It is important that the commutator is
ordered $[S_1,H_0]=S_1H_0-H_0S_1$; reversing it reverses the
energy-denominator sign.

Write
$S_1=\bigl(\begin{smallmatrix}0&X\\-X^\dagger&0\end{smallmatrix}\bigr)$,
where $X:\mathcal Q\to\mathcal P$.  Multiplying the block matrices in
\eqref{eq:ch21-block-matrices} gives
\begin{equation}
 [S_1,H_0]_{PQ}=XD-AX=-B
 \quad\Longleftrightarrow\quad AX-XD=B.
\label{eq:ch21-sylvester}
\end{equation}
This is a \emph{Sylvester equation}: a linear equation for an unknown
operator $X$ in which fixed operators multiply it on the left and right.
It is not a division by the matrix $A-D$, since the matrices can act on
different spaces and need not commute.

\begin{theorem}[Finite non-resonant Schrieffer--Wolff generator]
\label{thm:ch21-sw-generator}
Under \eqref{eq:ch21-H-perturbation} and
Definition~\ref{def:ch21-coupling-nonresonance}, define the
cross-block matrix entries
\begin{equation}
\langle\lambda\lvert X\rvert\mu\rangle:=
\begin{cases}
\displaystyle\frac{\langle\lambda\lvert V\rvert\mu\rangle}
 {E_\lambda-E_\mu}, &
       \langle\lambda\lvert V\rvert\mu\rangle\ne0,\\[6pt]
0,&\langle\lambda\lvert V\rvert\mu\rangle=0,
\end{cases}
\qquad \lambda\in I_P,\ \mu\in I_Q.
\label{eq:ch21-X-entries}
\end{equation}
Then $S_1=\bigl(\begin{smallmatrix}0&X\\-X^\dagger&0\end{smallmatrix}\bigr)$
is anti-Hermitian, block-off-diagonal, and solves
$[S_1,H_0]=-V_{\mathrm o}$.  The zero choice on uncoupled pairs remains
well defined even when their energies coincide.
\end{theorem}
\begin{proof}
The denominator is nonzero in the first case by
\eqref{eq:ch21-nonresonance}; the second case involves no division.
In the chosen eigenbases,
\begin{align*}
\langle\lambda\lvert AX-XD\rvert\mu\rangle
 &=(E_\lambda-E_\mu)
     \langle\lambda\lvert X\rvert\mu\rangle\\
 &=\langle\lambda\lvert V\rvert\mu\rangle,
\end{align*}
including when the matrix element of $V$ is zero.  Thus $AX-XD=B$
entrywise.  The opposite off-diagonal block of
$[S_1,H_0]$ is the adjoint of the first block, because
$S_1^\dagger=-S_1$ and $H_0^\dagger=H_0$ make $[S_1,H_0]$ Hermitian.
More explicitly, the opposite entry is
\begin{equation*}
 \langle\mu\lvert S_1\rvert\lambda\rangle
 =-\overline{\langle\lambda\lvert S_1\rvert\mu\rangle}
 =\frac{\langle\mu\lvert V\rvert\lambda\rangle}{E_\mu-E_\lambda}
\end{equation*}
when coupled, and zero otherwise.  Both diagonal blocks of
$[S_1,H_0]$ vanish because $H_0$ is block diagonal and $S_1$ is
off diagonal.  Collecting the four blocks proves the commutator
equation; the displayed block form proves anti-Hermiticity directly.
\end{proof}

\begin{remark}[Why the solution need not be unique]
If $E_\lambda=E_\mu$ for an uncoupled cross-block pair, the
corresponding entry of the homogeneous equation
$AX-XD=0$ is unconstrained.  Setting it to zero is a convenient,
explicit choice, not a uniqueness theorem.  If the spectra of $A$ and
$D$ are disjoint, the homogeneous equation has only the zero solution,
and the constructed $X$ is unique.
\end{remark}

\subsection{The ordered second-order transformation}
\label{subsec:ch21-sw-bch}

For any finite matrices $S$ and $T$, write
$\operatorname{ad}_S(T):=[S,T]$ and
$\operatorname{ad}_S^j(T)$ for its $j$-fold iteration.  We rotate the
Hamiltonian according to
\begin{equation}
 \widetilde H(t):=e^{tS_1}(H_0+tV)e^{-tS_1}.
\label{eq:ch21-sw-rotated}
\end{equation}
The orientation is fixed: we conjugate by $e^{+tS_1}$ on the left.
The finite-matrix exponential exists for every real $t$ and the
rotated Hamiltonian is self-adjoint and isospectral to $H(t)$, even
though the compression $P\widetilde H(t)P$ need not reproduce the exact
low-energy spectrum at a finite, non-small $t$.

To obtain the perturbative expansion, differentiate
$e^{tS_1}Te^{-tS_1}$ at $t=0$.  Its successive derivatives are
$\operatorname{ad}_{S_1}^j(T)$.  Hence the ordered polynomial through
degree two is
\begin{align}
B_2(t):={}&H_0+t\bigl(V+[S_1,H_0]\bigr)\notag\\
 &+t^2\left([S_1,V]
        +\frac12[S_1,[S_1,H_0]]\right).
\label{eq:ch21-bch-two-polynomial}
\end{align}
The equality between $B_2$ and the Taylor polynomial is a finite
algebraic fact, irrespective of whether $t$ is small.  The claim that
it approximates $\widetilde H(t)$ needs a separate remainder statement.

\begin{theorem}[Exact compressed degree-two identity]
\label{thm:ch21-sw-compressed}
Assume that $P,Q$ are complementary orthogonal projections on
$\mathcal H$, $H_0$ commutes with $P$, and $S_1$ obeys
\eqref{eq:ch21-generator-equation}.  Then, as an \emph{exact polynomial
identity} on $\mathcal P=\operatorname{ran}P$,
\begin{equation}
 P B_2(t)P
 =PH_0P+tPVP+\frac{t^2}{2}P[S_1,V]P.
\label{eq:ch21-sw-polynomial}
\end{equation}
The right-hand side is the second-order Schrieffer--Wolff effective
Hamiltonian, denoted $H_{\mathrm{eff}}^{[2]}(t)$.  No exponential
appears in the asserted equality.
\end{theorem}
\begin{proof}
Substitute $[S_1,H_0]=-V_{\mathrm o}$ into
\eqref{eq:ch21-bch-two-polynomial}.  The first-order coefficient is
$V-V_{\mathrm o}=V_{\mathrm d}$; in particular,
$PV_{\mathrm d}P=PVP$.  For the second-order term,
\begin{align*}
[S_1,V]+\tfrac12[S_1,[S_1,H_0]]
 &= [S_1,V]-\tfrac12[S_1,V_{\mathrm o}]\\
 &= [S_1,V_{\mathrm d}]+\tfrac12[S_1,V_{\mathrm o}].
\end{align*}
A product of an off-diagonal and a diagonal block operator is
off-diagonal, so $P[S_1,V_{\mathrm d}]P=0$.  Therefore the compressed
coefficient becomes
$\tfrac12P[S_1,V_{\mathrm o}]P$.  Moreover,
$P[S_1,V]P=P[S_1,V_{\mathrm o}]P$ by the same block argument.
The constant term is $PH_0P$, yielding exactly
\eqref{eq:ch21-sw-polynomial}.  Every equality used here is an
ordinary finite-matrix identity, with no perturbative limiting
argument.
\end{proof}

\begin{proposition}[Intermediate-state formula]
\label{prop:ch21-sw-intermediate}
For $\lambda,\lambda'\in I_P$ the second-order correction has matrix
entries
\begin{equation}
\begin{aligned}
\langle\lambda\lvert W^{(2)}\rvert\lambda'\rangle
  &=\frac12\sum_{\substack{\mu\in I_Q\\
       V_{\lambda\mu}V_{\mu\lambda'}\ne0}}
  V_{\lambda\mu}V_{\mu\lambda'}
  \left(\frac{1}{E_\lambda-E_\mu}
        +\frac{1}{E_{\lambda'}-E_\mu}\right),\\
W^{(2)}&:=\frac12P[S_1,V]P,
\qquad V_{ab}:=\langle a\lvert V\rvert b\rangle.
\end{aligned}
\label{eq:ch21-sw-intermediate-sum}
\end{equation}
Only nonzero coupled terms are included in the displayed sum, so its
denominators are defined under \eqref{eq:ch21-nonresonance}.  In
particular,
\begin{equation}
\langle\lambda\lvert W^{(2)}\rvert\lambda\rangle
 =\sum_{\substack{\mu\in I_Q\\V_{\lambda\mu}\ne0}}
 \frac{|V_{\lambda\mu}|^2}{E_\lambda-E_\mu}.
\label{eq:ch21-sw-diagonal-shift}
\end{equation}
\end{proposition}
\begin{proof}
Since $S_1$ is off diagonal, the $PP$ block of $S_1V$ is
$S_{1,PQ}V_{QP}$ and that of $VS_1$ is $V_{PQ}S_{1,QP}$.
Insert the orthonormal resolution $I_{\mathcal Q}=\sum_{\mu\in I_Q}
\lvert\mu\rangle\langle\mu\rvert$ to get
\begin{equation*}
\langle\lambda\lvert[S_1,V]\rvert\lambda'\rangle
=\sum_{\mu\in I_Q}\bigl((S_1)_{\lambda\mu}
 V_{\mu\lambda'}-V_{\lambda\mu}(S_1)_{\mu\lambda'}\bigr).
\end{equation*}
Theorem~\ref{thm:ch21-sw-generator} gives
$(S_1)_{\lambda\mu}=V_{\lambda\mu}/(E_\lambda-E_\mu)$
on coupled pairs and
$(S_1)_{\mu\lambda'}=-V_{\mu\lambda'}/(E_{\lambda'}-E_\mu)$.
Terms with vanishing couplings contribute zero by definition.  This
yields \eqref{eq:ch21-sw-intermediate-sum}.  Taking
$\lambda'=\lambda$ and using Hermiticity of $V$ proves
\eqref{eq:ch21-sw-diagonal-shift}.
\end{proof}

The denominator in \eqref{eq:ch21-sw-diagonal-shift} explains the
physical interpretation: the retained state makes a virtual excursion
to a discarded eigenstate and returns.  In a strictly ordered baseline
spectrum $E_\lambda<E_\mu$ for every coupled pair, each diagonal
second-order shift is nonpositive.  Without that ordering, cancellations
are possible.  Thus even the \emph{sign} of the correction is not
determined by the mere existence of a $P/Q$ splitting.

\subsection{What is the remainder, exactly?}
\label{subsec:ch21-sw-remainder}

There are two different uses of the notation $O(t^3)$.  In a formal
power series ring $\operatorname{End}(\mathcal H)[[t]]$, the statement
\begin{equation}
P\widetilde H(t)P\equiv H_{\mathrm{eff}}^{[2]}(t)
 \pmod{t^3}
\label{eq:ch21-sw-formal-mod}
\end{equation}
means \emph{coefficientwise} equality in degrees zero, one and two.
It does not evaluate an operator at a numerical $t$, and it does not
by itself bound the discarded terms.  For the finite analytic
exponential we can additionally obtain an explicit bound; this is
an elementary finite-matrix consequence, stated here separately from
the supplied A10 Lean theorem.

\begin{proposition}[Finite analytic Taylor remainder]
\label{prop:ch21-sw-analytic-remainder}
Under the hypotheses of Theorem~\ref{thm:ch21-sw-compressed}, for real
$t$ let
\begin{equation}
 R_3(t):=\widetilde H(t)-B_2(t).
\label{eq:ch21-rem-definition}
\end{equation}
Then the exact compressed analytic Hamiltonian obeys
\begin{equation}
 P\widetilde H(t)P=H_{\mathrm{eff}}^{[2]}(t)+PR_3(t)P,
\label{eq:ch21-sw-exact-plus-remainder}
\end{equation}
where, writing $\mathcal U_r(T)=e^{rS_1}Te^{-rS_1}$,
\begin{align}
R_3(t)={}&\frac{t^3}{2}\int_0^1(1-u)^2
  \mathcal U_{ut}\bigl(\operatorname{ad}_{S_1}^3(H_0)\bigr)\,du\notag\\
 &+t^3\int_0^1(1-u)
  \mathcal U_{ut}\bigl(\operatorname{ad}_{S_1}^2(V)\bigr)\,du.
\label{eq:ch21-sw-integral-remainder}
\end{align}
For the operator norm, using $S_1^\dagger=-S_1$,
\begin{equation}
\|PR_3(t)P\|\leq\|R_3(t)\|
\leq |t|^3\left(
\frac{\|\operatorname{ad}_{S_1}^3(H_0)\|}{6}
+\frac{\|\operatorname{ad}_{S_1}^2(V)\|}{2}\right).
\label{eq:ch21-sw-norm-bound}
\end{equation}
The bound is finite-volume and its constant depends on the chosen
carrier, $H_0,V$ and the selected generator $S_1$.
\end{proposition}
\begin{proof}
For a fixed $T$, put $F_T(t)=e^{tS_1}Te^{-tS_1}$.
Differentiating the matrix exponential gives
$F_T^{(j)}(r)=\mathcal U_r(\operatorname{ad}_{S_1}^jT)$.
Taylor's theorem with integral remainder through order two for
$F_{H_0}$ gives
\begin{equation*}
F_{H_0}(t)=H_0+t\operatorname{ad}_{S_1}(H_0)
+\frac{t^2}{2}\operatorname{ad}_{S_1}^2(H_0)
+\frac{t^3}{2}\int_0^1(1-u)^2
  \mathcal U_{ut}(\operatorname{ad}_{S_1}^3H_0)\,du.
\end{equation*}
Taylor's theorem through order one for $F_V$ gives
\begin{equation*}
F_V(t)=V+t\operatorname{ad}_{S_1}(V)
+t^2\int_0^1(1-u)
 \mathcal U_{ut}(\operatorname{ad}_{S_1}^2V)\,du.
\end{equation*}
Add $F_{H_0}(t)+tF_V(t)$ and compare with
\eqref{eq:ch21-bch-two-polynomial}; this proves
\eqref{eq:ch21-sw-integral-remainder} and then
\eqref{eq:ch21-sw-exact-plus-remainder}.
For real $t$ and $u$, $e^{utS_1}$ is unitary, so
$\|\mathcal U_{ut}(T)\|=\|T\|$.  Apply the triangle inequality and
$\int_0^1(1-u)^2du=1/3$ and
$\int_0^1(1-u)du=1/2$ to obtain
\eqref{eq:ch21-sw-norm-bound}.  Finally, orthogonal projections are
contractions, hence $\|PR_3P\|\leq\|R_3\|$.
\end{proof}

\begin{remark}[A finite error is not a uniform scaling estimate]
The right side of \eqref{eq:ch21-sw-norm-bound} may increase with
lattice size or as the non-resonant denominators become small.  It is
a genuine estimate at fixed finite carrier, but not a proof of an
$L\to\infty$ limit or a scale-independent perturbative expansion.
The coefficientwise polynomial result of the supplied Core A10 module
should not be reported as formal verification of this analytic bound.
\end{remark}

\subsection{A two-level calculation, including what fails at finite coupling}
\label{subsec:ch21-two-level}

\begin{example}[Explicit Schrieffer--Wolff rotation]
\label{ex:ch21-sw-two-level}
Let $\mathcal H=\mathbb C^2$, $P=\operatorname{diag}(1,0)$ and
$Q=\operatorname{diag}(0,1)$.  Take
\begin{equation}
 H_0=\begin{pmatrix}0&0\\0&1\end{pmatrix},\qquad
 V=\begin{pmatrix}0&1\\1&0\end{pmatrix}.
\label{eq:ch21-sw-2x2-input}
\end{equation}
The low state has energy zero and the high state energy one.  The
Sylvester denominator is $0-1=-1$, so
\begin{equation}
 S_1=\begin{pmatrix}0&-1\\1&0\end{pmatrix},\qquad
 [S_1,H_0]=-V,
 \quad[S_1,V]=\begin{pmatrix}-2&0\\0&2\end{pmatrix}.
\label{eq:ch21-sw-2x2-generator}
\end{equation}
The effective Hamiltonian is therefore the scalar operator
\begin{equation}
 H_{\mathrm{eff}}^{[2]}(t)=-t^2 P
 \quad\text{on }\operatorname{ran}P.
\label{eq:ch21-sw-2x2-effective}
\end{equation}
We can test exactly which part of this assertion is an approximation.
Since $S_1^2=-I$, the rotation is
\begin{equation}
 e^{tS_1}=\begin{pmatrix}\cos t&-\sin t\\
                                   \sin t&\cos t\end{pmatrix}.
\end{equation}
Direct multiplication gives the \emph{exact} compressed entry
\begin{align}
 (P\widetilde H(t)P)_{11}
 &=\sin^2t-t\sin(2t)
   =-t^2+t^4+O(t^6),
\label{eq:ch21-sw-2x2-rotated}\end{align}
and the off-diagonal entry
\begin{equation}
 (\widetilde H(t))_{12}
 =t\cos(2t)-\tfrac12\sin(2t)
 =-\tfrac43t^3+O(t^5).
\label{eq:ch21-sw-2x2-off}
\end{equation}
Thus the first-order mixing is canceled but the chosen rotation is
not an exact block diagonalization for arbitrary $t$.  Independently,
the exact eigenvalues of $H_0+tV$ are
\begin{equation}
 E_\pm(t)=\frac{1\pm\sqrt{1+4t^2}}2.
\label{eq:ch21-sw-2x2-eigvals}
\end{equation}
At $t=2/3$, the lower eigenvalue is $-1/3$ (with eigenvector
$(-2,1)$), whereas \eqref{eq:ch21-sw-2x2-effective} gives $-4/9$.
The two values differ because a second-order effective operator is
not an exact finite-coupling spectral identification.
\end{example}

The same caution applies when the quadratic correction happens to
vanish.  A nonzero perturbation can couple retained and discarded
states through \emph{different} intermediate energies with cancelling
denominators.

\begin{example}[Zero second-order correction with nonzero leakage]
\label{ex:ch21-sw-three-level}
Take $\mathcal P=\operatorname{span}\{e_1\}$,
$\mathcal Q=\operatorname{span}\{e_2,e_3\}$ and
\begin{equation}
 H_0=\operatorname{diag}(0,-1,1),\quad
 V=\begin{pmatrix}0&1&1\\1&0&0\\1&0&0\end{pmatrix},\quad
 S_1=\begin{pmatrix}0&1&-1\\-1&0&0\\1&0&0\end{pmatrix}.
\label{eq:ch21-sw-three-input}
\end{equation}
Here $PVQ\ne0$ and $[S_1,H_0]=-V$, but
\begin{equation}
\langle e_1\lvert W^{(2)}\rvert e_1\rangle
 =\frac{|V_{12}|^2}{0-(-1)}
  +\frac{|V_{13}|^2}{0-1}=1-1=0.
\label{eq:ch21-sw-three-cancel}
\end{equation}
A vanishing compressed second-order commutator therefore does
\emph{not} imply $PVQ=QVP=0$.  Conversely, if both off-block
couplings vanish, choosing $S_1=0$ gives the exact identity-rotation
compression $P(H_0+tV)P=PH_0P+tPVP$ and zero correction.  That
forward implication does not require a non-resonance hypothesis.
\end{example}

\subsection{Scope of the finite construction and the remaining chapter obligations}
\label{subsec:ch21-effective-scope}

The discrete procedure has three mathematically independent inputs:
(i) a precisely selected coordinate budget and its orthogonal shell,
(ii) an invariant baseline $H_0$ and a cross-block perturbation $V$,
and (iii) control of coupled energy denominators.  Earlier sections
provided exact finite CAR words, charge-selection rules, and conditional
leakage witnesses for spinful perturbations.  Those statements do not
by themselves furnish the spectral decomposition, denominator bounds,
or block invariance needed to insert a particular interacting
Hamiltonian into the Schrieffer--Wolff theorem.

The separately verified finite support results cover the non-resonant
commutator equation and skew-Hermiticity of the constructed generator,
two- and three-level matrix witnesses, the compressed \emph{degree-two
BCH polynomial} identity, the decomposition of shell projections and
their orthogonality and adjoint properties, and a nonzero repeated-mode-one
shell witness. One projector identity is $P+Q=\Pi_{\boldsymbol N,K}$ on the
ambient Fock carrier, whereas the generic compressed BCH theorem assumes
$P+Q=I$ on its whole carrier. Restricting to
$\mathcal B_{\boldsymbol N,K}$ and transferring an actual model's
operators and hypotheses is a separate bridge, not established merely
by placing the two statements next to each other.  These finite support
proofs do not constitute formal verification of the present Chapter~21
manuscript.

The remaining constructive tasks include proving, for any chosen
spinful Hamiltonian, that its finite blocks act on the proposed
carriers; checking charge-sector and energy-margin constraints at
each intermediate application; constructing a valid $H_0$ and its
adapted eigenbasis or another certified Sylvester solver; and proving
any claimed uniform denominator or remainder bounds.  The
coupling-wise condition is conditional, and a vanishing second-order
term is not a substitute for checking it.  Furthermore, a Mott or
spin gap would require an independently defined spectral gap and a
proof of its behavior under relevant size or coupling limits.
A continuum RG flow requires a scale-changing construction and a
comparison between effective theories at successive scales; energy
shell projection at one finite $K$ is not such a construction.
Finally, an isolated finite Hermitian Hamiltonian has a real discrete
spectrum, so scattering lifetimes or dissipative flow cannot be
inferred solely from the matrix identities above.

What has been established here is therefore deliberately narrower
and more precise: a well-defined \emph{finite energy-shell filter}, a
non-resonant \emph{first-order unitary generator}, an exact
\emph{second-order compressed polynomial identity}, and a finite
analytic remainder estimate with explicit hypotheses.  These results
provide the operator-level starting point for a future controlled
effective-Hamiltonian analysis, without prematurely identifying its
outcome with a charge gap or a continuum renormalization-group law.


\appendix
\part{Supporting Mathematical Appendices}

\chapter{Fourier Characters and Algebraic Normalization}
\label{app:fourier_characters}

\section{Introduction and Scope}

The discrete Fourier transform on a finite lattice $\Lambda = \mathbb{Z}/L\mathbb{Z}$ forms the kinematic bridge between position-space fermions and momentum-space modes. In Chapter~\ref{ch:fourier}, the Fourier transform was formulated over the complex field $\mathbb{C}$ using the canonical primitive root $\zeta = e^{2\pi i / L}$.

However, from an algebraic and foundational perspective, several critical questions arise:
\begin{enumerate}
    \item \emph{What is the minimal algebraic structure required for discrete Fourier duality?} Does character orthogonality hold over general commutative rings, or are integral domains strictly necessary?
    \item \emph{How do we rigorously decouple the algebraic character identities from the metric normalization $1/\sqrt{L}$?} In pure algebra, square roots of integers may not exist, yet exact two-sided Fourier inversion remains valid whenever $L$ is invertible.
    \item \emph{How is representative independence reconciled with the centered momentum band $\Lambda^*$?} The momentum band $\Lambda^* = \{-h+1, \dots, h\}$ is a set of signed integers, not a group under ordinary integer addition.
\end{enumerate}

This appendix provides the complete algebraic answers to these questions. We develop the three-tiered character architecture, prove geometric sum cancellation over arbitrary fields, establish the unscaled operator identities $T S = L I$ and $S T = L I$, and compute the exact lattice dispersion algebraically without relying on transcendental trigonometric functions.

\section{The Minimal Ring Constraint and a Zero-Divisor Counterexample}

It is tempting to develop discrete Fourier analysis over an arbitrary commutative ring $R$ equipped with an element $\zeta \in R$ of multiplicative order $L$. However, this generality fails in the presence of zero divisors.

\begin{example}[Failure of Character Orthogonality over Zero Divisors]
Consider the commutative ring $R = \mathbb{Z}/15\mathbb{Z}$ with lattice size $L = 2$.
The element $\zeta = 4 \in \mathbb{Z}/15\mathbb{Z}$ satisfies:
\begin{equation}
\zeta^1 = 4 \not\equiv 1 \pmod{15}, \qquad \zeta^2 = 16 \equiv 1 \pmod{15}.
\end{equation}
Thus $\zeta = 4$ has exact multiplicative order $L = 2$. Furthermore, the lattice size $L = 2$ is an invertible unit in $R$, since $2 \times 8 = 16 \equiv 1 \pmod{15}$.
However, evaluating the geometric character sum for the non-trivial frequency $k = 1$ yields:
\begin{equation}
\sum_{x=0}^{1} \zeta^{1 \cdot x} = \zeta^0 + \zeta^1 = 1 + 4 = 5 \not\equiv 0 \pmod{15}.
\end{equation}
Even though $L$ is invertible and $\zeta$ has order $L$, the character sum fails to vanish because $(\zeta - 1) = 3$ is a zero divisor in $\mathbb{Z}/15\mathbb{Z}$, and $3 \times 5 = 15 \equiv 0$.
\end{example}

Consequently, character orthogonality cannot be established over general commutative rings. Throughout this appendix, we work over an arbitrary \emph{field} $K$ (or an integral domain) of characteristic not dividing $L$, containing a primitive $L$-th root of unity $\zeta$.

\section{The Three-Tier Character Architecture}

To handle the signed momentum band $\Lambda^* = \{-h+1, \dots, h\}$ without imposing an artificial group law, we define Fourier characters across three distinct algebraic tiers.

\begin{definition}[Tier 1: Integer Character Pairing]
Let $K$ be a field and let $\zeta \in K^\times$ be an element of multiplicative order $L$. The \emph{integer character pairing} is the map $\chi_{\mathbb{Z}}: \mathbb{Z} \times \mathbb{Z} \to K^\times$ defined by:
\begin{equation}
\chi_{\mathbb{Z}}(\zeta, k, x) \coloneqq \zeta^{k x}.
\end{equation}
Because $\zeta \neq 0$, the power $\zeta^{k x}$ is well-defined for all signed integers $k, x \in \mathbb{Z}$ via field inversion $\zeta^{-1}$ for negative exponents.
\end{definition}

\begin{lemma}[Representative Independence]
\label{lem:a01_rep_independence}
Let $\zeta \in K$ be a primitive $L$-th root of unity (so $\zeta^L = 1$). For all integers $k, x, s, t \in \mathbb{Z}$:
\begin{equation}
\chi_{\mathbb{Z}}(\zeta, k + s L, \; x + t L) = \chi_{\mathbb{Z}}(\zeta, k, x).
\end{equation}
\end{lemma}

\begin{proof}
Expanding the exponent of $\zeta$:
\begin{align*}
(k + s L)(x + t L) &= k x + k t L + s x L + s t L^2 \\
&= k x + L (k t + s x + s t L).
\end{align*}
Since $\zeta \neq 0$, exponentiation respects integer addition:
\[
\zeta^{(k + s L)(x + t L)} = \zeta^{k x} \left(\zeta^L\right)^{k t + s x + s t L} = \zeta^{k x} (1)^{k t + s x + s t L} = \zeta^{k x}.
\]
\end{proof}

\begin{definition}[Tier 2: Residue Character]
For residue classes $k, x \in \mathbb{Z}/L\mathbb{Z}$ with canonical unsigned representatives $k_{\mathrm{val}}, x_{\mathrm{val}} \in \{0, 1, \dots, L-1\}$, the \emph{residue character} is:
\begin{equation}
\chi_{\mathrm{res}}(L, \zeta, k, x) \coloneqq \chi_{\mathbb{Z}}(\zeta, k_{\mathrm{val}}, x_{\mathrm{val}}).
\end{equation}
By Lemma~\ref{lem:a01_rep_independence}, $\chi_{\mathrm{res}}(L, \zeta, k, x)$ is independent of the choice of integer representatives for $k$ and $x$.
\end{definition}

\begin{definition}[Tier 3: Band Character]
For a signed momentum mode $k \in \Lambda^*$ and a lattice site $x \in \Lambda = \mathbb{Z}/L\mathbb{Z}$, the \emph{band character} is:
\begin{equation}
\chi_{\mathrm{band}}(L, \zeta, k, x) \coloneqq \chi_{\mathbb{Z}}(\zeta, k, x_{\mathrm{val}}) = \chi_{\mathrm{res}}(L, \zeta, \pi(k), x),
\end{equation}
where $\pi: \Lambda^* \to \mathbb{Z}/L\mathbb{Z}$ is the canonical projection $\pi(k) = k \bmod L$.
\end{definition}

\begin{theorem}[Character Homomorphism Laws]
\label{thm:character_homomorphisms}
Let $\zeta \in K$ be a primitive $L$-th root of unity. The character pairing satisfies:
\begin{enumerate}
    \item \textbf{Spatial Homomorphism:} For all $k \in \mathbb{Z}/L\mathbb{Z}$ and $x, y \in \mathbb{Z}/L\mathbb{Z}$,
    \begin{equation}
    \chi_{\mathrm{res}}(k, x + y) = \chi_{\mathrm{res}}(k, x) \chi_{\mathrm{res}}(k, y).
    \end{equation}
    \item \textbf{Momentum Homomorphism:} For all $k, p \in \mathbb{Z}/L\mathbb{Z}$ and $x \in \mathbb{Z}/L\mathbb{Z}$,
    \begin{equation}
    \chi_{\mathrm{res}}(k + p, x) = \chi_{\mathrm{res}}(k, x) \chi_{\mathrm{res}}(p, x).
    \end{equation}
    \item \textbf{Inversion:} $\chi_{\mathrm{res}}(-k, x) = \chi_{\mathrm{res}}(k, x)^{-1} = \chi_{\mathrm{res}}(k, -x)$.
    \item \textbf{Transported Band Addition:} For $k, p \in \Lambda^*$, let $k \oplus p \in \Lambda^*$ denote the transported band addition defined in Chapter~\ref{ch:lattice_band}. Then:
    \begin{equation}
    \chi_{\mathrm{band}}(k \oplus p, x) = \chi_{\mathrm{band}}(k, x) \chi_{\mathrm{band}}(p, x).
    \end{equation}
    \item \textbf{Complex Conjugation:} If $K = \mathbb{C}$ and $|\zeta| = 1$, then $\overline{\chi(k, x)} = \chi(-k, x)$.
\end{enumerate}
\end{theorem}

\begin{proof}
Items 1--3 follow directly from Lemma~\ref{lem:a01_rep_independence} and the index laws $\zeta^{k(x+y)} = \zeta^{kx}\zeta^{ky}$.
For item 4: by Theorem~\ref{thm:wrap_lemma}, $k \oplus p = k + p - w L$ where $w \in \{-1, 0, 1\}$.
By Lemma~\ref{lem:a01_rep_independence}:
\[
\chi_{\mathrm{band}}(k \oplus p, x) = \zeta^{(k + p - w L) x_{\mathrm{val}}} = \zeta^{(k+p) x_{\mathrm{val}}} = \zeta^{k x_{\mathrm{val}}} \zeta^{p x_{\mathrm{val}}} = \chi_{\mathrm{band}}(k, x) \chi_{\mathrm{band}}(p, x).
\]
For item 5: if $|\zeta| = 1$, then $\overline{\zeta} = \zeta^{-1}$, so $\overline{\zeta^{kx}} = (\zeta^{-1})^{kx} = \zeta^{-kx} = \chi(-k, x)$.
\end{proof}

\section{Geometric Sum Cancellation and Dual Orthogonality}

\begin{lemma}[Non-Trivial Character Cancellation]
\label{lem:nontrivial_character_sum}
Let $K$ be a field, and let $\zeta \in K$ be a primitive $L$-th root of unity. For any residue $k \in \mathbb{Z}/L\mathbb{Z}$:
\begin{equation}
\sum_{x \in \mathbb{Z}/L\mathbb{Z}} \chi_{\mathrm{res}}(k, x) = \begin{cases} L \cdot 1_K, & \text{if } k = 0, \\ 0, & \text{if } k \neq 0. \end{cases}
\end{equation}
\end{lemma}

\begin{proof}
If $k = 0$, then $\chi_{\mathrm{res}}(0, x) = \zeta^0 = 1$ for all $x$. The sum of $L$ copies of $1_K$ is $L \cdot 1_K$.

Now assume $k \neq 0$. Since $\zeta$ is a primitive $L$-th root of unity and $k \not\equiv 0 \pmod L$, $\zeta^{k_{\mathrm{val}}} \neq 1$.
Consider the finite sum $S = \sum_{x=0}^{L-1} (\zeta^{k_{\mathrm{val}}})^x$. Multiplying by $(\zeta^{k_{\mathrm{val}}} - 1)$:
\[
(\zeta^{k_{\mathrm{val}}} - 1) S = \sum_{x=0}^{L-1} (\zeta^{k_{\mathrm{val}}})^{x+1} - \sum_{x=0}^{L-1} (\zeta^{k_{\mathrm{val}}})^x = (\zeta^{k_{\mathrm{val}}})^L - 1 = (\zeta^L)^{k_{\mathrm{val}}} - 1 = 1^{k_{\mathrm{val}}} - 1 = 0.
\]
Because $K$ is a field and $\zeta^{k_{\mathrm{val}}} - 1 \neq 0$, we can divide by $(\zeta^{k_{\mathrm{val}}} - 1)$ to conclude $S = 0$.
\end{proof}

\begin{theorem}[Spatial and Momentum Orthogonality]
\label{thm:app_orthogonality}
Let $K$ be a field, $\zeta \in K$ a primitive $L$-th root of unity, and $h = L/2 \ge 1$.
\begin{enumerate}
    \item \textbf{Spatial Orthogonality on $\Lambda^*$:} For any two momentum modes $k, p \in \Lambda^*$:
    \begin{equation}
    \sum_{x \in \Lambda} \chi_{\mathbb{Z}}(\zeta, k - p, x_{\mathrm{val}}) = L \delta_{k, p} \cdot 1_K.
    \end{equation}
    \item \textbf{Momentum Orthogonality on $\Lambda$:} For any two lattice sites $x, y \in \Lambda$:
    \begin{equation}
    \sum_{k \in \Lambda^*} \chi_{\mathbb{Z}}(\zeta, k, (x - y)_{\mathrm{val}}) = L \delta_{x, y} \cdot 1_K.
    \end{equation}
\end{enumerate}
\end{theorem}

\begin{proof}
\begin{enumerate}
    \item By Theorem~\ref{thm:character_homomorphisms}, $\chi_{\mathbb{Z}}(\zeta, k - p, x_{\mathrm{val}}) = \chi_{\mathrm{res}}(\pi(k) - \pi(p), x)$.
    By Lemma~\ref{lem:nontrivial_character_sum}, the sum over $x$ equals $L \cdot 1_K$ if $\pi(k) = \pi(p)$, and $0$ otherwise.
    Since the projection $\pi: \Lambda^* \to \mathbb{Z}/L\mathbb{Z}$ is a bijection with inverse $r$ (Lemma~\ref{lem:section_identities}), $\pi(k) = \pi(p)$ if and only if $k = p$.
    \item Using the bijection $\pi: \Lambda^* \to \mathbb{Z}/L\mathbb{Z}$, summing $k$ over $\Lambda^*$ is identical to summing the residue $\kappa$ over $\mathbb{Z}/L\mathbb{Z}$.
    Applying Lemma~\ref{lem:nontrivial_character_sum} with role of $k$ and $x$ exchanged yields $L \delta_{x, y} \cdot 1_K$.
\end{enumerate}
\end{proof}

\section{Unscaled Transforms and Two-Sided Inversion}

\begin{definition}[Unscaled Forward and Backward Transforms]
Let $K^\Lambda$ and $K^{\Lambda^*}$ denote the $K$-vector spaces of functions on the lattice and band, respectively.
\begin{enumerate}
    \item The \emph{unscaled analysis map} $S: K^\Lambda \to K^{\Lambda^*}$ is defined by:
    \begin{equation}
    (S f)(k) \coloneqq \sum_{x \in \Lambda} f(x) \chi_{\mathbb{Z}}(\zeta, -k, x_{\mathrm{val}}).
    \end{equation}
    \item The \emph{unscaled synthesis map} $T: K^{\Lambda^*} \to K^\Lambda$ is defined by:
    \begin{equation}
    (T g)(x) \coloneqq \sum_{k \in \Lambda^*} g(k) \chi_{\mathbb{Z}}(\zeta, k, x_{\mathrm{val}}).
    \end{equation}
\end{enumerate}
\end{definition}

\begin{theorem}[Two-Sided Unscaled Inversion]
\label{thm:a01_unscaled_inversion}
The linear maps $S$ and $T$ satisfy:
\begin{equation}
T \circ S = L \cdot I_{K^\Lambda}, \qquad S \circ T = L \cdot I_{K^{\Lambda^*}}.
\end{equation}
Whenever $L \cdot 1_K \neq 0$ in $K$, $S$ is a linear isomorphism with inverse $S^{-1} = L^{-1} T$.
\end{theorem}

\begin{proof}
Let $f \in K^\Lambda$ and evaluate $(T(Sf))(y)$ for $y \in \Lambda$:
\begin{align*}
(T(Sf))(y) &= \sum_{k \in \Lambda^*} (Sf)(k) \chi(k, y) = \sum_{k \in \Lambda^*} \left( \sum_{x \in \Lambda} f(x) \chi(-k, x) \right) \chi(k, y) \\
&= \sum_{x \in \Lambda} f(x) \left( \sum_{k \in \Lambda^*} \chi(k, y - x) \right).
\end{align*}
By Theorem~\ref{thm:app_orthogonality}(2), the inner sum is $L \delta_{x, y} \cdot 1_K$.
Therefore:
\[
(T(Sf))(y) = \sum_{x \in \Lambda} f(x) (L \delta_{x, y} \cdot 1_K) = L f(y).
\]
Hence $T \circ S = L I$. A completely symmetric calculation exchanging the order of summation and using Theorem~\ref{thm:app_orthogonality}(1) proves $S \circ T = L I$.
\end{proof}

\section{Exact Dispersion Algebra in Arbitrary Fields}

In physical lattices, the nearest-neighbor discrete Laplacian is given by $(\Delta f)(x) = f(x+1) + f(x-1) - 2 f(x)$.
Its eigenvalue on plane-wave modes is determined purely algebraically without trigonometry.

\begin{lemma}[Exact Second-Difference Eigenvalue]
\label{lem:laplacian_algebraic}
Let $K$ be a field and $\zeta \in K^\times$. For any integer $k \in \mathbb{Z}$:
\begin{equation}
(\zeta^k - 1)(1 - \zeta^{-k}) = \zeta^k + \zeta^{-k} - 2.
\end{equation}
\end{lemma}

\begin{proof}
Expanding the left-hand side:
\[
(\zeta^k - 1)(1 - \zeta^{-k}) = \zeta^k \cdot 1 - \zeta^k \zeta^{-k} - 1 \cdot 1 + 1 \cdot \zeta^{-k} = \zeta^k - 1 - 1 + \zeta^{-k} = \zeta^k + \zeta^{-k} - 2.
\]
\end{proof}

\begin{proposition}[Connection to Standard Trigonometric Dispersion]
Specializing to $K = \mathbb{C}$ with the canonical root $\zeta = e^{2\pi i / L}$:
\begin{equation}
\zeta^k + \zeta^{-k} - 2 = - 4 \sin^2\left(\frac{\pi k}{L}\right).
\end{equation}
\end{proposition}

\begin{proof}
Using Euler's formula:
\[
\zeta^k + \zeta^{-k} - 2 = 2 \cos\left(\frac{2\pi k}{L}\right) - 2 = -2 \left(1 - \cos\left(\frac{2\pi k}{L}\right)\right).
\]
Using the standard double-angle identity $1 - \cos(2\theta) = 2 \sin^2(\theta)$ with $\theta = \frac{\pi k}{L}$ yields:
\[
-2 \left(2 \sin^2\left(\frac{\pi k}{L}\right)\right) = - 4 \sin^2\left(\frac{\pi k}{L}\right).
\]
\end{proof}

\begin{remark}[Formalization Note]
In the verified Lean~4 core (Appendix~A01), all character homomorphisms, the representative-independence theorem, the geometric sum cancellation, and the unscaled operator identities are proved without noncomputable choice over generic fields. The physical normalization $a = 1/\sqrt{L}$ is introduced only in the complex Hilbert tier to construct the unitary isometry $U: \ell^2(\Lambda) \to \ell^2(\Lambda^*)$.
\end{remark}

\chapter{CAR Hilbert Carriers and Sign Involutions}
\label{app:car_hilbert}

\section{Introduction and the Two-Tiered Architecture}

The canonical anticommutation relations (CAR) govern the algebra of fermionic creation and annihilation operators. In physics textbooks, algebraic anticommutation relations and Hilbert-space adjoints are frequently introduced simultaneously. However, in constructive mathematics and formal verification, conflating these two aspects introduces subtle vulnerabilities:
\begin{enumerate}
    \item \textbf{Non-Unitary Similarities:} Let $c_i, c_i^\dagger$ satisfy the algebraic CAR on a vector space $V$. For any invertible linear transformation $S \in \operatorname{GL}(V)$, the similarity-transformed operators:
    \begin{equation}
    \tilde{c}_i \coloneqq S c_i S^{-1}, \qquad \tilde{c}_i^\dagger \coloneqq S c_i^\dagger S^{-1}
    \end{equation}
    satisfy all algebraic anticommutation relations identically:
    \[
    \{\tilde{c}_i, \tilde{c}_j\} = 0, \qquad \{\tilde{c}_i^\dagger, \tilde{c}_j^\dagger\} = 0, \qquad \{\tilde{c}_i, \tilde{c}_j^\dagger\} = \delta_{ij} I.
    \]
    However, unless $S$ is unitary, $\tilde{c}_i^\dagger \neq (\tilde{c}_i)^\dagger$. Thus, algebraic CAR does not imply adjoint compatibility.
    \item \textbf{Carrier Metric Distinctions:} The function carrier $X \to \mathbb{C}$ carries a natural product norm, but not an intrinsic Euclidean $\ell^2$ inner product. Conflating function spaces with Euclidean spaces risks installing incompatible metric instances.
\end{enumerate}

To resolve these issues, we establish a strict two-tiered algebraic hierarchy:
\begin{itemize}
    \item \textbf{Tier 1: Abstract Algebraic CAR,} defined over an arbitrary complex vector space $V$ without any metric or topological structure.
    \item \textbf{Tier 2: Hilbert-Space CAR,} defined over finite-dimensional complex Euclidean spaces, adding the true Hilbert adjoint condition $c_i^\dagger = (c_i)^\dagger$.
\end{itemize}

\section{Tier 1: Abstract Algebraic CAR Relations}

\begin{definition}[Algebraic CAR System]
Let $V$ be a complex vector space and let $\iota$ be a finite index set. An \emph{algebraic CAR system} on $(V, \iota)$ consists of two operator families $c, c^\dagger : \iota \to \operatorname{End}_{\mathbb{C}}(V)$ satisfying:
\begin{enumerate}
    \item $\{c_i, c_j\} = c_i c_j + c_j c_i = 0$ for all $i, j \in \iota$.
    \item $\{c_i^\dagger, c_j^\dagger\} = c_i^\dagger c_j^\dagger + c_j^\dagger c_i^\dagger = 0$ for all $i, j \in \iota$.
    \item $\{c_i, c_j^\dagger\} = c_i c_j^\dagger + c_j^\dagger c_i = \delta_{ij} I_V$ for all $i, j \in \iota$.
\end{enumerate}
The local number operator for mode $i \in \iota$ is defined algebraically by:
\begin{equation}
n_i \coloneqq c_i^\dagger c_i.
\end{equation}
\end{definition}

Without assuming any inner product, the algebraic relations alone imply fundamental commutation theorems.

\begin{theorem}[CAR Bilinear Commutator Identity]
\label{thm:app_bilinear_commutator}
Let $(c, c^\dagger)$ be an algebraic CAR system on $(V, \iota)$. For any four indices $a, b, c, d \in \iota$:
\begin{equation}
\left[ c_a^\dagger c_b, \; c_c^\dagger c_d \right] = \delta_{bc} c_a^\dagger c_d - \delta_{ad} c_c^\dagger c_b.
\end{equation}
\end{theorem}

\begin{proof}
We evaluate the commutator $[A B, C D]$ for operators satisfying $\{B, C\} = \delta_{bc} I$, $\{D, A\} = \delta_{ad} I$, $\{A, C\} = 0$, and $\{B, D\} = 0$.
By non-commutative ring algebra:
\begin{align*}
[c_a^\dagger c_b, \; c_c^\dagger c_d] &= c_a^\dagger c_b c_c^\dagger c_d - c_c^\dagger c_d c_a^\dagger c_b \\
&= c_a^\dagger (\delta_{bc} I - c_c^\dagger c_b) c_d - c_c^\dagger (\delta_{ad} I - c_a^\dagger c_d) c_b \\
&= \delta_{bc} c_a^\dagger c_d - \delta_{ad} c_c^\dagger c_b - c_a^\dagger c_c^\dagger c_b c_d + c_c^\dagger c_a^\dagger c_d c_b.
\end{align*}
Using $\{c_a^\dagger, c_c^\dagger\} = 0$ and $\{c_b, c_d\} = 0$:
\[
c_c^\dagger c_a^\dagger c_d c_b = (- c_a^\dagger c_c^\dagger) (- c_b c_d) = c_a^\dagger c_c^\dagger c_b c_d.
\]
The last two terms cancel identically:
\[
- c_a^\dagger c_c^\dagger c_b c_d + c_a^\dagger c_c^\dagger c_b c_d = 0.
\]
This leaves exactly $\delta_{bc} c_a^\dagger c_d - \delta_{ad} c_c^\dagger c_b$.
\end{proof}

\begin{corollary}[Number Operator Commutators and Idempotence]
\label{cor:app_number_comm}
For all modes $i, j \in \iota$:
\begin{enumerate}
    \item $[n_i, c_j^\dagger] = \delta_{ij} c_j^\dagger$.
    \item $[n_i, c_j] = - \delta_{ij} c_j$.
    \item $[n_i, n_j] = 0$.
    \item $n_i^2 = n_i$ (idempotence).
\end{enumerate}
\end{corollary}

\begin{proof}
\begin{enumerate}
    \item Set $a = i, b = i, c = j$ in the algebraic expansion:
    \[
    [c_i^\dagger c_i, c_j^\dagger] = c_i^\dagger \{c_i, c_j^\dagger\} - \{c_i^\dagger, c_j^\dagger\} c_i = c_i^\dagger (\delta_{ij} I) - 0 = \delta_{ij} c_j^\dagger.
    \]
    \item Similarly, $[c_i^\dagger c_i, c_j] = c_i^\dagger \{c_i, c_j\} - \{c_i^\dagger, c_j\} c_i = 0 - (\delta_{ij} I) c_i = - \delta_{ij} c_i$.
    \item Setting $a = b = i$ and $c = d = j$ in Theorem~\ref{thm:app_bilinear_commutator}:
    \[
    [n_i, n_j] = [c_i^\dagger c_i, c_j^\dagger c_j] = \delta_{ij} c_i^\dagger c_j - \delta_{ij} c_j^\dagger c_i = \delta_{ij} (n_i - n_i) = 0.
    \]
    \item Setting $i = j$ in the anticommutator $\{c_i, c_i^\dagger\} = I$:
    \[
    n_i^2 = (c_i^\dagger c_i)(c_i^\dagger c_i) = c_i^\dagger (I - c_i^\dagger c_i) c_i = c_i^\dagger c_i - (c_i^\dagger)^2 c_i^2.
    \]
    Since $2 (c_i^\dagger)^2 = \{c_i^\dagger, c_i^\dagger\} = 0$, we have $(c_i^\dagger)^2 = 0$, so $n_i^2 = c_i^\dagger c_i = n_i$.
\end{enumerate}
\end{proof}

\section{Tier 2: Concrete Occupation Hilbert Spaces}

We now construct the concrete representation of CAR on the fermionic Fock space.

\begin{definition}[Fermionic Fock Space]
Let $\iota$ be a finite index set of cardinality $|\iota| = d$. The \emph{fermionic Fock space} $\mathcal{F}(\iota)$ is the complex Euclidean space:
\begin{equation}
\mathcal{F}(\iota) \coloneqq \ell^2(\mathcal{P}(\iota)),
\end{equation}
where $\mathcal{P}(\iota)$ is the power set of all subsets of $\iota$.
The canonical orthonormal basis vectors are the occupation kets:
\begin{equation}
|S\rangle \in \mathcal{F}(\iota) \quad \text{for each } S \subseteq \iota.
\end{equation}
\end{definition}

\begin{proposition}[Carrier Properties]
\label{prop:fock_carrier_properties}
\begin{enumerate}
    \item \textbf{Finite Dimension:} $\dim_{\mathbb{C}} \mathcal{F}(\iota) = 2^{|\iota|}$.
    \item \textbf{Inner Product:} For $u, v \in \mathcal{F}(\iota)$, $\langle u \mid v \rangle = \sum_{S \subseteq \iota} \overline{u(S)} v(S)$.
    \item \textbf{Orthonormality:} $\langle S \mid T \rangle = \delta_{S, T}$ for all $S, T \subseteq \iota$.
    \item \textbf{Basis Resolution:} Every vector $v \in \mathcal{F}(\iota)$ expands uniquely as $v = \sum_{S \subseteq \iota} v(S) |S\rangle$.
\end{enumerate}
\end{proposition}

\begin{lemma}[Adjoint Verification via Basis Pairings]
\label{lem:adjoint_pairing}
Let $A, B \in \operatorname{End}_{\mathbb{C}}(\mathcal{F}(\iota))$ be linear operators. If the inner product pairing satisfies:
\begin{equation}
\langle S \mid A T \rangle = \langle B S \mid T \rangle \quad \text{for all basis pairs } S, T \subseteq \iota,
\end{equation}
then $B$ is identically the Hilbert adjoint of $A$:
\begin{equation}
B = A^\dagger.
\end{equation}
\end{lemma}

\begin{proof}
Let $u = \sum_S u(S) |S\rangle$ and $v = \sum_T v(T) |T\rangle$ be arbitrary vectors in $\mathcal{F}(\iota)$. By sesquilinearity:
\begin{align*}
\langle u \mid A v \rangle &= \sum_{S, T \subseteq \iota} \overline{u(S)} v(T) \langle S \mid A T \rangle \\
&= \sum_{S, T \subseteq \iota} \overline{u(S)} v(T) \langle B S \mid T \rangle = \langle B u \mid v \rangle.
\end{align*}
By the Riesz representation characterization of the adjoint, this proves $B = A^\dagger$.
\end{proof}

\section{Fermionic Sign Mechanics and Jordan--Wigner Ordering}

To construct the operators $c_i, c_i^\dagger$, we fix a strict linear order $<$ on $\iota$.

\begin{definition}[Preceding Occupation Count and Sign]
For any mode $i \in \iota$ and configuration $S \subseteq \iota$:
\begin{enumerate}
    \item The \emph{preceding occupied count} is the number of occupied modes strictly below $i$:
    \begin{equation}
    \sigma(i, S) \coloneqq \#\{ j \in S \mid j < i \}.
    \end{equation}
    \item The \emph{fermionic sign factor} is:
    \begin{equation}
    \epsilon(i, S) \coloneqq (-1)^{\sigma(i, S)} \in \{+1, -1\}.
    \end{equation}
\end{enumerate}
\end{definition}

To prevent natural number underflow in formal induction, the preceding count is governed by strictly additive recursion formulas.

\begin{lemma}[Additive Preceding Count Laws]
\label{lem:additive_count}
Let $i, j \in \iota$ and $S \subseteq \iota$.
\begin{enumerate}
    \item \textbf{Insertion Formula:} For $j \notin S$:
    \begin{equation}
    \sigma(i, S \cup \{j\}) = \sigma(i, S) + \begin{cases} 1, & \text{if } j < i, \\ 0, & \text{if } j \ge i. \end{cases}
    \end{equation}
    \item \textbf{Erasure Formula:} For $j \in S$:
    \begin{equation}
    \sigma(i, S) = \sigma(i, S \setminus \{j\}) + \begin{cases} 1, & \text{if } j < i, \\ 0, & \text{if } j \ge i. \end{cases}
    \end{equation}
    \item \textbf{Self Invariance:} $\sigma(i, S \cup \{i\}) = \sigma(i, S)$ and $\sigma(i, S \setminus \{i\}) = \sigma(i, S)$.
\end{enumerate}
\end{lemma}

\begin{proof}
The set $\{k \in S \cup \{j\} \mid k < i\}$ partitions into $\{k \in S \mid k < i\}$ and $\{j\}$ if $j < i$, or just $\{k \in S \mid k < i\}$ if $j \ge i$. Taking cardinalities yields the additive formula. The erasure formula follows by writing $S = (S \setminus \{j\}) \cup \{j\}$ and applying insertion.
\end{proof}

\begin{theorem}[Permutation Sign Exchange Identities]
\label{thm:app_sign_exchange}
For any two distinct modes $i \neq j$ in $\iota$:
\begin{enumerate}
    \item If $i, j \notin S$:
    \begin{equation}
    \epsilon(j, S) \epsilon(i, S \cup \{j\}) = - \epsilon(i, S) \epsilon(j, S \cup \{i\}).
    \end{equation}
    \item If $i, j \in S$:
    \begin{equation}
    \epsilon(j, S) \epsilon(i, S \setminus \{j\}) = - \epsilon(i, S) \epsilon(j, S \setminus \{i\}).
    \end{equation}
    \item If $i \notin S$ and $j \in S$:
    \begin{equation}
    \epsilon(j, S) \epsilon(i, S \setminus \{j\}) = - \epsilon(i, S) \epsilon(j, S \cup \{i\}).
    \end{equation}
\end{enumerate}
\end{theorem}

\begin{proof}
Without loss of generality, assume $j < i$.
\begin{enumerate}
    \item By Lemma~\ref{lem:additive_count}(1), since $j < i$, $\sigma(i, S \cup \{j\}) = \sigma(i, S) + 1$.
    On the other hand, since $i > j$, $\sigma(j, S \cup \{i\}) = \sigma(j, S) + 0 = \sigma(j, S)$.
    Therefore:
    \[
    \epsilon(j, S) \epsilon(i, S \cup \{j\}) = (-1)^{\sigma(j, S)} (-1)^{\sigma(i, S) + 1} = - (-1)^{\sigma(j, S) + \sigma(i, S)}.
    \]
    While the right-hand side has:
    \[
    \epsilon(i, S) \epsilon(j, S \cup \{i\}) = (-1)^{\sigma(i, S)} (-1)^{\sigma(j, S)} = (-1)^{\sigma(i, S) + \sigma(j, S)}.
    \]
    The ratio of the two products is $(-1)^1 = -1$.
    \item The proofs for items 2 and 3 follow identically from Lemma~\ref{lem:additive_count} by tracking the single relative shift $[j < i] = 1$.
\end{enumerate}
\end{proof}

\begin{definition}[Concrete Creation and Annihilation Operators]
The linear operators $c_i, c_i^\dagger \in \operatorname{End}_{\mathbb{C}}(\mathcal{F}(\iota))$ are defined on basis kets $|S\rangle$ by:
\begin{align}
c_i |S\rangle &\coloneqq \begin{cases} \epsilon(i, S) |S \setminus \{i\}\rangle, & \text{if } i \in S, \\ 0, & \text{if } i \notin S, \end{cases} \\
c_i^\dagger |S\rangle &\coloneqq \begin{cases} \epsilon(i, S) |S \cup \{i\}\rangle, & \text{if } i \notin S, \\ 0, & \text{if } i \in S. \end{cases}
\end{align}
\end{definition}

\begin{theorem}[Concrete CAR and Adjoint Verification]
\label{thm:app_concrete_car}
The operators $(c, c^\dagger)$ satisfy:
\begin{enumerate}
    \item $\{c_i, c_j\} = 0$ for all $i, j \in \iota$.
    \item $\{c_i^\dagger, c_j^\dagger\} = 0$ for all $i, j \in \iota$.
    \item $\{c_i, c_j^\dagger\} = \delta_{ij} I$ for all $i, j \in \iota$.
    \item $c_i^\dagger = (c_i)^\dagger$ for all $i \in \iota$.
\end{enumerate}
\end{theorem}

\begin{proof}
\begin{enumerate}
    \item For $i = j$: $c_i^2 |S\rangle = 0$ because $i \notin S \setminus \{i\}$.
    For $i \neq j$: if both $i, j \in S$, applying $c_i c_j$ yields $\epsilon(j, S) \epsilon(i, S \setminus \{j\}) |S \setminus \{i, j\}\rangle$.
    Applying $c_j c_i$ yields $\epsilon(i, S) \epsilon(j, S \setminus \{i\}) |S \setminus \{i, j\}\rangle$.
    By Theorem~\ref{thm:app_sign_exchange}(2), the two coefficients sum to zero. If either $i \notin S$ or $j \notin S$, both compositions vanish.
    \item Identical to item 1 using Theorem~\ref{thm:app_sign_exchange}(1).
    \item For $i \neq j$: follows from Theorem~\ref{thm:app_sign_exchange}(3).
    For $i = j$: if $i \in S$, $c_i^\dagger c_i |S\rangle = 1 |S\rangle$ and $c_i c_i^\dagger |S\rangle = 0$.
    If $i \notin S$, $c_i^\dagger c_i |S\rangle = 0$ and $c_i c_i^\dagger |S\rangle = 1 |S\rangle$.
    In both cases, $(c_i c_i^\dagger + c_i^\dagger c_i)|S\rangle = |S\rangle$.
    \item To verify adjointness, let $S, T \subseteq \iota$.
    $\langle S \mid c_i T \rangle$ is non-zero only when $i \in T$ and $S = T \setminus \{i\}$, in which case it equals $\epsilon(i, T)$.
    Conversely, $\langle c_i^\dagger S \mid T \rangle$ is non-zero only when $i \notin S$ and $S \cup \{i\} = T$, in which case it equals $\epsilon(i, S)$.
    By Lemma~\ref{lem:additive_count}(3), $\epsilon(i, T) = \epsilon(i, S \cup \{i\}) = \epsilon(i, S)$.
    Therefore $\langle S \mid c_i T \rangle = \langle c_i^\dagger S \mid T \rangle$.
    By Lemma~\ref{lem:adjoint_pairing}, $c_i^\dagger = (c_i)^\dagger$.
\end{enumerate}
\end{proof}

\section{Global Parity and Ordered Product Mechanics}

Because operator composition on $\operatorname{End}_{\mathbb{C}}(\mathcal{F}(\iota))$ is non-commutative, infinite or finite products of operators cannot be written without specifying a linear order.

\begin{definition}[Global Parity Operator]
Let $\operatorname{sort}(\le, \iota) = (i_1, i_2, \dots, i_d)$ be the sorted list of all modes in $\iota$. The \emph{global parity operator} $\Gamma$ is:
\begin{equation}
\Gamma \coloneqq \prod_{k=1}^d (I - 2 n_{i_k}).
\end{equation}
\end{definition}

\begin{theorem}[Parity Properties]
\label{thm:app_parity_properties}
\begin{enumerate}
    \item On basis kets: $\Gamma |S\rangle = (-1)^{|S|} |S\rangle$.
    \item Involution and Hermiticity: $\Gamma^2 = I$ and $\Gamma^\dagger = \Gamma$.
    \item Anticommutation with Fermions: $\{\Gamma, c_i\} = 0$ and $\{\Gamma, c_i^\dagger\} = 0$ for all $i \in \iota$.
\end{enumerate}
\end{theorem}

\begin{proof}
\begin{enumerate}
    \item Since $[n_i, n_j] = 0$, the product acts as the product of eigenvalues $(1 - 2 [i \in S])$. Each occupied mode contributes $(1 - 2) = -1$, while unoccupied modes contribute $(1 - 0) = 1$. The product is $(-1)^{|S|}$.
    \item $(-1)^{|S|} (-1)^{|S|} = 1$, so $\Gamma^2 = I$. Since eigenvalues are real, $\Gamma^\dagger = \Gamma$.
    \item $c_i$ maps $|S\rangle$ (parity $(-1)^{|S|}$) to $|S \setminus \{i\}\rangle$ (parity $(-1)^{|S|-1} = -(-1)^{|S|}$). Thus $\Gamma c_i = - c_i \Gamma$.
\end{enumerate}
\end{proof}

\begin{remark}[Formalization Note]
In the verified Lean~4 development (Appendix~A02), all sign exchange lemmas, the basis extension constructor, and the adjoint pairing theorems are verified constructively. In particular, the absence of natural underflow ensures that sign proofs do not require classical casing over integer signs.
\end{remark}

\chapter{Energy Budgets and Filtered Maps}
\label{app:energy_budgets}

\section{Overview and Purpose}
This appendix provides the technical framework for the energy budget regularization scheme utilized in Chapters~\ref{ch:vacuum_budget} through~\ref{ch:haldane_completeness}. It establishes integer excitation energy formulas, sector charge stability, and operator margin filtrations.

\section{Core Properties}
\begin{itemize}
    \item \textbf{Integer Energy Grading:} Proof that excitation energy is a strictly non-negative integer for all configurations in admissible sectors.
    \item \textbf{Sorted Particle Rank Formula:} Proof that sorted momentum indices $s_1 < s_2 < \dots < s_{h+N}$ uniquely determine the excitation energy relative to the sector sea ground.
    \item \textbf{Filtration Preservation:} Proof that operators with maximum energy shift $M$ map $B(N, K)$ into $B(N, K+M)$.
\end{itemize}

\begin{remark}[Formalization Note]
In the verified Lean~4 formalization (Appendix A03, Energy Budgets), all 114 lemmas on energy gradings, projections, and ground states are verified in the Lean 4 kernel.
\end{remark}

\chapter{Current Margins, Density Kinematics, and Partition Combinatorics}
\label{app:density_partitions}

\section{Overview and Purpose}
This appendix supplies the detailed combinatorial and algebraic lemmas underlying the Heisenberg algebra (Chapter~\ref{ch:heisenberg_algebra}) and Haldane completeness (Chapter~\ref{ch:haldane_completeness}). It covers margin tracking, non-wrapping density kinematics, and the particle-hole partition bijection.

\section{Core Properties}
\begin{itemize}
    \item \textbf{Non-Wrapping Kinematics:} Proof that density shifts $k \mapsto k \ominus m$ do not wrap around the Brillouin zone boundaries when operating on states satisfying the margin conditions.
    \item \textbf{Exact Edge Domains:} Explicit computation of the top and bottom edge domains for commutators.
    \item \textbf{Partition Bijection:} The bijective mapping between particle-hole excitation sequences and integer partitions $\lambda \vdash E$, establishing $\dim \mathcal{H}(N, E) = p(E)$.
\end{itemize}

\begin{remark}[Formalization Note]
In the verified Lean~4 formalization (Appendix A04, including Current Margins, Density Kinematics, and Partition Combinatorics), all margin bounds and partition bijections are verified without placeholders or nonstandard axioms.
\end{remark}


\backmatter

\begin{thebibliography}{99}
\addcontentsline{toc}{chapter}{Bibliography}

\bibitem{giamarchi2004}
T.~Giamarchi, \emph{Quantum Physics in One Dimension}, International Series of Monographs on Physics, Vol.~121, Oxford University Press, Oxford, 2004.

\bibitem{gogolin1998}
A.~O.~Gogolin, A.~A.~Nersesyan, and A.~M.~Tsvelik, \emph{Bosonization and Strongly Correlated Systems}, Cambridge University Press, Cambridge, 1998.

\bibitem{vondelft1998}
J.~von Delft and H.~Schoeller, \emph{Bosonization for Beginners --- Refermionization for Experts}, Annalen der Physik (Leipzig) \textbf{7} (1998), no.~4, 225--305.

\bibitem{haldane1981}
F.~D.~M.~Haldane, \emph{`Luttinger liquid theory' of one-dimensional quantum fluids. I. Properties of the Luttinger model and their extension to the general 1D interacting spinless Fermi gas}, Journal of Physics C: Solid State Physics \textbf{14} (1981), no.~19, 2585--2609.

\bibitem{haag1996}
R.~Haag, \emph{Local Quantum Physics: Fields, Particles, Algebras}, 2nd ed., Texts and Monographs in Physics, Springer-Verlag, Berlin, 1996.

\bibitem{luttinger1963}
J.~M.~Luttinger, \emph{An exactly soluble model of a many-fermion system}, Journal of Mathematical Physics \textbf{4} (1963), no.~9, 1154--1162.

\bibitem{mattis1965}
D.~C.~Mattis and E.~H.~Lieb, \emph{Exact solution of a many-fermion system and its associated boson field}, Journal of Mathematical Physics \textbf{6} (1965), no.~2, 304--312.

\bibitem{lieb1961}
E.~Lieb, T.~Schultz, and D.~Mattis, \emph{Two soluble models of an antiferromagnetic chain}, Annals of Physics \textbf{16} (1961), no.~3, 407--466.

\end{thebibliography}


\end{document}
